
% new file names from date 20171119: 
% block3_171119.tex
\renewcommand{\thisfile}{block3\_190328}   

\cleardoublepage
\mysection{DG and Triangulated Bifunctors} \label{sec:DG-tri-bifun}

\AYcopyright 
 
In this section we extend the theory of triangulated derived functors, that was 
studied in Section \ref{sec:der-funcs}, to  
{\em triangulated bifunctors}. As stated in Convention \ref{conv:2490}, all 
linear structures and operations (such as categories and functors) are  
$\K$-linear, where $\K$ is a fixed commutative base 
ring. The symbol $\ot$ means $\ot_{\K}$. 



\mysubsection{DG Bifunctors} \label{subsec:dg-bifun}
We already talked about bifunctors in Subsection \ref{subsec:bifunc}. That 
was for categories without further structure. Here we will consider 
DG categories, and matters become a bit more complicated. 

As a warmup, let us begin with linear bifunctors. 

\begin{dfn} \label{dfn:5040}
Let $\cat{C}_1$ and $\cat{C}_2$ be linear categories. 
We define the linear category $\cat{C}_1 \ot \cat{C}_2$ as follows. The 
set of objects is 
\[ \opn{Ob}(\cat{C}_1 \ot \cat{C}_2) :=
\opn{Ob}(\cat{C}_1) \times \opn{Ob}(\cat{C}_2) . \]
For each pair of objects 
$(M_1, M_2), (N_1, N_2)$ in $\cat{C}_1 \ot \cat{C}_2$,
i.e.\ $M_i, N_i \in \opn{Ob}(\cat{C}_i)$,
we let 
\[ \opn{Hom}_{\cat{C}_1 \ot \cat{C}_2}
\bigl( (M_1, M_2), (N_1, N_2) \bigr) := 
\opn{Hom}_{\cat{C}_1} (M_1, N_1) \ot 
\opn{Hom}_{\cat{C}_2} (M_2, N_2) . \]

Given morphisms 
$\phi_i \in \opn{Hom}_{\cat{C}_i}(L_i, M_i)$
and
$\psi_i \in \opn{Hom}_{\cat{C}_i}(M_i, N_i)$, 
for $i = 1, 2$, their tensors are morphisms 
\[ \phi_1 \ot \phi_2 \in \opn{Hom}_{\cat{C}_1 \ot \cat{C}_2}
\bigl( (L_1, L_2), (M_1, M_2) \bigr) \]
and 
\[ \psi_1 \ot \psi_2 \in 
\opn{Hom}_{\cat{C}_1 \ot \cat{C}_2}
\bigl( (M_1, M_2), (N_1, N_2) \bigr) . \]
Every morphism in $\cat{C}_1 \ot \cat{C}_2$ is a sum of such pure tensors. We 
define the composition to be 
\[ \begin{aligned}
& (\psi_1 \ot \psi_2) \circ (\phi_1 \ot \phi_2) := 
(\psi_1 \circ \phi_1) \ot (\psi_2 \circ \phi_2)
\\ & \quad 
\in \opn{Hom}_{\cat{C}_1 \ot \cat{C}_2}
\bigl( (L_1, L_2), (N_1, N_2) \bigr) .
\end{aligned} \]
\end{dfn}

\begin{dfn} \label{dfn:1965}
Let $\cat{C}_1$, $\cat{C}_2$ and $\cat{D}$ be linear categories. A 
{\em linear bifunctor}% 
\index{Bifunctor! linear}
\[ F : \cat{C}_1 \times \cat{C}_2 \to \cat{D} \]
is, by definition, a linear functor 
$F : \cat{C}_1 \ot \cat{C}_2 \to \cat{D}$.
\end{dfn}

Now to the DG situation. 

\begin{dfn} \label{dfn:2010}
Let $\cat{C}_1$ and $\cat{C}_2$ be DG categories. 
We define the DG category $\cat{C}_1 \ot \cat{C}_2$ as follows: the 
set of objects is 
\[ \opn{Ob}(\cat{C}_1 \ot \cat{C}_2) :=
\opn{Ob}(\cat{C}_1) \times \opn{Ob}(\cat{C}_2) . \]
For each pair of objects 
$(M_1, M_2), (N_1, N_2)$ in $\cat{C}_1 \ot \cat{C}_2$
we let 
\[ \opn{Hom}_{\cat{C}_1 \ot \cat{C}_2}
\bigl( (M_1, M_2), (N_1, N_2) \bigr) := 
\opn{Hom}_{\cat{C}_1} (M_1, N_1) \ot 
\opn{Hom}_{\cat{C}_2} (M_2, N_2) . \]
The formula for the composition involves the Koszul sign rule. Given 
morphisms 
$\phi_i \in \opn{Hom}_{\cat{C}_i}(L_i, M_i)^{d_i}$
and
$\psi_i \in \opn{Hom}_{\cat{C}_i}(M_i, N_i)^{e_i}$, 
for $i = 1, 2$, their tensors are morphisms 
\[ \phi_1 \ot \phi_2 \in 
\opn{Hom}_{\cat{C}_1 \ot \cat{C}_2}
\bigl( (L_1, L_2), (M_1, M_2) \bigr)^{d_1 + d_2} \]
and 
\[ \psi_1 \ot \psi_2 \in 
\opn{Hom}_{\cat{C}_1 \ot \cat{C}_2}
\bigl( (M_1, M_2), (N_1, N_2) \bigr)^{e_1 + e_2} . \]
We define the composition to be 
\[ \begin{aligned}
& (\psi_1 \ot \psi_2) \circ (\phi_1 \ot \phi_2) := 
(-1)^{d_1 \cd e_2} \cd (\psi_1 \circ \phi_1) \ot
(\psi_2 \circ \phi_2)
\\ & \qquad 
\in \opn{Hom}_{\cat{C}_1 \ot \cat{C}_2}
\bigl( (L_1, L_2), (N_1, N_2) \bigr)^{d_1 + d_2 + e_1 + e_2} . 
\end{aligned} \]
\end{dfn}

\begin{exa}  \label{exa:2010}
Suppose $\cat{C}_1$ and $\cat{C}_2$ are single-object DG 
categories, say $\opn{Ob}(\cat{C}_i) = \{ x_i \}$. 
Then $\cat{C}_1 \ot \cat{C}_2$ also has a single-object, say 
$y := (x_1, x_2)$. 
The endomorphism DG rings satisfy
\[ \opn{End}_{\cat{C}_1 \ot \cat{C}_2}(y) = 
\opn{End}_{\cat{C}_1}(x_1) \ot \opn{End}_{\cat{C}_2}(x_2) . \]
See formula \ref{eqn:4106} and Example \ref{exa:1100}. 
\end{exa}

DG functors between DG categories were introduced in Definition \ref{dfn:1065}.

\begin{dfn} \label{dfn:1966}
Let $\cat{C}_1$, $\cat{C}_2$ and $\cat{D}$ be DG categories. A 
{\em DG bifunctor}%
\index{Bifunctor! differential graded}
\[ F : \cat{C}_1 \times \cat{C}_2 \to \cat{D} \]
is, by definition, a DG functor
$F : \cat{C}_1 \ot \cat{C}_2 \to \cat{D}$,
where $\cat{C}_1 \ot \cat{C}_2$ is the DG category from Definition
\ref{dfn:2010}. 
\end{dfn}

Warning: due to the signs that the composition of odd morphisms acquires, 
a DG bifunctor $F$ does not become, after we forget the DG structure,
a linear bifunctor in the sense of Definition \ref{dfn:1965}.

\begin{prop} \label{prop:5040}
Given a DG bifunctor
$F : \cat{C}_1 \times \cat{C}_2 \to \cat{D}$, there are induced linear 
bifunctors on the strict subcategories
\[ \opn{Str}(F) : \opn{Str}(\cat{C}_{1})  \times \opn{Str}(\cat{C}_{2}) 
\to  \opn{Str}(\cat{D}) \]
and on the homotopy categories
\[ \opn{Ho}(F) : \opn{Ho}(\cat{C}_{1})  \times \opn{Ho}(\cat{C}_{2}) 
\to  \opn{Ho}(\cat{D}) . \]
\end{prop}

\begin{exer} \label{exer:5040}
Prove this proposition. 
\end{exer}

Let $F$ be a DG bifunctor as in Definition \ref{dfn:1966}. If we fix an object 
$M_1 \in \cat{C}_{1}$, then 
$F(M_1, -) :  \cat{C}_2 \to \cat{D}$
is a DG functor. Similarly if we fix the second argument of $F$. 

For DG bifunctors there are several options for contravariance. 
In Subsection \ref{subsec:contrvar-dg-func} we talked about contravariant DG 
functors, and the opposite DG category $\cat{C}^{\mrm{op}}$ of a given DG 
category $\cat{C}$.  

\begin{dfn} \label{dfn:3125}
Let $\cat{C}_1$, $\cat{C}_2$ and $\cat{D}$ be DG categories. A 
{\em DG bifunctor that is contravariant in the first or the second 
argument}% 
\index{Bifunctor! contravariant}
is, by definition, a DG bifunctor
\[ F : \cat{C}_1^{\diamondsuit_1} \times \cat{C}_2^{\diamondsuit_2} 
\to \cat{D} \]
as in Definition \ref{dfn:1966}, where the symbols $\diamondsuit_1$ and 
$\diamondsuit_2$ are either $\bra{\mrm{empty}}$ or $\opn{op}$, as the case may 
be. 
\end{dfn}

Here are the two main examples of DG bifunctors. We give each of them in 
the commutative version and the noncommutative version. 

\begin{exa} \label{exa:1966}
Consider a commutative ring $A$. The category of complexes of $A$-modules is 
the DG category $\dcat{C}(A)$, and we take 
$\cat{C}_1 = \cat{C}_2 = \cat{D} := \dcat{C}(A)$. 
For each pair of objects 
$M_1, M_2 \in \dcat{C}(A)$ there is an object 
$F(M_1, M_2) := M_1 \ot_{A} M_2 \in \dcat{C}(A)$.
This is the usual tensor product of complexes. 
We define the action of $F$ on morphisms as follows: given 
$\phi_i \in \opn{Hom}_{\dcat{C}(A)}(M_i, N_i)^{k_i}$, 
we let 
\[ F(\phi_1, \phi_2) := \phi_1 \ot \phi_2 \in
\opn{Hom}_{\dcat{C}(A)} \bigl( F(M_1, M_2), F(N_1, N_2) \bigr)^{k_1 + k_2} . \]
The result is a DG bifunctor
\[ F = (- \ot_A -) : \dcat{C}(A) \times \dcat{C}(A) \to \dcat{C}(A) . \]
\end{exa}

\begin{exa} \label{exa:1965}
Consider DG rings $A_0, A_1, A_2$ (possibly noncommutative, but 
$\K$-central). 
Define the DG categories 
$\cat{C}_i := \dcat{C}(A_{i - 1} \ot A_i^{\mrm{op}})$
and
$\cat{D} := \dcat{C}(A_{0} \ot A_{2}^{\mrm{op}})$.
For each pair of objects 
$M_1 \in \cat{C}_1$ and $M_2 \in \cat{C}_2$ there is a DG $\K$-module 
$F(M_1, M_2) := M_1 \ot_{A_1} M_2$;
see Definition \ref{dfn:1100}.
This has a canonical DG $(A_{0} \ot A_{2}^{\mrm{op}})$-module structure:
\[  (a_0 \ot a_2) \cd  (m_1 \ot m_2) := (-1)^{\lms j_2 \cd (k_1 + k_2)} \cd
(a_0 \cd m_1) \ot (m_2 \cd a_2) \]
for elements 
$a_i \in A_i^{j_i}$  and $m_i \in M_i^{k_i}$. 
In this way $F(M_1, M_2)$ becomes an object of $\cat{D}$. 
We define the action of $F$ on morphisms as follows: given 
$\phi_i \in \opn{Hom}_{\cat{C}_i}(M_i, N_i)^{k_i}$,
we let 
\[ F(\phi_1, \phi_2) := \phi_1 \ot \phi_2 \in 
\opn{Hom}_{\cat{D}} \bigl( F(M_1, M_2), F(N_1, N_2) \bigr)^{k_1 + k_2} . \]
The result is a DG bifunctor
\[ F = (- \ot_{A_1} - ) : 
\dcat{C}(A_{0} \ot A_1^{\mrm{op}}) \times 
\dcat{C}(A_{1} \ot A_2^{\mrm{op}}) \to 
\dcat{C}(A_{0} \ot A_2^{\mrm{op}}) . \]
Compare this example to the one-sided construction in Example \ref{exa:1118}. 
\end{exa}

\begin{exa} \label{exa:1967}
Again we take a commutative ring $A$, but now our bifunctor $F$ arises from 
Hom, and so there is contravariance in the first argument. 
We define the DG categories 
$\cat{C}_1 = \cat{C}_2 = \cat{D} := \dcat{C}(A)$,
the same as in Example \ref{exa:1966}.
For every pair of objects 
$M_1, M_2 \in \dcat{C}(A)$ there is an object 
$F(M_1, M_2) := \opn{Hom}_{A}(M_1, M_2) \in \dcat{C}(A)$.
This is the usual Hom complex. 
We define the action of $F$ on morphisms as follows: given 
\[ \phi_1 \in \opn{Hom}_{\cat{C}_1^{\mrm{op}}}(M_1, N_1)^{k_1} =
\opn{Hom}_{\dcat{C}(A)^{\mrm{op}}}(M_1, N_1)^{k_1} =
\opn{Hom}_{A}(N_1, M_1)^{k_1} \]
and
\[ \phi_2 \in \opn{Hom}_{\cat{C}_2}(M_2, N_2)^{k_2} =
\opn{Hom}_{\dcat{C}(A)}(M_2, N_2)^{k_2} =
\opn{Hom}_{A}(M_2, N_2)^{k_2} \]
we let 
\[ \begin{aligned}
& F(\phi_1, \phi_2) := \opn{Hom}(\phi_1, \phi_2) \in
\opn{Hom}_{A} \bigl( \opn{Hom}_{A}(M_1, M_2), \opn{Hom}_{A}(N_1, N_2)
\bigr)^{k_1 + k_2}
\\
& \qquad = 
\opn{Hom}_{\cat{D}} \bigl( F(M_1, M_2), F(N_1, N_2) \bigr)^{k_1 + k_2} .
\end{aligned} \]
The result is a DG bifunctor
\[ F = \opn{Hom}_A(-, -) : \dcat{C}(A)^{\mrm{op}} \times \dcat{C}(A) \to 
\dcat{C}(A) . \]
\end{exa}

\begin{exa} \label{exa:1968}
Consider DG rings $A_0, A_1, A_2$ (possibly noncommutative, but 
$\K$-central). There is a DG bifunctor 
\[ F := \opn{Hom}_{A_1}(-, -) :
\dcat{C}(A_1 \ot A_0^{\mrm{op}})^{\mrm{op}} \times 
\dcat{C}(A_1 \ot A_2^{\mrm{op}}) \to 
\dcat{C}(A_0 \ot A_2^{\mrm{op}}) . \]
The details here are so confusing that we just leave them out.
(We shall come back to this bifunctor in Sections 
\ref{sec:perf-tilt-NC}-\ref{sec:rigid-DC-NC}, 
and then we shall have to deal with the details.)
\end{exa}

 

\mysubsection{Triangulated Bifunctors} \label{subsec:tri-bifun}

T-additive categories and triangulated categories were defined in Section 
\ref{sec:triangulated}. 

Suppose $(\cat{K}_1, \opn{T}_1)$ and $(\cat{K}_2, \opn{T}_2)$ are 
T-additive categories (linear over the base ring $\K$). There are two induced 
translation automorphisms on the category 
$\cat{K}_1 \times \cat{K}_2$~: 
\begin{equation} \label{eqn:4050}
\opn{T}_1(M_1, M_2) := \bigl( \opn{T}_1(M_1), M_2 \bigr) 
\end{equation}
and 
\begin{equation} \label{eqn:4051}
\opn{T}_2(M_1, M_2) := \bigl( M_1, \opn{T}_2(M_2) \bigr) . 
\end{equation}
These two functors commute: 
$\opn{T}_2 \circ \opn{T}_1 = \opn{T}_1 \circ \opn{T}_2$.

\begin{dfn} \label{dfn:1970}
Let $(\cat{K}_1, \opn{T}_1)$, $(\cat{K}_2, \opn{T}_2)$ 
and $(\cat{L}, \opn{T})$ be T-additive categories. 
A {\em T-additive bifunctor} 
\[ (F, \tau_1, \tau_2) : (\cat{K}_1, \opn{T}_1) \times (\cat{K}_2, \opn{T}_2) 
\to (\cat{L}, \opn{T}) \]
is made up of an additive bifunctor 
$F : \cat{K}_1 \times \cat{K}_2 \to \cat{L}$,
as in Definition \ref{dfn:1965}, together with isomorphisms 
$\tau_i : F \circ \opn{T}_i \iso \opn{T} \circ \, F$ 
of bifunctors 
$\cat{K}_1 \times \cat{K}_2 \to \cat{L}$. 
The condition is that 
\begin{equation*} \tag{$\dag$}
\tau_1 \circ \tau_2 = - \tau_2 \circ \tau_1 
\end{equation*}
as isomorphisms of bifunctors 
\[ F \circ \opn{T}_2 \circ \opn{T}_1 = F \circ \opn{T}_1 \circ \opn{T}_2 
\twoiso \opn{T} \circ \opn{T} \circ \, F . \]
\end{dfn}

The reason for the sign appearing in condition ($\dag$) will become clear 
in the proof of Lemma \ref{lem:4225}. 

Let $(F, \tau_1, \tau_2)$ be a T-additive bifunctor as in Definition 
\ref{dfn:1970}. If we fix an object $M_1 \in \cat{K}_{1}$, then 
\[ \bigl( F(M_1, -), \tau_2 \bigr) :  (\cat{K}_2, \opn{T}_2) \to 
(\cat{L}, \opn{T}) \]
is a T-additive functor. Similarly if we fix the second argument of $F$. 

In the next exercises we let the reader establish several operations on 
T-additive bifunctors. 

\begin{exer} \label{exer:1975}
In the situation of Definition \ref{dfn:1970}, suppose 
$(G, \tau) : (\cat{L}, \opn{T}) \to (\cat{L}', \opn{T}')$
is a T-additive functor into a fourth T-additive category 
$(\cat{L}', \opn{T}')$. Write the explicit formula for the T-additive bifunctor 
\[ (G, \tau) \circ (F, \tau_1, \tau_2) : 
(\cat{K}_1, \opn{T}_1) \times (\cat{K}_2, \opn{T}_2) 
\to (\cat{L}', \opn{T}') . \]
This should be compared to Definition \ref{dfn:1840}. 
\end{exer}

\begin{exer} \label{exer:1976}
In the situation of Definition \ref{dfn:1970}, suppose 
\[ (F', \tau'_1, \tau'_2) : (\cat{K}_1, \opn{T}_1) \times 
(\cat{K}_2, \opn{T}_2) \to (\cat{L}, \opn{T}) \]
is another T-additive bifunctor. Write the definition of a morphism of 
T-additive bifunctors 
$\ze : (F, \tau_1, \tau_2) \to (F', \tau'_1, \tau'_2)$.
Use Definition \ref{dfn:1840} as a template. 
\end{exer}

\begin{exer} \label{exer:1971}
Give a definition of a T-additive trifunctor. Show that if $F$ and $G$ are 
T-additive bifunctors, then $G(-, F(-, -))$ and $G(F(-, -), -)$
are T-additive trifunctors (whenever these compositions makes sense). 
\end{exer}

We now move to triangulated categories. 

\begin{dfn} \label{dfn:1971}
Let $(\cat{K}_1, \opn{T}_1)$, $(\cat{K}_2, \opn{T}_2)$ 
and $(\cat{L}, \opn{T})$ be triangulated categories. 
A {\em triangulated bifunctor}% 
\index{Bifunctor! triangulated}
\[ (F, \tau_1, \tau_2) : (\cat{K}_1, \opn{T}_1) \times (\cat{K}_2, \opn{T}_2) 
\to (\cat{L}, \opn{T}) \]
is a T-additive bifunctor that respects the triangulated structure in each 
argument. Namely, for every distinguished triangle 
\[ L_1 \xar{\al_1} M_1 \xar{\be_1} N_1 \xar{\ga_1} \opn{T}_1(L_1) \]
in $\cat{K}_1$, and every object $L_2 \in \cat{K}_2$,
the triangle 
\[ F(L_1, L_2) \xar{F(\al_1, \opn{id})} 
F(M_1, L_2) \xar{F(\be_1, \opn{id})} 
F(N_1, L_2) \xar{\tau_1 \circ F(\ga_1, \opn{id})} 
\opn{T}(F(L_1, L_2)) \]
in $\cat{L}$ is distinguished; and the same for distinguished triangles in the 
second argument. 
\end{dfn}

The operations on triangulated bifunctors are the same as those on T-additive 
bifunctors (see exercises above). 

We now discuss triangulated bifunctors in our favorite setup:
DG modules in abelian categories. This is done both in the covariant and in 
the contravariant direction, in either argument.
Recall that the opposite homotopy category
$\dcat{K}(A, \cat{M})^{\mrm{op}}$
was given a canonical triangulated structure in Subsection 
\ref{subsec:opp-hom-triang}.

\begin{setup} \label{set:3231}
We are given this data:
\begin{itemize}
\item DG rings $A_1$ and $A_2$, and abelian categories $\cat{M}_1$ and 
$\cat{M}_2$.

\item Direction indicators $\diamondsuit_1$ and $\diamondsuit_2$, that are 
either $\bra{\mrm{empty}}$ or $\mrm{op}$. 

\item Full subcategories 
$\cat{C}_i \sub \dcat{C}(A_i, \cat{M}_i)$, whose homotopy categories 
$\cat{K}_i := \lb \opn{Ho}(\cat{C}_i) \sub \dcat{K}(A_i, \cat{M}_i)$
are such that $\cat{K}_i^{\diamondsuit_i}$ is a triangulated subcategory of 
$\dcat{K}(A_i, \cat{M}_i)^{\diamondsuit_i}$.
\end{itemize} 
\end{setup}

The translation functor of the category 
$\dcat{K}(A_i, \cat{M}_i)$ is denoted by $\opn{T}_i$, and the  
translation functor of the category 
$\dcat{K}(A_i, \cat{M}_i)^{\diamondsuit_i}$ is denoted 
by $\opn{T}_i^{\diamondsuit_i}$. 

\begin{dfn} \label{dfn:2020}
Under Setup \ref{set:3231}, let $\cat{L}$ be another triangulated 
category. A
{\em triangulated bifunctor from $(\cat{K}_1, \cat{K}_2)$ to $\cat{L}$ that is 
contravariant in the first or the second argument}%
\index{Bifunctor! contravariant}%
\index{Bifunctor! triangulated}
is, by definition, a triangulated bifunctor
\[ (F, \tau_1, \tau_2) : 
(\cat{K}_1^{\diamondsuit_1}, \opn{T}_1^{\diamondsuit_1}) \times 
(\cat{K}_2^{\diamondsuit_2}, \opn{T}_2^{\diamondsuit_2}) 
\to (\cat{L}, \opn{T}) \]
as in Definition \ref{dfn:1971}, where the contravariance is according to the 
direction indicators $\diamondsuit_1$ and $\diamondsuit_2$.
\end{dfn}

Next we explain how to obtain triangulated bifunctors from DG bifunctors. 
This is done in the following setup:

\begin{setup} \label{set:3230}
In addition to the data from Setup \ref{set:3231}, we are given: 
\begin{itemize}
\item A DG ring $B$ and an abelian category $\cat{N}$.

\item A DG bifunctor
$F : \cat{C}_1^{\diamondsuit_1} \times \cat{C}_2^{\diamondsuit_2}
\to \dcat{C}(B, \cat{N})$.
\end{itemize}
\end{setup}

The translation functor of the category 
$\dcat{K}(B, \cat{N})$ is denoted by $\opn{T}$.  
For every object
$(M_1, M_2) \in \cat{C}_1 \times \cat{C}_2$ 
there are isomorphisms
\begin{equation} \label{eqn:2015}
\tau_{i, M_1, M_2} : 
F(\opn{T}_i(M_1, M_2))  \iso \opn{T}(F(M_1, M_2))
\end{equation}
in $\dcat{C}(B, \cat{N})$, arising from Definition \tup{\ref{dfn:1150}}.
Let us make this explicit (only for $i = 2$, since the case $i = 1$ is  
similar). Fixing the object $M_1$ we obtain a DG functor
$G :  \cat{C}_2 \to \dcat{C}(B, \cat{N})$,  
$G(M_2) := F(M_1, M_2)$. 
The isomorphism 
\[ \tau_{2, M_1, M_2} : G(\opn{T}_2(M_2)) \iso \opn{T}(G(M_2)) \]
is then 
\begin{equation} \label{eqn:5042}
\tau_{2, M_1, M_2} = \opn{t}_{G(M_2)} \circ \, G(\opn{t}_{M_2})^{-1} . 
\end{equation}
We are using the little t operator here. 

\begin{lem} \label{lem:2015}
Under Setups \tup{\ref{set:3231}} and \tup{\ref{set:3230}}, assume that 
the direction indicators $\diamondsuit_1$ and $\diamondsuit_2$ are both
$\bra{\mrm{empty}}$. Fix $i \in \{ 1, 2 \}$.
Letting the object $(M_1, M_2)$ in \tup{(\ref{eqn:2015})} vary, we get an 
isomorphism 
$\tau_i : F \circ \opn{T}_i \iso \opn{T} \circ \, F$
of additive bifunctors 
$\opn{Str}(\cat{C}_{1}) \times \opn{Str}(\cat{C}_{2}) \to
\dcat{C}_{\mrm{str}}(B, \cat{N})$.
\end{lem}

\begin{proof}
This is an almost immediate consequence of the fact that the little 
t operators are morphisms of functors (see Theorem \ref{thm:1260}(2)),
\end{proof}

According to Proposition \ref{prop:5040}, the DG bifunctor $F$ induces an 
additive bifunctor 
$F : \cat{K}_{1}^{\diamondsuit_1} \times \cat{K}_{2}^{\diamondsuit_2} \to 
\dcat{K}(B, \cat{N})$ 
on the homotopy categories. 

\begin{lem} \label{lem:4225}    
Under Setups \tup{\ref{set:3231}} and \tup{\ref{set:3230}}, assume that 
the direction indicators $\diamondsuit_1$ and $\diamondsuit_2$ are both
$\bra{\mrm{empty}}$. Then 
\[ (F, \tau_1, \tau_2) :  \cat{K}_1 \times \cat{K}_2 \to 
\dcat{K}(B, \cat{N}) \]
is a triangulated bifunctor. 
\end{lem}

\begin{proof}
The only challenge is to prove that $(F, \tau_1, \tau_2)$ is a T-additive 
bifunctor; and in that, all we have to prove is that condition 
($\dag$) in Definition \ref{eqn:4051} holds, namely that 
$\tau_1 \circ \tau_2 = - \tau_2 \circ \tau_1$.
The rest hinges on single-argument considerations, that were already handled in 
Theorems \ref{thm:1150} and \ref{thm:1265}. 

So let us prove the equality ($\dag$). Choose a pair of objects 
$(M_1, M_2)$. We have diagram (\ref{eqn:5043}) in the category 
$\dcat{C}(B, \cat{N})$. 
\begin{figure}
\begin{equation} \label{eqn:5043}
\UseTips  \xymatrix @C=5ex @R=6.5ex {
&
F(\opn{T}_1(M_1), \opn{T}_2(M_2))
\ar[dr]^{F(\opn{t}_{M_1}^{-1}, \opn{id})}
\ar[dl]_{F(\opn{id}, \opn{t}_{M_2}^{-1})}
\\
F(\opn{T}_1(M_1), M_2)
\ar[d]_{\opn{t}_{F(\opn{T}_1(M_1), M_2)}}
\ar[dr]^{F(\opn{t}_{M_1}^{-1}, \opn{id})}
&
&
F(M_1, \opn{T}_2(M_2))
\ar[d]^{\opn{t}_{F(M_1, \opn{T}_2(M_2))}}
\ar[dl]_{F(\opn{id}, \opn{t}_{M_2}^{-1})}
\\
\opn{T}(F(\opn{T}_1(M_1), M_2))
\ar[d]_{\opn{T}(F(\opn{t}_{M_1}^{-1}, \opn{id}))}
&
F(M_1, M_2)
\ar[dl]^{\ \ \opn{t}_{F(M_1, M_2)}}
\ar[dr]_{\opn{t}_{F(M_1, M_2)}\ }
&
\opn{T}(F(M_1, \opn{T}_2(M_2))
\ar[d]^{\opn{T}(F(\opn{id}, \opn{t}_{M_2}^{-1}))}
\\
\opn{T}(F(M_1, M_2))
\ar[dr]^{\opn{T}(\opn{t}_{F(M_1, M_2)})}
&
&
\opn{T}(F(M_1, M_2))
\ar[dl]_{\opn{T}(\opn{t}_{F(M_1, M_2)}) \ \ }
\\
&
\opn{T}(\opn{T}(F(M_1, M_2)))
} 
\end{equation}
\end{figure}
Going from top to bottom on the left edge is the 
morphism $\tau_1 \circ \tau_2$, and going on the right edge is the 
morphism $\tau_2 \circ \tau_1$.
The bottom diamond is trivially commutative. 
The two triangles, with common vertex at $F(M_1, M_2)$, are $(-1)$-commutative, 
because $\opn{t} : \opn{Id} \to \opn{T}$ is a degree $-1$ morphism of DG 
functors. Since they occur on both sides, these signs cancel each other. 
Finally, the top diamond is $(-1)$-commutative, because
\[ (\opn{t}_{M_1}^{-1}, \opn{id}) \circ (\opn{id}, \opn{t}_{M_2}^{-1}) = 
(\opn{t}_{M_1}^{-1}, \opn{t}_{M_2}^{-1}) = 
- (\opn{id}, \opn{t}_{M_2}^{-1}) \circ (\opn{t}_{M_1}^{-1}, \opn{id}) . 
\qedhere \]
\end{proof}

\begin{thm} \label{thm:4225} 
Under Setups \tup{\ref{set:3231}} and \tup{\ref{set:3230}},
there are canonical translation isomorphisms $\tau_1$ and $\tau_2$ such that 
\[ (F, \tau_1, \tau_2) : 
(\cat{K}_1^{\diamondsuit_1}, \opn{T}_1^{\diamondsuit_1}) \times 
(\cat{K}_2^{\diamondsuit_2}, \opn{T}_2^{\diamondsuit_2})
\to (\dcat{K}(B, \cat{N}), \opn{T})  \]
is a triangulated bifunctor. 
\end{thm}

\begin{proof}
For $i = 1, 2$ let us define the indicator 
\[ \heartsuit_i :=
\begin{cases}
\mrm{flip} & \tup{if} \hspace{1ex} \diamondsuit_i = \mrm{op}, 
\\
\bra{\mrm{empty}} & \tup{if} \hspace{1ex} \diamondsuit_i =  \bra{\mrm{empty}} .
\end{cases} \]
We get DG categories $\cat{C}_i^{\heartsuit_i}$: 
if $\diamondsuit_i =  \bra{\mrm{empty}}$ then
$\cat{C}_i^{\heartsuit_i} = \cat{C}_i \sub \dcat{C}(A_i, \cat{M}_i)$,
and if $\diamondsuit_i = \mrm{op}$ then 
\[ \cat{C}_i^{\heartsuit_i} = \cat{C}_i^{\mrm{flip}} 
= \opn{Flip}(\cat{C}_i^{\mrm{op}}) \sub 
\dcat{C}(A_i, \cat{M}_i)^{\mrm{flip}} = 
\dcat{C}(A_i^{\mrm{op}}, \cat{M}_i^{\mrm{op}}) , \]
as in the proof of Theorem \ref{thm:2995}. Likewise there are triangulated 
categories $\cat{K}_i^{\heartsuit_i}$. 

We start by composing $F$ with the isomorphism of DG categories 
$\opn{Flip}^{-1}$ from Theorem \ref{thm:2495}, in each coordinate $i$ for which 
$\diamondsuit_i = \opn{op}$. This gives a new DG bifunctor 
\[ F' : \cat{C}_1^{\heartsuit_1} \times \cat{C}_2^{\heartsuit_2} \to
\dcat{C}(B, \cat{N}) . \]
Now we can apply Lemma \ref{lem:4225}, to get a triangulated bifunctor 
\[ F' : \cat{K}_1^{\heartsuit_1} \times \cat{K}_2^{\heartsuit_2} \to
\dcat{K}(B, \cat{N}) . \]
Specifically, this means that we have a pair of translations 
$(\tau'_1, \tau'_2)$ such that \lb $(F', \tau'_1, \tau'_2)$ is 
a triangulated bifunctor. 

Finally we compose $F'$ with the isomorphism of triangulated categories 
$\opn{Flip}$ from formula (\ref{eqn:3002}), in each coordinate $i$ for which 
$\diamondsuit_i = \opn{op}$. This recovers the original bifunctor $F$ on
$\cat{K}_1^{\diamondsuit_1} \times \cat{K}_2^{\diamondsuit_2}$.
There are translations $\tau_i$, that are gotten from the 
translations $\tau'_i$ like in the proof of Theorem \ref{thm:2995}. 
\end{proof}



\mysubsection{Derived Bifunctors}  \label{subsec:der-bifun}

Here we explain what are right and left derived bifunctors of triangulated 
bifunctors. The definitions and results are very similar to the single-argument 
case. For the sake of simplicity, we shall mostly ignore the translation 
functors and the translation isomorphisms; enough was said about them in the 
previous subsection. 

The next setup will be used throughout this subsection.

\begin{setup} \label{set:1976}
The following are given: 
\begin{itemize}
\item Triangulated categories $\cat{K}_1$, $\cat{K}_2$ and $\cat{E}$.

\item A triangulated bifunctor 
$F : \cat{K}_1 \times \cat{K}_2 \to \cat{E}$.

\item Denominator sets of cohomological origin
$\cat{S}_1 \sub \cat{K}_1$ and $\cat{S}_2 \sub \cat{K}_2$.
\end{itemize}
We write $\cat{D}_i := (\cat{K}_i)_{\cat{S}_i}$, 
and $\opn{Q}_i : \cat{K}_i \to \cat{D}_i$ are the localization functors.
\end{setup}

The localized category 
$\cat{D}_i := (\cat{K}_i)_{\cat{S}_i}$
is triangulated, and the localization functor 
$\opn{Q}_i : \cat{K}_i \to \cat{D}_i$
is triangulated. On the product categories we get a functor 
\[ \opn{Q}_1 \times \opn{Q}_2 : 
\cat{K}_1 \times \cat{K}_2 \to \cat{D}_1 \times \cat{D}_2 . \]

Before embarking on the definitions of the derived bifunctors, we want to 
introduce the relevant categories of functors. These are the bifunctor variants 
of what we had in Subsection \ref{subsec:fun-cats}. 

\begin{dfn} \label{dfn:3130} 
We denote by 
$\cat{LinBiFun}(\cat{K}_1 \times \cat{K}_2, \cat{E})$
the category whose objects are the linear bifunctors 
$F : \cat{K}_1 \times \cat{K}_2 \to \cat{E}$.
The morphisms are the obvious ones. 
\end{dfn}

This is a linear category. Actually, this category is the same as the category 
$\cat{LinFun}(\cat{K}_1 \ot \cat{K}_2, \cat{E})$ 
of linear functors
$F : \cat{K}_1 \ot \cat{K}_2 \to \cat{E}$, 
only with the extra information that the source is a product category.
     
\begin{dfn} \label{dfn:2036}
We denote by 
$\cat{TrBiFun}(\cat{K}_1 \times \cat{K}_2, \cat{E})$
the category whose objects are the $\K$-linear triangulated bifunctors 
$(F, \tau_1, \tau_2) : \cat{K}_1 \times \cat{K}_2 \to \cat{E}$.
The morphisms are those of T-additive bifunctors. 
\end{dfn}

There is a functor 
\begin{equation} \label{eqn:3130}
\cat{TrBiFun}(\cat{K}_1 \times \cat{K}_2, \cat{E}) \to 
\cat{LinBiFun}(\cat{K}_1 \times \cat{K}_2, \cat{E}) 
\end{equation}
that forgets $(\tau_1, \tau_2)$. It is a faithful additive functor. 

Suppose $U_i : \cat{K}'_i \to \cat{K}_i$ are triangulated functors between 
triangulated categories. We get an induced additive functor
\begin{equation} \label{eqn:2036}
\cat{Fun}(U_1 \times U_2, \opn{Id}_{\cat{E}}) :
\cat{TrBiFun}(\cat{K}_1 \times \cat{K}_2, \cat{E}) \to 
\cat{TrBiFun}(\cat{K}'_1 \times \cat{K}'_2, \cat{E})
\end{equation}
The formula is 
$F \mapsto F \circ (U_1 \times U_2)$.

\begin{lem} \label{lem:2035}
If the functors $U_1$ and $U_2$ are equivalences, then the functor 
in \tup{(\ref{eqn:2036})} is an equivalence.
\end{lem}

\begin{proof}
This is basically the same as the proof of Lemma \ref{lem:2335}.
\end{proof}

As in Definition \ref{dfn:2037} we denote by 
\[ \cat{TrBiFun}_{\cat{S}_1 \times \cat{S}_2}
(\cat{K}_1 \times \cat{K}_2, \cat{E}) \sub 
\cat{TrBiFun}(\cat{K}_1 \times \cat{K}_2, \cat{E}) \]
the full subcategory on the triangulated bifunctors $F$ such that 
$F(\cat{S}_1 \times \cat{S}_2) \sub \cat{E}^{\times}$. 

\begin{lem} \label{lem:2036} 
The functor 
\[ \cat{Fun}(\opn{Q}_1 \times \opn{Q}_2, \opn{Id}_{\cat{E}}) :
\cat{TrBiFun}(\cat{D}_1 \times \cat{D}_2, \cat{E}) \to 
\cat{TrBiFun}_{\cat{S}_1 \times \cat{S}_2} (\cat{K}_1 \times \cat{K}_2, \cat{E})
\]
is an isomorphism of categories. 
\end{lem}

\begin{proof}
This is basically the same as the proof of Lemma \ref{lem:2336}, when 
combined with the isomorphism of triangulated categories 
$(\cat{K}_1 \times \cat{K}_2)_{\cat{S}_1 \times \cat{S}_2} \to 
\cat{D}_1 \times \cat{D}_2$ 
from Proposition \ref{prop:2028}. 
\end{proof}

We are ready to talk about derived bifunctors.  

\begin{dfn} \label{dfn:1975}
Under Setup \ref{set:1976}, 
a {\em triangulated right derived bifunctor of $F$ with respect to 
$\cat{S}_1 \times \cat{S}_2$}%
\index{Bifunctor! derived}
is a pair $(\mrm{R} F, \eta^{\mrm{R}})$, where 
\[ \mrm{R} F : \cat{D}_1 \times \cat{D}_2 \to \cat{E} \]
is a triangulated bifunctor, and 
\[ \eta^{\mrm{R}}: F \twoto \mrm{R} F \circ (\opn{Q}_1 \times \opn{Q}_2) \] 
is a morphism of triangulated bifunctors, 
such that the following universal property holds:
\begin{itemize}
\item[(R)] Given a pair $(G, \th)$, consisting of a triangulated 
bifunctor
$G : \cat{D}_1 \times \cat{D}_2 \to \cat{E}$ 
and a morphism of triangulated bifunctors 
$\th : F \twoto G \circ (\opn{Q}_1 \times \opn{Q}_2)$,
there is a unique morphism of triangulated bifunctors 
$\mu : \mrm{R}^{} F \twoto G$
such that 
$\th = (\mu \circ \opn{id}_{\opn{Q}_1 \times \opn{Q}_2}) * \eta^{\mrm{R}}$.
\end{itemize}
\end{dfn}

Here is a diagram showing property (R):
\begin{equation} \label{eqn:2076}
\UseTips  \xymatrix @C=14ex @R=10ex {
\cat{K}_1 \times \cat{K}_2
\ar[r]^{F} _(0.45){}="q"  ^(0.55){}="q2" 
\ar[d]_{\opn{Q}_1 \times \opn{Q}_2}
& 
\cat{E} 
\\
\cat{D}_1 \times \cat{D}_2
\ar[ur]^(0.32){\mrm{R} F} ^(0.45){}="f" ^(0.5){}="f1" 
\ar@{=>}  "q";"f" _{\eta^{\mrm{R}}}
\ar@{->}@(r,d)[ru]_{G} ^(0.55){}="g"  ^(0.48){}="g1"
\ar@{=>}  "q2";"g" ^(0.7){\th}
\ar@{=>}  "f";"g1" _(0.45){\mu}
} 
\end{equation}
 
\begin{prop} \label{prop:2092}
If a triangulated right derived bifunctor exists, then it is unique up to a 
unique isomorphism. 
\end{prop}

\begin{proof}
This is just like the proof of Proposition \ref{prop:1350}. We leave the small 
changes to the reader. 
\end{proof}

Existence in general is like Theorem \ref{thm:1470}, but a bit more complicated.

\begin{thm} \label{thm:2000}
Under Setup \tup{\ref{set:1976}}, 
assume there are full triangulated subcategories 
$\cat{J}_1 \subseteq \cat{K}_1$ and 
$\cat{J}_2 \subseteq \cat{K}_2$
with these two properties\tup{:}
\begin{enumerate}
\rmitem{a} 
If $\phi_1 : I_1 \to J_1$ is a morphism in $\cat{J}_1 \cap \cat{S}_1$
and $\phi_2 : I_2 \to J_2$ is a morphism in $\cat{J}_2 \cap \cat{S}_2$, 
then 
\[ F(\phi_1, \phi_2) : F(I_1, I_2) \to F(J_1, J_2) \]
is an isomorphism in $\cat{E}$.
 
\rmitem{b} For every $p \in \{1, 2\}$ and every object $M_p \in \cat{K}_p$,
there exists a morphism $\rho_p : M_p \to I_p$ in $\cat{S}_p$ with target 
$I_p \in \cat{J}_p$. 
\end{enumerate}
Then the triangulated right derived bifunctor%
\index{Bifunctor! derived}
\[ (\mrm{R} F, \eta^{\mrm{R}}) : \cat{D}_1 \times \cat{D}_2 \to \cat{E} \]
of $F$ with respect to $\cat{S}_1 \times \cat{S}_2$
exists. Moreover, for every pair of objects $I_1 \in \cat{J}_1$ and 
$I_2 \in \cat{J}_2$ the morphism  
\[ \eta^{\mrm{R}}_{I_1, I_2} : F(I_1, I_2) \to \mrm{R} F (I_1, I_2) \]
in $\cat{E}$ is an isomorphism.
\end{thm}

In applications we will see that either $\cat{J}_1 = \cat{K}_1$ or 
$\cat{J}_2 = \cat{K}_2$; namely we will only need to resolve in the second or 
in the first argument, respectively. 

\begin{proof}
It will be convenient to change notation.
For $p = 1, 2$ let's define 
$\cat{K}'_p := \cat{J}_p$, $\cat{S}'_p := \cat{K}'_p \cap \cat{S}_p$ and
$\cat{D}'_p := (\cat{K}'_p)_{\cat{S}'_p}$. 
The localization functors are 
$\opn{Q}'_p : \cat{K}'_p \to \cat{D}'_p$. 
The inclusions are 
$U_p : \cat{K}'_p \to \cat{K}_p$, and their localizations are the functors 
$V_p : \cat{D}'_p \to \cat{D}_p$.
By Lemma \ref{lem:2086} the functors $V_p$ are equivalences.

The situation is depicted in these diagrams. 
We have this commutative diagram of 
products of triangulated functors between products of triangulated 
categories:
\begin{equation} \label{eqn:2050}
\UseTips \xymatrix @C=10ex @R=6ex {
\cat{K}'_1 \times \cat{K}'_2
\ar[r]^{U_1 \times U_2}
\ar[d]_{\opn{Q}'_1 \times \opn{Q}'_2}
&
\cat{K}_1 \times \cat{K}_2
\ar[d]^{\opn{Q}_1 \times \opn{Q}_2}
\\
\cat{D}'_1 \times \cat{D}'_2
\ar[r]^{V_1 \times V_2}
&
\cat{D}_1 \times \cat{D}_2
} 
\end{equation}

\noindent 
The arrow $V_1 \times V_2$ is an equivalence. 
Diagram (\ref{eqn:2050}) induces a commutative diagram of linear categories:
\begin{equation} \label{eqn:2051}
\UseTips \xymatrix @C=16ex @R=5.5ex {
\cat{TrBiFun}(\cat{K}'_1 \times \cat{K}'_2, \cat{E})
&
\cat{TrBiFun}(\cat{K}_1 \times \cat{K}_2, \cat{E})
\ar[l]_{\cat{Fun}(U_1 \times U_2, \opn{Id})}
\\
\cat{TrBiFun}_{\cat{S}'_1 \times \cat{S}'_2}
(\cat{K}'_1 \times \cat{K}'_2, \cat{E})
\ar[u]_{\mrm{f.f.\, emb}}
&
\cat{TrBiFun}_{\cat{S}_1 \times \cat{S}_2}
(\cat{K}_1 \times \cat{K}_2, \cat{E})
\ar[u]^{\mrm{f.f.\, emb}}
\ar[l]^{\mrm{equiv}}_{\cat{Fun}(U_1 \times U_2, \opn{Id})}
\\
\cat{TrBiFun}(\cat{D}'_1 \times \cat{D}'_2, \cat{E})
\ar[u]^(0.47){\cat{Fun}(\opn{Q}'_1 \times \opn{Q}'_2, \opn{Id})}_{\mrm{isom}}
&
\cat{TrBiFun}(\cat{D}_1 \times \cat{D}_2, \cat{E})
\ar[l]_{\cat{Fun}(V_1 \times V_2, \opn{Id})}^{\mrm{equiv}}
\ar[u]_(0.47){\cat{Fun}(\opn{Q}_1 \times \opn{Q}_2, \opn{Id})}^{\mrm{isom}}
} 
\end{equation}
According to Lemmas \ref{lem:2035} and \ref{lem:2036}, the arrows in the 
diagram above that are marked ``isom'' or ``equiv'' are isomorphisms or 
equivalences, respectively. By definition the arrows marked ``f.f.\ emb'' are 
fully faithful embeddings. 

We know that $\cat{S}_p \sub \cat{K}_p$ are left denominator sets. Therefore 
(see Proposition \ref{prop:2100})
$\cat{S}_1 \times \cat{S}_2 \sub \cat{K}_1 \times \cat{K}_2$
is a  left denominator set.
Condition (a) of Theorem \ref{thm:2000} says that $F$ sends morphisms in 
$\cat{S}'_1 \times \cat{S}'_2$ to isomorphisms in $\cat{E}$. 
Condition (b) there says that there are enough right 
$(\cat{K}'_1 \times \cat{K}'_2)$-resolutions in 
$\cat{K}_1 \times \cat{K}_2$.

Thus we are in a position to use the abstract Theorem \ref{thm:2080}.
It says that there is an abstract right derived functor 
$(\mrm{R} F, \eta^{\mrm{R}}) :  \cat{D}_1 \times \cat{D}_2 \to \cat{E}$.
However, going over the proof of Theorem \ref{thm:2080},
we see that all constructions there can be made within the triangulated 
setting, namely in diagram (\ref{eqn:2051}) instead of in  diagram 
(\ref{eqn:2084}). Therefore $\mrm{R} F$ is an object of the category in the 
bottom right corner of (\ref{eqn:2051}), and the morphism 
$\eta^{\mrm{R}} : F \twoto \mrm{R} F \circ (\opn{Q}_1 \times \opn{Q}_2)$
is in the category in the top right corner of (\ref{eqn:2051}).
\end{proof}

We now treat left derived bifunctors. Due to the similarity to the right 
derived functors, we give only a few details here. 

\begin{dfn} \label{dfn:2105}
Under Setup \ref{set:1976}, 
a {\em triangulated left derived bifunctor of $F$ with respect to 
$\cat{S}_1 \times \cat{S}_2$}%
\index{Bifunctor! derived}
is a pair $(\mrm{L} F, \eta^{\mrm{L}})$, where 
\[ \mrm{L} F : \cat{D}_1 \times \cat{D}_2 \to \cat{E} \]
is a triangulated bifunctor, and 
\[ \eta^{\mrm{L}} :  \mrm{L} F \circ (\opn{Q}_1 \times \opn{Q}_2) \twoto F \] 
is a morphism of triangulated bifunctors, 
such that the following universal property holds:
\begin{itemize}
\item[(L)] Given a pair $(G, \th)$, consisting of a triangulated 
bifunctor
$G : \cat{D}_1 \times \cat{D}_2 \to \cat{E}$ 
and a morphism of triangulated bifunctors 
$\th :  G \circ (\opn{Q}_1 \times \opn{Q}_2) \twoto F$,
there is a unique morphism of triangulated bifunctors 
$\mu : G \twoto \mrm{L}^{} F$
such that 
$\th =  \eta^{\mrm{L}} * (\mu \circ \opn{id}_{\opn{Q}_1 \times \opn{Q}_2})$.
\end{itemize}
\end{dfn}

\begin{prop} \label{prop:2105}
If a triangulated left derived bifunctor exists, then it is unique up to a 
unique isomorphism. 
\end{prop}

\begin{proof}
This is the opposite of Proposition \ref{prop:2092}, and we leave it to the 
reader to make the adjustments.
\end{proof}

\begin{thm} \label{thm:2106}
In the situation of Definition \tup{\ref{dfn:2105}}, 
assume there are full triangulated subcategories 
$\cat{P}_1 \subseteq \cat{K}_1$ and 
$\cat{P}_2 \subseteq \cat{K}_2$
with these two properties\tup{:}
\begin{enumerate}
\rmitem{a} If $\phi_1 : P_1 \to Q_1$ is a
morphism in $\cat{P}_1 \cap \cat{S}_1$
and $\phi_2 : P_2 \to Q_2$ is a 
morphism in $\cat{P}_2 \cap \cat{S}_2$, then 
\[ F(\phi_1, \phi_2) : F(P_1, P_2) \to F(Q_1, Q_2) \]
is an isomorphism in $\cat{E}$.
 
\rmitem{b} 
For every $i \in \{1, 2\}$ and every object $M_i \in \cat{K}_i$,
there exists a morphism $\rho_i : P_i \to M_i$ in $\cat{S}_i$ with target 
$P_i \in \cat{P}_i$. 
\end{enumerate}
Then the triangulated left derived bifunctor%
\index{Bifunctor! derived}
\[ (\mrm{L} F, \eta^{\mrm{L}}) : \cat{D}_1 \times \cat{D}_2 \to \cat{E} \]
of $F$ with respect to $\cat{S}_1 \times \cat{S}_2$
exists. Moreover, for every pair of objects $P_1 \in \cat{P}_1$ and 
$P_2 \in \cat{P}_2$ the morphism  
\[ \eta^{\mrm{L}}_{P_1, P_2} :  \mrm{L} F(P_1, P_2) \to F(P_1, P_2) \]
in $\cat{E}$ is an isomorphism.
\end{thm}

\begin{exer} \label{exer:4635}
Prove Theorem \ref{thm:2106}.
(Hint: This is the opposite of Theorem \ref{thm:2000}. The proof is analogous.)
\end{exer}


In applications we will see that either $\cat{P}_1 = \cat{K}_1$ or 
$\cat{P}_2 = \cat{K}_2$; namely we will only need to resolve in the second or 
in the first argument, respectively. 

There are more specialized definitions of triangulated derived bifunctors 
whose sources are categories of DG modules, like Definitions \ref{dfn:3236} and 
\ref{dfn:3237}. We leave them to the reader. 

\begin{rem} \label{5050}
Deriving triangulated bifunctors that are contravariant in one or two arguments 
involves passage to opposite categories; see Definition \ref{dfn:2020}. 
Proposition \ref{prop:2356} tells us that when 
$\cat{K}_i^{\mrm{op}} \sub \dcat{K}(A_i, \cat{M}_i)^{\mrm{op}}$
are full triangulated subcategories, the sets 
$\cat{S}_i^{\mrm{op}}$
of quasi-isomorphisms in them are of cohomological origin. So we are still 
within the conditions of Setup \ref{set:1976}, and Theorems \ref{thm:2000} and 
\ref{thm:2106} can be applied. 
We leave it to the reader to work out the details in the various cases. 
However, the important case of the bifunctor $\opn{Hom}(-, -)$ 
will be given a full treatment in Subsection \ref{subsec:RHom}.
\end{rem}
 




\cleardoublepage
\mysection{Resolving Subcategories of 
\texorpdfstring{$\dcat{K}(A, \cat{M})$}{K(A,M)}} 
\label{sec:resol} 

\AYcopyright

In this section we are back in the more concrete setting: $A$ is a DG ring, 
and $\cat{M}$ is an abelian category, both over the commutative base ring 
$\K$; see Convention \ref{conv:2490}. We will define {\em K-projective} and 
{\em K-injective} DG modules in
$\bcat{K}(A, \cat{M})$. These DG modules form full triangulated 
subcategories of $\bcat{K}(A, \cat{M})$, and are concrete versions of the 
abstract categories $\cat{P}$ and $\cat{J}$, that played important roles in 
Subsection \ref{subsec:tri-der-funs}.
For $\bcat{K}(A)$ we also define {\em K-flat DG modules}.

\mysubsection{K-Injective DG Modules} \label{subsec:K-inj}

For every integer $p$ we have the $p$-th cohomology functor 
$\opn{H}^p : \bcat{C}_{\mrm{str}}(A, \cat{M}) \to \cat{M}$.
There is equality 
$\opn{H}^p = \opn{H}^0 \circ \opn{T}^p$,
where $\opn{T}$ is the translation functor of $\bcat{C}(A, \cat{M})$. 
The functors $\opn{H}^p$ pass to the homotopy category, and 
$\opn{H}^0 : \bcat{K}(A, \cat{M}) \to \cat{M}$
is a cohomological functor, in the sense of Definition \ref{dfn:1510}.

\begin{dfn} \label{dfn:1505}
A DG module $N \in \dcat{C}(A, \cat{M})$ is called {\em acyclic}% 
\index{Differential graded module! acyclic}
if $\mrm{H}^p(N) = 0$ for all $p$. 
\end{dfn}

\begin{dfn} \label{dfn:1506}
A DG module $I \in \dcat{C}(A, \cat{M})$ is called 
{\em K-injective}%
\index{Differential graded module! K-injective}
if for every acyclic DG module $N \in \dcat{C}(A, \cat{M})$, the DG $\K$-module 
$\opn{Hom}_{A, \cat{M}}(N, I)$ is acyclic.
\end{dfn}

The definition above characterizes K-injectives as objects of 
$\dcat{C}(A, \cat{M})$. 
The next proposition shows that being K-injective is intrinsic to the 
triangulated category $\dcat{K}(A, \cat{M})$, with the cohomological functor 
$\opn{H}^0$ (that tells us which are the acyclic objects). 

\begin{prop} \label{prop:1515}
A DG module $I \in \dcat{K}^{}(A, \cat{M})$ is K-injective if and only if 
$\opn{Hom}_{\dcat{K}(A, \cat{M})}(N, I) = 0$ for every acyclic DG module 
$N \in \dcat{K}^{}(A, \cat{M})$.
\end{prop}

\begin{proof}
This is because for every integer $p$ we have 
\[ \opn{H}^p \bigl( \opn{Hom}_{A, \cat{M}}(N, I) \bigr) \cong
\opn{H}^0 \bigl( \opn{Hom}_{A, \cat{M}}(\opn{T}^{-p}(N), I) \bigr) 
= \opn{Hom}_{\dcat{K}(A, \cat{M})}(\opn{T}^{-p}(N), I) , \]
and $N$ is acyclic iff $\opn{T}^{-p}(N)$ is acyclic.
\end{proof}

\begin{dfn} \label{dfn:1507}
Let $M \in \dcat{C}(A, \cat{M})$. A 
{\em K-injective resolution}% 
\index{Resolution! K-injective}
of $M$ is a quasi-isomorphism $\rho : M \to I$ in 
$\dcat{C}_{\mrm{str}}(A, \cat{M})$, where 
$I$ is a K-injective DG module. 
\end{dfn}

\begin{rem} \label{rem:5051}
Let $\rho : M \to I$ be a K-injective resolution in 
$\dcat{C}_{\mrm{str}}(A, \cat{M})$. 
By a slight abuse of notation, we shall often refer to the induced 
quasi-isomorphism 
$\opn{P}(\rho) : M \to I$ in $\dcat{K}(A, \cat{M})$
as a K-injective resolution of $M$. Cf.\ Definitions \ref{dfn:1508} and
\ref{dfn:1512} below. 
\end{rem}

In the next section we will prove existence of K-injective resolutions in 
several contexts. Here is an easy example of a K-injective complex. 

\begin{exer} \label{exer:1505}
Let $I \in \dcat{K}(\cat{M})$ be a complex of injective objects of 
$\cat{M}$, with zero differential. Prove that $I$ is K-injective. 
\end{exer}

\begin{dfn} \label{dfn:1509}
Let $\cat{K}$ be a full subcategory of $\dcat{K}^{}(A, \cat{M})$. 
The full subcategory of  $\cat{K}$ on the K-injective DG modules in it is 
denoted by $\cat{K}_{\mrm{inj}}$. In other words, 
$\cat{K}_{\mrm{inj}} = \dcat{K}^{}(A, \cat{M})_{\mrm{inj}} \cap \cat{K}$.
\end{dfn}

\begin{rem} \label{rem:5050}
Warning: the property of a DG module $I \in \cat{K}$ being K-injective is in 
general not intrinsic to the subcategory 
$\cat{K} \sub \dcat{K}(A, \cat{M})$. This is because the 
condition in Proposition \ref{prop:1515} (and in Definition \ref{dfn:1506}) 
has to be tested against all acyclic DG modules 
$N \in \dcat{K}(A, \cat{M})$.
\end{rem}

\begin{prop}  \label{prop:130}
If  $\cat{K}$ is a full triangulated subcategory of
$\dcat{K}^{}(A, \cat{M})$, then $\cat{K}_{\mrm{inj}}$ is a full triangulated 
subcategory of $\cat{K}$.
\end{prop}

\begin{proof}
It suffices to prove that 
$\dcat{K}^{}(A, \cat{M})_{\mrm{inj}}$ is a 
triangulated subcategory of $\dcat{K}^{}(A, \cat{M})$.
It is easy to see that $\dcat{K}^{}(A, \cat{M})_{\tup{inj}}$ is closed under
translations. Suppose 
$I \to J \to K \xar{\, \triangle\, }$
is a distinguished triangle in $\dcat{K}^{}(A, \cat{M})$ such that  
$I, J$ are K-injective DG modules. We have to show that $K$ is also
K-injective. Take any acyclic DG module $N \in \dcat{K}^{}(A, \cat{M})$. 
There is an exact sequence
\[ \opn{Hom}_{\dcat{K}(A, \cat{M})}(N, J) \to 
\opn{Hom}_{\dcat{K}(A, \cat{M})}(N, K)
\to \opn{Hom}_{\dcat{K}(A, \cat{M})}(N, \opn{T}(I))  \]
in $\cat{Mod} \K$. Because $J$ and $\opn{T}(I)$ are 
K-injectives, Proposition \ref{prop:1515} says that
\[ \opn{Hom}_{\dcat{K}(A, \cat{M})}(N, J) = 0 = 
\opn{Hom}_{\dcat{K}(A, \cat{M})}(N, \opn{T}(I)) . \]
Therefore 
$\opn{Hom}_{\dcat{K}(A, \cat{M})}(N, K) = 0$. 
But $N$ is an arbitrary acyclic DG module, so $K$ is K-injective.
\end{proof}

\begin{exa} \label{exa:1520}
Let $\star$ be some boundedness condition (namely $\mrm{b}$, $+$ or $-$). 
We know that $\dcat{K}^{\star}(A, \cat{M})$ 
is a full triangulated subcategory of $\dcat{K}^{}(A, \cat{M})$.
Hence $\dcat{K}^{\star}(A, \cat{M})_{\mrm{inj}}$
is a triangulated subcategory too.  
\end{exa}

\begin{dfn} \label{dfn:1508}
Let $\cat{K}$ be a full triangulated subcategory of 
$\dcat{K}(A, \cat{M})$. We say that $\cat{K}$ 
{\em has enough K-injectives}% 
\index{Resolution! K-injective}
if every DG module $M \in \cat{K}$ admits a K-injective resolution inside 
$\cat{K}$. I.e.\ there is a quasi-isomorphism $\rho : M \to I$
where $I \in \cat{K}_{\mrm{inj}}$.
\end{dfn}

Here is the crucial fact regarding K-injectives. 

\begin{lem} \label{lem:1510}
Let  $s : I \to M$ be a quasi-isomorphism in $\dcat{K}(A, \cat{M})$,
and assume $I$ is K-injective. Then $s$ has a left inverse, namely 
there is a morphism $t : M \to I$ in $\dcat{K}(A, \cat{M})$ such that 
$t \circ s = \opn{id}_{I}$. 
Moreover, this morphism $t$ is unique. 
\end{lem}

\begin{proof}
Consider a distinguished triangle 
$I \xar{s} M \to N \xar{\, \triangle \, }$
in $\dcat{K}(A, \cat{M})$ that's built on $s$. The long exact 
cohomology sequence tells us that $N$ is an acyclic DG module. 
Therefore
$\opn{Hom}_{\dcat{K}^{}(A, \cat{M})}(\opn{T}^p(N), I) = 0$ 
for all $p$. 
The exact sequence 
\[ \begin{aligned}
& \opn{Hom}_{\dcat{K}(A, \cat{M})}(N, I) \to 
\opn{Hom}_{\dcat{K}(A, \cat{M})}(M, I)
\\
& \quad 
\to \opn{Hom}_{\dcat{K}(A, \cat{M})}(I, I) 
\to \opn{Hom}_{\dcat{K}(A, \cat{M})}(\opn{T}^{-1}(N), I) 
\end{aligned} \]
shows that the homomorphism
\[ (-) \circ s : \opn{Hom}_{\dcat{K}(A, \cat{M})}(M, I) \iso 
\opn{Hom}_{\dcat{K}(A, \cat{M})}(I, I) \]
is a bijection.
We take $t : M \to I$ to be the unique morphism in $\dcat{K}(A, \cat{M})$
such that $t \circ s = \opn{id}_I$. 
\end{proof}

\begin{thm} \label{thm:3135}
\index{Differential graded module! K-injective}
Let $A$ be a DG ring, let $\cat{M}$ be an abelian category, and let
$\cat{K}$ be a full triangulated subcategory of $\dcat{K}(A, \cat{M})$.
Write $\cat{S} := \cat{K} \cap \, \dcat{S}(A, \cat{M})$,
and let $\cat{D} := \cat{K}_{\cat{S}}$ be the localization, with localization 
functor $\opn{Q} : \cat{K} \to \cat{D}$. 
Then for every $M \in \cat{K}$ and $I \in \cat{K}_{\mrm{inj}}$
the homomorphism  
\[ \opn{Q}_{M, I} : \opn{Hom}_{\cat{K}}(M, I) \to \opn{Hom}_{\cat{D}}(M, I) \]
is bijective.  
\end{thm}

\begin{proof}
We know that $\opn{Q} : \cat{K} \to \cat{D}$ is a left Ore localization with 
respect to $\cat{S}$. 
Suppose $q : M \to I$ is a morphism in $\cat{D} = \cat{K}_{\cat{S}}$.
Let us present $q$ as a left fraction: 
$q = \opn{Q}(s)^{-1} \circ \opn{Q}(a)$, 
where $a : M \to N$ and $s : I \to N$ are morphisms in 
$\cat{K}$, and $s$ is a quasi-isomorphism. 
By Lemma \ref{lem:1510}, $s$ has a left inverse 
$t : M \to I$ in $\cat{K}$. We get a morphism 
$t \circ a : M \to I$ in $\cat{K}$, and an easy calculation shows that
$\opn{Q}(t \circ a) = q$ in $\cat{D}$.
This proves surjectivity of the function $\opn{Q}_{M, I}$.

Now let's prove injectivity of $\opn{Q}_{M, I}$.  
If $a : M \to I$ is a morphism in $\cat{K}$ such that 
$\opn{Q}_{M, I}(a) = 0$, then by axiom (LO4) of left Ore localization (see 
Definition \ref{dfn:2310}) there is a quasi-isomorphism $s : I \to L$ in 
$\cat{K}$ such that $s \circ a = 0$ in $\cat{K}$. Let $t$ be the left inverse 
of 
$s$. Then $a = t \circ s \circ a= 0$ in $\cat{K}$. 
\end{proof}

\begin{cor} \label{cor:3135}
In the situation of Theorem \tup{\ref{thm:3135}}, 
the localization functor 
$\opn{Q} : \cat{K}_{\mrm{inj}} \to \cat{D}$ 
is fully faithful.
\end{cor}

\begin{proof}
Since $\cat{K}_{\mrm{inj}}$ is a full subcategory of 
$\cat{K}$, for every $J, I \in \cat{K}_{\mrm{inj}}$
we have bijections
\[ \opn{Hom}_{\cat{K}_{\mrm{inj}}}(J, I) \iso \opn{Hom}_{\cat{K}}(J, I) 
\xar{\opn{Q}_{J, I}} \opn{Hom}_{\cat{D}}(J, I) , \]
where the bijection $\opn{Q}_{J, I}$ is by the theorem above. 
\end{proof}

\begin{cor} \label{cor:1515}
In the situation of Theorem \tup{\ref{thm:3135}}, if
$\cat{K}$ has enough K-injectives, then the localization functor
$\opn{Q} : \cat{K}_{\mrm{inj}} \to \cat{D}$
is an equivalence.
\end{cor}

\begin{proof}
By Corollary \ref{cor:3135} the functor $\opn{Q}$ is fully faithful. The 
extra condition guarantees that $\opn{Q}$ is essentially surjective on objects. 
\end{proof}

\begin{cor} \label{cor:1509}
Let $\star$ be any boundedness condition. If $\dcat{K}^{\star}(A, \cat{M})$ has 
enough K-injectives, then the triangulated functor
\[ \opn{Q} : \dcat{K}^{\star}(A, \cat{M})_{\mrm{inj}} \to
\dcat{D}^{\star}(A, \cat{M}) \]
is an equivalence.
\end{cor}

\begin{proof}
Since $\dcat{K}^{\star}(A, \cat{M})$ is a full 
triangulated subcategory of $\dcat{K}(A, \cat{M})$,
this is a special case of the Corollary \ref{cor:1515}.
\end{proof}

\begin{rem} \label{rem:1511}
These last results are of tremendous importance, both theoretically and 
practically. In the theory, Corollary \ref{cor:1509}
shows that the localized category 
$\dcat{D}^{\star}(A, \cat{M})$, 
which is too big to lie inside the original universe $\cat{U}$ (see Remark 
\ref{rem:1405}), is equivalent to a  $\cat{U}$-category. 
On the practical side, Corollary \ref{cor:3135} means that among K-injective 
objects we do not need fractions to represent morphisms. 
\end{rem}

\begin{cor} \label{cor:1511}
Let $\cat{K} \sub \cat{K}'$ be full triangulated subcategories of 
$\dcat{K}^{}(A, \cat{M})$.
Define 
$\cat{S} := \cat{K} \cap \, \dcat{S}(A, \cat{M})$,
$\cat{S}' := \cat{K}' \cap \, \dcat{S}(A, \cat{M})$,
$\cat{D} := \cat{K}_{\cat{S}}$
and
$\cat{D}' := \cat{K}'_{\cat{S}'}$.
If $\cat{K}$ and $\cat{K}'$ have enough K-injectives, then the canonical 
functor $\cat{D} \to \cat{D}'$ is fully faithful. 
\end{cor}

\begin{proof}
Combine Corollary \ref{cor:1515} with the fact that the canonical functor
$\cat{K} \to \cat{K}'$ is fully faithful.
\end{proof}

\begin{cor} \label{cor:1922}
Let $\phi : I \to J$ be a morphism in $\dcat{C}_{\mrm{str}}(A, \cat{M})$ 
between K-injective objects. Then $\phi$ is a homotopy equivalence if and only 
if it is a quasi-isomorphism. 
\end{cor}

\begin{proof}
One implication is trivial. For the reverse implication, if $\phi$ is a 
quasi-isomorphism then it is an isomorphism in 
$\dcat{D}(A, \cat{M})$, and by Corollary \ref{cor:3135} for 
$\cat{K} = \dcat{K}(A, \cat{M})$
we see that $\phi$ is an isomorphism in $\dcat{K}(A, \cat{M})$. 
\end{proof}
   
\begin{thm}  \label{thm:5050}   
\index{Derived functor! triangulated right}
In the situation of Theorem \tup{\ref{thm:3135}}, assume that $\cat{K}$ has 
enough K-injectives. Let $\cat{E}$ be a triangulated category, and let 
$F : \cat{K} \to \cat{E}$
be a triangulated functor. Then $F$ has a right derived functor 
$(\mrm{R} F, \eta^{\mrm{R}}) : \cat{D} \to \cat{E}$. 
Furthermore, for every $I \in \cat{K}_{\mrm{inj}}$ the morphism 
$\eta^{\mrm{R}}_I : F(I) \to \mrm{R} F(I)$
in $\cat{E}$ is an isomorphism. 
\end{thm}

\begin{proof}
We will use Theorem \ref{thm:1470}. In the notation of that theorem, let 
$\cat{J} := \cat{K}_{\mrm{inj}}$.
Condition (a) of that theorem holds (this is the ``enough K-injectives'' 
assumption). Next, Corollary \ref{cor:1922} implies that every 
quasi-isomorphism $\phi : I \to J$ in $\cat{K}_{\mrm{inj}}$ is actually an 
isomorphism. Therefore $F(\phi)$ is an isomorphism in $\cat{E}$, and this is 
condition (b) of  Theorem \ref{thm:1470}. 
\end{proof}

Here is another useful definition. It is a variant of Definition \ref{dfn:2082}.

\begin{dfn} \label{dfn:1512}
Let $\cat{K}$ be a full triangulated subcategory of 
$\dcat{K}(A, \cat{M})$, and assume $\cat{K}$ has enough K-injectives.
A  {\em system of K-injective resolutions} in $\cat{K}$ is a pair $(I, \rho)$, 
where $I : \opn{Ob}(\cat{K}) \to \opn{Ob}(\cat{K}_{\mrm{inj}})$
is a function, and 
$\rho = \{ \rho_M \}_{M \in \opn{Ob}(\cat{K})}$ 
is a collection of quasi-isomorphisms
$\rho_M : M \to I(M)$ in $\cat{K}$.
Moreover, if $M \in \opn{Ob}(\cat{K}_{\mrm{inj}})$, then 
$I(M) = M$ and $\rho_M = \opn{id}_M$. 
\end{dfn}

\begin{exa} \label{exa:1530}
Let $A$ be any DG ring. We will prove later that 
$\dcat{K}^{}(A)$ has enough K-injectives. Therefore, given a triangulated 
functor $F : \dcat{K}^{}(A) \to \cat{E}$ into a triangulated category 
$\cat{E}$, the right derived functor  
$(\mrm{R} F, \eta^{\mrm{R}}) : \dcat{D}(A) \to \cat{E}$ 
exists. Suppose we choose a system  of K-injective resolutions $(I, \rho)$ in 
$\dcat{K}(A)$. Then we get a presentation of $(\mrm{R} F, \eta^{\mrm{R}})$ as 
follows: $\mrm{R} F(M) := F(I(M))$ and $\eta^{\mrm{R}}_M := F(\rho_M)$. 
\end{exa}

The proposition below is a variant of Lemma \ref{lem:2081}.

\begin{prop} \label{prop:1510}
In the situation of Theorem \tup{\ref{thm:3135}}, suppose $\cat{K}$ has enough 
K-injectives, and let $(I, \rho)$ be a system of K-injective resolutions in 
$\cat{K}$. Then the function $I$ extends uniquely to a triangulated functor
$I : \cat{D} \to \cat{K}_{\mrm{inj}}$, such that 
$\opn{Id}_{\cat{K}_{\mrm{inj}}} = I \circ \opn{Q}|_{\cat{K}_{\mrm{inj}}}$, 
and 
$\rho : \opn{Id}_{\cat{D}} \twoto \opn{Q} \circ \, I$
is an isomorphism of triangulated functors.
\end{prop}

\begin{proof}
The proof is the same as that of Lemma \ref{lem:2081}, 
except that here we use Corollary \ref{cor:1515}. 
\end{proof}

The next corollary is a categorical interpretation of the last proposition.

\begin{cor}[Functorial K-Injective Resolutions] \label{cor:1920}
Let $\cat{K}$ be a full triangulated subcategory of 
$\dcat{K}(A, \cat{M})$, and assume $\cat{K}$ has enough K-injectives.
\begin{enumerate}
\item There is a triangulated functor 
$I : \cat{K} \to \cat{K}$ and a morphism of triangulated functors 
$\rho : \opn{Id}_{\cat{K}} \twoto I$, such that for every object $M \in \cat{K}$
the object $I(M)$ is K-injective, and the morphism 
$\rho_M : M \to I(M)$ is a quasi-isomorphism. 

\item If $(I', \rho')$ is another such pair, then there is a unique isomorphism 
of triangulated functors $\ze : I \twoiso I'$ such that 
$\rho' = \ze \circ \rho$. 
\end{enumerate}
\end{cor}

\begin{exer} \label{exer:1920}
Prove Corollary \ref{cor:1920}.
\end{exer}

Infinite direct sums in $\dcat{C}_{\mrm{str}}(A, \cat{M})$ were discussed in 
Proposition \ref{prop:3140}. Their existence depends only on the existence of 
infinite direct sums in $\cat{M}$. 

\begin{thm} \label{thm:3140}
\index{Direct sums! infinite {\indash} in $\dcat{D}(A, \cat{M})$}
Let $A$ be a DG ring, let $\cat{M}$ be an abelian category, and let 
$\{ M_x \}_{x \in X}$ be a collection of objects of 
$\dcat{C}(A, \cat{M})$. Assume that $\dcat{C}(A, \cat{M})$ has enough 
K-injectives, and that the direct sum 
$M := \bigoplus_{x \in X} M_x$
exists in $\dcat{C}_{\mrm{str}}(A, \cat{M})$. 
Then $M$ is the direct sum of the collection $\{ M_x \}_{x \in X}$ in 
$\dcat{D}(A, \cat{M})$. 
\end{thm}

\begin{proof}
Let's use these shorthands: 
$\cat{C}_{\mrm{str}} := \dcat{C}_{\mrm{str}}(A, \cat{M})$,
$\cat{K} := \dcat{K}(A, \cat{M})$ and 
$\cat{D} := \dcat{D}(A, \cat{M})$.
For each index $x$ there is a morphism $e_x : M_x \to M$ in 
$\cat{C}_{\mrm{str}}$,
that passes to a morphism 
$\opn{Q}(e_x) : M_x \to M$ in $\cat{D}$. We need to verify that the collection 
of morphisms $\{ \opn{Q}(e_x) \}_{x \in X}$ has the universal property of a 
coproduct. 

Take any object $N \in \cat{D}$, and choose a K-injective 
resolution $N \to J$ in $\cat{C}_{\mrm{str}}$.
There is an isomorphism of DG $\K$-modules 
\begin{equation} \label{eqn:3140}
\begin{aligned}
&
\opn{Hom}_{A, \cat{M}}(M, J) = 
\opn{Hom}_{A, \cat{M}} \Bigl( \bigoplus\nolimits_{x \in X} M_x, \, J \Bigr)
\\
& \quad 
\cong \prod\nolimits_{x \in X} \opn{Hom}_{A, \cat{M}}(M_x, J) .
\end{aligned}
\end{equation}
From it we obtain isomorphisms of $\K$-modules 
\begin{equation} \label{eqn:3145}
\begin{aligned}
&
\opn{Hom}_{\cat{D}}(M, N) \cong^{\dag} \opn{Hom}_{\cat{K}}(M, J)
= \opn{H}^0 \bigl( \opn{Hom}_{A, \cat{M}}(M, J) \bigr) 
\\
& \quad 
\cong^{\ddag} \opn{H}^0 \Bigl( 
\prod\nolimits_{x \in X} \opn{Hom}_{A, \cat{M}}(M_x, J) \Bigr) 
\\
& \quad 
\cong^{\sharp} \prod\nolimits_{x \in X}
\opn{H}^0 \bigl( \opn{Hom}_{A, \cat{M}}(M_x, J) \bigr) 
\\
& \quad 
= \prod\nolimits_{x \in X} \opn{Hom}_{\cat{K}}(M_x, J)
% \\
% & \qquad 
\cong^{\dag} \prod\nolimits_{x \in X} \opn{Hom}_{\cat{D}}(M_x, N) . 
\end{aligned} 
\end{equation}
The isomorphisms $\cong^{\dag}$ are due to Theorem \ref{thm:3135}; the 
isomorphism $\cong^{\ddag}$ is from formula (\ref{eqn:3140}); and the 
isomorphism $\cong^{\sharp}$ is because the functor $\opn{H}^0$ commutes 
with products. Since the composed isomorphism in (\ref{eqn:3145}) is induced by 
the collection of morphisms $\{ \opn{Q}(e_x) \}_{x \in X}$, we see that the 
universal property of a coproduct holds. 
\end{proof}

\begin{rem} \label{rem:5055}
The concept of a K-injective complex (i.e.\ a K-injective object of 
$\dcat{C}(\cat{M})$), as well as the similar concepts of K-projective and 
K-flat complexes that will be discussed below, were introduced by N. 
Spaltenstein \cite{Spa} in 1988. At about the same time other authors (B. Keller 
\cite{Kel1}, M. Bockstedt and A. Neeman \cite{BoNe}, J. Bernstein and V. 
Lunts \cite{BeLu}) discovered this concept 
independently, with other names (such as {\em homotopically injective complex}).
The texts \cite{BeLu} and \cite{Kel1} already talk about DG modules over DG 
rings. 
\end{rem}






\mysubsection{K-Projective DG Modules} \label{subsec:K-proj}

This subsection is dual to the previous one, and so we will be brief. 

\begin{dfn} \label{dfn:1520}
A DG module $P \in \dcat{C}(A, \cat{M})$ is called 
{\em K-projective}% 
\index{Differential graded module! K-projective}
if for every acyclic DG module $N \in \dcat{C}(A, \cat{M})$, the DG $\K$-module 
$\opn{Hom}_{A, \cat{M}}(P, N)$ is acyclic.
\end{dfn}

\begin{dfn} \label{dfn:1521}
Let $M \in \dcat{C}(A, \cat{M})$. A 
{\em K-projective resolution}% 
\index{Resolution! K-projective}
of $M$ is a quasi-isomorphism 
$\rho : P \to M$ in $\dcat{C}_{\mrm{str}}(A, \cat{M})$, where 
$P$ is a K-projective DG module. 
\end{dfn}

Remark \ref{rem:5051} on terminology applies here too. 

\begin{prop} \label{prop:1520}
A DG module $P \in \dcat{K}^{}(A, \cat{M})$ is K-projective if and only if 
$\opn{Hom}_{\dcat{K}(A, \cat{M})}(P, N) = 0$ for every acyclic DG module 
$N \in \dcat{K}^{}(A, \cat{M})$.
\end{prop}

The proof is like that of Proposition \ref{prop:1515}.



\begin{dfn} \label{dfn:1522}
Let $\cat{K}$ be a full subcategory of $\dcat{K}^{}(A, \cat{M})$. 
The full subcategory of  $\cat{K}$ on the K-projective DG modules in it is 
denoted by $\cat{K}_{\mrm{prj}}$. In other words, 
$\cat{K}_{\mrm{prj}} = \dcat{K}^{}(A, \cat{M})_{\mrm{prj}} \cap \cat{K}$.
\end{dfn}

The warning in Remark  \ref{rem:5050} applies here too. 

\begin{prop}  \label{prop:1521}
If  $\cat{K}$ is a full triangulated subcategory of $\dcat{K}^{}(A, 
\cat{M})$, then $\cat{K}_{\mrm{prj}}$ is a full triangulated subcategory of 
$\cat{K}$.
\end{prop}

The proof is like that of Proposition \ref{prop:130}.

\begin{exa} \label{exa:1521}
Let $\star$ be some boundedness condition (namely $\mrm{b}$, $+$ or $-$). 
Since $\dcat{K}^{\star}(A, \cat{M})$ 
is a full triangulated subcategory of $\dcat{K}^{}(A, \cat{M})$,
we see that $\dcat{K}^{\star}(A, \cat{M})_{\mrm{prj}}$
is a full triangulated subcategory too.  
\end{exa}

\begin{dfn} \label{dfn:1523}
Let $\cat{K}$ be a full triangulated subcategory of 
$\dcat{K}(A, \cat{M})$. We say that $\cat{K}$ 
{\em has enough K-projectives}%
\index{Resolution! K-projective}
if every DG module $M \in \cat{K}$ admits a K-projective resolution inside 
$\cat{K}$. I.e.\ there is a quasi-isomorphism $\rho : P \to M$
where $P \in \cat{K}_{\mrm{prj}}$.
\end{dfn}

\begin{lem} \label{lem:1521}
Let  $s : M \to P$ be a quasi-isomorphism in 
$\dcat{K}(A, \cat{M})$, and assume $P$ is K-projective. Then $s$ has 
a right inverse; namely there is a morphism $t : P \to M$ in
$\dcat{K}(A, \cat{M})$ such that $s \circ t = \opn{id}_P$.
Moreover, this morphism $t$ is unique.  
\end{lem}

The proof is almost the same as that of Lemma \ref{lem:1510}.

\begin{thm} \label{thm:3145}
\index{Differential graded module! K-projective}
Let $A$ be a DG ring, let $\cat{M}$ be an abelian category, and let
$\cat{K}$ be a full triangulated subcategory of $\dcat{K}(A, \cat{M})$.
Write $\cat{S} := \cat{K} \cap \, \dcat{S}(A, \cat{M})$,
and let $\cat{D} := \cat{K}_{\cat{S}}$ be the localization, with localization 
functor $\opn{Q} : \cat{K} \to \cat{D}$. 
Then for every $M \in \cat{K}$ and $P \in \cat{K}_{\mrm{prj}}$
the homomorphism  
\[ \opn{Q}_{P, M} : \opn{Hom}_{\cat{K}}(P, M) \to \opn{Hom}_{\cat{D}}(P, M) 
\]
is bijective.  
\end{thm}

The proof is like that of Theorem \ref{thm:3135}, only now we use the fact that
the functor $\opn{Q}$ is a right Ore localization. 

The results from here to the end of the subsection are also proved like their 
K-injective counterparts in Subsection \ref{subsec:K-inj}. 

\begin{cor} \label{cor:3145} 
In the situation of Theorem \tup{\ref{thm:3145}}, 
the localization functor 
$\opn{Q} : \cat{K}_{\mrm{prj}} \to \cat{D}$
is fully faithful.
\end{cor}
        
\begin{cor} \label{cor:1520}
In the situation of Theorem \tup{\ref{thm:3145}}, 
if $\cat{K}$ has enough K-pro\-jectives, then the localization functor
$\opn{Q} : \cat{K}_{\mrm{prj}} \to \cat{D}$
is an equivalence.
\end{cor}

\begin{cor} \label{cor:3146}
Let $\star$ be any boundedness condition. If $\dcat{K}^{\star}(A, \cat{M})$ has 
enough K-projectives, then the triangulated functor
\[ \opn{Q} : \dcat{K}^{\star}(A, \cat{M})_{\mrm{prj}} \to
\dcat{D}^{\star}(A, \cat{M}) \]
is an equivalence.
\end{cor}

\begin{cor} \label{cor:1521}
Let $\cat{K} \sub \cat{K}'$ be full triangulated subcategories of 
$\dcat{K}^{}(A, \cat{M})$.
Define 
$\cat{S} := \cat{K} \cap \, \dcat{S}(A, \cat{M})$,
$\cat{S}' := \cat{K}' \cap \, \dcat{S}(A, \cat{M})$,
$\cat{D} := \cat{K}_{\cat{S}}$
and
$\cat{D}' := \cat{K}'_{\cat{S}'}$.
If $\cat{K}$ and $\cat{K}'$ have enough K-projectives, then the canonical 
functor $\cat{D} \to \cat{D}'$ is fully faithful. 
\end{cor}

\begin{cor} \label{cor:1923} 
Let $\phi : P \to Q$ be a morphism in $\dcat{C}_{\mrm{str}}(A, \cat{M})$ 
between K-projective objects. Then $\phi$ is a homotopy equivalence if and only 
if it is a quasi-isomorphism. 
\end{cor}

\begin{thm}  \label{thm:5051}  
\index{Derived functor! triangulated left}
In the situation of Theorem \tup{\ref{thm:3145}}, assume that $\cat{K}$ has 
enough K-projectives. Let $\cat{E}$ be a triangulated category, and let 
$F : \cat{K} \to \cat{E}$
be a triangulated functor. Then $F$ has a left derived functor 
$(\mrm{L} F, \eta^{\mrm{L}}) : \cat{D} \to \cat{E}$.
Furthermore, for every $P \in \cat{K}_{\mrm{prj}}$ the morphism 
$\eta^{\mrm{L}}_P :\mrm{L} F(P) \to F(P)$
in $\cat{E}$ is an isomorphism. 
\end{thm}

\begin{dfn} \label{dfn:1524}
Let $\cat{K}$ be a full triangulated subcategory of 
$\dcat{K}(A, \cat{M})$, and assume $\cat{K}$ has enough K-projectives.
A  {\em system of K-projective resolutions} in $\cat{K}$ is a pair $(P, \rho)$, 
where $P : \opn{Ob}(\cat{K}) \to \opn{Ob}(\cat{K}_{\mrm{prj}})$
is a function, and 
$\rho = \{ \rho_M \}_{M \in \opn{Ob}(\cat{K})}$ 
is a collection of quasi-isomorphisms
$\rho_M : P(M) \to M$ in $\cat{K}$.
Moreover, if $M \in \opn{Ob}(\cat{K}_{\mrm{prj}})$, then 
$P(M) = M$ and $\rho_M = \opn{id}_M$. 
\end{dfn}

\begin{exa} \label{exa:1531}
Let $A$ be any DG ring. We will prove later that $\bcat{K}^{}(A)$
has enough K-projectives. Therefore, given a triangulated functor 
$F : \bcat{K}^{}(A) \to \cat{E}$ into a triangulated category $\cat{E}$, 
the left derived functor  
$(\mrm{L} F, \eta^{\mrm{L}}) : \bcat{D}^{}(A) \to \cat{E}$
exists. Suppose we choose a system  of K-projective resolutions $(P, \rho)$ in 
$\bcat{K}(A)$. Then we get a presentation of $(\mrm{L} F, \eta^{\mrm{R}})$ as 
follows: $\mrm{L} F(M) := F(P(M))$ and $\eta^{\mrm{L}}_M := F(\rho_M)$. 
\end{exa}

\begin{prop} \label{prop:1522}
In the situation of Theorem \tup{\ref{thm:3145}}, suppose $\cat{K}$ has enough 
K-projectives, and let $(P, \rho)$ be a system  of K-projective resolutions.
Then the function $P$ extends uniquely to a triangulated functor
$P : \cat{D} \to \cat{K}_{\mrm{prj}}$, such that 
$\opn{Id}_{\cat{K}_{\mrm{prj}}} = P \circ \opn{Q}|_{\cat{K}_{\mrm{prj}}}$, 
and 
$\rho :  \opn{Q} \circ \, P  \twoto \opn{Id}_{\cat{D}}$
is an isomorphism of triangulated functors.
\end{prop}

\begin{cor}[Functorial K-Projective Resolutions] \label{cor:1921}
Let $\cat{K}$ be a full triangulated subcategory of 
$\dcat{K}(A, \cat{M})$, and assume $\cat{K}$ has enough K-projectives.
\begin{enumerate}
\item There is a triangulated functor 
$P : \cat{K} \to \cat{K}$ and a morphism of triangulated functors 
$\rho : P \twoto \opn{Id}_{\cat{K}}$, such that for every object 
$M \in \cat{K}$
the object $P(M)$ is K-projective, and the morphism 
$\rho_M : P(M) \to M$ is a quasi-isomorphism. 

\item If $(P', \rho')$ is another such pair, then there is a unique isomorphism 
of triangulated functors $\ze : P' \twoiso P$ such that 
$\rho' =  \rho \circ \ze$. 
\end{enumerate}
\end{cor}




\mysubsection{K-Flat DG Modules} \label{subsec:K-flat}

Recall that for a DG ring $A$, its opposite DG ring is $A^{\mrm{op}}$.
The objects of $\dcat{C}(A^{\mrm{op}})$ are the right DG $A$-modules. 

\begin{dfn} \label{dfn:1525}
A DG module $P \in \dcat{C}(A)$ is called 
{\em K-flat}% 
\index{Differential graded module! K-flat}
if for every acyclic DG module $N \in \dcat{C}(A^{\mrm{op}})$, the DG 
$\K$-module $N \ot_A P$ is acyclic.
\end{dfn}

\begin{dfn} \label{dfn:4720}
Let $M \in \dcat{C}(A)$. A 
{\em K-flat resolution}% 
\index{Resolution! K-flat}
of $M$ is a quasi-iso\-morphism $\rho : P \to M$ in 
$\dcat{C}_{\mrm{str}}(A)$, where $P$ is a K-flat DG module. 
\end{dfn}

Remark \ref{rem:5051} on terminology applies here as well. 

\begin{dfn} \label{dfn:3145}
Let $\cat{K}$ be a full triangulated subcategory of $\bcat{K}(A)$.
We denote by $\cat{K}_{\mrm{flat}}$ the full subcategory of 
$\bcat{K}(A)$ on the K-flat complexes in it. 
Thus $\cat{K}_{\mrm{flat}} = \cat{K} \cap \bcat{K}(A)_{\mrm{flat}}$.
\end{dfn}

The warning in Remark  \ref{rem:5050} applies here too. 

\begin{prop} \label{prop:1525}
If $P \in \dcat{C}(A)$ is K-projective then it is K-flat. 
\end{prop}

\begin{proof}
Let $\K^*$ be an injective cogenerator of $\dcat{M}(\K) = \cat{Mod} \K$.
This means that $\K^*$ is an injective $\K$-module, such that every nonzero 
$\K$-module $W$ admits a nonzero homomorphism $W \to \K^*$. 
See Examples \ref{exa:1665} and \ref{exa:1665}. 
It is not hard to see that a DG $\K$-module $W$ is acyclic if and only if 
$\opn{Hom}_{\K}(W, \K^*)$ is acyclic. 

Take an acyclic complex $N \in \dcat{C}(A^{\mrm{op}})$. Then by Hom-tensor 
adjunction there is an isomorphism of DG $\K$-modules
\[ \opn{Hom}_{\K}(N \ot_A P, \K^*) \cong 
\opn{Hom}_{A} \bigl(P, \opn{Hom}_{\K}(N, \K^*) \bigr) . \]
The right side is acyclic by our assumptions. Hence so is the left side. It 
follows that $N \ot_A P$ is acyclic. 
\end{proof}

\begin{prop} \label{prop:1526}
A DG module $P \in \dcat{K}(A)$ is K-flat iff 
\[ \opn{Hom}_{\dcat{K}(A)} \bigl(P, \opn{Hom}_{\K}(N, J) \bigr) = 0 \]
for every acyclic $N \in \dcat{C}(A^{\mrm{op}})$ and every injective 
$J \in \cat{Mod} \K$.
\end{prop}

\begin{exer}
Prove Proposition \ref{prop:1526}. (Hint: look at the proof of Proposition 
\ref{prop:1525}.)
\end{exer}

\begin{prop} \label{prop:3145}
Let $\cat{K}$ be a full triangulated subcategory of $\bcat{K}(A)$.
Then $\cat{K}_{\mrm{flat}}$ is a full triangulated subcategory of
$\cat{K}$. 
\end{prop}

\begin{exer} \label{exer:3145}
Prove Proposition \ref{prop:3145}. (Hint: The challenge is to show that 
$\cat{K}_{\mrm{flat}}$ is closed under distinguished triangles. 
Use Proposition \ref{prop:1526}.)
\end{exer}

In Theorem \ref{thm:2107} on the left derived bifunctor 
$(- \ot^{\mrm{L}}_{A} -)$ of the bifunctor $(- \ot_A -)$ we will use K-flat 
resolutions. 

\begin{rem} \label{rem:1525} 
In view of Proposition \ref{prop:1525}, the reader might wonder why we bother 
with K-flat DG modules. The reason is that there are situations in which there 
are enough K-flat objects but not enough K-projective objects. One such 
situation is that of the category 
$\dcat{C}(\mcal{A}) = \dcat{C}(\cat{Mod} \mcal{A})$,
where $(X, \mcal{A})$ is a ringed space and $\cat{Mod} \mcal{A}$
is the category of sheaves of $\AA$-modules; see Example \ref{exa:4725}. 
Another situation is when K-flatness is needed is for DG bimodules over 
noncommutative DG rings, as in Sections \ref{sec:perf-tilt-NC} and later.
\end{rem}



\mysubsection{Opposite Resolving Subcategories}
\label{subsec:resol-op}

In this subsection we explain how K-injective and K-projective DG modules in 
$\dcat{K}(A, \cat{M})$ behave in the opposite category 
$\dcat{K}(A, \cat{M})^{\mrm{op}}$, and then use this 
to prove existence of contravariant triangulated derived functors.

Recall the flipping operation from Subsections \ref{subsec:contrvar-dg-func}, 
\ref{subsec:opp-hom-triang} and \ref{subsec:opp-dercat-triang}.
There is an isomorphism of DG categories 
\begin{equation} \label{eqn:3150}
\opn{Flip} : \dcat{C}(A, \cat{M})^{\mrm{op}} \iso 
\dcat{C}(A, \cat{M})^{\mrm{flip}} = 
\dcat{C}(A^{\mrm{op}}, \cat{M}^{\mrm{op}}) , 
\end{equation}
and the corresponding isomorphisms of the homotopy categories
\begin{equation} \label{eqn:3151}
\opn{Flip} : \dcat{K}(A, \cat{M})^{\mrm{op}} \iso 
\dcat{K}(A, \cat{M})^{\mrm{flip}} = 
\dcat{K}(A^{\mrm{op}}, \cat{M}^{\mrm{op}})
\end{equation}
and on the derived categories
\begin{equation} \label{eqn:3152}
\opn{Flip} : \dcat{D}(A, \cat{M})^{\mrm{op}} \iso 
\dcat{D}(A, \cat{M})^{\mrm{flip}} = 
\dcat{D}(A^{\mrm{op}}, \cat{M}^{\mrm{op}}) . 
\end{equation}
By definition of the triangulated structure on 
$\dcat{K}(A, \cat{M})^{\mrm{op}}$,
the isomorphisms (\ref{eqn:3151}) and  (\ref{eqn:3152})
are of triangulated categories. 

The objects of $\dcat{C}(A, \cat{M})$ and
$\dcat{C}(A, \cat{M})^{\mrm{op}}$ are the same, and the property of a DG module 
$M$ being acyclic is the same in both categories. 

\begin{lem} \label{lem:3150}
A DG module $M \in \dcat{C}(A, \cat{M})$ is acyclic iff the DG module \lb 
$\opn{Flip}(M) \in \dcat{C}(A^{\mrm{op}}, \cat{M}^{\mrm{op}})$
is acyclic. 
\end{lem}

\begin{proof}
This follows from Theorem \ref{thm:2495}(4). 
\end{proof}

\begin{lem} \label{lem:3151}
A DG module $P \in \dcat{K}(A, \cat{M})$ is K-projective iff
the DG module 
$I := \opn{Flip}(P) \in \dcat{K}(A^{\mrm{op}}, \cat{M}^{\mrm{op}})$
is K-injective. 
\end{lem}

\begin{proof} 
According to Lemma \ref{lem:3150}, the functor 
$\opn{Flip}$ gives a bijection between the set of acyclic DG modules 
$M \in \dcat{K}(A, \cat{M})$
and the set of acyclic DG modules 
$ N := \opn{Flip}(M) \in \dcat{K}(A^{\mrm{op}}, \cat{M}^{\mrm{op}})$.
Since the functor (\ref{eqn:3151}) is an equivalence, we 
have an equality and an isomorphism
\begin{equation} \label{eqn:5055}
\opn{Hom}_{\dcat{K}(A, \cat{M})} (P, M) = 
\opn{Hom}_{\dcat{K}(A, \cat{M})^{\mrm{op}}} (M, P) \cong 
\opn{Hom}_{\dcat{K}(A^{\mrm{op}}, \cat{M}^{\mrm{op}})} (N, I)  
\end{equation}
in $\dcat{M}(\K)$.

By Proposition \ref{prop:1515}, the DG 
module $I$ is K-injective in 
$\dcat{K}(A^{\mrm{op}}, \cat{M}^{\mrm{op}})$
iff for every acyclic DG module 
$N \in \dcat{K}(A^{\mrm{op}}, \cat{M}^{\mrm{op}})$
we have 
$\opn{Hom}_{\dcat{K}(A^{\mrm{op}}, \cat{M}^{\mrm{op}})}(N, I) = 0$.
Similarly, by Proposition \ref{prop:1520}, the DG 
module $P$ is K-projective in $\dcat{K}(A, \cat{M})$
iff for every acyclic DG module 
$M \in \dcat{K}(A, \cat{M})$ we have 
$\opn{Hom}_{\dcat{K}(A, \cat{M})}(P, M) = 0$.
In view of (\ref{eqn:5055}) we conclude that $I$ is K-injective iff $P$ is 
K-projective.
\end{proof}

Here are the contravariant modifications of Theorems \ref{thm:5050} and 
\ref{thm:5051}. 

\begin{setup} \label{set:3150}
We are given:
\begin{itemize}
\item A full additive subcategory
$\cat{K} \sub \dcat{K}(A, \cat{M})$, 
s.t.\ $\cat{K}^{\mrm{op}} \sub \dcat{K}(A, \cat{M})^{\mrm{op}}$
is triangulated. The localization of $\cat{K}^{\mrm{op}}$ at the 
quasi-isomorphisms in it is denoted by $\cat{D}^{\mrm{op}}$.  

\item A triangulated category $\cat{E}$ and  a triangulated functor
$F : \cat{K}^{\mrm{op}} \to \cat{E}$.
\end{itemize}
\end{setup}

\begin{thm} \label{thm:3150}
Under Setup \tup{\ref{set:3150}}, 
assume that $\cat{K}$ has enough K-projectives. 
Then the triangulated right derived functor 
$(\mrm{R} F, \eta^{\mrm{R}}) : \cat{D}^{\mrm{op}} \to \cat{E}$
of $F$ exists. Furthermore, for every K-projective DG module
$P \in \cat{K}$, the morphism 
$\eta^{\mrm{R}}_P : F(P) \to \mrm{R} F(P)$ 
in $\cat{E}$ is an isomorphism. 
\end{thm}

\begin{proof}
Let 
$\cat{K}^{\mrm{flip}} := \opn{Flip}(\cat{K}^{\mrm{op}}) \sub
\dcat{K}(A^{\mrm{op}}, \cat{M}^{\mrm{op}})$.
This is a triangulated category. Its localization w.r.t.\ the 
quasi-isomorphisms is denoted by $\cat{D}^{\mrm{flip}}$. We have isomorphisms of 
triangulated categories
$\opn{Flip} : \cat{K}^{\mrm{op}} \iso \cat{K}^{\mrm{flip}}$
and 
$\opn{Flip} : \cat{D}^{\mrm{op}} \iso \cat{D}^{\mrm{flip}}$.
Let 
$F^{\mrm{flip}} := F \circ \opn{Flip}^{-1} : \cat{K}^{\mrm{flip}} \to \cat{E}$.
This is a triangulated functor. 

Since $\opn{Flip}$ sends quasi-isomorphisms to 
quasi-isomorphisms, and by Lemma \ref{lem:3151}, each K-projective 
resolution
$\rho : P \to M$ of $M$ in $\cat{K}$ goes to a K-injective resolution 
$\opn{Flip}(\rho) : \opn{Flip}(M) \to \opn{Flip}(P)$
of $N := \opn{Flip}(M)$ in $\cat{K}^{\mrm{flip}}$. 
We are given that $\cat{K}$ has enough K-projectives. Therefore 
$\cat{K}^{\mrm{flip}}$ has enough K-injectives. 

Now we can apply Theorem \ref{thm:5050} to deduce that the right 
derived functor
$(\mrm{R} F^{\mrm{flip}}, \eta^{\mrm{flip}}) : \cat{D}^{\mrm{flip}} 
\to \cat{E}$
of $F^{\mrm{flip}}$ exists. Then 
\[ (\mrm{R} F, \eta^{\mrm{R}}) := 
(\mrm{R} F^{\mrm{flip}}, \eta^{\mrm{flip}}) \circ \opn{Flip} : 
\cat{D}^{\mrm{op}} \to \cat{E} \]
is a right derived functor of $F$.

The assertion about the morphism $\eta^{\mrm{R}}_P$ is clear from the 
discussion 
above. 
\end{proof}

\begin{thm} \label{thm:3151}
Under Setup \tup{\ref{set:3150}}, 
assume that $\cat{K}$ has enough K-injectives. 
Then the left derived functor 
$(\mrm{L} F, \eta^{\mrm{L}}) : \cat{D}^{\mrm{op}} \to \cat{E}$
of $F$ exists. Furthermore, for every K-injective DG module
$I \in \cat{K}$, the morphism 
$\eta^{\mrm{L}}_I : \mrm{L} F(I) \to F(I)$ 
in $\cat{E}$ is an isomorphism. 
\end{thm}

\begin{proof}
The same as that of the previous theorem, except that now we use Lemma
\ref{lem:3151} and Theorem \ref{thm:5051}. 
\end{proof}








\cleardoublepage
\mysection{Existence of Resolutions} \label{sec:exist-resol}     

\AYcopyright
 
In this section we continue in the more concrete setting: 
$\K$ is a nonzero commutative base ring, $A$ is a central DG $\K$-ring, 
and $\cat{M}$ is a $\K$-linear abelian category; see Convention \ref{conv:2490}.
The DG category $\dcat{C}^{\star}(A, \cat{M})$
was introduced in Section \ref{subsec:DGModinM}. 
We will prove existence of  K-projective and K-injective resolutions in 
$\dcat{C}^{\star}(A, \cat{M})$,
under suitable conditions.  


\mysubsection{Direct and Inverse Limits of Complexes}
\label{subsec:lims-cplx}
We shall have to work with limits in this section. Limits in abstract abelian 
and DG categories (not to mention triangulated categories) are a very 
delicate issue. We will try to be as concrete as possible, in order to avoid 
pitfalls and confusion. The notation we use for direct and inverse limits
of systems indexed by $\N$ was presented in Subsection \ref{subsec:inv-dir-lim}.
 
\begin{prop} \label{prop:1660} \mbox{}
\begin{enumerate}
\item Let $\{ M_k \}_{k \in \N}$ be a direct system in 
$\dcat{C}_{\mrm{str}}(A, \cat{M})$. Assume that for every $i$ the direct limit 
$\lim_{k \to} M^i_k$ exists in $\cat{M}$. Then 
the direct limit 
$M = \lim_{k \to} M_k$ exists in $\dcat{C}_{\mrm{str}}(A, \cat{M})$, and in 
degree $i$ it is $M^i = \lim_{k \to} M^i_k$.

\item Let $\{ M_k \}_{k \in \N}$ be an inverse system in 
$\dcat{C}_{\mrm{str}}(A, \cat{M})$. Assume that for every $i$ the inverse limit 
$\lim_{\lar k} M^i_k$ exists in $\cat{M}$. Then 
the inverse limit 
$M = \lim_{\lar k} M_k$ exists in $\dcat{C}_{\mrm{str}}(A, \cat{M})$, and in 
degree $i$ it is $M^i = \lim_{\lar k} M^i_k$.
\end{enumerate}
\end{prop}

\begin{proof}
We will only prove item (1); the proof of item (2) is similar. 
For each integer $i$ define 
$M^i := \lim_{k \to} M^i_k \in \cat{M}$. 
By the universal property of the direct limit, the differentials 
$\d : M^i_k \to M^{i + 1}_k$ induce differentials 
$\d : M^i \to M^{i + 1}$, and in this way we obtain a complex 
$M := \{ M^i \}_{i \in \Z} \in \dcat{C}_{}(\cat{M})$.
Similarly, every element $a \in A^j$ induces morphisms 
$a : M^i \to M^{i + j}$ in $\cat{M}$, and thus $M$ becomes an object of 
$\dcat{C}(A, \cat{M})$. There are morphisms $M_k \to M$ in 
$\dcat{C}_{\mrm{str}}(A, \cat{M})$, and it is easy to see that these make $M$ 
into a direct limit of the system $\{ M_k \}_{k \in \N}$.
\end{proof}

Since limits exist in $\cat{M} = \cat{Mod} \K$, the proposition above says 
that they exist in $\bcat{C}_{\mrm{str}}(A)$. Likewise they exist in the 
category $\dcat{G}_{\mrm{str}}(\K)$ of graded $\K$-modules.

We say that a direct system $\{ M_k \}_{k \in \N}$ in $\cat{M}$ is {\em 
eventually stationary} if $\mu_k : M_k \to M_{k + 1}$ are isomorphisms for 
large $k$. Similarly we can talk about an eventually stationary inverse system. 
The limit of an eventually stationary system (direct or inverse) always exists: 
it is $M_k$ for large enough $k$. 

\begin{prop} \label{prop:1575} \mbox{}
\begin{enumerate}
\item Let $\{ M_k \}_{k \in \N}$ be a direct system in 
$\bcat{C}_{\mrm{str}}(A, \cat{M})$.  
Assume that for each $i$ the direct system $\{ M^i_k \}_{k \in \N}$ in 
$\cat{M}$ is eventually stationary. Then the direct limit 
$M = \lim_{k \to} M_k$ exists in $\dcat{C}_{\mrm{str}}(A, \cat{M})$, 
the  direct limit 
$\lim_{k \to} \opn{H}(M_k)$ exists in $\dcat{G}_{\mrm{str}}(\cat{M})$, 
and the canonical morphism 
$\lim_{k \to} \, \opn{H}(M_k) \to \opn{H}(M)$ 
in $\dcat{G}_{\mrm{str}}(\cat{M})$ is an isomorphism. 

\item Let $\{ M_k \}_{k \in \N}$ be an inverse system in 
$\bcat{C}_{\mrm{str}}(A, \cat{M})$. Assume that for each $i$ the inverse system 
$\{ M^i_k \}_{k \in \N}$ in $\cat{M}$ is eventually stationary. Then the 
inverse limit $M = \lim_{\leftarrow k} M_k$ exists in 
$\dcat{C}_{\mrm{str}}(A, \cat{M})$, the  inverse limit 
$\lim_{\leftarrow k} \opn{H}(M_k)$ exists in $\dcat{G}_{\mrm{str}}(\cat{M})$, 
and the canonical morphism 
$\opn{H}(M) \to \lim_{\leftarrow k} \, \opn{H}(M_k)$
in $\dcat{G}_{\mrm{str}}(\cat{M})$ is an isomorphism. 
\end{enumerate}
\end{prop}

\begin{proof}
Again, we only prove item (1). 
As mentioned above, for each $i$ the limit
$M^i = \lim_{k \to} M^i_k$ exists in $\cat{M}$.
By Proposition \ref{prop:1660} the limit 
$M = \lim_{k \to} M_k$ exists in $\dcat{C}_{\mrm{str}}(A, \cat{M})$.

Regarding the cohomology: fix an integer $i$. 
Take $k$ large enough such that 
$M^{i'}_k \to M^{i'}_{k'}$ are isomorphisms for all $k \leq k'$ and 
$i - 1 \leq  i' \leq i + 1$. Then 
$M^{i'}_{k'} \to M^{i'}$ 
are isomorphisms in this range, and therefore 
$\opn{H}^i(M_{k'}) \to \opn{H}^i(M)$ are isomorphisms for all $k \leq k'$. 
We see that the direct system 
$\{ \opn{H}^i(M_k) \}_{k \in \N}$
is eventually stationary, and its direct limit is 
$\opn{H}^i(M)$.
\end{proof}

When we drop the abstract abelian category $\cat{M}$, i.e.\ when we work 
with $\cat{M} = \cat{Mod} \K =  \dcat{M}(\K)$ and 
$\dcat{C}_{\mrm{str}}(A, \cat{M}) = \dcat{C}_{\mrm{str}}(A)$, 
there is no problem of 
existence of limits. The next proposition says that furthermore ``direct limits 
are exact'' in $\bcat{C}_{\mrm{str}}(A)$. 

\begin{prop} \label{prop:1585}
Let $\{ M_k \}_{k \in \N}$ be a direct system in $\bcat{C}_{\mrm{str}}(A)$,
with limit $M = \lim_{k \to} M_k$. 
Then the canonical homomorphism 
$\lim_{k \to} \opn{H}(M_k) \to \opn{H}(M)$ 
in $\dcat{G}_{\mrm{str}}(\K)$ is bijective. 
\end{prop}

\begin{exer}  \label{exer:1741}
Prove Proposition \ref{prop:1585}. (Hint: forget the action of $A$, and work 
with complexes of $\K$-modules.)
\end{exer}

Exactness of inverse limits tends to be much more complicated than that of 
direct limits, even for $\K$-modules. We always have to impose some condition 
on the inverse system to have exactness in the limit. 

Let 
$\bigl( \{ M_k \}_{k \in \N}, \{ \mu_k \}_{k \in \N} \bigr)$
be an inverse system in $\bcat{M}(\K)$.
For every $k \leq l$ let 
$M_{k, l} \sub M_k$ be the image of the homomorphism
$\mu_{k, l} : M_l \to M_k$.
Note that there are inclusions 
$M_{k, l + 1} \sub M_{k, l}$, so for fixed $k$ we have an inverse system 
$\{ M_{k, l} \}_{l \geq k}$
of submodules of $M_k$. 

\begin{dfn} \label{dfn:1590}
We say that an inverse system $\{ M_k \}_{k \in \N}$ in $\bcat{M}(\K)$
has the {\em Mittag-Leffler property}% 
\index{Mittag-Leffler property}
if for every index $k$, the 
inverse system $\{ M_{k, l} \}_{l \geq k}$ 
is eventually stationary.
\end{dfn}

\begin{exa} \label{exa:1590}
If the inverse system 
$\bigl( \{ M_k \}_{k \in \N}, \{ \mu_k \}_{k \in \N} \bigr)$
in $\bcat{M}(\K)$ satisfies at least one of the following conditions, then it 
has the  Mittag-Leffler property:
\begin{itemize}
\rmitem{a} The system has surjective transitions. 

\rmitem{b} The system is eventually stationary. 

\rmitem{c} For every $k \in \N$ there exists some $l \geq k$ such that
$M_{k, l} = 0$. This is called the {\em trivial Mittag-Leffler property}, and 
one says that the system is {\em pro-zero}. 
\end{itemize}
\end{exa}

\begin{thm}[Mittag-Leffler Argument] \label{thm:1535}
\index{Mittag-Leffler property}
Let $\{ M_k \}_{k \in \N}$ be an inverse system in $\dcat{C}_{\mrm{str}}(A)$, 
with inverse limit 
$M = \lim_{\leftarrow k} M_k$. Assume the system satisfies these two 
conditions\tup{:}
\begin{enumerate}
\rmitem{a} For every $i \in \Z$ the inverse system
$\{ M^i_k \}_{k \in \N}$ in $\dcat{M}(\K)$
has the Mittag-Leffler property.

\rmitem{b} For every $i \in \Z$ the inverse system
$\{ \opn{H}^i(M_k) \}_{k \in \N}$ in $\dcat{M}(\K)$
has the Mittag-Leffler property.
\end{enumerate}
Then the canonical homomorphisms 
$\opn{H}^i(M) \to \lim_{\leftarrow k} \opn{H}^i(M_k)$ 
are bijective for all $i$.
\end{thm}

\begin{proof}
We can forget all about the graded $A$-module structure, and just view this as 
an inverse system in $\dcat{C}_{\mrm{str}}(\Z)$, i.e.\ an inverse system of 
complexes of abelian groups. Now this is a special case of  
\cite[Proposition 1.12.4]{KaSc1} or 
\cite[Chapitre $0_{\mrm{III}}$, Proposition 13.2.3]{{EGA-III}}.
\end{proof}

The most useful instance of the ML argument is this:

\begin{cor} \label{cor:1590}
Let $\{ M_k \}_{k \in \N}$ be an inverse system in $\dcat{C}_{\mrm{str}}(A)$, 
with inverse limit 
$M = \lim_{\leftarrow k} M_k$. Assume the system satisfies these two 
conditions\tup{:}
\begin{enumerate}
\rmitem{a} For every $i \in \Z$ the inverse system
$\{ M^i_k \}_{k \in \N}$ has surjective transitions. 

\rmitem{b} For every $k$ the DG module $M_k$ is acyclic.  
\end{enumerate}
Then $M$ is acyclic. 
\end{cor}

\begin{proof}
Conditions (a) and (b) here imply conditions (a) and (b) of Theorem 
\ref{thm:1535}, respectively. 
\end{proof}

\begin{exer} \label{exer:1904}
Prove Corollary \ref{cor:1590} directly, without resorting to Theorem 
\ref{thm:1535}. 
\end{exer}

By {\em generalized integers}% 
\index{Generalized integer}
we mean elements of the ordered set 
$\Z \cup \{ \pm \infty \}$.  Recall that for a subset $S \subseteq \Z$, its 
infimum is $\opn{inf}(S) \in \Z \cup \{ \pm \infty \}$,
where $\opn{inf}(S) = + \infty$ iff $S = \varnothing$. 
Likewise the supremum is $\opn{sup}(S) \in \Z \cup \{ \pm \infty \}$,
where $\opn{sup}(S) = - \infty$ iff $S = \varnothing$.
For $i \in \Z \cup \{ \pm \infty \}$ 
the expression $-i$ has a value in $\Z \cup \{ \pm \infty \}$.
If $i_0, i_1 \in \Z \cup \{ \pm \infty \}$ satisfy 
$i_0 \leq i_1$, then $-i_1 \leq -i_0$. 
For $i_0, i_1 \in \Z \cup \{ \infty \}$, the expressions 
$i_0 + i_1$ and $-i_0 - i_1$ have values in 
$\Z \cup \{ \pm \infty \}$. 
A generalized integer $i$ is called {\em finite} if $i < \infty$,
i.e.\ if $i \in  \Z \cup \{ -\infty \}$. 

\begin{dfn} \label{dfn:4765}
Given $i_0 \in \Z \cup \{ -\infty \}$ and 
$i_1 \in \Z \cup \{ \infty \}$, with $i_0 \leq i_1$,
the {\em integer interval}%
\index{Integer interval}
with these endpoints is the set 
\[ [i_0, i_1] := \{ i \in \Z  \mid i_0 \leq i \leq i_1 \} \sub \Z . \]
There is also the empty integer interval $\varnothing$.
When the elements of an integer interval represent degrees, we 
sometimes call it a {\em degree interval}%
\index{Degree interval}.
\end{dfn}

We exclude the possibilities $i_0 = i_1 = \infty$ and $i_0 = i_1 = -\infty$
for endpoints of intervals.
In calculus notation we have 
$[i_0, \infty] = [i_0, \infty) \cap \Z$, etc.; 
but we avoid the open interval notation on purpose. 
Whenever we refer to an integer interval $[i_0, i_1]$, it is always 
assumed that the generalized integers $i_0$ and $i_1$ satisfy the conditions in 
Definition \ref{dfn:4765}.
Further properties of integer intervals will be studied in Subsection 
\ref{subsec:bounded-revis}. 

\begin{dfn} \label{dfn:3155}
Let $M = \{ M^i \}_{i \in \Z}$ be a graded object in $\cat{M}$,
and let 
$S := \{ i \in \Z \mid  M^i \neq 0 \} \sub \Z$. 
\begin{enumerate}
\item The {\em supremum} of $M$ is
$\sup(M) := \opn{sup}(S) \in \Z \cup \{ \pm \infty \}$. 

\item The {\em infimum} of $M$ is 
$\inf(M) := \opn{inf}(S) \in \Z \cup \{ \pm \infty \}$. 

\item The {\em amplitude} of $M$ is 
$\opn{amp}(M) := \sup (M) - \inf (M) \in \N \cup \{ \pm \infty \}$.
\end{enumerate}
\end{dfn}

Note that for $M = 0$ this reads $S = \varnothing$, 
$\inf (M) = \infty$, $\sup (M) = -\infty$ 
and $\opn{amp} (M) = -\infty$.
Thus $M$ is bounded (resp.\ bounded above, resp.\ bounded below) iff 
$\opn{amp} (M) < \infty$ (resp.\ $\opn{sup} (M) < \infty$, resp.\
$\opn{inf} (M) > -\infty$).



\mysubsection{Totalizations}
Consider a complex 
\begin{equation} \label{eqn:3300}
(M, \pa) = \bigl( \cdots \to M^{-1} \xar{\pa^{-1}}
M^{0} \xar{\pa^{0}} M^{1} \to \cdots \bigr) 
\end{equation}
with entries in the abelian category $\dcat{C}_{\mrm{str}}(A, \cat{M})$. 
This means that for each $i \in \Z$ there is a DG module 
$(M^{i}, \d_{M^i}) \in \dcat{C}(A, \cat{M})$, 
and a morphism 
$\pa^i : M^i \to M^{i + 1}$
in $\dcat{C}_{\mrm{str}}(A, \cat{M})$, 
such that $\pa^{i + 1} \circ \pa^i = 0$.
Recall that the strictness of $\pa^i$ says that it has degree $0$, and that
\begin{equation} \label{eqn:3365}
\pa^i \circ \d_{M^i}  = \d_{M^{i + 1}} \circ \pa^i .
\end{equation}
In terms of our symbolic notation, $(M, \pa)$ is an object of the DG category 
\lb $\dcat{C}(\dcat{C}_{\mrm{str}}(A, \cat{M}))$. 

Each DG module $(M^i, \d_{M^i})$ has its own internal structure:
\begin{equation} \label{eqn:3305}
(M^i, \d_{M^i}) = \bigl( \{ M^{i, j} \}_{j \in \Z}, 
\{ \d_{M^{i}}^{\lms j} \}_{j \in \Z} \bigr) . 
\end{equation}
Here $M^{i, j} \in \cat{M}$, 
$\d_{M^{i}}^{\lms j} : M^{i, j} \to M^{i, j + 1}$ 
is a morphism in $\cat{M}$, and the relation 
$\d_{M^{i}}^{\lms j + 1} \circ \d_{M^{i}}^{\lms j} = 0$ holds.
There is an action of the DG ring $A$ on $M^i$, prescribed by a DG $\K$-ring 
homomorphism 
$f_{M^i} : A \to \opn{End}_{\cat{M}}(M^i)$.

Sometimes we view $M$ as a double complex, and then we refer to its rows and 
columns: $M^i$ is the $i$-th column of $M$, and
$\{ M^{i, j} \}_{i \in \Z}$ is its $j$-th row.
See diagram (\ref{eqn:5057}). 
\begin{equation} \label{eqn:5057}
\UseTips \xymatrix @C=3ex @R=3ex {
&
\vdots
& &
\vdots
& &
\vdots
\\
\cdots
\ar[r]
& 
M^{-1, 1}
\ar[rr]^{\pa^{-1}}
\ar[u]
& &
M^{0, 1}
\ar[rr]^{\pa^{0}}
\ar[u]
& &
M^{1, 1}
\ar[r]
\ar[u]
& 
\cdots
\\ \\
\cdots
\ar[r]
& 
M^{-1, 0}
\ar[rr]^{\pa^{-1}}
\ar[uu]^{\d_{M^{-1}}^{0}}
& &
M^{0, 0}
\ar[rr]^{\pa^{0}}
\ar[uu]^{\d_{M^{0}}^{0}}
& &
M^{1, 0}
\ar[r]
\ar[uu]^{\d_{M^{1}}^{0}}
& 
\cdots
\\
&
*[u]{\vdots}
\ar[u]
& &
*[u]{\vdots}
\ar[u]
& &
*[u]{\vdots}
\ar[u]
} 
\end{equation}

As explained in Proposition \ref{prop:1660},
if the abelian category $\cat{M}$ has 
countable direct sums, then so does $\dcat{C}_{\mrm{str}}(A, \cat{M})$. 
Recall that $A^{\natural}$ is the underlying graded ring of the DG ring $A$. 

\begin{dfn} \label{dfn:3300}
Assume that $\cat{M}$ has countable direct sums. Given a complex 
$(M, \pa)$ with entries in $\dcat{C}_{\mrm{str}}(A, \cat{M})$, with notation as 
in \tup{(\ref{eqn:3300})}, its {\em direct sum totalization}% 
\index{Totalization! direct sum}
is the object 
\[ \opn{Tot}^{\oplus}(M, \pa)  = 
\bigl( \opn{Tot}^{\oplus}(M), \d_{\mrm{Tot}} \bigr) \in  
\dcat{C}(A, \cat{M}) \]
defined as follows, in four stages. 
\begin{enumerate}
\item There is a graded object 
\[ \opn{Tot}^{\oplus}(M) := \bigoplus_{i \in \Z} \, 
\opn{T}^{-i}(M^i) \in \dcat{G}(A^{\natural}, \cat{M}) . \]

\item On each summand $\opn{T}^{-i}(M^i)$ of $\opn{Tot}^{\oplus}(M)$
there is a differential 
$\d_{\opn{T}^{-i}(M^i)}$, 
and we define the degree $1$ operator 
$\d_M := \bigoplus_{i \in \Z} \d_{\opn{T}^{-i}(M^i)}$ 
on $\opn{Tot}^{\oplus}(M)$.

\item For each $i$ let 
\[ \opn{tot}(\pa)^i := \opn{t}^{-1}_{\lms \opn{T}^{-(i + 1)}(M^{i + 1})} \circ 
\opn{T}^{-i}(\pa^i) : 
\opn{T}^{-i}(M^i) \to \opn{T}^{-(i + 1)}(M^{i + 1}) . \]
We define the degree $1$ operator 
$\opn{tot}(\pa) :=  \bigoplus_{i \in \Z} \, \opn{tot}(\pa)^i$
on the graded object $\opn{Tot}^{\oplus}(M)$.

\item The differential $\d_{\mrm{Tot}}$ on the graded object 
$\opn{Tot}^{\oplus}(M)$ is 
$\d_{\mrm{Tot}} := \d_M + \opn{tot}(\pa)$.
\end{enumerate}
\end{dfn}

For the definition to be valid we need to verify something: 

\begin{lem} \label{lem:3030}
The degree $1$ operator $\d_{\mrm{Tot}}$ makes the pair 
$\bigl( \opn{Tot}^{\oplus}(M), \d_{\mrm{Tot}} \bigr)$
into an object of $\dcat{C}(A, \cat{M})$. 
\end{lem}

\begin{proof}
We have to check that 
$\d_{\mrm{Tot}} \circ \d_{\mrm{Tot}} = 0$
and that $\opn{tot}(\pa)$ is $A$-linear. 
Below we shall verify the first condition; the second verification is similar, 
and is left to the reader. 

For a given $i \in \Z$ consider the composed morphism starting at the summand
$\opn{T}^{-i}(M^i)$ of the total DG module:
\[ \begin{aligned}
&
\opn{T}^{-i}(M^i) \xar{\d_{\mrm{Tot}}}
\opn{T}^{-i}(M^i) \oplus \opn{T}^{-(i + 1)}(M^{i + 1}) 
\\
& \quad 
\xar{\d_{\mrm{Tot}}}
\opn{T}^{-i}(M^i) \oplus \opn{T}^{-(i + 1)}(M^{i + 1})
\oplus \opn{T}^{-(i + 2)}(M^{i + 2}) . 
\end{aligned} \]
Let us view these direct sums as columns, with matrices of operators acting on 
them from the left. So we have to calculate the matrix product 
\[ 
\bmat{ \d_{\opn{T}^{-i}(M^i)} & 0
\\[0.3em]
\opn{t}^{-1} \circ \opn{T}^{-i}(\pa^i) & \d_{\opn{T}^{-(i + 1)}(M^{i + 1})}
\\[0.3em]
0 & \opn{t}^{-1} \circ \opn{T}^{-(i + 1)}(\pa^{i + 1}) }
\circ
\bmat{ \d_{\opn{T}^{-i}(M^i)} 
\\[0.3em]
\opn{t}^{-1} \circ \opn{T}^{-i}(\pa^i) } 
\]
The resulting column of operators has  
$\d_{\opn{T}^{-i}(M^i)} \circ \d_{\opn{T}^{-i}(M^i)} = 0$
in the first position. In the second position we have
\begin{equation} \label{eqn:3366}
\opn{t}^{-1} \circ \opn{T}^{-i}(\pa^i) \circ \d_{\opn{T}^{-i}(M^i)} +
\d_{\opn{T}^{-(i + 1)}(M^{i + 1})} \circ \opn{t}^{-1} 
\circ \opn{T}^{-i}(\pa^i) .
\end{equation}
Using Proposition \ref{prop:4155} we have 
\[ \d_{\opn{T}^{-(i + 1)}(M^{i + 1})} \circ  \opn{t}^{-1} = 
- \opn{t}^{-1} \circ \, \d_{\opn{T}^{-i}(M^{i + 1})} = 
- \opn{t}^{-1} \circ \opn{T}^{-i}(\d_{M^{i + 1}}) . \]
Therefore (\ref{eqn:3366}) equals 
\[ \opn{t}^{-1} \circ \opn{T}^{-i} \bigl( \pa^i \circ \d_{M^i} -
\d_{M^{i + 1}} \circ \pa^i \bigr) , \]
which is zero by (\ref{eqn:3365}). 
In the bottom of the column we have
\begin{equation} \label{eqn:3367}
\opn{t}^{-1} \circ \opn{T}^{-(i + 1)}(\pa^{i + 1}) \circ 
\opn{t}^{-1} \circ \opn{T}^{-i}(\pa^{i}) . 
\end{equation}
According to Definition \ref{dfn:1171} we know that 
\[ \opn{T}^{-(i + 1)}(\pa^{i + 1}) \circ \opn{t}^{-1} = 
\opn{t}^{-1} \circ \opn{T}(\opn{T}^{-(i + 1)}(\pa^{i + 1})) = 
\opn{t}^{-1} \circ \opn{T}^{-i}(\pa^{i + 1})) . \]
Therefore (\ref{eqn:3367}) equals 
$\opn{t}^{-1} \circ \opn{T}^{-i} \bigl( \pa^{i + 1} \circ \pa^{i} \bigr) = 0$. 
\end{proof}

\begin{exer} \label{exer:4750}
Finish the proof of the lemma (the $A$-linearity of $\opn{tot}(\pa)$). 
\end{exer}

\begin{exa} \label{exa:3300}
If $M^i = 0$ for $i \neq -1, 0$, then 
\[ \opn{Tot}^{\oplus}(M, \pa) = 
\opn{Cone}(\pa^{-1} : M^{-1} \to M^0) , \]
the standard cone from Definition \ref{dfn:1172}.
In Proposition \ref{prop:3305} there is  a more general statement. 
\end{exa}

In Definition \ref{dfn:2320} we talked about smart truncations of DG 
modules. Here is another sort of truncations. 

\begin{dfn} \label{dfn:3066}
Let $\cat{N}$ be an abelian category. 
For a complex $N \in \dcat{C}(\cat{N})$ its 
{\em stupid truncations}%
\index{Truncation of a complex! stupid}%
\index{1-Sttleq@$\opn{stt}^{\leq q}(N)$}
at an integer $q$ are 
\[ \opn{stt}^{\leq q}(N) := 
\bigl( \cdots \to N^{q - 1} \to N^{q} \to 0 \to 0 \to \cdots \bigr) \] 
and% 
\index{1-Sttgeq@$\opn{stt}^{\geq q}(N)$}
\[ \opn{stt}^{\geq q}(N) := 
\bigl( \cdots \to 0 \to 0 \to N^{q} \to N^{q + 1} \to \cdots \bigr)  . \]
\end{dfn}

These truncations fit into a short exact sequence 
\begin{equation} \label{eqn:2147}
0 \to \opn{stt}^{\geq q}(N) \to N \to \opn{stt}^{\leq q - 1}(N) \to  0
\end{equation}
in $\dcat{C}_{\mrm{str}}(\cat{N})$. 

Warning: The stupid truncations do not apply to DG modules 
$N \in \dcat{C}(B, \cat{N})$, unless the DG ring $B$ is a ring.   
However, in what follows we will use stupid truncations in 
$\dcat{C}(\cat{N})$, where 
$\cat{N}$ is the abelian category $\dcat{C}_{\mrm{str}}(A, \cat{M})$. 
This could be a bit confusing. 

\begin{prop} \label{prop:3305}
Let 
\[ (M, \pa) =  \bigl( \cdots \to M^{-1} \xar{\pa^{-1}}
M^{0} \xar{\pa^{0}} M^{1} \to 0 \to \cdots \bigr) \]
be a complex with entries in the abelian category 
$\dcat{C}_{\mrm{str}}(A, \cat{M})$, let 
\[ (M', \pa) :=  \bigl( \cdots \to M^{-1} \xar{\pa^{-1}}
M^{0} \to 0 \to \cdots \bigr) , \]
namely the stupid truncation of $M$ below $0$, and let 
$\rho : \opn{Tot}^{\oplus}(M', \pa) \to M^1$
be the morphism in $\dcat{C}_{\mrm{str}}(A, \cat{M})$ induced by
$\pa^{0} : M^{0} \to M^1$. 
Then there is an isomorphism 
\[ \tag{\dag} \opn{T} \bigl( \opn{Tot}^{\oplus}(M, \pa) \bigr) \cong 
\opn{Cone} \bigl( - \rho : \opn{Tot}^{\oplus}(M', \pa) \to M^1 \bigr) \]
in $\dcat{C}_{\mrm{str}}(A, \cat{M})$. 
\end{prop}

\begin{proof}
By comparing translations we see that the objects on both sides of (\dag) match 
in each degree. Likewise for all the differentials that do not involve 
$\pa^0$. A calculation similar to that in the proof of Lemma 
\ref{lem:3030} shows that there is a sign change there, and this is why we need 
$-\rho$ on the right side of (\dag). 
\end{proof}

\begin{cor} \label{cor:3310}
If, in the situation of Proposition \tup{\ref{prop:3305}}, the DG module  
$\opn{Tot}^{\oplus}(M, \pa)$ is acyclic, then 
$\rho : \opn{Tot}^{\oplus}(M', \pa) \to M^1$ 
is a quasi-isomorphism. 
\end{cor}

\begin{proof}
If $\opn{Tot}^{\oplus}(M, \pa)$ is acyclic then so is 
$\opn{T} \bigl( \opn{Tot}^{\oplus}(M, \pa) \bigr)$. 
The proposition implies that $-\rho$ is a quasi-isomorphism. Therefore $\rho$ 
is a quasi-isomorphism. 
\end{proof}

\begin{dfn} \label{dfn:3335}
Let $M \in \dcat{C}(A, \cat{M})$.
\begin{enumerate}
\item A {\em filtration}% 
\index{Filtration on a DG module}
on $M$ is a direct system
$\{ F_j(M) \}_{j \geq -1}$ in $\dcat{C}_{\mrm{str}}(A, \cat{M})$
consisting of subobjects $F_j(M) \sub M$, indexed by the interval 
$[-1, \infty] \sub \Z$. 

\item We say that $M = \lim_{j \to} F_j(M)$
if the direct limit exists, and the 
canonical morphism $\lim_{j \to} F_j(M) \to M$ in 
$\dcat{C}_{\mrm{str}}(A, \cat{M})$ is an isomorphism. 

\item The filtration $F$ induces, for each $j \geq 0$, the subquotient 
\[ \opn{Gr}^F_j(M) := F_{j}(M) / F_{j - 1}(M) \in \dcat{C}(A, \cat{M}) . \] 
\end{enumerate}
\end{dfn}

Our filtrations will always be ascending. But sometimes they will be indexed 
by a subinterval $[j_0, j_1]$ of $[-1, \infty]$; and in this case 
$\opn{Gr}^F_j(M)$ will be defined only for 
$j \in [j_0 + 1, j_1]$.
 
In the next proposition we take $\cat{M} = \dcat{M}(\K)$; but really 
all we need is that $\cat{M}$ has exact countable direct limits. 

Recall that for a DG $A$-module $M$ we denote by $\opn{Z}(M)$
the object of cocycles of $M$. It is an object of 
$\dcat{G}(\K)$.

\begin{prop} \label{prop:3300}
Let $(M, \pa)$ be a complex with entries in the abelian category 
$\dcat{C}_{\mrm{str}}(A)$, 
with internal structure \tup{(\ref{eqn:3305})}. Assume that the two conditions 
below hold\tup{:}
\begin{itemize}
\rmitem{i} For every $j \in \Z$ the complex  
\[ \bigl( \cdots \to M^{-1, j} \xar{\pa^{-1}}
M^{0, j} \xar{\pa^{0}} M^{1, j} \to \cdots \bigr) \]
with entries in $\dcat{M}(\K)$ is acyclic. 

\rmitem{ii} There is some $j_1 \in \Z$ such that for all 
$j \geq j_1$ the complex 
\[ \bigl( \cdots \to \opn{Z}^j(M^{-1}) \xar{\pa^{-1}}
\opn{Z}^j(M^{0}) \xar{\pa^{0}} \opn{Z}^j(M^{1}) \to \cdots \bigr) \]
with entries in $\dcat{M}(\K)$ is acyclic. 
\end{itemize}
Then the DG $A$-module $\opn{Tot}^{\oplus}(M, \pa)$ is acyclic. 
\end{prop}

\begin{proof} The proof is done in two steps. 

\smallskip \noindent
Step 1. Let us replace condition (ii) by the stronger condition: 
\begin{itemize}
\rmitem{ii$'$} There is some $j'_1 \in \Z$ such that 
$M^{i, j} = 0$ for all $j > j'_1$. 
\end{itemize}
For every $i$ we introduce a filtration 
$\{ F_q(M^i) \}_{q \geq -1}$
on the DG $\K$-module $M^i$ as follows: 
\[ F_q (M^i) := \opn{stt}^{\geq j'_1 - q}(M^i)
= \bigl( \cdots \to 0 \to M^{i, j_1' - q} \xar{\d_{M^i}} M^{i, j'_1 - q + 1} 
\to \cdots \bigr) , \]
the stupid truncation above $j'_1 - q$. 
These filtrations induces a filtration on the DG $\K$-module 
$\opn{Tot}^{\oplus}(M, \pa)$~: 
\[ F_q \bigl( \opn{Tot}^{\oplus}(M) \bigr) := 
\opn{Tot}^{\oplus} 
\bigl( \cdots \to F_q(M^{-1}) \xar{\pa^{-1}}
F_q(M^{0}) \xar{\pa^{0}} F_q(M^{1}) \to \cdots \bigr) . \]
We note that $F_{-1} \bigl( \opn{Tot}^{\oplus}(M) \bigr) = 0$, and 
\begin{equation} \label{eqn:3032}
\bigcup_{q \geq -1} \, F_q \bigl( \opn{Tot}^{\oplus}(M) \bigr) = 
\opn{Tot}^{\oplus}(M) . 
\end{equation}

For every $q \geq 0$ the DG $\K$-module 
$\opn{Gr}^F_q \bigl( \opn{Tot}^{\oplus}(M) \bigr)$
is isomorphic in $\dcat{C}_{\mrm{str}}(\K)$, up to signs of differentials, to 
the complex in condition (i) 
with index $j = j'_1 - q$, so it is acyclic. For $q = -1$ 
we trivially have an acyclic DG $\K$-module 
$F_{-1} \bigl( \opn{Tot}^{\oplus}(M) \bigr)$. 
For every $q \geq 0$ there is an exact sequence 
\[ 0 \to F_{q - 1} \bigl( \opn{Tot}^{\oplus}(M) \bigr) \to 
F_q \bigl( \opn{Tot}^{\oplus}(M) \bigr) \to 
\opn{Gr}^F_q \bigl( \opn{Tot}^{\oplus}(M) \bigr) \to 0 \]
in $\dcat{C}_{\mrm{str}}(\K)$. By induction on $q$ we conclude that 
$F_q \bigl( \opn{Tot}^{\oplus}(M) \bigr)$
is acyclic. Finally, by $(\ref{eqn:3032})$, and the fact that cohomology 
commutes with direct limits in $\dcat{C}_{\mrm{str}}(\K)$, we see that the DG 
module $\opn{Tot}^{\oplus}(M, \pa)$ is acyclic. 

\medskip \noindent 
Step 2. Now we assume conditions (i) and (ii). 
For each $i$ we introduce a filtration 
$\{ G_q(M^i) \}_{q \geq 0}$
on the DG $\K$-module $M^i$ as follows: 
\[ G_q(M^i) := \opn{smt}^{\leq j_1 + q}(M^i) = 
\bigl( \cdots \to M^{i, j_1 + q - 1} \xar{\d_{M^i}} 
\opn{Z}^{j_1 + q}(M^i) \to 0 \to  \cdots \bigr) , \]
the smart truncation below $j_1 + q$. 

For every value of $q$ the complex 
\[ G_q(M) := \bigl( \cdots \to G_q(M^{-1}) \xar{\pa^{-1}}
G_q(M^{0}) \xar{\pa^{0}} G_q(M^{1}) \to \cdots \bigr) \]
satisfies conditions (i) and (ii$'$), with 
$j'_1 := j_1 + q$. Indeed, for $j < j_1 + q$ the $j$-th row of $G_q(M)$ is 
\[ \bigl( \cdots \to M^{-1, j} \xar{\pa^{-1}}
M^{0, j} \xar{\pa^{0}} M^{1, j} \to \cdots \bigr) , \]
and it is exact by condition (i). 
For  $j = j_1 + q$ its $j$-th row is 
\[ \bigl( \cdots \to \opn{Z}^j(M^{-1}) \xar{\pa^{-1}}
\opn{Z}(M^{0}) \xar{\pa^{0}} \opn{Z}^j(M^{1}) \to \cdots \bigr) , \]
and it is exact by condition (ii). 
And for $j > j_1 + q$ the $j$-th row of $G_q(M)$ is zero. 
Therefore, by step 1, the DG module 
$G_q \bigl( \opn{Tot}^{\oplus}(M) \bigr) := 
\opn{Tot}^{\oplus} \bigl( G_q(M) \bigr)$ 
is acyclic. 

Finally, since $M = \lim_{q \to} G_q(M)$, 
it follows that 
\[ \opn{Tot}^{\oplus}(M) = 
\lim_{q \to} \, G_q \bigl( \opn{Tot}^{\oplus}(M) \bigr) . \]
So $\opn{Tot}^{\oplus}(M)$ is acyclic.
\end{proof}

The second half of this subsection deals with products instead of direct sums. 
As explained in Proposition \ref{prop:1660}, 
if the abelian category $\cat{M}$ has 
countable products, then so does $\dcat{C}_{\mrm{str}}(A, \cat{M})$. 

\begin{dfn} \label{dfn:3320}
Assume that $\cat{M}$ has countable products. Given a complex 
$(M, \pa)$ with entries in $\dcat{C}_{\mrm{str}}(A, \cat{M})$, with notation as 
in \tup{(\ref{eqn:3300})}, its 
{\em product totalization}%
\index{Totalization! product}
is the object 
\[ \opn{Tot}^{\Pi}(M, \pa)  = 
\bigl( \opn{Tot}^{\Pi}(M), \d_{\mrm{Tot}} \bigr) \in  
\dcat{C}(A, \cat{M}) \]
defined as follows, in four stages. 
\begin{enumerate}
\item There is a graded object 
\[ \opn{Tot}^{\Pi}(M) := \prod_{i \in \Z} \, 
\opn{T}^{-i}(M^i) \in \dcat{G}(A^{\natural}, \cat{M}) . \]

\item On each factor $\opn{T}^{-i}(M^i)$ of $\opn{Tot}^{\Pi}(M)$
there is a differential 
$\d_{\opn{T}^{-i}(M^i)}$, 
and we let 
$\d_M := \prod_{i \in \Z} \, \d_{\opn{T}^{-i}(M^i)}$,
which is a differential on $\opn{Tot}^{\Pi}(M)$. 

\item For each $i$ let 
\[ \opn{tot}(\pa)^i := \opn{t}^{-1}_{\lms \opn{T}^{-(i + 1)}(M^{i + 1})} \circ 
\opn{T}^{-i}(\pa^i) : 
\opn{T}^{-i}(M^i) \to \opn{T}^{-(i + 1)}(M^{i + 1}) . \]
We define the degree $1$ operator 
$\opn{tot}(\pa) := \prod_{i \in \Z} \, \opn{tot}(\pa)^i$
on the graded object $\opn{Tot}^{\Pi}(M)$.

\item The differential $\d_{\mrm{Tot}}$ on the graded object 
$\opn{Tot}^{\Pi}(M)$ is 
$\d_{\mrm{Tot}} := \d_M + \opn{tot}(\pa)$.
\end{enumerate}
\end{dfn}

Note that (if $\cat{M}$ also has countable direct sums) there is a canonical 
embedding 
\begin{equation} \label{eqn:3370}
\opn{Tot}^{\oplus}(M) \sub \opn{Tot}^{\Pi}(M) 
\end{equation}
in $\dcat{C}_{\mrm{str}}(A, \cat{M})$. 
In case $M^i = 0$ for $\abs{i} \gg 0$, then this is an equality.

For Definition \ref{dfn:3320} to be valid we need to verify something: 

\begin{lem} \label{lem:3031}
The degree $1$ operator $\d_{\mrm{Tot}}$ makes the pair 
$\bigl( \opn{Tot}^{\Pi}(M), \d_{\mrm{Tot}} \bigr)$
into an object of $\dcat{C}(A, \cat{M})$. 
\end{lem}

\begin{proof}
This is the same calculation as in the proof of Lemma \ref{lem:3030},
heuristically using the inclusion (\ref{eqn:3370}).
\end{proof}

\begin{prop} \label{prop:3320}
Let 
\[ (M, \pa) =  \bigl( \cdots \to 0 \to M^{-1} \xar{\pa^{-1}}
M^{0} \xar{\pa^{0}} M^{1} \to \cdots \bigr) \]
be a complex with entries in the abelian category 
$\dcat{C}_{\mrm{str}}(A, \cat{M})$, let 
\[ (M', \pa) :=  \bigl( \cdots \to 0 \to M^{0} \xar{\pa^{0}} M^{1} \to \cdots 
\bigr) , \]
namely the stupid truncation of $M$ above $0$, and let 
$\rho : M^{-1} \to \opn{Tot}^{\Pi}(M', \pa)$
be the morphism in $\dcat{C}_{\mrm{str}}(A, \cat{M})$ induced by
$\pa^{-1} : M^{-1} \to M^0$. Then there is an isomorphism 
\[ \opn{Tot}^{\Pi}(M, \pa) \cong 
\opn{Cone} \bigl( \rho : M^{-1} \to \opn{Tot}^{\Pi}(M', \pa) \bigr) \]
in $\dcat{C}_{\mrm{str}}(A, \cat{M})$. 
\end{prop}

\begin{proof}
Like the proof of Proposition \ref{prop:3305}, only now there are no sign 
issues. 
\end{proof}

This directly implies:

\begin{cor} \label{cor:3320}
If, in the situation of Proposition \tup{\ref{prop:3320}}, the DG module \lb 
$\opn{Tot}^{\Pi}(M, \pa)$ is acyclic, then 
$\rho : M^{-1} \to \opn{Tot}^{\Pi}(M', \pa)$
is a quasi-isomorphism. 
\end{cor}

In Subsection \ref{subsec:zero} 
we discussed quotients in categories. They are dual to subobjects. Similarly, 
we now introduce a notion dual to a filtration of an object. 

\begin{dfn} \label{dfn:3336}
Let $M \in \dcat{C}(A, \cat{M})$.
\begin{enumerate}
\item A {\em cofiltration}% 
\index{Cofiltration on a DG module}
on $M$ is an inverse system
$\{ F_j(M) \}_{j \geq -1}$ in $\dcat{C}_{\mrm{str}}(A, \cat{M})$
consisting of quotients $M \surj F_j(M)$, indexed by the interval 
$[-1, \infty] \sub \Z$. 

\item We say that $M = \lim_{\lar j} F_j(M)$
if the inverse limit exists, and the 
canonical morphism $M \to \lim_{\lar j} F_j(M)$
in $\dcat{C}_{\mrm{str}}(A, \cat{M})$ is an isomorphism.

\item The cofiltration $F$ gives rise to the subquotients 
\[ \opn{Gr}_j^F(M) := \opn{Ker} \bigl( F_j(M) \to F_{j - 1}(M) \bigr) 
\in \dcat{C}(A, \cat{M}) . \]
\end{enumerate}
\end{dfn}

Recall that for a DG $A$-module $M$ we denote by 
$\opn{Y}(M) := M / \opn{B}(M)$; this is the object of decocycles of $M$, and 
it belongs to $\dcat{G}(\K)$.
See Definition \ref{dfn:2993} and equation \ref{eqn:4145}. 

\begin{prop} \label{prop:3321}
Let $(M, \pa)$ be a complex with entries in the abelian category 
$\dcat{C}_{\mrm{str}}(A)$, 
with internal structure \tup{(\ref{eqn:3305})}. Assume that the two conditions 
below hold\tup{:}
\begin{itemize}
\rmitem{i} For every $j \in \Z$ the complex  
\[ \bigl( \cdots \to M^{-1, j} \xar{\pa^{-1}}
M^{0, j} \xar{\pa^{0}} M^{1, j} \to \cdots \bigr) \]
with entries in $\dcat{M}(\K)$ is acyclic. 

\rmitem{ii} There is some $j_0 \in \Z$ such that for all 
$j \leq j_0$ the complex 
\[ \bigl( \cdots \to \opn{Y}^j(M^{-1}) \xar{\opn{Y}^j(\pa^{-1})}
\opn{Y}^j(M^{0}) \xar{\opn{Y}^j(\pa^{0})} \opn{Y}^j(M^{1}) \to \cdots \bigr) \]
with entries in $\dcat{M}(\K)$ is acyclic. 
\end{itemize}
Then the DG $A$-module $\opn{Tot}^{\Pi}(M, \pa)$ is acyclic. 
\end{prop}

\begin{proof} The proof is divided into two steps. 

\smallskip \noindent
Step 1. We replace condition (ii) by the stronger condition: 
\begin{itemize}
\rmitem{ii$'$} There is some $j'_{0} \in \Z$ such that 
$M^{i, j} = 0$ for all $j < j'_{0}$. 
\end{itemize}
For every $i$ we introduce a cofiltration 
$\{ F_q(M^i) \}_{q \geq -1}$
on the DG $\K$-module $M^i$ as follows: 
$F_q (M^i) := \opn{stt}^{\leq j'_{0} + q}(M^i)$, 
the stupid truncation below $j'_{0} + q$. 
These cofiltrations induce a cofiltration on the DG $\K$-module 
$\opn{Tot}^{\Pi}(M, \pa)$~: 
\[ F_q \bigl( \opn{Tot}^{\Pi}(M) \bigr) := 
\opn{Tot}^{\Pi} 
\bigl( \cdots \to F_q(M^{-1}) \xar{\pa^{-1}}
F_q(M^{0}) \xar{\pa^{0}} F_q(M^{1}) \to \cdots \bigr) . \]
We note that $F_{-1} \bigl( \opn{Tot}^{\Pi}(M) \bigr) = 0$, and 
\begin{equation} \label{eqn:3320}
\lim_{\lar q} \, F_q \bigl( \opn{Tot}^{\Pi}(M) \bigr) = 
\opn{Tot}^{\Pi}(M) . 
\end{equation}

For every $q \geq 0$ the DG $\K$-module 
$\opn{Gr}^F_q \bigl( \opn{Tot}^{\Pi}(M) \bigr)$
is isomorphic in $\dcat{C}_{\mrm{str}}(\K)$ to the complex in condition (i) 
with index $j = j'_{0} + q$ (up to a sign change of the differentials), so it 
is acyclic. For $q = -1$ we trivially have an acyclic DG $\K$-module 
$F_{-1} \bigl( \opn{Tot}^{\Pi}(M) \bigr)$. 
For every $q \geq 0$ there is an exact sequence 
\[ 0 \to \opn{Gr}^F_q \bigl( \opn{Tot}^{\Pi}(M) \bigr) \to 
F_q \bigl( \opn{Tot}^{\Pi}(M) \bigr) \to
F_{q - 1} \bigl( \opn{Tot}^{\Pi}(M) \bigr) \to 0 \]
in $\dcat{C}_{\mrm{str}}(\K)$. By induction on $q$ we conclude that 
$F_q \bigl( \opn{Tot}^{\Pi}(M) \bigr)$
is acyclic. Finally, because the inverse system of complexes of $\K$-modules
$(\ref{eqn:3320})$ has surjective transitions, the Mittag-Leffler argument 
(see Corollary \ref{cor:1590}) says that the DG module 
$\opn{Tot}^{\Pi}(M, \pa)$ is acyclic. 

\medskip \noindent 
Step 2. Here we assume conditions (i) and (ii). 
For each $i$ we introduce a cofiltration 
$\{ G_q(M^i) \}_{q \geq 0}$
on the DG $\K$-module $M^i$ as follows: 
\[ G_q(M^i) := \opn{smt}^{\geq j_0 - q}(M^i) = 
\bigl( \cdots \to 0 \to \opn{Y}^{j_0 - q}(M^i) 
\xar{\d_{M^i}} M^{i, j_0 - q + 1} \to \cdots \bigr) , \]
the smart truncation above $j_0 - q$. 

For each $q$ the complex 
\[ G_q(M) := \bigl( \cdots \to G_q(M^{-1}) \xar{G_q(\pa^{-1})}
G_q(M^{0}) \xar{G_q(\pa^{0})} G_q(M^{1}) \to \cdots \bigr) \]
satisfies conditions (i) and (ii$'$), with $j'_{0} := j_0 - q$.
Indeed, its $j$-th row is 
\[ \bigl( \cdots \to M^{-1, j} \xar{\pa^{-1}}
M^{0, j} \xar{\pa^{0}} M^{1, j} \to \cdots \bigr) \]
for $j > j_0 - q$; it is 
\[ \bigl( \cdots \to \opn{Y}^j(M^{-1}) \xar{\opn{Y}^j(\pa^{-1})}
\opn{Y}(M^{0}) \xar{\opn{Y}^j(\pa^{0})} \opn{Y}^j(M^{1}) \to \cdots \bigr) \]
for $j = j_0 - q$; and it is zero for $j < j_0 - q$. All these rows are exact. 
Therefore, by step 1, the DG module 
$G_q \bigl( \opn{Tot}^{\Pi}(M) \bigr) := \opn{Tot}^{\Pi} \bigl( G_q(M) \bigr)$
is acyclic. 

Finally, since $M = \lim_{\lar q} \, G_q(M)$, 
it follows that 
\[ \opn{Tot}^{\Pi}(M) = \lb \lim_{\lar q} \,
G_q \bigl( \opn{Tot}^{\Pi}(M) \bigr) . \]
This is an inverse system of complexes of $\K$-modules with surjective 
transitions. According to the Mittag-Leffler argument the limit
$\opn{Tot}^{\Pi}(M)$ is acyclic.
\end{proof}



\mysubsection{K-Projective Resolutions in 
\texorpdfstring{$\bcat{C}^-(\cat{M})$}{C-(M)}} 
\label{subsec:exis-K-prj}


Recall that $\cat{M}$ is a $\K$-linear abelian category, and 
$\dcat{C}(\cat{M})$ is the DG category of complexes in $\cat{M}$. The strict 
category $\dcat{C}_{\mrm{str}}(\cat{M})$ is abelian. 

The next definition is inspired by the work of B. Keller 
\cite[Section 3.1]{Kel1}. 

\begin{dfn} \label{dfn:1535}
Let $P$ be an object of $\dcat{C}(\cat{M})$. 
\begin{enumerate}
\item A {\em semi-projective filtration}% 
\index{Filtration on a DG module! semi-projective}
on $P$ is a filtration 
$F = \{ F_j(P) \}_{j \geq -1}$
on $P$ as an object of $\dcat{C}_{\mrm{str}}(\cat{M})$, such that:
\begin{itemize}
\item $F_{-1}(P) = 0$. 
\item  Each $\opn{Gr}^F_j(P)$ 
is a complex of projective objects of $\cat{M}$ with zero differential.
\item $P = \lim_{j \to} F_j(P)$ in $\dcat{C}_{\mrm{str}}(\cat{M})$.
\end{itemize}

\item The complex $P$ is called a 
{\em semi-projective complex}% 
\index{Complex in abelian category! semi-projective}
if it admits some semi-projective filtration. 
\end{enumerate}
\end{dfn}

K-projective objects in $\dcat{C}(\cat{M})$ were introduced in Definition 
\ref{dfn:1520}.

\begin{thm} \label{thm:1536}
\index{Complex in abelian category! semi-projective}
\index{Complex in abelian category! K-projective}
Let $\cat{M}$ be an abelian category, and let $P$ be a 
semi-projective complex in $\dcat{C}(\cat{M})$. Then $P$ is 
K-projective. 
\end{thm}

\begin{proof} The proof is in four steps. 

\medskip \noindent
Step 1. We start by proving that if $P = \opn{T}^k(Q)$, the translation of a 
projective object $Q \in \cat{M}$, then $P$ is K-projective. This is easy: 
given an acyclic complex $N \in \dcat{C}(\cat{M})$, we have 
\[ \opn{Hom}_{\cat{M}}(P, N) = 
\opn{Hom}_{\cat{M}} \bigl( \opn{T}^k(Q), N \bigr) \cong
\opn{T}^{-k} \bigl( \opn{Hom}_{\cat{M}}(Q, N) \bigr) \]
in $\dcat{C}_{\mrm{str}}(\K)$. But 
$\opn{Hom}_{\cat{M}}(Q, -)$ is an exact functor 
$\cat{M} \to \dcat{M}(\K)$, so $\opn{Hom}_{\cat{M}}(Q, N)$ is an acyclic 
complex. 

\medskip \noindent
Step 2. Now $P$ is a complex of projective objects of $\cat{M}$ with zero 
differential. This means that 
$P \cong \bigoplus_{k \in \Z} \, \opn{T}^k(Q_k)$ 
in $\dcat{C}_{\mrm{str}}(\cat{M})$, where each $Q_k$ is a projective object in 
$\cat{M}$. But then 
\[ \opn{Hom}_{\cat{M}}(P, N) \cong 
\prod_{k \in \Z} \, \opn{Hom}_{\cat{M}} \bigl( \opn{T}^k(Q_k), N \bigr)  \]
in $\dcat{C}_{\mrm{str}}(\K)$. 
This is an easy case of Proposition \ref{prop:1535}. By step 1 and the fact 
that a product of acyclic complexes in $\dcat{C}_{\mrm{str}}(\K)$ is acyclic 
(itself an easy case of the Mittag-Leffler argument), we conclude that 
$\opn{Hom}_{\cat{M}}(P, N)$ is acyclic. 

\medskip \noindent
Step 3. Fix a semi-projective filtration $F = \{ F_j(P) \}_{j \geq -1}$
on $P$. Here we prove that for every $j$ the complex  
$F_j(P)$ is K-projective. This is done by induction on $j \geq -1$. 
For $j = -1$ it is trivial. For $j \geq 0$ there is an exact sequence 
\begin{equation} \label{eqn:1555}
0 \to F_{j - 1}(P) \to F_j(P) \to \opn{Gr}^F_j(P) \to 0 
\end{equation}
in $\dcat{C}_{\mrm{str}}(\cat{M})$. 
In each degree $i \in \Z$ the exact sequence 
\[ 0 \to F_{j - 1}(P)^i \to F_j(P)^i \to \opn{Gr}^F_j(P)^i \to 0 \]
in $\cat{M}$ splits, because $\opn{Gr}^F_j(P)^i$ is a projective object.
Thus the exact sequence (\ref{eqn:1555}) is split exact in the abelian category 
$\bcat{G}_{\mrm{str}}(\cat{M})$ of graded objects in $\cat{M}$. 

Let $N \in \dcat{C}(\cat{M})$ be an acyclic complex. Applying the 
functor $\opn{Hom}_{\cat{M}}(-, N)$ 
to the sequence of complexes (\ref{eqn:1555}) we obtain a sequence 
\begin{equation} \label{eqn:1556}
0 \to \opn{Hom}_{\cat{M}} \bigl( \opn{Gr}^F_j(P), N \bigr) \to 
\opn{Hom}_{\cat{M}} \bigl( F_j(P), N \bigr) \to
\opn{Hom}_{\cat{M}} \bigl( F_{j - 1}(P), N \bigr) \to  0 
\end{equation}
in $\dcat{C}_{\mrm{str}}(\K)$. Because (\ref{eqn:1555}) is split exact in 
$\bcat{G}_{\mrm{str}}(\cat{M})$, the sequence (\ref{eqn:1556}) is split exact 
in $\bcat{G}_{\mrm{str}}(\K)$. Therefore (\ref{eqn:1556}) is exact in 
$\dcat{C}_{\mrm{str}}(\K)$.

By the induction hypothesis the complex 
$\opn{Hom}_{\cat{M}} \bigl( F_{j - 1}(P), N \bigr)$
is acyclic. By step 2 the complex
$\opn{Hom}_{\cat{M}} \bigl( \opn{Gr}^F_{j}(P), N \bigr)$
is acyclic. The long exact cohomology sequence associated to 
(\ref{eqn:1556}) shows that the complex 
$\opn{Hom}_{\cat{M}} \bigl( F_j(P), N \bigr)$
is acyclic too.

\medskip \noindent
Step 4. We keep the semi-projective filtration 
$F = \{ F_j(P) \}_{j \geq -1}$ from step 3. 
Take any acyclic complex $N \in \bcat{C}(\cat{M})$. By 
Proposition \ref{prop:1535} we know that 
\[ \opn{Hom}_{\cat{M}}(P, N) \cong 
\lim_{\leftarrow j} \, \opn{Hom}_{\cat{M}} \bigl( F_j(P), N \bigr) \]
in $\dcat{C}_{\mrm{str}}(\K)$. According to step 3 the complexes 
$\opn{Hom}_{\cat{M}} \bigl( F_j(P), N \bigr)$ are all acyclic. 
The exactness of the sequences (\ref{eqn:1556}) implies that 
the inverse system \lb 
$\bigl\{ \opn{Hom}_{\cat{M}} \bigl( F_j(P), N \bigr) \bigr\}_{j \geq -1}$
in $\dcat{C}_{\mrm{str}}(\K)$ has surjective transitions. Now the 
Mittag-Leffler 
argument (Corollary \ref{cor:1590}) says that the inverse limit complex
$\opn{Hom}_{\cat{M}}(P, N)$ is acyclic. 
\end{proof}

\begin{prop} \label{prop:1900}
Let $\cat{M}$ be an abelian category. If $P \in \dcat{C}(\cat{M})$ is a bounded 
above complex of projectives, then $P$ is a semi-projective complex.
\end{prop}

\begin{proof}
Say $P$ is nonzero and $\opn{sup}(P) = i_1 \in \Z$. 
For $j \geq -1$ define 
$F_j(P) := \opn{stt}^{\geq i_1 - j}(P)$,
the stupid truncation above $i_1 - j$ from Definition \ref{dfn:3066}.
Then $\{ F_j(P) \}_{j \geq -1}$ is a semi-projective filtration on $P$. 
\end{proof}

The next theorem is dual to \cite[Lemma I.4.6(1)]{RD}, in the sense of 
changing direction of arrows. We give a detailed proof. 

\begin{thm} \label{thm:1540}
Let $\cat{M}$ be an abelian category, and let $\cat{P} \sub \cat{M}$ be a
full subcategory such that every object $M \in \cat{M}$ admits an 
epimorphism $P \surj M$ from some object $P \in \cat{P}$. Then every complex 
$M \in \bcat{C}^{-}(\cat{M})$ admits a quasi-isomorphism 
$\rho : P \to  M$ in $\dcat{C}^-_{\mrm{str}}(\cat{M})$,
such that $\opn{sup}(P) = \opn{sup}(M)$, and 
each $P^i$ is an object of $\cat{P}$.
\end{thm}

\begin{proof}
This is divided into steps. 

\smallskip \noindent
Step 1. We can assume that $M \neq 0$. 
After translating $M$, we can assume that $\opn{sup}(M) = 0$. 

\medskip \noindent
Step 2. Say the differential of the complex $M$ is 
$\d^{i}_M : M^i \to M^{i + 1}$.
Let us choose an epimorphism $\rho^{0} : P^{0} \surj M^{0}$ in $\cat{M}$ 
from some object $P^{0} \in \cat{P}$. We get a morphism 
$\de^{0} :  M^{-1} \oplus P^0 \to M^0$
whose components are $\d^{-1}_M$ and $-\rho^{0}$. Thus 
$\opn{Ker}(\de^{0}) = M^{-1} \times_{M^{0}} P^0$,
the fibered product. 
Next we choose an epimorphism
$\psi^{-1} : P^{-1} \surj \opn{Ker}(\de^{0})$  
from some object $P^{-1} \in \cat{P}$.
So there is an exact sequence 
\begin{equation} \label{eqn:5060}
P^{-1} \xar{\psi^{-1}} M^{-1} \oplus P^0 \xar{\de^0} M^0 \to 0  .
\end{equation}
The components of $\psi^{-1}$ are denoted by 
$\rho^{-1} : P^{-1} \to  M^{-1}$ and 
$\d_P^{-1} : P^{-1} \to  P^{0}$.
We have this commutative diagram 
\begin{equation} \label{eqn:5058}
\UseTips \xymatrix @C=6ex @R=6ex {
&
P^{-1}
\ar[d]^{\rho^{-1}}
\ar[r]^{\d_P^{-1}}
&
P^0 
\ar[d]^{\rho^0}
\ar[r]
&
0
\\
M^{-2}
\ar[r]^{\d_M^{-2}}
&
M^{-1}
\ar[r]^{\d_M^{-1}}
&
M^0 
\ar[r]
&
0
} 
\end{equation}
in $\cat{M}$.

\medskip \noindent
Step 3. This is the inductive step. Here $i \leq -1$, and we already have 
objects $P^{i}, \ldots, P^0$ in $\cat{P}$, and morphisms 
$\rho^{i}, \ldots, \rho^0$ and $\d_P^{i}, \ldots, \d_P^{-1}$,
that fit into this commutative diagram
\begin{equation} \label{eqn:1745}
\UseTips \xymatrix @C=6ex @R=6ex {
&
P^i
\ar[d]^{\rho^i}
\ar[r]^{\d_P^i}
&
P^{i + 1}
\ar[d]^{\rho^{i + 1}}
\ar[r] 
&
\cdots 
\ar[r]^{\d_P^{-1}}
& 
P^0 
\ar[d]^{\rho^0}
\ar[r]
&
0
\\
M^{i - 1}
\ar[r]^{\d_M^{i - 1}}
&
M^i
\ar[r]^{\d_M^{i}}
&
M^{i + 1}
\ar[r] 
&
\cdots 
\ar[r]^{\d_M^{-1}}
& 
M^0 
\ar[r]
&
0
} 
\end{equation}
in $\cat{M}$. Also 
$\d_P^{j + 1} \circ \d_P^{j} = 0$ for all $j$ in the interval $[i, -1]$

Define the morphism 
$\de^{i}  :  M^{i - 1} \oplus P^{i} \to  M^{i} \oplus P^{i + 1}$
to be the one with components 
$\d_M^{i - 1}$, $- \rho^i$ and $- \d_P^{i}$.
Expressing direct sums of objects as columns, and letting matrices of morphisms 
act on them from the left, we have this matrix representation of $\de^i$~:
\begin{equation} \label{eqn:1905}
\de^i = \bmat{ \d_M^{i - 1} & - \rho^i \\[0.3em] 0 & - \d_P^i } .
\end{equation}
So 
$\opn{Ker}(\de^{i}) = M^{i - 1} \times_{(M^i \oplus P^{i + 1})} P^{i}$,
the fibered product. 
 
Let us choose an epimorphism
$\psi^{i - 1} : P^{i - 1} \surj  \opn{Ker}(\de^{i})$ 
from an object $P^{i - 1} \in \cat{P}$. 
We get an exact sequence 
\begin{equation} \label{eqn:1750}
P^{i - 1} \xar{\psi^{i - 1}} M^{i - 1} \oplus P^{i} \xar{\de^i}
M^{i} \oplus P^{i + 1} .
\end{equation}
The components of the morphism $\psi^{i - 1}$ are denoted by 
$\rho^{i - 1} : P^{i - 1} \to  M^{i - 1}$ and 
$\d_P^{i - 1} : P^{i - 1} \to  P^{i}$.
In matrix representation is 
\begin{equation} \label{eqn:5061}
\psi^{i - 1} = \bmat{ \rho^{i - 1} \\[0.3em] \d_P^{i - 1} } . 
\end{equation}
In this way we obtain the slightly bigger diagram than (\ref{eqn:1745}):
\begin{equation} \label{eqn:1751}
\UseTips \xymatrix @C=6ex @R=6ex {
P^{i - 1}
\ar[d]^{\rho^{i - 1}}
\ar[r]^{\d_P^{i - 1}}
&
P^i
\ar[d]^{\rho^i}
\ar[r]^{\d_P^i}
&
P^{i + 1}
\ar[d]^{\rho^{i + 1}}
\ar[r] 
&
\cdots 
\ar[r]
& 
P^0 
\ar[d]^{\rho^0}
\ar[r]
&
0
\\
M^{i - 1}
\ar[r]^{\d_M^{i - 1}}
&
M^i
\ar[r]^{\d_M^{i}}
&
M^{i + 1}
\ar[r] 
&
\cdots 
\ar[r]
& 
M^0 
\ar[r]
&
0
} 
\end{equation}
Because $\de^i \circ \psi^{i - 1} = 0$ in (\ref{eqn:1750}), it follows that
$\d_P^{i} \circ \d_P^{i - 1} = 0$, and also 
\begin{equation} \label{eqn:1753}
\rho^{i} \circ \d_P^{i - 1} = \d_M^{i - 1}  \circ \rho^{i - 1} ,
\end{equation}
showing that diagram (\ref{eqn:1751}) is commutative. 

\medskip \noindent
Step 4. We carry out the construction in step 3 inductively for all $i \leq -1$, 
thus obtaining a diagram like (\ref{eqn:1751}) that goes infinitely to the 
left. 

Letting $P^i := 0$ for positive $i$, the 
collection $P := \{ P^i \}_{i \in \Z}$ becomes a complex, with 
differential $\d_P := \{ \d_P^{i} \}_{i \in \Z}$.
By equation (\ref{eqn:1753}), the collection 
$\rho := \{ \rho^{i} \}_{i \in \Z}$ is a strict morphism 
of complexes $\rho : P \to M$. 

\medskip \noindent
Step 5. It remains to prove that $\rho$ is a quasi-isomorphism.

Take any $i \leq 0$. Let us examine this diagram: 
\begin{equation} \label{eqn:1752}
\UseTips \xymatrix @C=6ex @R=6ex {
P^{i - 1}
\ar[r]^{- \psi^{i - 1}}
\ar[d]^{(0, \opn{id})}
&
M^{i - 1} \oplus P^{i}
\ar[r]^{\de^{i}}
\ar[d]^{\opn{id}}
&
M^{i} \oplus P^{i + 1}
\ar[d]^{\opn{id}}
\\
M^{i - 2} \oplus P^{i - 1}
\ar[r]^{\de^{i - 1}}
&
M^{i - 1} \oplus P^{i}
\ar[r]^{\de^{i}}
&
M^i \oplus P^{i+ 1}
}
\end{equation}
Comparing the formula for $\de^{i - 1}$ in (\ref{eqn:1905}) to the 
formula for $- \psi^{i - 1}$ in (\ref{eqn:5061}), we see that this diagram is 
commutative. 
An easy calculation using (\ref{eqn:1753}) shows that 
$\de^{i} \circ \de^{i - 1} = 0$. 
The top row is exact, because (up to signs) for $i = 0$ it is part of the 
exact sequence (\ref{eqn:5060}), and  for $i \leq -1$ it is the exact sequence 
(\ref{eqn:1750}). 
It follows that the bottom row is also exact.

Let $N =  \{ N^i \}_{i \in \Z}$ be the complex with  components 
$N^i := M^{i - 1} \oplus P^{i}$ for all $i$. 
The differential
$\d_N= \{ \d_N^{i} \}_{i \in \Z}$
is $\d^i_N := -\de^{i} : N^i \to N^{i + 1}$.
We have $N^i = 0$ for $i \geq 2$. 
The exactness of the sequence (\ref{eqn:5060}) says that 
$\opn{H}^1(N) = 0$, and the exactness of the second row in 
(\ref{eqn:1752}) says that $\opn{H}^i(N) = 0$ for $i \leq 0$. 
Hence the complex $N$ is acyclic. 
On the other hand, by the definition of the morphisms 
$\de^i$ in (\ref{eqn:1905}), we see that $N$ is just the standard cone on the 
strict morphism of complexes
$\opn{T}^{-1}(\rho) : \opn{T}^{-1}(P) \to \opn{T}^{-1}(M)$.
% where $(M', \d_{M'})$ is the complex $(M, -\d_{M})$. 
See Definition \ref{dfn:1172}. 
Therefore $\opn{T}^{-1}(\rho)$ is a quasi-isomorphism, and so is $\rho$.
\end{proof}

\begin{exa} \label{exa:4725}
Let $(X, \AA)$ be a ringed space. Usually there aren't enough projective 
objects in the abelian category $\dcat{M}(\AA) = \cat{Mod} \AA$ of sheaves of 
$\AA$-modules on $X$; but there are always enough flat objects. A general method 
to produce flat sheaves is as follows. Given an open set $U \sub X$ with 
inclusion map $g : U \to X$,  
the extension by zero sheaf $g_{!}(\AA|_{U})$ is a flat $\AA$-module.  
Let us call $g_{!}(\AA|_{U})$ a pseudo-free $\AA$-module of pseudo-rank 
$1$ and pseudo-support $U$ . A {\em pseudo-free $\AA$-module} is, by 
definition, a sheaf $\PP$ which is a direct sum of rank $1$ pseudo-free 
$\AA$-modules (with varying pseudo-supports). A pseudo-free $\AA$-module is 
flat, and every $\AA$-module $\MM$ admits an epimorphism
$\PP \surj \MM$ from some pseudo-free $\AA$-module $\PP$. 
More on this in \cite[Proposition II.1.2]{RD} or 
\cite[Proposition 2.4.12]{KaSc1}. 
(Warning: in case $(X, \OO_X)$ is a scheme, the pseudo-free $\OO_X$-modules are 
usually not quasi-coherent.) 

Now let $\dcat{C}(\AA) =  \dcat{C}(\cat{Mod} \AA)$ be the category of complexes 
of $\AA$-modules. Because $\dcat{M}(\AA)$ has enough flat objects, 
Theorem \ref{thm:1540} tells us that every complex 
$\MM \in \dcat{C}^{-}(\AA)$ admits a quasi-isomorphism 
$\PP \to \MM$ from a bounded above complex of flat 
$\AA$-modules $\PP$. Similarly to Proposition \ref{prop:1900} and Theorem 
\ref{thm:1536}, one can show that such a complex $\PP$ is K-flat. 
The conclusion is that the subcategory 
$\dcat{C}^{-}(\AA)$ has enough K-flat objects. 

But in fact more is true: the category of unbounded complexes $\dcat{C}^{}(\AA)$ 
also has enough K-flat objects; see \cite{Spa}. 
For an even more general statement see Remark \ref{rem:4675} below. 
\end{exa}

\begin{dfn} \label{dfn:2895}
Let $\cat{M}$ be an abelian category, and let $\cat{M}' \sub \cat{M}$ be a 
full abelian subcategory. 
We say that {\em $\cat{M}'$ has enough projectives relative to $\cat{M}$} 
if every object $M \in \cat{M}'$ admits an epimorphism 
$P \surj M$, where $P$ is an object of $\cat{M}'$ that is projective in the 
bigger category $\cat{M}$.
\end{dfn}

Of course, if $\cat{M}'$ has enough projectives relative to $\cat{M}$
then $\cat{M}'$ itself has enough projectives. 

Thick abelian categories were defined in Definition \ref {dfn:2323}.
The next theorem is the opposite of \cite[Lemma I.4.6(3)]{RD}. 

\begin{thm} \label{thm:2895}
Let $\cat{M}$ be an abelian category, and let $\cat{M}' \sub \cat{M}$ be a 
thick abelian subcategory that has enough projectives relative to $\cat{M}$. 
Let $M \in \dcat{C}(\cat{M})$ be a complex with bounded above cohomology, 
such that $\opn{H}^i(M) \in \cat{M}'$ for all $i$.
Then there is a quasi-isomorphism $\rho : P \to M$ in 
$\dcat{C}_{\mrm{str}}(\cat{M})$, where 
$P \in \dcat{C}^-(\cat{M}')$, each $P^i$ is projective in $\cat{M}$,
and $\opn{sup}(P) = \opn{sup}(\opn{H}(M))$. 
\end{thm}

\begin{proof}
The proof of Theorem \ref{thm:2880}, reversed, works here. 
(Despite being located later in the book, Theorem \ref{thm:2880} is 
not logically dependent on the current theorem.)
To be explicit, let us take $\cat{N} := \cat{M}^{\mrm{op}}$ and 
$\cat{N}' := \cat{M}'^{\, \mrm{op}}$.
Since monomorphisms in $\cat{M}$ become epimorphisms in $\cat{N}$, 
and projective objects in $\cat{M}$ become injective objects in $\cat{N}$,
the full abelian subcategory $\cat{N}' \sub \cat{N}$ satisfies 
assumptions of Theorem \ref{thm:2880}. By Theorem \ref{thm:2495} we have a 
canonical isomorphism of categories 
$\bcat{C}^{+}_{\mrm{str}}(\cat{N}) \iso 
\dcat{C}^-_{\mrm{str}}(\cat{M})^{\mrm{op}}$
that respects the cohomology functor. 
Thus a quasi-isomorphism $N \to J$ in 
$\bcat{C}^{+}_{\mrm{str}}(\cat{N})$
gives rise to a quasi-isomorphism $P \to M$ in 
$\bcat{C}^{-}_{\mrm{str}}(\cat{M})$.
\end{proof}

\begin{cor} \label{cor:1900}
\index{Resolution! K-projective}
If $\cat{M}$ is an abelian category with enough projectives, then  
$\bcat{C}^{-}(\cat{M})$ has enough K-projectives. 
\end{cor}

\begin{proof}
According to either Theorem \ref{thm:1540} or Theorem \ref{thm:2895},
every $M \in \bcat{C}^{-}(\cat{M})$ admits a quasi-isomorphism $P \to M$
from a bounded above complex of projectives $P$. Now use Proposition 
\ref{prop:1900} and Theorem \ref{thm:1536}. 
\end{proof}

\begin{cor} \label{cor:1750}
Let $\cat{M}$ be an abelian category with enough projectives, 
and let $M \in \dcat{C}(\cat{M})$ be a complex with bounded above cohomology. 
Then $M$ has a K-projective resolution $P \to M$,
such that $\opn{sup}(P) = \opn{sup}(\opn{H}(M))$,
and every $P^i$ is an projective object of $\cat{M}$. 
\end{cor}

\begin{proof}
We may assume that $\opn{H}(M)$ is not zero. 
Let $i := \opn{sup}(\opn{H}(M)) \in \Z$, and take $N := \opn{smt}^{\leq i}(M)$,
the smart truncation from  Definition \ref{dfn:2320}. Then 
$N \to M$ is a quasi-isomorphism and $\opn{sup}(N) = i$.
According to either Theorem \ref{thm:1540} or Theorem \ref{thm:2895},
there is a quasi-isomorphism $P \to N$, where $P$ is a complex 
of projectives and $\opn{sup}(P) = i$. 
By Proposition \ref{prop:1900} and Theorem \ref{thm:1536} the complex $P$ is 
K-projective. The composed quasi-isomorphism $P \to M$ is what we are looking 
for.
\end{proof}

\begin{cor} \label{cor:3175}
Under the assumptions of Theorem \tup{\ref{thm:2895}}, the canonical functor 
$\dcat{D}^-(\cat{M}') \to \dcat{D}^-_{\cat{M}'}(\cat{M})$
is an equivalence. 
\end{cor}

\begin{proof}
Consider the commutative diagram 
\[ \UseTips \xymatrix @C=6ex @R=6ex {
\dcat{K}^-(\cat{M}')_{\catt{M}\tup{-}\mrm{prj}}
\ar@(l,u)[drr]^{F}
\ar[d]_{\opn{Q}}
&
\\
\dcat{D}^-(\cat{M}')
\ar[r]^{G}
& 
\dcat{D}^-_{\cat{M}'}(\cat{M})
\ar[r]^{\tup{f.f.}}
& 
\dcat{D}(\cat{M})
} \]
where $\dcat{K}^-(\cat{M}')_{\catt{M}\tup{-}\mrm{prj}}$ is the homotopy 
category of bounded above complexes of objects of $\cat{M}'$ that are 
projective in $\cat{M}$. Since these are K-projective complexes in 
$\dcat{K}(\cat{M})$, the functors $F$ and $\opn{Q}$ are fully faithful, by 
Corollary \ref{cor:3145}. The functor marked ``f.f.'' is the fully faithful 
inclusion of this full subcategory. We conclude that the functor $G$ is also 
fully faithful. The theorem tells us that the essential 
image of $G$ is $\dcat{D}^-_{\cat{M}'}(\cat{M})$. 
\end{proof}

Here is an important instance where Theorem \ref{thm:2895} and  
Corollary \ref{cor:3175} apply. 

\begin{exa} \label{exa:2900}
Let $A$ be a left noetherian ring. Consider the abelian category 
$\cat{M} := \dcat{M}(A) = \cat{Mod} A$, and the thick abelian subcategory 
$\cat{M}' := \dcat{M}_{\mrm{f}}(A) = \cat{Mod}_{\mrm{f}} A$ of finitely 
generated modules. Then $\cat{M}'$ has enough projectives relative to 
$\cat{M}$. 
Theorem \ref{thm:2895} tells us that if 
$M \in \dcat{C}(\cat{M})$ is a complex such that the modules 
$\opn{H}^i(M)$ are all finitely generated, and $\opn{H}^i(M) = 0$ for 
$i \gg 0$, then there is a resolution $P \to M$, where $P$ is a bounded above 
complex of finitely generated projective modules. 
Corollary \ref{cor:3175} says that the canonical functor 
\[ \dcat{D}^{-}(\dcat{M}_{\mrm{f}}(A)) = \dcat{D}^{-}(\cat{M}') \to 
\dcat{D}^{-}_{\cat{M}'}(\cat{M}) = \dcat{D}^{-}_{\mrm{f}}(\dcat{M}(A))
= \dcat{D}^{-}_{\mrm{f}}(A) \]
is an equivalence. 
See Example \ref{exa:1930} for another approach to this problem. 
\end{exa}

We end this subsection with a result on K-flat complexes over rings. 

\begin{prop} \label{prop:4535}
Let $A$ be a ring and $P$ a bounded above complex of flat $A$-modules. 
Then $P$ is a K-flat complex. 
\end{prop}

\begin{exer} \label{exer:4535}
Prove this proposition. (Hint: see the proofs of Proposition \ref{prop:1900} 
and Theorem \ref{thm:1536}.)
\end{exer}




\newpage
\mysubsection{K-Projective Resolutions in 
\texorpdfstring{$\bcat{C}(A)$}{C(A)}}
\label{subsec:exis-K-prj-dgmods}

In this subsection $A$ is a DG $\K$-ring (without any vanishing assumption). 

Let $A^{\natural}$ be the graded ring that we have after forgetting the 
differential of $A$. In Subsection \ref{subsec:DGModinM} we introduced the 
functor 
$\opn{Und} : \dcat{C}(A) \to \dcat{G}(A^{\natural})$
that forgets the differentials of DG modules. We shall now use the abbreviated 
form $M^{\natural}  := \opn{Und}(M)$. 

Recall that the translation $\opn{T}^{-i}(A)$ is a DG $A$-module in which the 
element $\opn{t}^{-i}(1_A)$ is in degree $i$. 
This element is a cocycle, and when we forget the differentials, the graded 
module $\opn{T}^{-i}(A)^{\natural}$ is free over the graded ring 
$A^{\natural}$, with basis $\opn{t}^{-i}(1_A)$. Therefore, for every DG 
$A$-module $M$ there is a canonical isomorphism
\begin{equation} \label{eqn:1755}
\opn{Hom}_{A} \bigl( \opn{T}^{-i}(A), M \bigr) \cong \opn{T}^{i}(M)
\end{equation}
in $\dcat{C}_{\mrm{str}}(\K)$, and canonical isomorphisms 
\begin{equation} \label{eqn:1600}
\opn{Hom}_{\dcat{C}_{\mrm{str}}(A)} \bigl( \opn{T}^{-i}(A), M \bigr) \cong
\opn{Z}^0  \bigl( \opn{Hom}_{A} ( \opn{T}^{-i}(A), M ) \bigr) \cong \opn{Z}^i(M)
\end{equation}
in $\dcat{M}(\K)$. 
(Actually, (\ref{eqn:1755}) is an isomorphism in $\dcat{C}_{\mrm{str}}(A)$,
but this uses the DG $A$-bimodule structure of $\opn{T}^{-i}(A)$.)

We begin with a definition that is very similar to Definition \ref{dfn:1535}. 

\begin{dfn} \label{dfn:1575}
Let $P$ be an object of $\dcat{C}(\cat{A})$. 
\begin{enumerate}
\item We say that $P$ is a {\em free DG $A$-module}% 
\index{Differential graded module! free}
if there is an isomorphism
$P \cong \bigoplus_{s \in S} \opn{T}^{-i_s}(A)$
in $\dcat{C}_{\mrm{str}}(A)$, 
for some indexing set $S$ and some collection of integers 
$\{ i_s \}_{s \in S}$. 

\item A {\em semi-free filtration}% 
\index{Filtration on a DG module! semi-free}
on $P$ is a filtration 
$F = \{ F_j(P) \}_{j \geq -1}$
of $P$ in $\dcat{C}_{\mrm{str}}(A)$, such that:
\begin{itemize}
\item $F_{-1}(P) = 0$. 
\item  Each $\opn{Gr}^F_j(P)$ is a free DG $A$-module.
\item $P = \lim_{j \to} F_j(P)$.
\end{itemize}

\item The DG module $P$ is called {\em semi-free}%
\index{Differential graded module! semi-free}
if it admits some semi-free filtration. 
\end{enumerate}
\end{dfn}

Note that the direct limit in item (2) is just the union of DG submodules.

There is an alternative description of semi-free filtrations, that is useful 
for understanding their structure. We will give it in Definition 
\ref{dfn:4732} below, after some combinatorial preliminaries. 

The next definition has already appeared within Example \ref{exa:4625}.

\begin{dfn} \label{dfn:4730}
A {\em graded set}%
\index{Graded set}
is a set $S$ that is partitioned into subsets: 
$S = \coprod_{i \in \Z} \, S^i$. 
The elements of $S^i$ are said to have degree $i$,
so that $S^i = \{ s \in S \mid \opn{deg}(s) = i \}$.
\end{dfn}

A free DG $A$-module $P$ as in Definition \ref{dfn:1575}(1) can be described as 
follows. The indexing set $S$ can be made into a graded set by defining
$S^k := \{ s \in S \mid i_s = k \}$.
In other words, $\opn{deg}(s) = i_s$. We denote by $A \cd s$ the free DG 
$A$-module with basis $s$, so that $\d(s) = 0$, and there is an isomorphism 
$A \cd s \iso \opn{T}^{-i_s}(A)$ in $\dcat{C}_{\mrm{str}}(A)$
sending $s \mapsto \opn{t}^{-i_s}(1_A)$. Writing 
\begin{equation} \label{eqn:5062}
A \cd S := \bigoplus_{s \in S} \ A \cd s
\end{equation}
we obtain $A \cd S \cong P$ as DG $A$-modules. 
Note that $A^{\natural} \cd S \cong P^{\natural}$
as graded $A^{\natural}$-modules. 

\begin{dfn} \label{dfn:4731}
Let $S$ be a graded set. A {\em filtration} on $S$ is a direct system
$F = \bigl\{ F_{j}(S) \bigr\}_{j \geq -1}$
of subsets of $S$, such that 
$F_{-1}(S) = \varnothing$ and $\bigcup_{j} F_{j}(S) = S$.  
The pair $(S, F)$ is called a {\em filtered graded set}% 
\index{Graded set! filtered}.
\end{dfn}

Note that each $F_{j}(S)$ is itself a graded set, with degree $i$ component
$F_{j}(S)^i := F_{j}(S) \cap S^i$.

\begin{dfn} \label{dfn:4732}
Let $P$ be a DG $A$-module. A 
{\em semi-basis}%
\index{Semi-basis! of DG module}
of $P$ is a filtered graded set 
$(S, F)$, together with an isomorphism 
$P^{\natural}  \cong A^{\natural} \cd S$
of graded $A^{\natural}$-modules, such that under this isomorphism we have 
$\d_P(F_j(S)) \sub A^{\natural} \cd F_{j - 1}(S)$
for every $j \geq 0$. 
\end{dfn}

\begin{prop} \label{prop:4730}
Let $P$ be a DG $A$-module. 
\begin{enumerate}
\item If $P$ has a semi-basis $(S, F)$, then $P$ has an induced semi-free 
filtration $\{ F_j(P) \}_{j \geq -1}$ such that 
$F_j(P)^{\natural} = A^{\natural} \cd F_{j}(S)$
as graded $A^{\natural}$-modules for every $j \geq -1$.

\item Every semi-free filtration of $P$ is induced by some semi-basis. 
\end{enumerate}
\end{prop}

\begin{exer} \label{exer:4725}
Prove Proposition \ref{prop:4730}.
\end{exer}

\begin{exa} \label{exa:1575}
If $A$ is a ring, then a free DG $A$-module $P$ is a complex of free $A$-modules
with zero differential. A semi-free DG $A$-module $P$ is also a  complex of 
free $A$-modules, but there is a differential on it, and there is a subtle 
condition on $P$ imposed by the existence of a semi-free filtration. If the 
complex $P$ happens to be bounded above, then it is automatically semi-free, 
with a filtration like the one in the proof of Proposition \ref{prop:1900}. 
\end{exa}

\begin{exer} \label{exa:1577}
Find a ring $A$, and a complex $P$ of free $A$-modules, that is not semi-free. 
(Hint: Take the ring $A = \K[\ep]$ of dual numbers. Find a complex of free 
$A$-modules $P$ that is acyclic but not null-homotopic. Now use Theorem 
\ref{thm:1575} and Corollary \ref{cor:1923}  to deduce that $P$ is not 
semi-free.)
\end{exer}

\begin{prop} \label{prop:4636}
Let $A$ and $B$ be DG rings, and let $P \in \dcat{C}(A)$ and 
$Q \in \dcat{C}(B)$ be semi-free DG modules. Then 
$P \ot Q \in \dcat{C}(A \ot B)$
is a semi-free DG module.
\end{prop}

\begin{exer} \label{exer:4636}
Prove this proposition. (Hint: let $\bigl\{ F_j(P) \bigr\}_{j \geq -1}$
and \lb $\bigl\{ G_k(Q) \bigr\}_{k \geq -1}$ be semi-free filtrations. 
Define 
\[ E_{j}(P \ot Q) := \sum_{k = 0}^{j} \, 
\bigl( F_k(P) \ot G_{j - k}(Q) \bigr) \sub P \ot Q . \]
Show that this is a semi-free filtration.) 
\end{exer}

\begin{thm} \label{thm:1575}
\index{Differential graded module! K-projective}
\index{Differential graded module! semi-free}
Let $P$ be an object of $\dcat{C}(A)$. If $P$ is semi-free, then it is 
K-projective. 
\end{thm}

\begin{proof}
This is similar to the proof of Theorem \ref{thm:1536}. 

\smallskip \noindent
Step 1. We start by proving that if $P = \opn{T}^{-i}(A)$, a translation of 
$A$, 
then $P$ is K-projective. This is easy: 
given an acyclic $N \in \dcat{C}(A)$, we have 
\[ \opn{Hom}_{A}(P, N) = 
\opn{Hom}_{A} \bigl( \opn{T}^{-i}(A), N \bigr) \cong
\opn{T}^{i} \bigl( \opn{Hom}_{A}(A, N) \bigr)
\cong \opn{T}^{i}(N)\]
in $\dcat{C}_{\mrm{str}}(\K)$, and this is acyclic. 

\medskip \noindent
Step 2. Now $P$ is free, say 
$P \cong \bigoplus_{s \in S} \, \opn{T}^{-i_s}(A)$.
Then 
\[ \opn{Hom}_{A}(P, N) \cong 
\prod_{s \in S} \, \opn{Hom}_{A} \bigl( \opn{T}^{-i_s}(A), N \bigr) . \]
By step 1 and the fact that a product of acyclic complexes in 
$\dcat{C}_{\mrm{str}}(\K)$ is 
acyclic, we conclude that $\opn{Hom}_{\cat{M}}(P, N)$ is acyclic. 

\medskip \noindent
Step 3. Fix a semi-free filtration $F = \{ F_j(P) \}_{j \geq -1}$
of $P$. Here we prove that for every $j \geq -1$ the DG module  
$F_j(P)$ is K-projective. This is done by induction on $j \geq -1$. 
For $j = -1$ it is trivial. For $j \geq 0$ there is an exact sequence 
\begin{equation} \label{eqn:1580}
0 \to F_{j - 1}(P) \to F_j(P) \to \opn{Gr}^F_j(P) \to 0 
\end{equation}
in the abelian category $\dcat{C}_{\mrm{str}}(A)$. 
Because 
$\opn{Gr}^F_j(P)$
is a free DG module, it is a projective object in the abelian category 
$\bcat{G}_{\mrm{str}}(A^{\natural})$ of graded modules over the graded ring 
$A^{\natural}$. 
Therefore the sequence (\ref{eqn:1580}) is split exact in 
$\bcat{G}_{\mrm{str}}(A^{\natural})$.

Let $N \in \dcat{C}(A)$ be an acyclic DG module. Applying the 
functor $\opn{Hom}_{A}(-, N)$ 
to the sequence (\ref{eqn:1580}) we obtain a sequence 
\begin{equation} \label{eqn:1581}
0 \to \opn{Hom}_{A} \bigl( \opn{Gr}^F_j(P), N \bigr) \to 
\opn{Hom}_{A} \bigl( F_j(P), N \bigr) \to
\opn{Hom}_{A} \bigl( F_{j - 1}(P), N \bigr) \to  0 
\end{equation}
in $\dcat{C}_{\mrm{str}}(\K)$. If we forget differentials this is a sequence in 
$\bcat{G}_{\mrm{str}}(\K)$. Because (\ref{eqn:1580}) is split exact in 
$\bcat{G}_{\mrm{str}}(A^{\natural})$, it follows that  (\ref{eqn:1581}) is 
split 
exact in $\bcat{G}_{\mrm{str}}(\K)$. Therefore (\ref{eqn:1581}) is exact in 
$\dcat{C}_{\mrm{str}}(\K)$.

By the induction hypothesis the DG $\K$-module  
$\opn{Hom}_{A} \bigl( F_{j - 1}(P), N \bigr)$
is acyclic. By step 2 the DG module
$\opn{Hom}_{A} \bigl( \opn{Gr}^F_{j}(P), N \bigr)$
is acyclic. The long exact cohomology sequence associated to 
(\ref{eqn:1581}) shows that the DG module 
$\opn{Hom}_{A} \bigl( F_j(P), N \bigr)$
is acyclic too.

\medskip \noindent
Step 4. We keep the semi-free filtration 
$F = \{ F_j(P) \}_{j \geq -1}$ from step 3. 
Take any acyclic $N \in \bcat{C}(\cat{M})$. By 
Proposition \ref{prop:1535} we know that 
\[ \opn{Hom}_{A}(P, N) \cong 
\lim_{\leftarrow j} \, \opn{Hom}_{A} \bigl( F_j(P), N \bigr) \]
in $\dcat{C}_{\mrm{str}}(\K)$. According to step 3 the complexes 
$\opn{Hom}_{A} \bigl( F_j(P), N \bigr)$ are all acyclic. 
The exactness of the sequences (\ref{eqn:1581}) implies that 
the inverse system \lb 
$\bigl\{ \opn{Hom}_{A} \bigl( F_j(P), N \bigr) \bigr\}_{j \geq -1}$
in $\dcat{C}_{\mrm{str}}(\K)$ has surjective transitions. Now the 
Mittag-Leffler argument (Corollary \ref{cor:1590}) says that the inverse limit 
complex $\opn{Hom}_{A}(P, N)$ is acyclic. 
\end{proof}

Here is a result similar to Theorem \ref{thm:1540}. 

\begin{thm} \label{thm:3305}
\index{Differential graded module! semi-free}
Let $A$ be a DG ring. Every $M \in \bcat{C}(A)$ admits a quasi-iso\-morphism
$\rho : P \to M$ in $\dcat{C}_{\mrm{str}}(A)$ from a semi-free DG $A$-module 
$P$.
\end{thm}

The proof of this theorem was communicated to us by B. Keller.
First we need a construction and a lemma. 

Let us define the DG $A$-module 
\begin{equation} \label{eqn:3307}
C := \opn{Cone} \bigl( \opn{id} : \opn{T}^{-1}(A) \to \opn{T}^{-1}(A) \bigr) , 
\end{equation}
the standard cone of the identity automorphism of $\opn{T}^{-1}(A)$
in $\dcat{C}_{\mrm{str}}(A)$. Since
\[ C = \opn{T}^{-1}(A) \oplus \opn{T}^{}(\opn{T}^{-1}(A)) = 
\opn{T}^{-1}(A) \oplus A \]
as graded objects, we have the elements 
$e_0 := 1_A \in A \sub C^0$ and 
$e_1 := \opn{t}^{-1}(1_A) \in \opn{T}^{-1}(A)^1 \sub C^1$.
They satisfy $\d_C(e_0) = e_1$. 
Note that the DG module $C$ is semi-free, with semi-free filtration 
\begin{equation} \label{eqn:3332}
F_j(C) := 
\begin{cases}
0 & \tup{if} \ \ j = -1
\\
\opn{T}^{-1}(A) &  \tup{if} \ \ j = 0 
\\
C &  \tup{if} \ \ j \geq 1 .
\end{cases}
\end{equation}
As a graded $A^{\natural}$-module, $C^{\natural}$ is free, with basis 
$(e_0, e_1)$. Of course the DG module $C$ is contractible. 

\begin{lem} \label{lem:3306}
Let $M \in \dcat{C}(A)$.
\begin{enumerate}
\item There is a homomorphism $\psi : Q \to M$ 
in $\dcat{C}_{\mrm{str}}(A)$, such that 
$\opn{Z}(\psi) : \opn{Z}(Q) \to \opn{Z}(M)$ is surjective,
and $Q = \bigoplus_{s \in S} \opn{T}^{-i_s}(A)$ 
for a collection of integers $\{ i_s \}_{s \in S}$. 

\item There is a surjective homomorphism
$\psi' : Q' \to M$ in $\dcat{C}_{\mrm{str}}(A)$,
such that 
$Q' := \bigoplus_{s \in S'} \opn{T}^{-i_s}(C)$
for some collection of integers $\{ i_s \}_{s \in S'}$,
where $C$ is the DG $A$-module from  formula \tup{(\ref{eqn:3307})}.
\end{enumerate}
\end{lem}

To clarify the notation in the lemma, the collections of integers  
$\{ i_s \}_{s \in S}$ and $\{ i_s \}_{s \in S'}$ are distinct; namely the
function $i : s \mapsto i_s$ is $i : S \sqcup S' \to \Z$. 

\begin{proof} \mbox{}

\smallskip \noindent
(1) For every cocycle $m \in \opn{Z}^i(M)$
there a homomorphism $\psi_m : \opn{T}^{-i}(A) \to M$ 
that sends the element 
$\opn{t}^{-i}(1_A) \in \opn{Z}^i(\opn{T}^{-i}(A))$
to $m$. Thus, if $\{ m_s \}_{s \in S}$ is a collection of 
homogeneous cocycles that generates $\opn{Z}(M)$ as a $\K$-module, we get a 
homomorphism $\psi$ as claimed. 

\medskip \noindent
(2) For every element $m \in M^i$ there is a homomorphism 
$\psi'_m : \opn{T}^{-i}(C) \to M$ 
in $\dcat{C}_{\mrm{str}}(A)$ that sends 
$\opn{t}^{-i}(e_0) \mapsto m$ and 
$\opn{t}^{-i}(e_1) \mapsto (-1)^i \cd \d_M(m)$. 
Hence, by taking a collection $\{ m_s \}_{s \in S'}$ of homogeneous elements 
that  generates $M$ as a $\K$-module, we get a homomorphism $\psi'$ as claimed. 
\end{proof}

The DG $A$-module $Q$ in item (1) of the lemma is free. The DG $A$-module 
$Q'$ in item (2) is semi-free, with semi-free filtration induced by that of 
$C$, namely 
\begin{equation} \label{eqn:3311}
F_{j}(Q') := \bigoplus_{s \in S'} \, \opn{T}^{-i_s}(F_j(C))
\end{equation}
for $j \geq -1$, where $F_j(C)$ is from (\ref{eqn:3332}). 

\begin{proof}[Proof of Theorem \tup{\ref{thm:3305}}]
The proof is in several steps. 

\smallskip \noindent
Step 1. We are going to produce an exact sequence 
\begin{equation} \label{eqn:3350}
\cdots \to Q^{-2} \xar{\pa^{-2}} Q^{-1} \xar{\pa^{-1}}
Q^0 \xar{\eta} M \to 0
\end{equation}
in $\dcat{C}_{\mrm{str}}(A)$ with these properties: 
\begin{itemize}
\rmitem{a} The sequence 
\[ \cdots \to \opn{Z}(Q^{-2}) \xar{\opn{Z}(\pa^{-2})} \opn{Z}(Q^{-1}) 
\xar{\opn{Z}(\pa^{-1})} \opn{Z}(Q^0) \xar{\opn{Z}(\eta)} \opn{Z}(M) \to 0 \]
in $\dcat{G}_{\mrm{str}}(\K)$ is exact.

\rmitem{b} For every $i \leq 0$ the DG $A$-module $Q^i$ has the following 
decomposition 
in $\dcat{C}_{\mrm{str}}(A)$~: 
$Q^i = Q_i \oplus Q'_i$, where
$Q_i = \bigoplus_{s \in S_i} \, \opn{T}^{-j_s}(A)$
and 
$Q'_i = \lb \bigoplus_{s \in S'_i} \, \opn{T}^{-j_s}(C)$
for some collection of integers 
$\{ j_s \}_{s \in S_i \sqcup S'_i}$. 
\end{itemize}
This will be done recursively on $i \leq 0$. 

By Lemma \ref{lem:3306} we can find DG $A$-modules $Q_0, Q'_0$ 
and homomorphisms 
$\psi_0 : Q_0 \to M$, $\psi'_0 : Q'_0 \to M$, such that both 
$\opn{Z}(\psi_0)$ and $\psi'_0$ are surjective. 
The DG modules $Q_0, Q'_0$ are of the required 
form (property (b) with $i = 0$). We let 
$Q^0 := Q_0 \oplus Q'_0$ and $\eta := \psi_0 \oplus \psi'_0$. 

Let $N^0 := \opn{Ker}(\eta)$. Then the sequence 
\begin{equation} \label{eqn:3352}
0 \to N^0 \to Q^0 \xar{\eta} M \to 0 
\end{equation}
in $\dcat{C}_{\mrm{str}}(A)$ is exact. The sequence 
\[ 0 \to \opn{Z}(N^0) \to \opn{Z}(Q^0) \xar{\opn{Z}(\eta)} \opn{Z}(M) \to 0 \]
in $\dcat{G}_{\mrm{str}}(\K)$ is also exact. Indeed, the exactness at 
$\opn{Z}(N^0)$ and $\opn{Z}(Q^0)$ is because the functor 
$\opn{Z} : \dcat{C}_{\mrm{str}}(\K) \to \dcat{G}_{\mrm{str}}(\K)$
is left exact (see Proposition \ref{prop:4145}), and the sequence 
(\ref{eqn:3352}) is exact. 
The exactness at $\opn{Z}(M)$ is by the surjectivity of 
$\opn{Z}(\psi_0) : \opn{Z}(Q_0) \to \opn{Z}(M)$.

Next we repeat this procedure with $N^0$ instead of $M$, to obtain an exact 
sequence 
\begin{equation} \label{eqn:3351}
0 \to N^{-1} \to Q^{-1} \xar{\pa^{-1}} N^0 \to 0 ,
\end{equation}
with $Q^{-1}$ of the required form (property (b) with $i = -1$), and such that 
\[ 0 \to \opn{Z}(N^{-1}) \to \opn{Z}(Q^{-1}) \xar{\opn{Z}(\pa^{-1})} 
\opn{Z}(N^0) \to 0 \]
is also exact. We then splice (\ref{eqn:3352}) and (\ref{eqn:3351})
to get the sequence 
\[ 0 \to N^{-1} \to Q^{-1} \xar{\pa^{-1}} Q^0 \xar{\eta} M . \]

Continuing recursively we obtain an exact sequence (\ref{eqn:3350}) that has 
properties (a) and (b). 

\medskip \noindent 
Step 2. We view the exact sequence (\ref{eqn:3350}) as an acyclic complex with 
entries in $\dcat{C}_{\mrm{str}}(A)$, that has $Q^0$ in degree $0$. 
Define the total DG $A$-modules 
\[ P :=\opn{Tot}^{\oplus} \bigl( 
\cdots \to Q^{-2} \xar{\pa^{-2}} Q^{-1} \xar{\pa^{-1}} Q^0 \to 0 \to \cdots
\bigr) \]
and 
\[ P^{\mrm{aug}} := \opn{Tot}^{\oplus} \bigl(
\cdots \to Q^{-2} \xar{\pa^{-2}} Q^{-1} \xar{\pa^{-1}}
Q^0 \xar{\eta} M \to 0 \to \cdots \bigr) . \]

Because the sequence (\ref{eqn:3350}) and the sequence appearing in property 
(a) are exact, we can use Proposition \ref{prop:3300} (with any 
$j_1 \in \Z$). The conclusion is that the DG $A$-module $P^{\mrm{aug}}$ is 
acyclic. Then, by Corollary \ref{cor:3310}, the homomorphism 
$\rho : P \to M$ in $\dcat{C}_{\mrm{str}}(A)$
that is induced by $\eta$ is a quasi-isomorphism. 

\medskip \noindent
Step 3. It remains to produce a semi-free filtration 
$\{ F_j(P) \}_{j \geq -1}$ on the DG $A$-module $P$. Our formula is this: 
$F_{-1}(P) := 0$ of course. For $k \geq 0$ we let 
\[ F_{2 \cd k}(P) := 
\Bigr( \bigoplus\nolimits_{l = 0}^{-(k - 1)} \, \opn{T}^{-l} (Q^l) \Bigl) 
\, \oplus \,   
\opn{T}^{k} \bigl( Q_{-k} \oplus F_0(Q'_{-k}) \bigr) , \]
where $F_0(Q'_{-k})$ comes from (\ref{eqn:3311}), and 
\[ F_{2 {\cdot} k + 1}(P)  := 
\bigoplus\nolimits_{l = 0}^{-k} \, \opn{T}^{-l} (Q^l) . \]
We leave it to the reader to verify that this is a semi-free filtration. 
\end{proof}

\begin{exer} \label{exer:4752}
Finish step 3 of the proof. 
\end{exer}

\begin{cor} \label{cor:1580}
Let $A$ be a DG ring. The category $\bcat{C}(A)$ 
has enough K-pro\-jectives%
\index{Resolution! K-projective}. 
\end{cor}

\begin{proof}
Combine Theorems \ref{thm:1575} and \ref{thm:3305}. 
\end{proof}

Recall that a DG ring $A$ is nonpositive if $A^i = 0$ for all $i > 0$.
According to Proposition \ref{prop:3170}, when $A$ is nonpositive the smart 
truncation functors exist. 

\begin{cor} \label{cor:3340}
If the DG ring $A$ is nonpositive, then every $M \in \bcat{C}(A)$ admits a 
quasi-iso\-morphism $\rho : P \to M$ in $\dcat{C}_{\mrm{str}}(A)$ from a 
semi-free DG $A$-module $P$, such that 
$\opn{sup}(P) = \opn{sup}(\opn{H}(M))$. 
\end{cor}

\begin{proof}
Let $i_1 := \opn{sup}(\opn{H}(M))$. If $i_1 = \pm \infty$ then there is nothing 
to prove beyond Theorem \ref{thm:3305}; so let us assume that $i_1 \in \Z$. 
Define $M' := \opn{smt}^{\leq i_1}(M)$. Then $M' \to M$ is a quasi-isomorphism, 
and $\opn{sup}(M') = i_1$. By replacing $M$ with $M'$ we can now assume that 
$\opn{sup}(M) = i_1$.

Since $\opn{sup}(M) \leq i_1$, in step 1 of the proof of the theorem we can 
take the DG module $Q_0$ to be concentrated in degrees $\leq i_1$.
Another consequence of the fact that $\opn{sup}(M) \leq i_1$ is that 
$M^{i_1} = \opn{Z}^{i_1}(M)$. This means that the homomorphism 
$\psi_0 : Q^{i_1}_0 \to M^{i_1}$ is already surjective; so we don't need to 
cover $M^{i_1}$ by $Q'_0$, and thus we can take $Q'_0$ to be concentrated in 
degrees $\leq i_1$. In this way we can arrange to have $\opn{sup}(Q^0) \leq 
i_1$. But then we also have $\opn{sup}(N^0) \leq i_1$. Thus 
recursively we can arrange to have $\opn{sup}(Q^l) \leq i_1$ for all $l$. 
Therefore in step 2 of the proof we get $\opn{sup}(P) \leq i_1$. 

Finally, because $\opn{sup}(\opn{H}(M)) = i_1$ we must have 
$\opn{sup}(P) = i_1$.
\end{proof}

\begin{dfn} \label{dfn:3355}
Let $A$ be a DG ring. A DG $A$-module $P$ is called a {\em 
pseudo-finite free DG $A$-module} if there are numbers $i_1 \in \Z$ and 
$r_i \in \N$ such that 
$ P \cong \bigoplus_{i \leq i_1} \opn{T}^{-i}(A)^{\oplus r_i}$
in $\dcat{C}_{\mrm{str}}(A)$. 
\end{dfn}

\begin{dfn} \label{dfn:1925}
Let $A$ be a DG ring and let $P$ be a DG $A$-module. 
\begin{enumerate}
\item A {\em pseudo-finite semi-free filtration}% 
\index{Filtration on a DG module! pseudo-finite semi-free}
on $P$ is a semi-free 
filtration $F = \lb \bigl\{ F_j(P) \bigr\}_{j \geq -1}$
satisfying this finiteness condition: there are numbers 
$i_1 \in \Z$ and $r_j \in \N$ such that 
$\opn{Gr}^F_j(P) \cong \opn{T}^{-i_1 + j}(A)^{\oplus r_j}$
in $\dcat{C}_{\mrm{str}}(A)$ for all $j \in \N$. 

\item We call $P$ a {\em pseudo-finite semi-free DG $A$-module}%
\index{Differential graded module! pseudo-finite semi-free}
if it admits some pseudo-finite semi-free filtration.
\end{enumerate}
\end{dfn}

Note that in item (1) of the definition above, for every $j \geq -1$ we have
\begin{equation} \label{eqn:4736}
F_j(P)^{\natural} \cong 
\bigoplus_{0 \leq k \leq j}  \opn{Gr}^F_j(P)^{\natural}
\cong \bigoplus_{0 \leq k \leq j}  
\opn{T}^{-i_1 + k}(A^{\natural})^{\oplus r_k}
\end{equation}
as graded $A^{\natural}$-modules. 

\begin{prop} \label{prop:4731}
Let $P$ be a DG $A$-module. The two conditions below are equivalent. 
\begin{enumerate}
\rmitem{i} $P$ is a pseudo-finite semi-free DG $A$-module.

\rmitem{ii} $P$ admits a semi-basis $(S, F)$ such that $S$ is bounded above, 
say by $i_1 \in \Z$\tup{;}
each of the sets $S^i$ is finite\tup{;} and 
$F_j(S) = \bigcup_{i \in [i_1 - j, i_1]}$
for every $j \geq -1$. 
\end{enumerate}
\end{prop}

\begin{exer} \label{exer:4731}
Prove Proposition \ref{prop:4731}. (Hint: the numbers $i_1$ in Definition 
\ref{dfn:1925}(1) and in condition (ii) are the same.)
\end{exer}
 
A graded $A^{\natural}$-module 
$P^{\natural}$ is called {\em pseudo-finite free} if there is an 
isomorphism of graded $A^{\natural}$-modules
$P^{\natural} \cong \bigoplus_{j \geq 0} 
\opn{T}^{-i_1 + j}(A^{\natural})^{\oplus r_j}$
for some numbers $i_1 \in \Z$ and $r_j \in \N$. 

\begin{prop} \label{prop:3175}
Let $A$ be a DG ring and let $P$ be a DG $A$-module. The following 
two conditions are equivalent\tup{:}
\begin{enumerate}
\rmitem{i} The DG $A$-module $P$ is pseudo-finite semi-free.

\rmitem{ii} The DG $A$-module $P$ admits a semi-free filtration 
$F = \{ F_j(P) \}_{j \geq -1}$ such that 
each $\opn{Gr}^F_j(P)$ is a pseudo-finite free DG $A$-module, and this 
asymptotic formula holds\tup{:} 
$\lim_{j \to \infty} \opn{sup} (\opn{Gr}^F_j(P)) = - \infty$.
\end{enumerate}
If $A$ is nonpositive, the conditions above are equivalent to\tup{:} 
\begin{enumerate}
\rmitem{iii} The graded $A^{\natural}$-module $P^{\natural}$ is pseudo-finite 
free.
\end{enumerate}
If $A$ is a ring \tup{(}i.e.\ $A^i = 0$ for all $i \neq 0$\tup{)}, the 
conditions above are equivalent to\tup{:} 
\begin{enumerate}
\rmitem{iv} $P$ is a bounded above complex of finitely generated 
free $A$-modules. 
\end{enumerate}
\end{prop}

\begin{exer} \label{exer:3175}
Prove the proposition above. (Hint: rephrase these conditions in terms of  
semi-bases.)
\end{exer}

\begin{rem} \label{rem:4635}
Condition (iv) of Proposition \ref{prop:3175} is a hint for the reason we chose 
the name ``pseudo-finite'' in Definition \ref{dfn:1925}, and also the name 
``pseudo-noetherian'' in Definitions \ref{dfn:4755} and \ref{dfn:3180} below 
(in conjunction with Theorem \ref{thm:3340}). Indeed, in \cite{SGA6} and 
\cite{SP}, a complex of modules over a ring $A$ is called {\em pseudo-coherent} 
if it is isomorphic, in $\dcat{D}(A)$, to a bounded above complex of finite 
free 
$A$-modules. 

For a generalization of the notion of pseudo-coherent complex to DG modules, 
see 
Definition \ref{dfn:4470}.
\end{rem}

\begin{dfn} \label{dfn:4755} 
A graded ring $A$ is called {\em left pseudo-noetherian} if it is
nonpositive, the ring $A^0$ is left noetherian, and each $A^i$ is a finitely 
generated (left) $A^0$-module. 
\end{dfn}

\begin{dfn} \label{dfn:3180} 
Let $A$ be a DG ring. 
\begin{enumerate}
\item We call $A$ a 
{\em cohomologically left pseudo-noetherian DG ring}%
\index{Differential graded ring! cohomologically pseudo-noetherian}
if its cohomology $\opn{H}(A)$ is a left pseudo-noetherian graded ring.

\item Suppose $A$ is cohomologically left pseudo-noetherian. We 
denote by $\dcat{D}_{\mrm{f}}(A)$ the full subcategory of 
$\dcat{D}(A)$ on the DG modules $M$ such that $\opn{H}^i(M)$ is a finitely 
generated $\opn{H}^0(A)$-module for every $i$. 
\end{enumerate}
\end{dfn}

Of course $\dcat{D}_{\mrm{f}}(A)$ is a full triangulated subcategory of 
$\dcat{D}(A)$. As usual we combine indicators: 
$\dcat{D}_{\mrm{f}}^{\star}(A) := \dcat{D}_{\mrm{f}}(A) \cap 
\dcat{D}^{\star}(A)$.

\begin{exa} \label{exa:3182}
Suppose $A^0$ is a nonzero noetherian commutative ring.
Let $X = \coprod_{i \leq 0} X^i$ be a 
nonpositive graded set, such that each graded piece $X^i$ is finite. 
Consider the strongly commutative graded polynomial ring
$A^0[X] := A^0 \ot \K [X]$
from Example \ref{exa:4625}.
The graded ring $A^0[X]$ is (left and right) pseudo-noetherian, but if $X$ is 
infinite then $A^0[X]$ is not noetherian (on either side).  

Likewise, let $Y = \coprod_{i < 0} Y^i$ be a negative graded set, such that 
each graded piece $Y^i$ is finite. Take the noncommutative graded 
polynomial ring
$A^0 \bra{Y} := A^0 \ot \K \bra{Y}$
from Example \ref{exa:4625} and Definition \ref{dfn:4306}. 
The graded ring $A^0\bra{Y}$ is (left and right) pseudo-noetherian, but if $Y$ 
has at least two elements, then $A^0\bra{Y}$ is not noetherian (on either side).
\end{exa}

\begin{rem} \label{rem:3205}
Let $A$ be a DG ring whose cohomology $\opn{H}(A)$ is a nonpositive graded 
ring. Define $A' := \opn{smt}^{\leq 0}(A)$, the smart truncation of 
$A$ below $0$. It is easy to see that $A'$ is 
a DG subring of $A$, and the inclusion $A' \to A$ is a DG ring 
quasi-isomorphism. By Theorem \ref{thm:2363} there is an equivalence of 
triangulated categories $\dcat{D}(A) \to \dcat{D}(A')$. 
This says that for many purposes we can assume that $A$ itself is a 
nonpositive DG ring. 
\end{rem}

\begin{thm} \label{thm:3340}
\index{Differential graded ring! cohomologically pseudo-noetherian}
\index{Differential graded module! pseudo-finite semi-free}
Let $A$ be a nonpositive cohomologically left pseudo-noether\-ian DG ring, and 
let $M \in \dcat{D}_{\mrm{f}}^{-}(A)$. 
Then there is a quasi-isomorphism 
$P \to M$ in $\dcat{C}_{\mrm{str}}(A)$, where $P$ is a pseudo-finite 
semi-free DG $A$-module, and $\opn{sup}(P) = \opn{sup}(\opn{H}(M))$. 
\end{thm}

\begin{proof}
The proof is divided into five parts. 

\smallskip \noindent 
Part 1. We may assume that $\opn{H}(M) \neq 0$, so that 
$i_1 := \opn{sup}(\opn{H}(M)) \in \Z$. 
We are going to construct a direct system 
\[ 0 = F_{-1}(P) \sub F_0(P) \sub F_1(P) \sub \cdots \]
of pseudo-finite semi-free DG $A$-modules, together with a direct system of DG 
$A$-module homomorphisms 
$F_j(\rho) : F_{j}(P) \to M$
for $j \geq 0$, having these properties:
\begin{enumerate}
\rmitem{a} Let $j \geq 0$. The homomorphism 
$\opn{H}^i(F_j(\rho)) : \opn{H}^i(F_{j}(P)) \to \opn{H}^i(M)$ 
is surjective for all $i$, and bijective for all $i > i_1 - j$. 

\rmitem{b} $F_0(P)$ is a pseudo-finite free DG $A$-module with 
$\opn{sup}(F_0(P)) = i_1$.

\rmitem{c} Let $j \geq 1$. There is an isomorphism 
$F_{j}(P) / F_{j - 1}(P) \cong \opn{T}^{-i_1 + j}(A)^{\oplus r_j}$
in $\dcat{C}_{\mrm{str}}(A)$, for some $r_j \in \N$.
\end{enumerate}
This will be done recursively on $j \geq 0$. . 

\medskip \noindent
Part 2. Here we deal with the cases $j = -1, 0$. 
For $j = -1$ we take $F_{-1}(P) := 0$ of course. 

Now we deal with $F_0(P)$. 
For every $i \leq i_1$ we choose finitely many generators of the 
$\opn{H}^0(A)$-module $\opn{H}^i(M)$, and then we lift them to elements of
$\opn{Z}^i(M)$. These choices give rise to a pseudo-finite free DG $A$-module 
$F_0(P)$ and a homomorphism 
$F_0(\rho) : F_0(P) \to M$ in 
$\dcat{C}_{\mrm{str}}(A)$, such that 
$\opn{sup}(F_0(P)) = i_1$, and the homomorphism 
$\opn{H}(F_0(\rho)) : \opn{H}(F_0(P)) \to \opn{H}(M)$ 
in $\dcat{G}_{\mrm{str}}(\K)$ is surjective.  
Condition (a) holds for $j = 0$, and so does condition (b). Condition (c) is 
not applicable here. 

\medskip \noindent
Part 3. 
Suppose that $j \geq 0$, and we have a direct system 
\[ 0 = F_{-1}(P) \sub F_0(P) \sub F_1(P) \sub \cdots \sub F_j(P) \]
of semi-free DG $A$-modules, together with a direct system of DG $A$-module 
homomorphisms
$F_{j'}(\rho) : F_{j'}(P) \to M$ 
for $j' \in [0, j]$, satisfying conditions (a)-(c) with index $j'$ 
instead of $j$. In this part we will construct the DG $A$-module 
$F_{j + 1}(P)$. 

Let 
\begin{equation} \label{eqn:5065}
\bar{N} := \opn{Ker} \bigl( \opn{H}^{i_1 - j}(F_j(\rho)) : 
\opn{H}^{i_1 - j}(F_{j}(P)) \to \opn{H}^{i_1 - j}(M) \bigr) .
\end{equation}
Because $F_{j}(P)$ is a pseudo-finite semi-free DG $A$-module, and $A$ is a 
cohomologically left pseudo-noetherian DG ring, it follows that $\bar{N}$ is a 
finitely generated $\opn{H}^0(A)$-module.
We choose elements 
$\bar{n}_1, \ldots, \bar{n}_{r_{j + 1}} \in \bar{N}$
that generate $\bar{N}$ as an $\opn{H}^0(A)$-module.
Then we lift them to elements
$n_1, \ldots, n_{r_{j + 1}} \in \opn{Z}^{i_1 - j}(F_{j}(P))$.

Consider the free DG $A$-module 
$Q' := \opn{T}^{-i_1 + j}(A)^{\oplus r_{j + 1}}$
with basis $(e'_1, \ldots, \lb e'_{r_{j + 1}})$ concentrated in degree 
$i_1 - j$. There is a homomorphism 
$\phi : Q' \to F_{j}(P)$ in $\dcat{C}_{\mrm{str}}(A)$
with formula $\phi(e'_k) := n_k$. Define 
\begin{equation} \label{eqn:4755}
F_{j + 1}(P) := \opn{Cone} \bigl( \phi : Q' \to F_{j}(P) \bigr) , 
\end{equation}
the standard cone of $\phi$. 
The DG $A$-module $F_{j + 1}(P)$ is pseudo-finite semi-free.
Condition (b) still holds, because $F_{0}(P)$ did not change. And 
$F_{j + 1}(P) / F_{j}(P) \cong Q$, where 
\begin{equation} \label{eqn:4760}
Q := \opn{T}(Q') \cong \opn{T}^{-i_1 + j + 1}(A)^{\oplus r_{j + 1}} , 
\end{equation}
so condition (c) holds also for $j + 1$. 

\medskip \noindent
Part 4. We continue part 3. Here we construct the 
homomorphism 
$F_{j + 1}(\rho) : F_{j + 1}(P) \to M$
in $\dcat{C}_{\mrm{str}}(A)$. 
For every $k \in [1, r_{j + 1}]$ let 
$e_k := \opn{t}_{Q'}(e'_k) \in Q^{i_1 - j - 1}$.
The sequence $(e_1, \ldots, e_{r_{j + 1}})$ is a basis of the free DG 
$A$-module $Q$. 
We have $\d_Q(e_k) = 0$ and $\opn{t}^{-1}_{Q}(e_k) = e'_k \in Q'$, 
so by the definition of the differential of the cone we have 
\[ \d_{F_{j + 1}(P)}(e_k) = \d_Q(e_k) + \phi(\opn{t}^{-1}_{Q}(e_k)) = 
\phi(e'_k) = n_k . \]
Now the cohomology class $\bar{n}_k$ of $n_k$ belongs to $\bar{N}$,
so 
$\opn{H}^{i_1 - j}(F_j(\rho))(\bar{n}_k) = 0$.
This means that the element 
$F_j(\rho)(n_k)$ is a coboundary in $M$, and we can lift it to an element 
$m_k \in M^{i_1 - j - 1}$.

By definition of the cone we have
$F_{j + 1}(P)^{\natural} = F_{j}(P)^{\natural} \oplus Q^{\natural}$ 
as graded $A^{\natural}$-modules.
We define the homomorphism 
$F_{j + 1}(\rho) : F_{j + 1}(P)^{\natural} \to M^{\natural}$ 
in $\dcat{G}_{\mrm{str}}(A^{\natural})$
to be the extension of $F_{j}(\rho)$ such that 
$F_{j + 1}(\rho)(e_k) := m_k$. Then 
$F_{j + 1}(\rho) : F_{j + 1}(P) \to M$ 
is actually a homomorphism in $\dcat{C}_{\mrm{str}}(A)$. 
         
There is equality $F_{j + 1}(P)^i = F_{j}(P)^i$
in all degrees $i \geq i_1 - j$.
Therefore  
$\opn{H}^i(F_{j + 1}(\rho)) = \opn{H}^i(F_j(\rho))$
for all $i > i_1 - j$, and these homomorphisms are bijective by assumption. 
The homomorphism 
$\opn{Z}^{i_1 - j}(F_{j}(P)) \to \opn{Z}^{i_1 - j}(F_{j + 1}(P))$
is surjective, and hence 
$\opn{H}^{i_1 - j}(F_{j}(P)) \to \opn{H}^{i_1 - j}(F_{j + 1}(P))$
is also surjective. Because 
the submodule $\bar{N} \sub \opn{H}^{i_1 - j}(F_j(P))$
goes to zero in $\opn{H}^{i_1 - j}(F_{j + 1}(P))$,
we have the following commutative diagram:
%\begin{equation} \label{eqn:5066}
\[ \UseTips \xymatrix @C=3ex @R=6ex {
0
\ar[r]
&
\bar{N}
\ar[r]
\ar[d]
&
\opn{H}^{i_1 - j}(F_{j}(P))
\ar[rrr]^(0.53){\opn{H}^{i_1 - j}(F_{j}(\rho))}
\ar[d]^{\mrm{surj}}
& & &
\opn{H}^{i_1 - j}(M)
\ar[r]
\ar[d]^{\opn{id}}
&
0
\\
0
\ar[r]
&
0
\ar[r]
&
\opn{H}^{i_1 - j}(F_{j + 1}(P))
\ar[rrr]^(0.53){\opn{H}^{i_1 - j}(F_{j + 1}(\rho))}
& & &
\opn{H}^{i_1 - j}(M)
\ar[r]
&
0
} \]
% \end{equation}
The top row is exact, by (\ref{eqn:5065}). Therefore the bottom row is exact, 
i.e.\ \lb $\opn{H}^{i_1 - j}(F_{j + 1}(\rho))$ is bijective.
We conclude that condition (a) holds with index $j + 1$. 

\medskip \noindent
Part 5. Define 
$P := \lim_{j \to} F_{j}(P)$ and 
$\rho := \lim_{j \to} F_{j}(\rho)$
in $\dcat{C}_{\mrm{str}}(A)$. 
Condition (c) implies that for every $i$ the direct system $\{ F_j(P^i) \}_{j 
\geq -1}$ is eventually stationary. Hence, by condition (a), the homomorphism 
$\opn{H}^i(\rho)$ is bijective. So $\rho$ is a quasi-isomorphism. 

By condition (b) and (c) the filtration $\{ F_j(P) \}_{j \geq -1}$
on $P$ satisfies condition (ii) of Proposition \ref{prop:3175}. Therefore $P$ 
is a pseudo-finite semi-free DG $A$-module.
\end{proof}

\begin{exa} \label{exa:1930} 
A special yet very important case of Theorem \ref{thm:3340} is this: 
$A$ is a left noetherian ring, and $M$ is a complex of $A$-modules 
with bounded above cohomology, such that each $\opn{H}^i(M)$ is a finitely 
generated $A$-module. Then $M$ has a resolution 
$P \to M$, where $P$ is a complex of finitely generated free $A$-modules, and 
$\opn{sup}(P) = \opn{sup}(\opn{H}(M))$. 
Compare this to Example \ref{exa:2900}.
\end{exa}

\begin{rem} \label{rem:4675}
In Example \ref{exa:4725} we discussed the existence of K-flat resolutions of 
unbounded complexes of sheaves of modules. Here is a more general assertion. 

Consider a {\em DG ringed space} $(X, \AA)$, namely $\AA$ is a 
sheaf of DG rings on a topological space $X$. There are no commutativity or 
boundedness conditions on $\AA$. By combining the concept of a 
semi-free DG module from Definition \ref{dfn:1575} with the concept of a 
pseudo-free module from Example \ref{exa:4725}, we get something new: a {\em 
semi-pseudo-free DG $\AA$-module}. This is a K-flat DG $\AA$-module. The 
arguments in \cite{Ye15} imply that every DG $\AA$-module $\MM$ admits a 
quasi-isomorphism $\PP \to \MM$ from a semi-pseudo-free DG $\AA$-module $\PP$.
Therefore $\dcat{C}(\AA)$ has enough K-flat objects.  
A detailed proof of this assertion will be published in the future. 
\end{rem}




\mysubsection{K-Injective Resolutions in 
\texorpdfstring{$\bcat{C}^+(\cat{M})$}{C+(M)}} 
\label{subsec:exis-K-inj}

In this subsection $\cat{M}$ is an abelian category, and $\dcat{C}(\cat{M})$ 
is the category of complexes in $\cat{M}$. Cofiltrations and their limits were 
introduced in Definition \ref{dfn:3336}.

\begin{dfn} \label{dfn:1620}
Let $I$ be a complex in $\dcat{C}(\cat{M})$. 
\begin{enumerate}
\item A {\em semi-injective cofiltration}% 
\index{Cofiltration on a DG module! semi-injective}
on $I$ is a cofiltration $G = \{ G_q(I) \}_{q \geq -1}$
in $\dcat{C}_{\mrm{str}}(\cat{M})$ such that:
\begin{itemize}
\item $G_{-1}(I) = 0$. 

\item  Each $\opn{Gr}_q^G(I)$ 
is a complex of injective objects of $\cat{M}$ with zero differential.

\item $I = \lim_{\lar q} \, G_q(I)$.
\end{itemize}

\item The complex $I$ is called a {\em semi-injective complex}% 
\index{Complex in abelian category! semi-injective}
if it admits some semi-injective cofiltration. 
\end{enumerate}
\end{dfn}


\begin{thm} \label{thm:1620}
\index{Complex in abelian category! semi-injective}
\index{Complex in abelian category! K-injective}
Let $\cat{M}$ be an abelian category, and let $I$ be a 
semi-injective complex in $\dcat{C}(\cat{M})$. Then $I$ is 
K-injective. 
\end{thm}


\begin{proof}
The proof is very similar to that of Theorem \ref{thm:1536}. 
 
\medskip \noindent
Step 1. We start by proving that if $I = \opn{T}^p(J)$, the translation of an 
injective object $J \in \cat{M}$, then $I$ is K-injective. This is easy: 
given an acyclic complex $N \in \dcat{C}(\cat{M})$, we have 
\[ \opn{Hom}_{\cat{M}}(N, I) = 
\opn{Hom}_{\cat{M}} \bigl( N, \opn{T}^p(J) \bigr) \cong
\opn{T}^{p} \bigl( \opn{Hom}_{\cat{M}}(N, J) \bigr) \]
in $\dcat{C}_{\mrm{str}}(\K)$. But 
$\opn{Hom}_{\cat{M}}(-, J)$ is an exact functor 
$\cat{M} \to \dcat{M}(\K)$, so $\opn{Hom}_{\cat{M}}(N, J)$ is an acyclic 
complex. 

\medskip \noindent
Step 2. Now $I$ is a complex of injective objects of $\cat{M}$ with zero 
differential. This means that 
$I \cong \prod_{p \in \Z} \opn{T}^p(J_p)$
in $\dcat{C}_{\mrm{str}}(\cat{M})$, where each $J_p$ is an injective object in 
$\cat{M}$. But then 
\[ \opn{Hom}_{\cat{M}}(N, I) \cong 
\prod_{p \in \Z} \, \opn{Hom}_{\cat{M}} \bigl( N, \opn{T}^p(J_p) \bigr) . \]
This is an easy case of Proposition \ref{prop:1535}(2). By step 1 and the fact 
that a product of acyclic complexes in $\dcat{C}_{\mrm{str}}(\K)$ is acyclic 
(itself an easy case of the Mittag-Leffler argument), we conclude that 
$\opn{Hom}_{\cat{M}}(N, I)$ is acyclic. 

\medskip \noindent
Step 3. Fix a semi-injective cofiltration $G = \{ G_q(I) \}_{q \geq -1}$
of $I$. Here we prove that for every $q$ the complex  
$G_q(I)$ is K-injective. This is done by induction on $q$. 
For $q = -1$ it is trivial. For $q \geq 0$ there is an exact sequence of 
complexes 
\begin{equation} \label{eqn:1620}
0 \to \opn{Gr}_q^{G}(I) \to G_q(I) \to G_{q - 1}(I) \to 0 
\end{equation}
in $\dcat{C}_{\mrm{str}}(\cat{M})$. 
In each degree $p \in \Z$ the exact sequence 
\[ 0 \to \opn{Gr}_q^{G}(I)^p \to G_q(I)^p \to G_{q - 1}(I)^p \to 0 \]
in $\cat{M}$ splits, because $\opn{Gr}^G_{q}(I)^p$ is an injective 
object. Thus the exact sequence (\ref{eqn:1620}) is split in the category 
$\bcat{G}_{\mrm{str}}(\cat{M})$ of graded objects in $\cat{M}$. 

Let $N \in \dcat{C}(\cat{M})$ be an acyclic complex. Applying the 
functor $\opn{Hom}_{\cat{M}}(N, -)$ 
to the sequence of complexes (\ref{eqn:1620}) we obtain a sequence 
\begin{equation} \label{eqn:1621}
0 \to \opn{Hom}_{\cat{M}} \bigl( N, \opn{Gr}^G_{q}(I) \bigr) \to  
\opn{Hom}_{\cat{M}} \bigl( N, G_q(I) \bigr) \to 
\opn{Hom}_{\cat{M}} \bigl( N, G_{q - 1}(I) \bigr) \to  0
\end{equation}
in $\dcat{C}_{\mrm{str}}(\K)$. Because (\ref{eqn:1620}) is split exact in 
$\bcat{G}_{\mrm{str}}(\cat{M})$, the sequence (\ref{eqn:1621}) is split exact 
in $\bcat{G}_{\mrm{str}}(\K)$. Therefore (\ref{eqn:1621}) is exact in 
$\dcat{C}_{\mrm{str}}(\K)$.

By the induction hypothesis the complex 
$\opn{Hom}_{\cat{M}} \bigl( N, G_{q - 1}(I) \bigr)$
is acyclic. By step 2 the complex
$\opn{Hom}_{\cat{M}} \bigl( N, \opn{Gr}^G_{q}(I) \bigr)$
is acyclic. The long exact cohomology sequence associated to 
(\ref{eqn:1621}) shows that the complex 
$\opn{Hom}_{\cat{M}} \bigl( N, G_{q}(I) \bigr)$
is acyclic too.

\medskip \noindent
Step 4. We keep the semi-injective cofiltration 
$G = \{ G_q(I) \}_{q \geq -1}$ from step 3. 
Take any acyclic complex $N \in \bcat{C}(\cat{M})$. By 
Proposition \ref{prop:1535} we know that 
\[ \opn{Hom}_{\cat{M}}(N, I) \cong 
\lim_{\leftarrow q} \, \opn{Hom}_{\cat{M}} 
\bigl( N, G_q(I) \bigr) \]
in $\dcat{C}_{\mrm{str}}(\K)$. According to step 3 the complexes 
$\opn{Hom}_{\cat{M}} \bigl( N, G_q(I) \bigr)$ are all acyclic. 
The exactness of the sequences (\ref{eqn:1621}) implies that 
the inverse system \lb 
$\bigl\{ \opn{Hom}_{\cat{M}} \bigl( N, G_q(I) \bigr) \bigr\}_{q \geq -1}$
in $\dcat{C}_{\mrm{str}}(\K)$ has surjective transitions. Now the 
Mittag-Leffler argument (Corollary \ref{cor:1590}) says that the inverse limit 
complex $\opn{Hom}_{\cat{M}}(N, I)$ is acyclic. 
\end{proof}


\begin{prop} \label{prop:1635}
Let $\cat{M}$ be an abelian category. If $I$ is a bounded below complex of 
injectives, then $I$ is a semi-injective complex.
\end{prop}

\begin{proof}
We can assume that $I \neq 0$. Let 
$p_0 := \opn{inf}(I) \in \Z$. For $q \geq -1$ let 
$G_q(I) := \opn{stt}^{\leq p_0 + q}(I)$,
the stupid truncation from Definition \ref{dfn:3066}.
The cofiltration $G = \{ G_q(I) \}_{q \geq -1}$ is semi-injective. 
\end{proof}

The next theorem is \cite[Lemma I.4.6(1)]{RD}. See also 
\cite[Proposition 1.7.7(i)]{KaSc1}. 

\begin{thm} \label{thm:1630}
Let $\cat{M}$ be an abelian category, and let $\cat{J} \sub \cat{M}$ be a
full subcategory such that every object $M \in \cat{M}$ admits a 
monomorphism $M \inj I$ to some object $I \in \cat{J}$. Then every complex 
$M \in \bcat{C}^{+}(\cat{M})$ 
admits a quasi-isomorphism $\rho : M \to  I$ in 
$\dcat{C}^+_{\mrm{str}}(\cat{M})$,
such that $\opn{inf}(I) = \opn{inf}(M)$,
and each $I^p$ is an object of $\cat{J}$.
\end{thm}

\begin{proof}
The proof is the same as that of Theorem \ref{thm:1540}, except for a 
mechanical reversal of arrows. To be more explicit, let us
take $\cat{N} := \cat{M}^{\mrm{op}}$ and $\cat{Q} := \cat{J}$.
Since monomorphisms in $\cat{M}$ become epimorphisms in $\cat{N}$, the 
full subcategory $\cat{Q} \sub \cat{N}$ satisfies the 
assumptions of Theorem \ref{thm:1540}. By Theorem \ref{thm:2495} we have a 
canonical isomorphism of categories 
$\bcat{C}^{-}_{\mrm{str}}(\cat{N}) \iso 
\dcat{C}^+_{\mrm{str}}(\cat{M})^{\mrm{op}}$. 
Thus a quasi-isomorphism $Q \to N$ in 
$\bcat{C}^{-}_{\mrm{str}}(\cat{N})$
gives rise to a quasi-isomorphism $M \to I$ in 
$\bcat{C}^{+}_{\mrm{str}}(\cat{M})$.
\end{proof}

\begin{dfn} \label{dfn:2880}
Let $\cat{M}$ be an abelian category, and let $\cat{M}' \sub \cat{M}$ be a 
full abelian subcategory. 
We say that {\em $\cat{M}'$ has enough injectives relative to $\cat{M}$} 
if every object $M \in \cat{M}'$ admits a monomorphism 
$M \inj I$, where $I$ is an object of $\cat{M}'$ that is injective in 
the bigger category $\cat{M}$.
\end{dfn}

Of course, in this situation the category $\cat{M}'$ itself has enough 
injectives. 

Thick abelian categories were defined in Definition \ref{dfn:2323}. 
The next theorem is \cite[Lemma I.4.6(3)]{RD}. See also 
\cite[Proposition 1.7.11]{KaSc1}. 

\begin{thm} \label{thm:2880}
Let $\cat{M}$ be an abelian category, and let $\cat{M}' \sub \cat{M}$ be a 
thick abelian subcategory that has enough injectives relative to $\cat{M}$.
Let $M \in \dcat{C}(\cat{M})$ be a complex with bounded below cohomology, 
such that $\opn{H}^i(M) \in \cat{M}'$ for all $i$.
Then there is a quasi-isomorphism $\rho : M \to I$ in 
$\dcat{C}_{\mrm{str}}(\cat{M})$, such that
$I \in \dcat{C}^+(\cat{M}')$, each $I^p$ is an injective object in 
$\cat{M}$, and $\opn{inf}(I) = \opn{inf}(\opn{H}(M))$. 
\end{thm}

Before the proof we need some auxiliary material. 

Suppose we are given morphisms $\psi_1 : K \to L_1$ and 
$\psi_2 : K \to L_2$ in $\cat{M}$. The {\em fibered coproduct} is the object 
\begin{equation} \label{eqn:2885}
L_1 \oplus_K L_2 := \opn{Coker} 
\bigl( (\psi_1, -\psi_2) : K \to L_1 \oplus L_2 \bigr) 
\end{equation}
in $\cat{M}$. 
It has an obvious universal property. The commutative diagram 
\begin{equation} \label{eqn:5072}
\UseTips \xymatrix @C=8ex @R=6ex {
K
\ar[r]^{\psi_1}
\ar[d]_{\psi_2}
&
L_1
\ar[d]^{\ep_1}
\\
L_2
\ar[r]^(0.37){\ep_2}
& 
L_1 \oplus_K L_2
} 
\end{equation}
in which $\ep_i$ are the morphisms induced by the embeddings 
$L_i \to L_1 \oplus L_2$, is sometimes called a {\em pushout diagram}. 

\begin{lem} \label{lem:2880}
In the situation above\tup{:}
\begin{enumerate}
\item The sequence
\[ \opn{Ker}(\psi_1) \xar{\psi_2} L_2 \xar{\ep_2} 
L_1 \oplus_K L_2 \xar{\pi_1} \opn{Coker}(\psi_1) \to 0 , \]
in which $\pi_1$ is the morphism induced by the projection 
$L_1 \oplus L_2 \to L_1$, is exact. 

\item Suppose $L_1 \to L'_1$ be a monomorphism. Then the induced morphism 
$L_1 \oplus_K L_2 \to L'_1 \oplus_K L_2$ 
is a monomorphism.
\end{enumerate}
\end{lem}

\begin{exer} \label{exer:2880}
Prove Lemma \ref{lem:2880}. (Hint: use the first sheaf trick,
Proposition \ref{prop:3635}.)  
\end{exer}

\begin{proof}[Proof of Theorem \tup{\ref{thm:2880}}]
The proof is an adaptation of the proof of \cite[Proposition 1.7.11]{KaSc1}. 
In the proof we use the objects of cocycles $\opn{Z}^p(L)$, 
co\-boundaries $\opn{B}^p(L)$ and decocycles $\opn{Y}^p(L)$, that are 
associated to a complex $L$ and an integer $p$; see Definition \ref{dfn:2993}.

\medskip \noindent
Step 1. We can assume that $\opn{H}(M) \neq 0$. 
By translating $M$, we may further assume that $\opn{inf}(\opn{H}(M)) = 0$. 
Then, after replacing $M$ with its smart truncation $\opn{smt}^{\geq 0}(M)$, 
we can assume that $\opn{inf}(M) = 0$. 

We are going to construct an inverse system 
\begin{equation} \label{eqn:5105}
F_p(I) = \bigl( \cdots \to 0 \to I^0 \xar{\d_I^0} I^1 
\xar{\d_I^{1}} \cdots \xar{\d_I^{p - 1}} I^p \to 0 \to \cdots \bigr)
\end{equation}
in $\dcat{C}_{\mrm{str}}(\cat{M}')$, indexed by $p \geq -1$, such that
the objects $I^p \in \cat{M}'$ are injective in $\cat{M}$. 
The morphism $F_{p + 1}(I) \to F_{p}(I)$ will send the object 
$I^{p + 1} = F_{p + 1}(I)^{p + 1}$ to 
$0$, and will be the identity in all other degrees. 
Simultaneously we will 
construct an inverse system of morphisms 
\begin{equation} \label{eqn:5106}
F_p(\phi) : M \to F_p(I)
\end{equation}
in $\dcat{C}_{\mrm{str}}(\cat{M})$.
The construction will be inductive. 
The morphisms $F_p(\phi)$ will satisfy this condition:
\begin{itemize}
\rmitem{$\heartsuit$} The morphism 
$\opn{H}^q(F_p(\phi)) : \opn{H}^q(M) \to \opn{H}^q(F_p(I))$
is a monomorphism for all $q \leq p$ and an isomorphism for all
$q \leq p - 1$. 
\end{itemize}

Then the complex 
$I := \lim_{\lar p} F_{p}(I)$ in $\dcat{C}(\cat{M}')$
is going to be of the required kind, and the morphism 
$\phi := \lim_{\lar p} F_{p}(\phi)$
in $\dcat{C}_{\mrm{str}}(\cat{M})$ is going to be a quasi-isomorphism.
Note that these inverse limits are innocent, since in each degree $q$ the 
inverse systems $\{ F_{p}(I)^q \}_{p \geq -1}$ and 
$\{ F_{p}(\phi)^q \}_{p \geq -1}$ are eventually stationary. 

\medskip \noindent
Step 2. In this step we begin the construction. 
Here $p \geq -1$, and we have a complex $F_p(I)$ as in (\ref{eqn:5105})
of the kind described above,
and a morphism $F_p(\phi) : M \to F_p(I)$ in 
$\dcat{C}_{\mrm{str}}(\cat{M})$ satisfying condition ($\heartsuit$). 
If  $p = -1$ then $F_{-1}(I) = 0$ and $F_{-1}(\phi) = 0$ of course.

We claim that the objects 
$F_p(I)^q$, $\opn{H}^q(F_p(I))$, 
$\opn{B}^{q}(F_p(I))$, $\opn{Y}^{q}(F_p(I))$ and 
$\opn{Z}^{q}(F_p(I))$ belong to $\cat{M}'$ for all $q \in [0, p + 1]$. 
For $q = p + 1$ all these objects are zero. 
For $F_p(I)^q = I^q$ it is given. For 
$\opn{H}^p(F_p(I)) = \opn{Coker}(\d_I^{p - 1})$ it is because $\cat{M}'$ is 
a full abelian subcategory of $\cat{M}$. For $\opn{H}^q(F_p(I))$ when 
$q < p$ we use the isomorphisms $\opn{H}^q(F_p(\phi))$ and the fact that 
$\opn{H}^q(M) \in \cat{M}'$. As for the rest of the objects listed, this is 
shown by descending induction on $q$, starting from $q = p$, using these short 
exact sequences 
\[ 0 \to \opn{Z}^{q}(F_p(I)) \to F_p(I)^{q} \to 
\opn{B}^{q + 1}(F_p(I)) \to  0 , \]
\[ 0 \to \opn{B}^{q}(F_p(I)) \to F_p(I)^q \to \opn{Y}^{q}(F_p(I)) \to  0 , \]
\[ 0 \to \opn{B}^{q}(F_p(I)) \to \opn{Z}^{q}(F_p(I)) \to 
\opn{H}^{q}(F_p(I)) \to  0 , \]
and the fact that $\cat{M}'$ is a full abelian subcategory of $\cat{M}$. 

\medskip \noindent
Step 3. The construction continues. Let us denote the components of the 
morphism $F_p(\phi)$ by
$\phi^q := F_p(\phi)^q : M^q \to F_p(I)^q = I^q$
for $q \in [0, p]$; outside this degree range we have $F_p(\phi)^q = 0$. 

The differential $\d_M^{p} : M^{p} \to M^{p + 1}$
induces a morphism $\psi_1 : \opn{Y}^p(M) \to \opn{Z}^{p + 1}(M)$.
There is also the morphism 
$\psi_2 := \opn{Y}^p(\phi) : \opn{Y}^p(M) \to \opn{Y}^{p}(F_p(I))$.
We define $N$ to be the fibered coproduct
\begin{equation} \label{eqn:5107}
N := \opn{Z}^{p + 1}(M) \oplus_{\opn{Y}^p(M)} \opn{Y}^{p}(F_p(I))
\end{equation}
relative to the morphisms $\psi_1$ and $\psi_2$. 
In the notation of Lemma \ref{lem:2880}, with the objects  $K := \opn{Y}^p(M)$, 
$L_1 := \opn{Z}^{p + 1}(M)$ and $L_2 := \opn{Y}^{p}(F_p(I))$,
and with the morphism $\psi_1, \psi_2$ that we have just defined, there are 
morphisms 
$\ep_1 : \opn{Z}^{p + 1}(M) \to N$,
$\ep_2 : \opn{Y}^{p}(F_p(I)) \to N$ and
$\pi_1 : N \to \opn{Coker}(\psi_1) = \opn{H}^{p + 1}(M)$.
According to Lemma \ref{lem:2880}(1) there is an exact sequence 
\begin{equation} \label{eqn:5109}
\opn{H}^{p}(M) \xar{\psi_2} \opn{Y}^{p}(F_p(I)) \xar{\ep_2} 
N \xar{\pi_1} \opn{H}^{p + 1}(M) \to 0 
\end{equation}
in $\cat{M}$. 
However, letting 
$\be : \opn{H}^{p}(F_p(I)) \to \opn{Y}^{p}(F_p(I))$ 
be the canonical monomorphism, there is equality 
$\psi_2 = \be \circ \opn{H}^p(F_p(\phi))$. By condition ($\heartsuit$) with 
indices $q = p$ we know that $\opn{H}^p(F_p(\phi))$ is a monomorphism. 
Therefore $\psi_2$ is a monomorphism, and (\ref{eqn:5109}) extends to a 
slightly longer exact sequence 
\begin{equation} \label{eqn:5115}
0 \to \opn{H}^{p}(M) \xar{\psi_2} \opn{Y}^{p}(F_p(I)) \xar{\ep_2} 
N \xar{\pi_1} \opn{H}^{p + 1}(M) \to 0 
\end{equation}
in $\cat{M}$. 

Let $\ga : \opn{Z}^{p + 1}(M) \inj M^{p + 1}$ be the canonical monomorphism,
and define 
$\psi_1^+ := \ga \circ \psi_1 : \opn{Y}^{p}(M) \to M^{p + 1}$. 
Let 
$\psi_2^+ := \psi_2 :  \opn{Y}^p(M) \to \opn{Y}^{p}(F_p(I))$.
We have the corresponding fibered coproduct
\begin{equation} \label{eqn:5108}
N^+ := M^{p + 1} \oplus_{\opn{Y}^p(M)} \opn{Y}^{p}(F_p(I)) . 
\end{equation}
The monomorphism $\ga : \opn{Z}^{p + 1}(M) \inj M^{p + 1}$ induces a 
morphism $\om : N \to N^+$ between these fibered coproducts, and Lemma 
\ref{lem:2880}(2) says that $\om$ is also a monomorphism. 

\medskip \noindent
Step 4. We now produce the object $I^{p + 1} \in \cat{M}'$. 
In step 2 we showed that $\opn{Y}^{p}(F_p(I)) \in \cat{M}'$. It is given that 
$\opn{H}^{q}(M) \in \cat{M}'$ for all $q$. 
The exact sequence (\ref{eqn:5109}) and Proposition \ref{prop:3171} 
say that $N \in \cat{M}'$ too. 
By assumption there is a monomorphism
$\chi : N \inj I^{p + 1}$
into some object  $I^{p + 1} \in \cat{M}'$ that's injective in $\cat{M}$. 
Because $\om : N \to N^+$ is a monomorphism, there is a morphism 
$\chi^+ : N^+ \to I^{p + 1}$
such that $\chi^+ \circ \om = \chi$. 

We have a canonical epimorphism $\al_I : I^{p} \surj \opn{Y}^{p}(F_p(I))$. 
Define the morphism 
$\d_{I}^{p} : I^{p} \to I^{p + 1}$
by the formula 
$\d_{I}^{p} := \chi \circ \ep_2 \circ \al_I$. Since 
$\al_I \cd \d_{I}^{p - 1} = 0$, we see that 
$\d_{I}^{p} \circ \d_{I}^{p - 1} = 0$. 
This is depicted in the commutative diagram below. 
\[ \UseTips \xymatrix @C=6ex @R=6ex {
I^{p - 1}
\ar[r]_(0.6){\d_I^{p - 1}}
\ar@(ur,ul)[rr]^{0}
&
I^p
\ar[r]^(0.4){\al_I}
\ar@(dr,dl)[rrr]_{\d_I^p}
&
\opn{Y}^{p}(F_p(I))
\ar[r]^(0.6){\ep_2}
&
N
\ar[r]^{\chi}
&
I^{p + 1}
} \]
Therefore we obtain a new complex 
\begin{equation} \label{eqn:5117}
F_{p + 1}(I) := \bigl( \cdots \to 0 \to I^0 \xar{\d_I^0} I^1 
\xar{\d_I^{1}} \cdots \xar{\d_I^{p - 1}} I^p 
\xar{\d_I^{p}} I^{p + 1} \to 0 \to \cdots \bigr) 
\end{equation}
in $\dcat{C}(\cat{M}')$, and it projects onto $F_{p}(I)$. 

\medskip \noindent
Step 5. In this step we construct the morphism 
$\phi^{p + 1} : M^{p + 1} \to I^{p + 1}$. 
Let $\ep^+_1 : M^{p + 1} \to N^+$ be the morphism that's part of the fibered 
coproduct (\ref{eqn:5108}). 
We define the morphism 
$\phi^{p + 1} : M^{p + 1} \to I^{p + 1}$
by the formula 
$\phi^{p + 1} := \chi^+ \circ \ep_1^+$. 

Let's examine the next diagram, where
$\al_M : M^p \surj \opn{Y}^{p}(M)$ is the canonical epimorphism, and 
$\psi_1^+$, $\psi_2^+ = \psi_2$, $\ep_1^+$, $\ep_2^+$ and $\ep_2$ are the 
morphisms involved in  the two fibered coproducts. 
\begin{equation} \label{eqn:5110}
\UseTips \xymatrix @C=6ex @R=6ex {
M^p
\ar[r]^(0.4){\al_M}
\ar[d]^{\phi^p}
&
\opn{Y}^{p}(M)
\ar[r]^(0.55){\psi_1^+}
\ar[d]^{\psi_2^+}
&
M^{p + 1}
\ar[r]^{\ep_1^+}
&
N^+
\ar[r]^{\chi^+}
&
I^{p + 1}
\ar[d]^{\opn{id}}
\\
I^p
\ar[r]^(0.4){\al_I}
&
\opn{Y}^{p}(F_p(I))
\ar[rr]^{\ep_2}
\ar[urr]^{\ep_2^+}
&
&
N
\ar[r]^{\chi}
\ar[u]_{\om}
&
I^{p + 1}
} 
\end{equation}
The left square is commutative because $\al : \opn{Id} \to \opn{Y}$ is a 
morphism of functors. The triangle that has $\opn{Y}^{p}(M)$ as a vertex is 
commutative because $N^+$ is the fibered coproduct. The other triangle, the 
one that has $N$ as a vertex, is commutative because $\om$ is a morphism 
between fibered coproducts. The right square is commutative by the defining 
property of $\chi^+$. So the whole diagram is commutative. 

We now take part of the commutative diagram (\ref{eqn:5110}), and add to it 
three curved arrows. 
\begin{equation} \label{eqn:5111}
\UseTips \xymatrix @C=6ex @R=5.5ex {
M^p
\ar[r]^(0.4){\al_M}
\ar[d]_{\phi^p}
\ar@(ur,ul)[rr]^{\d_{M}^{p}}
&
\opn{Y}^{p}(M)
\ar[r]^(0.55){\psi_1^+}
&
M^{p + 1}
\ar[r]^{\ep_1^+}
\ar@(ur,ul)[rr]^{\phi^{p + 1}}
&
N^+
\ar[r]^{\chi^+}
&
I^{p + 1}
\ar[d]^{\opn{id}}
\\
I^p
\ar[r]^(0.4){\al_I}
\ar@(dr,dl)[rrrr]_{\d_{I}^{p}}
&
\opn{Y}^{p}(F_p(I))
\ar[rr]^{\ep_2}
%\ar[urr]^{\ep_2^+}
&
&
N
\ar[r]^{\chi}
&
I^{p + 1}
} 
\end{equation}
The curved polygons are commutative, as can be seen from the definitions of the 
morphisms $\psi_1^+$, $\phi^{p + 1}$ and $\d_{I}^{p}$. We conclude that 
$\d_I^{p} \circ \phi^p = \phi^{p + 1} \circ \d_M^{p}$,
so there is a new morphism of complexes 
$F_{p + 1}(\phi) : M \to F_{p + 1}(I)$
whose degree $p + 1$ component is $\phi^{p + 1}$. 

\medskip \noindent
Step 6. Here we prove that $\opn{H}^{p + 1}(F_{p + 1}(\phi))$ is a 
monomorphism and $\opn{H}^{p}(F_{p + 1}(\phi))$ is an isomorphism.
This will verify condition ($\heartsuit$) with index $p + 1$.

Consider the fibered coproduct 
\begin{equation} \label{eqn:5116}
N^{\lozenge} := I^{p + 1} \oplus_{\opn{Y}^p(F_p(I))} \opn{Y}^{p}(F_p(I))
\end{equation} 
of the morphisms $\psi_1^{\lozenge}$ and $\psi_2^{\lozenge}$,
where $\psi_1^{\lozenge} : \opn{Y}^{p}(F_p(I)) \to I^{p + 1}$
is the morphism induced by $\d_I^p$, and 
$\psi_2^{\lozenge}$ is the identity of $\opn{Y}^{p}(F_p(I))$. 
Note that $\ep_1^{\lozenge} : I^{p + 1} \to N^{\lozenge}$ is an isomorphism, 
and 
$\opn{Coker}(\psi_1^{\lozenge}) = \opn{Y}^{p + 1}(F_{p + 1}(I)) = 
\opn{H}^{p + 1}(F_{p + 1}(I))$. 
The morphisms 
$\phi^{p + 1} : \opn{Z}^{p + 1}(M) \to I^{p + 1}$
and 
$\opn{Y}^{p}(F_{p}(\phi)) : \opn{Y}^{p}(M) \to \opn{Y}^{p}(F_{p}(I))$
induce a morphism 
$\om^{\lozenge} : N \to N^{\lozenge}$ between the fibered coproducts 
(\ref{eqn:5107}) 
and (\ref{eqn:5116}). According to Lemma \ref{lem:2880}(1) there is an 
induced commutative diagram
\begin{equation} \label{eqn:51168}
\UseTips \xymatrix @C=4ex @R=6ex {
0
\ar[r]
&
\opn{H}^{p}(M)
\ar[r]^(0.47){\psi_2}
\ar[d]_{\opn{H}^{p}(F_{p + 1}(\phi))}
&
\opn{Y}^{p}(F_{p}(I))
\ar[r]^(0.65){\ep_2}
\ar[d]_{\opn{id}}
&
N
\ar[r]^(0.4){\pi_2}
\ar[d]_{\om^{\lozenge}}
&
\opn{H}^{p + 1}(M)
\ar[r]
\ar[d]^{\opn{H}^{p + 1}(F_{p + 1}(\phi))}
&
0
\\
0
\ar[r]
&
\opn{H}^{p}(F_{p + 1}(I))
\ar[r]^{\psi_2^{\lozenge}}
&
\opn{Y}^{p}(F_{p}(I))
\ar[r]^(0.65){\ep_2^{\lozenge}}
&
N^{\lozenge}
\ar[r]^(0.3){\pi_2^{\lozenge}}
&
\opn{H}^{p + 1}(F_{p + 1}(I))
\ar[r]
&
0
} 
\end{equation}
The first row is an exact sequence, since it is (\ref{eqn:5115}).
The second row is also exact: exactness at the last three objects is by 
Lemma \ref{lem:2880}(1). As for exactness at $\opn{H}^{p}(F_{p + 1}(I))$,
we use the fact that the morphism 
$\opn{Y}^{p}(F_{p + 1}(I)) \to \opn{Y}^{p}(F_{p}(I))$
is an isomorphism, and its composition with  the canonical embedding 
$\opn{H}^{p}(F_{p + 1}(I)) \inj \opn{Y}^{p}(F_{p + 1}(I))$
is $\psi_2^{\lozenge}$, so the latter is a monomorphism.
A little calculation shows that $\om^{\lozenge} = \ep_1^{\lozenge} \circ \chi$, 
so it is a monomorphism. 

Finally, a diagram chase in (\ref{eqn:51168}), using the first sheaf trick, 
shows that $\opn{H}^{p + 1}(F_{p + 1}(\phi))$
is a monomorphism, and that $\opn{H}^{p}(F_{p + 1}(\phi))$
is an isomorphism. 
\end{proof}

\begin{cor} \label{cor:3185}
Under the assumptions of Theorem \tup{\ref{thm:2880}}, the canonical functor 
$\dcat{D}^+(\cat{M}') \to \dcat{D}^+_{\cat{M}'}(\cat{M})$
is an equivalence. 
\end{cor}

\begin{proof}
This is very similar to the proof of Corollary \ref{cor:3175}, and there is no 
need to repeat it. 
\end{proof}


Here is an important instance in which Theorem \ref{thm:2880} and Corollary 
\ref{cor:3185} apply. 

\begin{exa} \label{exa:2895}
Let $(X, \OO_X)$ be a noetherian scheme. Associated to it are these abelian 
categories: the category $\cat{M} := \cat{Mod} \OO_X$ of $\OO_X$-modules,
and the thick abelian subcategory $\cat{M}' := \cat{QCoh} \OO_X$
of quasi-coherent $\OO_X$-modules. 
According to \cite[Proposition II.7.6]{RD} the category $\cat{M}'$ has 
enough injectives relative to $\cat{M}$.
By Corollary \ref{cor:3185} the canonical functor
\[ \dcat{D}(\cat{QCoh} \OO_X) \to \dcat{D}^+_{\mrm{qc}}(\cat{Mod} \OO_X) \]
is an equivalence. 

For a more general statement, using other methods, see 
\cite[Corollary 5.5]{BoNe}. 
\end{exa}

\begin{cor} \label{cor:1630}
\index{Resolution! K-injective}
If $\cat{M}$ is an abelian category with enough injectives, then  
$\bcat{C}^{+}(\cat{M})$ has enough K-injectives. 
\end{cor}

\begin{proof}
According to either Theorem \ref{thm:1630} or Theorem \ref{thm:2880},
every $M \in \bcat{C}^{+}(\cat{M})$ admits a quasi-isomorphism $M \to I$
to bounded below complex of injectives $I$. Now use Proposition 
\ref{prop:1635} and Theorem \ref{thm:1620}. 
\end{proof}

\begin{cor} \label{cor:1902}
Let $\cat{M}$ be an abelian category with enough injectives, and let 
$M \in \dcat{C}(\cat{M})$ be a complex with bounded below cohomology. Then $M$ 
has a K-injective resolution $M \to I$,
such that $\opn{inf}(I) = \opn{inf}(\opn{H}(M))$, 
and every $I^p$ is an injective object of $\cat{M}$. 
\end{cor}

\begin{proof}
We may assume that $\opn{H}(M)$ is nonzero. Let 
$p := \opn{inf}(\opn{H}(M)) \in \Z$, and let 
$N := \opn{smt}^{\geq p}(M)$, the smart truncation from 
Definition \ref{dfn:2320}. So $M \to N$ is a quasi-isomorphism, and 
$\opn{inf}(N) = p$. 
According to either Theorem \ref{thm:1630} or Theorem \ref{thm:2880},
there is a quasi-isomorphism $N \to I$, where $I$ is a complex 
of injectives and $\opn{inf}(I) = p$. 
By Proposition \ref{prop:1635} and Theorem \ref{thm:1620} the complex $I$ is 
K-injective. The composed quasi-isomorphism $M \to I$ is what we are looking 
for. 
\end{proof}




\mysubsection{K-Injective Resolutions in 
\texorpdfstring{$\bcat{C}(A)$}{C(A)}} 
\label{subsec:exis-K-inj-dgmods}

Recall that we are working over a non\-zero commutative base ring $\K$, and 
$A$ is a central DG $\K$-ring. 

An {\em injective cogenerator} of the abelian category
$\dcat{M}(\K) = \cat{Mod} \K$ is an injective 
$\K$-module $\K^*$ with this property: if $M$ is a nonzero $\K$-module, then 
$\opn{Hom}_{\K}(M, \K^*)$ is nonzero. These always exist. Here are a few 
examples. 

\begin{exa} \label{exa:1665} 
For every $\K$ there is a canonical choice for an injective cogenerator: 
$\K^* := \opn{Hom}_{\Z}(\K, \Q / \Z)$.
See proof of Theorem \ref{thm:1115}.
Usually this is a very big module! 
\end{exa}

\begin{exa} \label{exa:1876}
Assume $\K$ is a complete noetherian local ring, with maximal ideal 
$\m$ and residue field $\mbb{k} = \K / \m$. In this case we would prefer to take
the smallest possible injective cogenerator $\K^*$, and this is the injective 
hull of $\mbb{k}$ as a $\K$-module. 

Here are some special cases. 
If $\K$ is a field, then $\K^* = \K = \mbb{k}$. 
If $\K = \what{\Z}_p$, the ring of $p$-adic integers, then 
$\mbb{k} = \mbb{F}_p$, and $\K^* \cong \what{\Q}_p / \what{\Z}_p$, which is the 
$p$-primary part of $\Q / \Z$. 
If $\K$ contains some field, then there exists a ring homomorphism 
$\mbb{k} \to \K$ that lifts the canonical surjection $\K \to \mbb{k}$.
(This is by the Cohen Structure Theorem, see \cite[Theorem 7.7]{Eis}.)
After choosing such a lifting, there is an isomorphism of $\K$-modules 
$\K^* \cong \opn{Hom}^{\mrm{cont}}_{\mbb{k}}(\K, \mbb{k})$,
where continuity is for the $\m$-adic topology on $\K$ and the discrete 
topology on $\mbb{k}$. 
\end{exa}

In this subsection we have the following standing convention.  

\begin{conv} \label{conv:5075}
There is a fixed injective cogenerator $\K^*$ of $\dcat{M}(\K)$.
\end{conv}

We view $\K^*$ as a DG $\K$-module concentrated in degree $0$ 
(with zero differential). 
For every $p \in \Z$ there is the DG $\K$-module $\opn{T}^{-p} (\K^*)$, 
which is concentrated in degree $p$.

It will be convenient to blur the distinction between DG $\K$-modules with 
zero differentials and graded $\K$-modules. Namely, if $N$ is a DG 
$\K$-module such that $\d_N = 0$, we will identify $N$ with the graded modules 
$N^{\natural}$, $\opn{H}(N)$, etc.

\begin{dfn} \label{dfn:5075}
For a DG $\K$-module $V$ we define the DG $\K$-module 
$V^* := \opn{Hom}_{\K}(V, \K^*)$. 
\end{dfn}

The operation $(-)^*$ is an exact contravariant functor from 
$\dcat{C}_{\mrm{str}}(\K)$ to itself, and also from 
$\dcat{G}_{\mrm{str}}(\K)$ to itself. Note that given 
$M \in \dcat{C}(A)$, its dual $M^*$ is an object of 
$\dcat{C}(A^{\mrm{op}})$, so that we have an exact functor 
\begin{equation} \label{eqn:5075}
(-)^* : \dcat{C}_{\mrm{str}}(A)^{\mrm{op}} \to 
\dcat{C}_{\mrm{str}}(A^{\mrm{op}}) . 
\end{equation}

\begin{dfn} \label{dfn:1805}
A DG $\K$-module $W$ is called {\em cofree} if there is an isomorphism
$W \cong V^*$ in $\dcat{C}_{\mrm{str}}(\K)$, for some free DG $\K$-module 
$V$. 
\end{dfn}

Since the differential of a free DG $\K$-module $V$ is zero, we see that 
a cofree DG module $W$ also has zero differential. 

\begin{lem} \label{lem:5075}
Let $V$ be a free DG $\K$-module, and let $W := V^*$. 
\begin{enumerate}
\item If $V \cong \boplus_{s \in S} \opn{T}^{p_s} (\K)$
for some indexing set $S$ and some collection of integers 
$\{ p_s \}_{s \in S}$, then 
$W \cong \prod_{s \in S} \opn{T}^{-p_s} (\K^*)$. 

\item As an object of the abelian category $\dcat{G}_{\mrm{str}}(\K)$, $W$ is 
injective. 
\end{enumerate}
\end{lem}

In item (1) the isomorphisms can be considered either in 
$\dcat{C}_{\mrm{str}}(\K)$ or in $\dcat{G}_{\mrm{str}}(\K)$; it doesn't matter, 
because $V$ and $W$ have zero differentials. However, item (2) is true only in 
the category $\dcat{G}_{\mrm{str}}(\K)$.

\begin{lem} \label{lem:3346}
Let $\phi : U \to V$ be a homomorphism in $\dcat{G}_{\mrm{str}}(\K)$. 
\begin{enumerate}
\item $\phi$ is injective iff $\phi^* : V^* \to U^*$ is surjective. 

\item $\phi$ is surjective iff $\phi^* : V^* \to U^*$ is injective.

\item The canonical homomorphism $U \to U^{**} = (U^*)^*$ in 
$\dcat{G}_{\mrm{str}}(\K)$ is injective. 

\item There exists an injective homomorphism 
$U \inj W$ in $\dcat{G}_{\mrm{str}}(\K)$
into some cofree DG $\K$-module $W$. 
\end{enumerate}
\end{lem}

\begin{exer}
Prove Lemmas \ref{lem:5075} and \ref{lem:3346}. 
\end{exer}

\begin{dfn} \label{dfn:3345}
Let $W$ be a cofree DG $\K$-module. The {\em cofree DG $A$-module coinduced 
from $W$} is the DG $A$-module
$I_W :=  \opn{Hom}_{\K}(A, W)$.
There is a homomorphism $\th_W : I_W \to W$
in $\dcat{C}_{\mrm{str}}(\K)$, with formula 
$\th_W(\chi) := \chi(1_A) \in W$. 
\end{dfn}

\begin{dfn} \label{dfn:3346}
A DG $A$-module $J$ is called {\em cofree}%
\index{Differential graded module! cofree}
if there is an isomorphism 
$J \cong I_W$ in $\dcat{C}_{\mrm{str}}(A)$ for some cofree DG $\K$-module 
$W$. 
\end{dfn}

\begin{lem} \label{lem:3345}
Consider the free DG $\K$-module 
$V := \boplus_{s \in S} \opn{T}^{p_s} (\K)$
and the cofree DG $\K$-module 
$W := V^*$. There are canonical isomorphisms 
\[ I_W \cong (A \ot V)^* \cong
\Bigl( \bigoplus\nolimits_{s \in S} \, \opn{T}^{p_s} (A) \Bigr)^*
\cong \prod\nolimits_{s \in S} \, \opn{T}^{-p_s} (A^*) \]
in $\dcat{C}_{\mrm{str}}(A)$.
\end{lem}

\begin{exer}
Prove the lemma above. 
\end{exer}

\begin{lem} \label{lem:3347}
Let $W$ be a cofree DG $\K$-module, and let $M$ be a DG $A$-module.
The homomorphism 
\[ \opn{Hom}(\opn{id}_M, \th_W) : \opn{Hom}_{A}(M, I_W) \to 
\opn{Hom}_{\K}(M, W) \]
in $\dcat{C}_{\mrm{str}}(\K)$ is an isomorphism. 
\end{lem}

\begin{proof}
The homomorphism
\[ \opn{Hom}(\opn{id}_M, \th_W) :
\opn{Hom}_{A} \bigl( M,  \opn{Hom}_{\K}(A, W) \bigr) \iso 
\opn{Hom}_{\K}(M, W) \]
is just adjunction for the DG ring homomorphism $\K \to A$, so it is bijective. 
\end{proof}

Recall that $\dcat{G}_{\mrm{str}}(A^{\natural})$ is the abelian category whose 
objects are the graded $A^{\natural}$-modules, and the morphisms are the 
$A$-linear homomorphisms of degree $0$. The forgetful functor 
$\opn{Und} : \dcat{C}_{\mrm{str}}(A) \to \dcat{G}_{\mrm{str}}(A^{\natural})$, 
$M \mapsto M^{\natural}$, is faithful. 

\begin{lem} \label{lem:1770}
Let $I$ be a cofree DG $A$-module. Then $I^{\natural}$ is an injective object 
of $\dcat{G}_{\mrm{str}}(A^{\natural})$.
\end{lem}

\begin{proof}
We can assume that $I = I_W$ for some cofree DG $\K$-module $W$. 
For every  $M \in \dcat{G}_{\mrm{str}}(A^{\natural})$ there are isomorphisms 
\[ \begin{aligned}
& \opn{Hom}_{\dcat{G}_{\mrm{str}}(A^{\natural})} \bigl( M, I_W^{\natural} 
\bigr) 
= \opn{Hom}_{A}(M, I_W)^0  
\\
& \qquad \cong^{\heartsuit} \opn{Hom}_{\K}(M, W)^0 =
\prod_{p \in \Z} \, \opn{Hom}_{\K}(M^{p}, W^p) 
\end{aligned} \]
in $\dcat{M}(\K)$.
The isomorphism $\cong^{\heartsuit}$ is by Lemma \ref{lem:3347}. 
For every $p$  the functor
$\dcat{G}_{\mrm{str}}(A^{\natural}) \to \dcat{M}(\K)$,  $M \mapsto M^p$, is 
exact. Because each $W^p$ is an injective object of $\dcat{M}(\K)$, 
the contravariant functor $\opn{Hom}_{\K}(-, W^p)$
from $\dcat{M}(\K)$ to itself is exact. 
And the product of exact functors into $\dcat{M}(\K)$ is exact. We conclude 
that 
the functor 
$\opn{Hom}_{\dcat{G}_{\mrm{str}}(A^{\natural})} \bigl( -, I_W^{\natural} \bigr)$
is exact. 
\end{proof}


Recall that for a DG $\K$-module $M$ we have the object of decocycles
\[ \opn{Y}(M) := \opn{Coker} \bigl( \d_M : \opn{T}^{-1}(M) \to M \bigr) =
M / \opn{B}(M) \in \dcat{G}(\K) . \]

\begin{lem} \label{lem:3348}
Let $W$ be a cofree DG $\K$-module, let $M$ be a DG $A$-module, and let 
$\chi : \opn{Y}(M) \to W$ be a homomorphism in $\dcat{G}_{\mrm{str}}(\K)$.
Then there is a unique homomorphism 
$\psi : M \to I_W$ in $\dcat{C}_{\mrm{str}}(A)$, such that the diagram 
\[ \UseTips \xymatrix @C=10ex @R=6ex {
\opn{Y}(M)
\ar[r]^{\opn{Y}(\psi)}
\ar@(u,u)[rr]^{\chi}
&
\opn{Y}(I_W)
\ar[r]^(0.45){\opn{Y}(\th_W)}
&
\opn{Y}(W) = W
} \]
in $\dcat{G}_{\mrm{str}}(\K)$ is commutative. 
\end{lem}

\begin{proof}
Since the differentials of $W$ and $\opn{Y}(M)$ are zero, the canonical 
homomorphism 
\[ \al : \opn{Hom}_{\K}(\opn{Y}(M), W)^0 = 
\opn{Z}^0 \bigl( \opn{Hom}_{\K}(\opn{Y}(M), W) \bigr) \to 
\opn{Z}^0 \bigl( \opn{Hom}_{\K}(M, W) \bigr) , \]
that's induced by the canonical surjection 
$M \surj \opn{Y}(M)$, is bijective. This gives us a unique homomorphism 
$\al(\chi) : M \to W$ in $\dcat{C}_{\mrm{str}}(\K)$. Next we use Lemma 
\ref{lem:3347} to obtain a unique 
$\psi : M \to I_W$ in $\dcat{C}_{\mrm{str}}(A)$ s.t.\ 
$\th_W \circ \psi = \al(\chi)$. This $\psi$ is what we are looking for. 
\end{proof}

The next definition is dual to Definition \ref{dfn:1575}. 

\begin{dfn} \label{dfn:1665}
Let $I$ be an object of $\dcat{C}(A)$.
\begin{enumerate}
\item A {\em semi-cofree cofiltration}% 
\index{Cofiltration on a DG module! semi-cofree}
on $I$ is a cofiltration
$G = \{ G_q(I) \}_{q \geq -1}$ on $I$ in $\dcat{C}_{\mrm{str}}(A)$ such that:
\begin{itemize}
\item $G_{-1}(I) = 0$. 

\item Each $\opn{Gr}_q^G(I)$ is a cofree DG $A$-module.

\item $I = \lim_{\lar q} \, G_q(I)$.
\end{itemize}

\item The DG $A$-module $I$ is called a {\em semi-cofree}% 
\index{Differential graded module! semi-cofree}
if it admits a semi-cofree cofiltration. 
\end{enumerate}
\end{dfn}

\begin{prop} \label{prop:3375}
If $I$ is a semi-cofree DG $A$-module, then $I^{\natural}$ is 
an injective object in the abelian category 
$\dcat{G}_{\mrm{str}}(A^{\natural})$. 
\end{prop}

\begin{prop} \label{prop:3185}
Assume $A$ is a ring. If $I$ is a semi-cofree DG $A$-module, then each $I^p$ is 
an injective $A$-module. 
\end{prop}

\begin{exer} \label{exer:3185}
Prove the two propositions above. 
\end{exer}

\begin{thm} \label{thm:1665}
\index{Differential graded module! semi-cofree}
\index{Differential graded module! K-injective}
Let $I$ be an object of $\dcat{C}(A)$. If $I$ is semi-cofree, then it is 
K-injective.
\end{thm}

\begin{proof}
The proof is very similar to those of Theorems \ref{thm:1536} and 
\ref{thm:1575}. But because the arguments involve limits, we shall give the 
full proof. 

\medskip \noindent
Step 1. Suppose $I$ is cofree; say 
$I \cong \prod_{s \in S} \, \opn{T}^{-p_s}(A^*)$. 
Lemma \ref{lem:3347} implies that for every DG $A$-module $N$ there is an 
isomorphism 
\[ \opn{Hom}_A(N, I) \cong 
\prod_{s \in S} \ \opn{Hom}_{\K} \bigl( \opn{T}^{p_s}(N), \K^* \bigr)  \]
of graded $\K$-modules.  It follows that if $N$ is acyclic, then so is 
$\opn{Hom}_A(N, I)$. 

\medskip \noindent
Step 2. Fix a semi-cofree cofiltration
$G = \{ G_q(I) \}_{q \geq -1}$ of $I$.
Here we prove that for every $q \geq -1$ the DG module  
$G_q(I)$ is K-injective. This is done by induction on $q \geq -1$. 
For $q= -1$ it is trivial. For $q \geq 0$ there is an exact sequence 
\begin{equation} \label{eqn:1766}
0 \to \opn{Gr}^G_q(I) \to G_q(I) \to G_{q - 1}(I) \to 0 
\end{equation}
in the category $\dcat{C}_{\mrm{str}}(A)$. 
Because $\opn{Gr}^G_q(I)$ is a cofree DG $A$-module, it is an injective object 
in the abelian category $\bcat{G}_{\mrm{str}}(A^{\natural})$; see Lemma 
\ref{lem:1770}. Therefore the sequence (\ref{eqn:1766}) is split exact in 
$\bcat{G}_{\mrm{str}}(A^{\natural})$.

Let $N \in \dcat{C}(A)$ be an acyclic DG module. Applying the 
functor $\opn{Hom}_{A}(N, -)$ 
to the sequence (\ref{eqn:1766}) we obtain a sequence 
\begin{equation} \label{eqn:1767}
0 \to \opn{Hom}_{A} \bigl( N, \opn{Gr}^G_q(I) \bigr) \to 
\opn{Hom}_{A} \bigl( N, G_q(I) \bigr) \to
\opn{Hom}_{A} \bigl( N, G_{q - 1}(I) \bigr) \to  0 
\end{equation}
in $\dcat{C}_{\mrm{str}}(\K)$. If we forget differentials this is a sequence in 
$\bcat{G}_{\mrm{str}}(\K)$. Because (\ref{eqn:1766}) is split exact in 
$\bcat{G}_{\mrm{str}}(A^{\natural})$, it follows that  (\ref{eqn:1767}) is 
split 
exact in $\bcat{G}_{\mrm{str}}(\K)$. Therefore (\ref{eqn:1767}) is exact in 
$\dcat{C}_{\mrm{str}}(\K)$.

By the induction hypothesis the DG $\K$-module  
$\opn{Hom}_{A} \bigl( N, G_{q - 1}(I) \bigr)$
is \lb acyclic. By step 1 the DG $\K$-module
$\opn{Hom}_{A} \bigl( N, \opn{Gr}^G_{q}(I) \bigr)$
is acyclic. The long exact cohomology sequence associated to 
(\ref{eqn:1767}) shows that the DG $\K$-module 
$\opn{Hom}_{A} \bigl( N, G_q(I) \bigr)$
is acyclic too.

\medskip \noindent
Step 3. We keep the semi-cofree cofiltration 
$G = \{ G_q(I) \}_{q \geq -1}$ from step 2. 
Take any acyclic $N \in \bcat{C}(A)$. By 
Proposition \ref{prop:1535} we know that 
\[ \opn{Hom}_{A}(N, I) \cong 
\lim_{\leftarrow j} \, \opn{Hom}_{A} \bigl( N, G_q(I) \bigr) \]
in $\dcat{C}_{\mrm{str}}(\K)$. According to step 2 the complexes 
$\opn{Hom}_{A} \bigl( N, G_q(I) \bigr)$ are all acyclic. 
The exactness of the sequences (\ref{eqn:1767}) implies that 
the inverse system \lb 
$\bigl\{ \opn{Hom}_{A} \bigl( N, G_q(I) \bigr) \bigr\}_{q \geq -1}$ 
in $\dcat{C}_{\mrm{str}}(\K)$ has surjective transitions. Now the 
Mittag-Leffler 
argument (Corollary \ref{cor:1590}) says that the inverse limit complex
$\opn{Hom}_{A}(N, I)$ is acyclic. 
\end{proof}

\begin{thm} \label{thm:3325}
\index{Differential graded module! semi-cofree}
Let $A$ be a DG ring. Every DG $A$-module $M$ admits a quasi-isomorphism
$\rho : M \to I$ in $\dcat{C}_{\mrm{str}}(A)$ to a semi-cofree DG $A$-module 
$I$.
\end{thm}

The proof of this theorem was communicated to us by B. Keller. We need a lemma 
first.

The semi-free DG $A$-module $C$ was defined in formula (\ref{eqn:3307}). 
Its dual DG $A$-module $C^* = \opn{Hom}_{\K}(C, \K^*)$  
is a semi-cofree DG $A$-module, with semi-cofree cofiltration 
$G_q(C^*) := F_q(C)^*$,
where $F_q(C)$ is from (\ref{eqn:3332}). 

\begin{lem} \label{lem:3325}
Let $M \in \dcat{C}(A)$.
\begin{enumerate}
\item There is a homomorphism $\psi : M \to J$ 
in $\dcat{C}_{\mrm{str}}(A)$, such that 
$\opn{Z}(\psi) : \opn{Z}(M) \to \opn{Z}(J)$ is injective,
and $J = \prod_{s \in S} \opn{T}^{-p_s}(A^*)$
for a collection of integers $\{ p_s \}_{s \in S}$. 

\item There is an injective homomorphism
$\psi' : M \to J'$ in $\dcat{C}_{\mrm{str}}(A)$,
such that 
$J' = \prod_{s \in S'} \, \opn{T}^{-p_s}(C^*)$
for some collection of integers $\{ p_s \}_{s \in S'}$.
\end{enumerate}
\end{lem}

To clarify the notation in the lemma: we have two distinct collections 
of integers, $\{ p_s \}_{s \in S}$ and 
$\{ p_s \}_{s \in S'}$. 

\begin{proof} \mbox{}

\smallskip \noindent
(1) By Lemma \ref{lem:3346} there is an injective homomorphism 
$\chi : \opn{Y}(M) \inj W$ 
for some cofree graded $\K$-module 
$W = \prod_{s \in S} \, \opn{T}^{-p_s}(\K^*)$.
Let $J := I_W$. 
According to Lemma \ref{lem:3348} there is a homomorphism 
$\psi : M \to J$ in $\dcat{C}_{\mrm{str}}(A)$ s.t.\ 
$\opn{Y}(\th_W) \circ \opn{Y}(\psi) = \chi$. Because $\chi$ is injective, so is
$\opn{Y}(\psi)$. 

\medskip \noindent
(2) Consider the DG $A^{\mrm{op}}$-module $M^*$. By Lemma \ref{lem:3306}
there is a collection of integers $\{ p_s \}_{s \in S'}$,
and a surjective homomorphism
$\phi : Q \to M^*$
in $\dcat{C}_{\mrm{str}}(A^{\mrm{op}})$,
where 
$Q := \bigoplus\nolimits_{s \in S'} \, \opn{T}^{p_s}(C)$. 
Dualizing we get a DG $A$-module 
$J' := Q^* \cong  \prod_{s \in S'} \opn{T}^{-p_s}(C^*)$, 
and an injective homomorphism
$\phi^* : M^{**} \to Q^* = J'$
in $\dcat{C}_{\mrm{str}}(A)$.
Composing $\phi^*$ with the canonical embedding 
$M \inj M^{**}$ gives us $\psi' : M \inj J'$. 
\end{proof}

The DG $A$-module $J$ in item (1) of the lemma is cofree. The DG $A$-module 
$J'$ in item (2) is semi-cofree, with semi-cofree cofiltration induced by that 
of $C^*$, namely 
\begin{equation} \label{eqn:3326}
G_{q}(J') :=  \prod_{s \in S'} \, \opn{T}^{-p_s}(G_q(C^*)) .
\end{equation}

\begin{proof}[Proof of Theorem \tup{\ref{thm:3325}}]
\mbox{}

\smallskip \noindent
Step 1. We are going to produce an exact sequence 
\begin{equation} \label{eqn:3338}
0 \to M \xar{\eta} J^0 \xar{\pa^0} J^1 \xar{\pa^1} J^2 \to \cdots  
\end{equation}
in $\dcat{C}_{\mrm{str}}(A)$ with these properties: 
\begin{itemize}
\rmitem{a} The sequence 
\[ 0 \to \opn{Y}(M) \xar{\opn{Y}(\eta)} \opn{Y}(J^0) \xar{\opn{Y}(\pa^0)} 
\opn{Y}(J^1) \xar{\opn{Y}(\pa^1)} \opn{Y}(J^2) \to \cdots \]
is exact in $\dcat{G}_{\mrm{str}}(\K)$.

\rmitem{b} For every $p \geq 0$ the DG $A$-module $J^p$ has the following 
decomposition 
in $\dcat{C}_{\mrm{str}}(A)$~: 
$J^p = J_p \oplus J'_p$, where
$J_p = \prod_{s \in S_p} \opn{T}^{-q_s}(A^*)$
and 
$J'_p = \lb \prod_{s \in S'_p} \opn{T}^{-q_s}(C^*)$ 
for some collection of integers 
$\{ q_s \}_{s \in S_p \sqcup S'_p}$. 
\end{itemize}
This will be done recursively on $p$. 

By Lemma \ref{lem:3325} we can find DG $A$-modules $J_0, J'_0$ 
and homomorphisms 
$\psi_0 : M \to J_0$, $\psi'_0 : M \to J'_0$, such that both $\opn{Y}(\psi_0)$
and $\psi'_0$ are injective. The DG modules $J_0, J'_0$ are of the required 
form (property (b) with $p = 0$). We let 
$J^0 := J_0 \oplus J'_0$ and $\eta := \psi_0 \oplus \psi'_0$. 

Let $N^0 := \opn{Coker}(\eta)$. Then the sequence 
\begin{equation} \label{eqn:3329}
0 \to M \xar{\eta} J^0 \to N^0 \to 0 
\end{equation}
in $\dcat{C}_{\mrm{str}}(A)$ is exact. The sequence 
\[ 0 \to \opn{Y}(M) \xar{\opn{Y}(\eta)} \opn{Y}(J^0) \to \opn{Y}(N^0) \to  0 \]
in $\dcat{G}_{\mrm{str}}(\K)$ is also exact. Indeed, the exactness at 
$\opn{Y}(J^0)$ and $\opn{Y}(N^0)$ is because the functor 
$\opn{Y} : \dcat{C}_{\mrm{str}}(\K) \to \dcat{G}_{\mrm{str}}(\K)$
is right exact (see Proposition \ref{prop:4145}), and the sequence 
(\ref{eqn:3329}) is exact. 
The exactness at $\opn{Y}(M)$ is by the injectivity of 
$\opn{Y}(\psi_0) : \opn{Y}(M) \to \opn{Y}(J_0)$.

Next we repeat this procedure with $N^0$ instead of $M$, to obtain an exact 
sequence 
\begin{equation} \label{eqn:3330}
0 \to N^0 \xar{\pa^0} J^1 \to N^1 \to 0 ,
\end{equation}
with $J^1$ of the required form (property (b) with $p = 1$), and such that 
\[ 0 \to \opn{Y}(N^0) \xar{\opn{Y}(\pa^0)} \opn{Y}(J^1) \to \opn{Y}(N^1) \to  
0 \]
is also exact. We then splice (\ref{eqn:3329}) and (\ref{eqn:3330})
to get the exact sequence 
\[ 0 \to M \xar{\eta} J^0 \xar{\pa^0} J^1 \to N^1 \to 0 . \]

Continuing recursively we obtain an exact sequence (\ref{eqn:3338}) that has 
properties (a) and (b). 

\medskip \noindent 
Step 2. We view the exact sequence (\ref{eqn:3338}) as an acyclic complex with 
entries in $\dcat{C}_{\mrm{str}}(A)$, that has $J^0$ in degree $0$. 
Define the total DG $A$-modules 
\[ I :=\opn{Tot}^{\Pi} \bigl( \cdots \to 0 \to
J^0 \xar{\pa^0} J^1 \xar{\pa^1} J^2 \to \cdots  \bigr) \]
and 
\[ I^{\mrm{aug}} := \opn{Tot}^{\Pi} \bigl( \cdots \to 0 \to M \xar{\eta} 
J^0 \xar{\pa^0} J^1 \xar{\pa^1} J^2 \to \cdots  \bigr) . \]

Because the sequence (\ref{eqn:3338}) and the sequence appearing in property 
(a) are exact, we can use Proposition \ref{prop:3321} (with any 
$j_0 \in \Z$). The conclusion is that the DG $A$-module $I^{\mrm{aug}}$ is 
acyclic. Then, by Corollary \ref{cor:3320}, the homomorphism 
$\rho : M \to I$ in $\dcat{C}_{\mrm{str}}(A)$ is a quasi-isomorphism. 

\medskip \noindent
Step 3. It remains to produce a semi-cofree cofiltration 
$\{ G_q(I) \}_{q \geq -1}$ on the DG $A$-module $I$. Our formula is this: 
$G_{-1}(I) := 0$ of course. For $r \geq 0$ we let 
\[ G_{2 \cd r}(I) := 
\Bigr( \bigoplus\nolimits_{p = 0}^{r - 1} \, \opn{T}^{-p} (J^p) \Bigl) 
\, \oplus \, \opn{T}^{-r} \bigl( J_r \oplus G_0(J'_r) \bigr) , \]
where $G_0(J'_r)$ comes from (\ref{eqn:3326}), and 
\[ G_{2 {\cdot} r + 1}(I)  := 
\bigoplus\nolimits_{p = 0}^{r} \, \opn{T}^{-p} (J^p) . \]
We leave it to the reader to verify that this is a semi-cofree cofiltration. 
\end{proof}

\begin{cor} \label{cor:1665}
Let $A$ be a DG ring. The category $\bcat{C}(A)$ has 
enough K-in\-jectives%
\index{Resolution! K-injective}.
\end{cor}

\begin{proof}
Combine Theorems \ref{thm:1665} and \ref{thm:3325}. 
\end{proof}

Recall that a DG ring $A$ is nonpositive if $A^i = 0$ for all $i > 0$. 
When $A$ is nonpositive, the DG module $A^*$ is concentrated in 
degrees $\geq 0$. 

\begin{cor} \label{cor:3345}
If the DG ring $A$ is nonpositive, then every $M \in \bcat{C}(A)$ admits a 
quasi-iso\-morphism $\rho : M \to I$ in $\dcat{C}_{\mrm{str}}(A)$ to a 
semi-cofree DG $A$-module $I$, such that 
$\opn{inf}(I) = \opn{inf}(\opn{H}(M))$. 
\end{cor}


\begin{proof}
Let $p_0 := \opn{inf}(\opn{H}(M))$. If $p_0 = \pm \infty$ then there is 
nothing to prove beyond Theorem \ref{thm:3325}; so let us assume 
that $p_0 \in \Z$. 
Define $M' := \opn{smt}^{\geq p_0}(M)$. Then $M \to M'$ is a quasi-isomorphism, 
and $\opn{inf}(M') = p_0$. By replacing $M$ with $M'$ we can now assume that 
$\opn{inf}(M) = p_0$.

Since $\opn{inf}(M) = p_0$ it follows that 
$M^{p_0} = \opn{Y}^{p_0}(M)$. This means that in step 1 of the proof of the 
theorem, the homomorphism 
$\psi_0 : M^{p_0} \to J^{p_0}_0$ is already injective. So, like in the proof of 
Corollary \ref{cor:3340}, we can arrange to have 
$\opn{inf}(J^0) \geq p_0$. Then, recursively, we can arrange to 
have $\opn{inf}(J^q) \geq p_0$ for all $q$. Therefore in step 2 of the proof we 
get $\opn{inf}(I) \geq p_0$. 

Finally, because $\opn{inf}(\opn{H}(M)) = p_0$ we must have 
$\opn{inf}(I) = p_0$.
\end{proof}

We end with the K-injective analogue of Remark \ref{rem:4675}.

\begin{rem} \label{rem:4676}
Suppose $(X, \AA)$ is a ringed space. We know that the abelian category 
$\dcat{M}(\AA) = \cat{Mod} \AA$ of sheaves of $\AA$-modules on $X$ has 
enough injective objects. Theorem \ref{thm:1630} tells us that the subcategory 
$\dcat{C}^{+}(\AA)$ has enough K-injective objects. But in fact more is true: 
the category of unbounded complexes $\dcat{C}^{}(\AA)$ also has enough 
K-injective objects; see \cite{Spa}. (There is another, more abstract proof of 
this statement: because $\dcat{M}(\AA)$ is a Grothendieck abelian category, it 
follows that $\dcat{C}^{}(\AA)$ has enough K-injectives. See 
\cite[Corollary 14.1.8]{KaSc2}.)

Now let's consider a DG ringed space $(X, \AA)$. Fix an injective cogenerator 
$\K^*$ of $\dcat{M}(\K)$. 
For a point $x \in X$, consider the cofree DG $\AA_x$-module 
$I_x := \opn{Hom}_{\K}(\AA_x, \K^*)$.
Pushing this module forward along the inclusion $\{ x \} \sub X$ we get a 
constant sheaf $\II_x$ with support the closed set 
$\ol{\{ x \}}$, that we refer to as the cofree DG $\AA$-module 
associated to $x$. By definition, a {\em cofree DG $\AA$-module} is a product 
of sheaves $\II_x$ (for varying points $x \in X$).  
Combining this notion with that of a semi-cofree DG module, we can talk about 
{\em semi-cofree DG  $\AA$-modules}. Like Theorem \ref{thm:1665}, one can show 
that a semi-cofree DG  $\AA$-module $\II$ is K-injective. 

Assume that the topological space $X$ is finite dimensional, in the sense that 
there is a natural number $d$ such that for every open set $U \sub X$, every 
sheaf of $\K$-module $\NN$ on $U$, and every $p > d$, the sheaf cohomology
$\opn{H}^p(U, \NN)$ vanishes. 
This condition is quite mild -- it is satisfied when $X$ is a finite 
dimensional noetherian scheme, and also when $X$ is a subspace of a finite 
dimensional topological manifold. 
Then a geometric variant of Theorem \ref{thm:3325} is 
true: every DG  $\AA$-module $\MM$ admits a quasi-isomorphism 
$\MM \to \II$ to a semi-cofree DG  $\AA$-module $\II$.
The finite dimensionality of $X$ is needed to invoke the ML argument. 
A detailed proof of this assertion will be published in the future. 
\end{rem}





\cleardoublepage
\mysection{Adjunctions, Equivalences and Cohomological Dimension} 
\label{sec:adj-equ-cohdim}

\AYcopyright

In this section we discuss the derived Hom and tensor bifunctors in several 
situations, and show how these bifunctors are related in adjunction 
formulas. Cohomological dimensions of functors and DG modules are introduced, 
and they are used to prove several theorems. 
Given a homomorphism between two DG rings, we study the derived restriction, 
induction and coinduction functors between the derived categories of these DG 
rings. We prove that for a quasi-isomorphism of DG rings the restriction 
functor is an equivalence. We also prove the existence of NC semi-free DG ring 
resolutions. 

As before, we work over a nonzero commutative base ring $\K$, and Convention 
\ref{conv:2490} is in place. 



\mysubsection{Boundedness Conditions Revisited} 
\label{subsec:bounded-revis}

This subsection is about notation. We start by introducing some combinatorial 
definitions that will enable a precise treatment of cohomological dimensions of 
objects and functors (in Subsection \ref{subsec:way-out}). 

We also change our notation for boundedness conditions in derived 
categories (Definition \ref{dfn:2122}) and for the translation functor 
(Definition \ref{dfn:2120}). These changes will simplify and improve our 
discussion (while hopefully not creating confusion). See Remark \ref{rem:4765} 
for an explanation. 

{\em Integer intervals} were introduced in Definition \ref{dfn:4765}. 
{\em Boundedness conditions} were introduced in Subsections 
\ref{subsec:DGModinM} and  \ref{subsec:bd-in-homot}.
We need to refine these notions. 

\begin{dfn} \label{dfn:4780}
\index{Boundedness! condition}
Let $S$ be an integer interval.
\begin{enumerate}
\item We call $S$ a {\em bounded above integer interval} if either 
$S = [i_0, i_1]$ with $i_1 \in \Z$, or if $S = \varnothing$. 

\item We call $S$ a {\em bounded below integer interval} if either 
$S = [i_0, i_1]$ with $i_0 \in \Z$, or if $S = \varnothing$.

\item We call $S$ a {\em bounded integer interval} if it is both bounded above 
and bounded below; i.e.\ either $S = [i_0, i_1]$ with $i_0, i_1 \in \Z$, or 
$S = \varnothing$.
\end{enumerate}
\end{dfn}

The next definition was already mentioned in passing in 
Subsection \ref{subsec:DGModinM}.

\begin{dfn} \label{dfn:4781}
The boundedness conditions in Definition \ref{dfn:4780} have abbreviations, 
that we call {\em boundedness indicators}%
\index{Boundedness! indicator}.
\begin{enumerate}
\item The indicator ``$-$'' stands for ``bounded above''.

\item The indicator ``$+$'' stands for ``bounded below''.

\item The indicator ``$\mrm{b}$'' stands for ``bounded''.

\item The indicator ``$\bra{\mrm{empty}}$'' (the empty symbol)
stands for the empty condition (no condition at all). 
\end{enumerate}
\end{dfn}

The symbol ``$\star$'' is used to denote an unspecified boundedness 
indicator, i.e.\ it is a placeholder for one of the boundedness indicators in 
Definition \ref{dfn:4781}.

\begin{dfn} \label{dfn:4782}
Let $S$ be an integer interval.
\begin{enumerate}
\item If $S =  [i_0, i_1]$, then its {\em length}% 
\index{Integer interval! length of}
is $i_1 - i_0 \in \N \cup \{ \infty \}$. If $S = \varnothing$, then then its 
length is $-\infty$.  

\item If $S =  [i_0, i_1]$, then the {\em reversed integer interval}%
\index{Integer interval! reversed}
is $-S := [-i_1, -i_0]$.
If $S = \varnothing$, then its reverse is also $\varnothing$.

\item If $S =  [i_0, i_1]$, and if $T = [j_0, j_1]$ is a second nonempty 
integer interval, then their sum% 
\index{Integer interval! sum of {\indash}s}
is $S + T := [i_0 + j_0, i_1 + j_1]$.
If $S = \varnothing$, and it $T$ is a second integer interval, then the sum is
$S + T := \varnothing$. 
\end{enumerate}
\end{dfn}

Thus an integer interval $S$ (whether empty or not) has finite length if and 
only if it is bounded. 

\begin{dfn}  \label{dfn:2125}
Let $M = \{ M^i \}_{i \in \Z}$ be a graded object in an abelian category 
$\cat{M}$, i.e.\ $M \in \dcat{G}(\cat{M})$. 
\begin{enumerate}
\item We say that $M$ is {\em concentrated} in an integer interval 
$S$ if the set \lb $\{ i \in \Z \mid M^i \neq 0 \}$ is contained in $S$.

\item The {\em concentration}%
\index{Concentration! of graded module}
\index{1-Con(M)@$\opn{con}(M)$}
of $M$ is the smallest integer interval 
$\opn{con}(M)$ in which $M$ is concentrated.

\item We say that $M$ is of {\em boundedness type $\star$},
for some boundedness indicator $\star$
%\index{Boundedness! condition on graded module}, 
if the integer interval $\opn{con}(M)$ satisfies the boundedness condition 
$\star$.
\end{enumerate}
\end{dfn}

Observe that  $\opn{con}(M) = \varnothing$ if and only if $M = 0$. 
Definition \ref{dfn:2125}(3) agrees with the boundedness 
conditions on a graded module $M$ from Subsections 
\ref{subsec:DGModinM} and  \ref{subsec:bd-in-homot}.

In Definition \ref{dfn:3155} we introduced the notions $\inf(M)$ and $\sup(M)$ 
for a graded object $M$. Using them we can describe the 
concentration $\opn{con}(M) = [i_0, i_1]$ when $M$ is nonzero:
$i_0 = \opn{inf} (M)$ and $i_1 = \opn{sup} (M)$.
The amplitude $\opn{amp}(M)$ is the length of the interval $\opn{con}(M)$.

Recall that to a DG module $M \in \dcat{C}^{}(A, \cat{M})$ we associate
the graded object $M^{\natural} \in \dcat{G}(\cat{M})$ that 
is gotten by forgetting the differential. The cohomology $\opn{H}(M)$
is another  object of $\dcat{G}(\cat{M})$. 

\begin{dfn} \label{dfn:4785}
Let  $A$ be a DG ring, $\cat{M}$ an abelian category, and 
$M \in \dcat{C}^{}(A, \cat{M})$.
\begin{enumerate}
\item The {\em concentration} of $M$% 
\index{Concentration! of DG module}
and the {\em boundedness type}% 
\index{Boundedness! type of DG module}
of $M$ are the same as in Definition \ref{dfn:2125}, when $M$ is seen as an 
object of $\dcat{G}^{}(\cat{M})$.

\item The {\em cohomological concentration}% 
\index{Concentration! cohomological {\indash} of DG module}
and the {\em cohomological boundedness type} of $M$%
\index{Boundedness! type of DG module}
are the concentration and boundedness type, respectively, of the 
graded object $\opn{H}(M) \in \dcat{G}^{}(\cat{M})$,
in the sense of Definition \ref{dfn:2125}.
\end{enumerate}
\end{dfn}

The next definition is in conflict with Definitions \ref{dfn:1460} and 
\ref{dfn:1461}; but we already warned (in Remark \ref{rem:1461}) that this 
change will take place. See Remark \ref{rem:4765} for further discussion. 

\begin{dfn} \label{dfn:2122}
Let $A$ be a DG ring, $\cat{M}$ an abelian category, and $\star$ a boundedness 
indicator. 
We denote by $\dcat{D}^{\star}(A, \cat{M})$% 
\index{1-D(A,M)@$\dcat{D}^{\star}(A, \cat{M})$}%
\index{1-D(M)@$\dcat{D}^{\star}(\cat{M})$}%
\index{Boundedness! condition}%
\index{Concentration! cohomological {\indash} of DG module}
the  full subcategory of $\dcat{D}(A, \cat{M})$ 
on the DG modules $M$ of cohomological boundedness type $\star$, as in 
Definition \ref{dfn:4785} above. 
\end{dfn}

Thus, for example,  a DG module $M$ belongs to 
$\dcat{D}^{\mrm{b}}(A, \cat{M})$ if and only if 
$\opn{con}(\opn{H}(M))$ is a bounded integer interval. 
As explained in subsection \ref{subsec:bd-in-homot}, 
$\dcat{D}^{\star}(A, \cat{M})$ is a full triangulated subcategory of 
$\dcat{D}^{}(A, \cat{M})$.

The notations $\dcat{C}^{\star}(A, \cat{M})$ and 
$\dcat{K}^{\star}(A, \cat{M})$
remain the same, as in Definitions \ref{dfn:3220} and \ref{dfn:3221}. 
Namely, these are the full subcategories of $\dcat{C}^{}(A, \cat{M})$ 
and $\dcat{K}^{}(A, \cat{M})$ respectively on the DG modules $M$ 
of boundedness type $\star$.%
\index{1-C(A,M)@$\dcat{C}^{\star}(A, \cat{M})$}%
\index{1-K(A,M)@$\dcat{K}^{\star}(A, \cat{M})$}%
\index{1-C(M)@$\dcat{C}^{\star}(\cat{M})$}%
\index{1-K(M)@$\dcat{K}^{\star}(\cat{M})$}

\begin{dfn} \label{dfn:2120}
Let $A$ be a DG ring and $\cat{M}$ an abelian category. 
For a DG module $M \in \dcat{C}(A, \cat{M})$ and an integer $i$, we write 
$M[i] := \opn{T}^i(M)$,
the $i$-th translation of $M$, 
as in Definition \ref{dfn:1170}.%
\index{1-M[i]@$M[i]$}%
\index{1-Ti(M)@$\opn{T}^i(M)$}%
\index{Translation! of DG module}

This notation applies also to the homotopy category $\dcat{K}(A, \cat{M})$,
the derived category $\dcat{D}(A, \cat{M})$, and every other T-additive 
category. 
\end{dfn}

The notation $M[i]$ makes it difficult to use the little t operator, and to 
talk about translation isomorphisms, but hopefully we won't require them 
anymore.

\begin{rem} \label{rem:4765}
We should say a few explanatory words on the change in Definition 
\ref{dfn:2122}, that conflicts with earlier definitions. 

Recall that in Definitions \ref{dfn:1460} and \ref{dfn:1461} we made a 
distinction between 
\begin{equation} \label{4770}
\dcat{D}^{\star}(A, \cat{M}) = 
\dcat{K}^{\star}(A, \cat{M})_{\dcat{S}^{\star}(A, \cat{M})} ,
\end{equation}
the localization of $\dcat{K}^{\star}(A, \cat{M})$ with respect to the 
quasi-isomorphisms in it, and $\dcat{D}(A, \cat{M})^{\star}$, 
the full subcategory of $\dcat{D}(A, \cat{M})$ on the complexes of 
cohomological boundedness type $\star$. This distinction was necessary in the 
early part of the book, when we were studying foundations and resolutions.
Now that the foundations are behind us, we no longer have a need for the 
subcategory $\dcat{D}^{\star}(-)$ in the sense of formula (\ref{4770}), i.e.\ 
the definition based on boundedness type. 

Still, one might ask: Why make the notational switch from 
$\dcat{D}(-)^{\star}$ to $\dcat{D}^{\star}(-)$, for the definition based on 
cohomological boundedness type? 
There are two reasons for that. The first is typographic: 
the notation $\dcat{D}(-)^{\star}$ becomes cumbersome when 
coupled with other indicators, such as
$\dcat{D}_{\mrm{f}}(-)^{\star}$
or $(\dcat{D}(-)^{\star})^{\mrm{op}}$.
After Definition \ref{dfn:2122} we can refer to these categories by the more 
elegant expressions $\dcat{D}^{\star}_{\mrm{f}}(-)$
or $\dcat{D}^{\star}(-)^{\mrm{op}}$, respectively. 

Second, the expression $\dcat{D}(-)^{\star}$ is a novelty, and nobody else uses 
it, whereas the notation $\dcat{D}^{\star}(-)$ is common. According 
to Proposition \ref{prop:1463}, in most situations the two distinct meanings of
$\dcat{D}^{\star}(-)$, the old one and the new one, refer to nested 
subcategories of $\dcat{D}(-)$ whose inclusion is an equivalence. Furthermore, 
in a few recent texts on derived categories (such as 
\cite[Definition tag=05RU]{SP} and \cite[Definition 1.11]{Ye8}) we find the 
same usage of the expression $\dcat{D}^{\star}(-)$ as in Definition 
\ref{dfn:2122}.
\end{rem}



\mysubsection{The Bifunctor \texorpdfstring{$\opn{RHom}$}{RHom}}   
\label{subsec:RHom}

Consider a DG ring $A$ and an abelian category $\cat{M}$.
Like in Example \ref{exa:1967} we get a DG bifunctor 
\[ F := \opn{Hom}_{A, \cat{M}}(-, -) : \dcat{C}(A, \cat{M})^{\mrm{op}} \times
\dcat{C}(A, \cat{M}) \to \dcat{C}(\K) . \]
In Definition \ref{dfn:2995} we put a triangulated structure on the opposite 
homotopy category $\dcat{K}(A, \cat{M})^{\mrm{op}}$. 
According to Theorem \ref{thm:4225}, there is an induced  
triangulated bifunctor 
\[ F = \opn{Hom}_{A, \cat{M}}(-, -) : \dcat{K}(A, \cat{M})^{\mrm{op}} \times
\dcat{K}(A, \cat{M}) \to \dcat{K}(\K) . \]
Postcomposing with the localization functor 
$\opn{Q} : \dcat{K}(\K) \to \dcat{D}(\K)$, 
we obtain a triangulated bifunctor 
\[ F = \opn{Hom}_{A, \cat{M}}(-, -) : \dcat{K}(A, \cat{M})^{\mrm{op}} \times
\dcat{K}(A, \cat{M}) \to \dcat{D}(\K) . \]
Next we pick full additive subcategories 
$\cat{K}_1, \cat{K}_2 \sub \dcat{K}(A, \cat{M})$
s.t.\ $\cat{K}_1^{\mrm{op}} \sub \dcat{K}(A, \cat{M})^{\mrm{op}}$
and
$\cat{K}_2 \sub \dcat{K}(A, \cat{M})$
are triangulated. In practice this choice would be by some boundedness 
conditions; for instance
$\cat{K}_1 := \dcat{K}^{-}(\cat{M})$
or $\cat{K}_2 := \dcat{K}^{+}(\cat{M})$,  
cf.\ Corollaries \ref{cor:1900} and \ref{cor:1630} respectively. 
We want to construct the right derived 
bifunctor of the triangulated bifunctor 
\[ F = \opn{Hom}_{A, \cat{M}}(-, -) : \cat{K}_1^{\mrm{op}} \times
\cat{K}_2 \to \dcat{D}(\K) . \]
This is done in the next theorem.

\begin{thm} \label{thm:2005}
Let $A$ be a DG $\K$-ring, let $\cat{M}$ be a $\K$-linear abelian 
category, and let $\cat{K}_1, \cat{K}_2 \sub \dcat{K}(A, \cat{M})$ 
be full additive subcategories, such that $\cat{K}_1^{\mrm{op}}$ is a 
full triangulated subcategory of $\dcat{K}(A, \cat{M})^{\mrm{op}}$, 
and $\cat{K}_2$ is a full triangulated subcategory of 
$\dcat{K}(A, \cat{M})$.
Assume either that $\cat{K}_1$ has enough K-projectives, or that $\cat{K}_2$ 
has 
enough K-injectives.
Let $\cat{D}_i$ denote the localization of $\cat{K}_i$ with respect to the 
quasi-isomorphisms in it. 
Then the triangulated bifunctor 
\[  \opn{Hom}_{A, \cat{M}}(-, -) : \cat{K}_1^{\mrm{op}} \times
\cat{K}_2 \to \dcat{D}(\K) \]
has a right derived bifunctor%
\index{1-RHom(A,M)@$\opn{RHom}_{A, \cat{M}}$}
\[  \opn{RHom}_{A, \cat{M}}(-, -) : \cat{D}_1^{\mrm{op}} \times
\cat{D}_2 \to \dcat{D}(\K) . \]
Moreover, if $P_1 \in \cat{K}_1$ is K-projective, or if 
$I_2 \in \cat{K}_2$ is K-injective, then the morphism 
\[ \eta^{\mrm{R}}_{P_1, I_2} : \opn{Hom}_{A, \cat{M}}(P_1, I_2) \to 
\opn{RHom}_{A, \cat{M}}(P_1, I_2) \]
in $\dcat{D}(\K)$ is an isomorphism. 
\end{thm}

\begin{proof}
If $\cat{K}_2$ has enough K-injectives, then we can take 
$\cat{J}_2 := \cat{K}_{2, \mrm{inj}}$, the full subcategory on the K-injectives 
inside $\cat{K}_2$. And we take 
$\cat{J}_1 := \cat{K}_{1}$. 
We claim that the conditions of Theorem \ref{thm:2000} are satisfied. Condition 
(b) is simply the assumption that $\cat{K}_2$ has enough K-injectives. As for 
condition (a): this is Lemma 
\ref{lem:2000} below, with its condition (ii). 

On the other hand, if $\cat{K}_1$ has enough K-projectives, then we can take 
$\cat{J}_1 := \cat{K}_{1, \mrm{prj}}$, 
the full subcategory on the K-projectives inside $\cat{K}_1$. 
And we take $\cat{J}_2 := \cat{K}_{2}$.
We claim that the conditions of Theorem \ref{thm:2000}  are satisfied for 
$\cat{J}_1^{\mrm{op}} \sub \cat{K}_1^{\mrm{op}}$. 
Condition (b) is simply the assumption that $\cat{K}_1$ has enough 
K-projectives: a quasi-isomorphism 
$\rho : P \to M$ in $\cat{K}_1$ becomes a quasi-isomorphism 
$\mrm{Op}(\rho) : \mrm{Op}(M) \to \mrm{Op}(P)$ in 
$\cat{K}_1^{\mrm{op}}$, and 
$\mrm{Op}(P) \in \cat{J}_1^{\mrm{op}}$.
As for condition (a): this is Lemma \ref{lem:2000} below,  with its condition 
(i). 

The last assertion also follows from \ref{lem:2000}.
\end{proof}

\begin{lem}  \label{lem:2000}
Suppose $\phi_1 : Q_1 \to P_1$ and $\phi_2 : I_2 \to J_2$
are quasi-iso\-morphisms in $\dcat{C}(A, \cat{M})$, 
and either of the conditions below holds\tup{:}
\begin{enumerate}
\rmitem{i} $Q_1$ and $P_1$ are both K-projective. 

\rmitem{ii} $I_2$ and $J_2$ are both K-injective.
\end{enumerate}
Then the homomorphism
\[  \opn{Hom}_{A, \cat{M}}(\phi_1, \phi_2) : 
\opn{Hom}_{A, \cat{M}}(P_1, I_2) \to \opn{Hom}_{A, \cat{M}}(Q_1, J_2) \]
in $\dcat{C}(\K)$ is a quasi-isomorphism. 
\end{lem}

\begin{proof}
We will only prove the case where $Q_1, P_1$ are both K-projective; the other 
case is very similar. 

The homomorphism in question factors as follows:
\[ \opn{Hom}_{A, \cat{M}}(\phi_1, \phi_2) = 
\opn{Hom}_{A, \cat{M}}(\phi_1, \opn{id}_{J_2})
\circ 
\opn{Hom}_{A, \cat{M}}(\opn{id}_{P_1}, \phi_2) . \]
It suffices to prove that each of the factors is a quasi-isomorphism. This can 
be done by a messy direct calculation, but we will provide an indirect proof 
that relies on properties of the homotopy category 
$\cat{K} := \dcat{K}(A, \cat{M})$ that 
were already established. 

Let $K_2$ be the standard cone on the homomorphism 
$\phi_2 : I_2 \to J_2$. So $K_2$ is acyclic. Because $P_1$ is K-projective it 
follows that $\opn{Hom}_{A, \cat{M}}(P_1, K_2)$ is acyclic. 
Thus for every integer $l$ we have 
\begin{equation} \label{eqn:2046}
 \opn{Hom}_{\cat{K}}(P_1[-l], K_2) \cong 
\opn{H}^{l} \bigl( \opn{Hom}_{A, \cat{M}}(P_1, K_2) \bigr) = 0 . 
\end{equation}
Next, there is a distinguished triangle 
\begin{equation} \label{eqn:2045}
I_2 \xar{\phi_2} J_2 \xar{\be_2} K_2 \xar{\ga_2} I_2[1]
\end{equation}
in $\cat{K}$. Applying the cohomological functor
$\opn{Hom}_{\cat{K}}(P_1[-l], -)$
to the distinguished triangle (\ref{eqn:2045}) yields a long exact sequence, as 
explained in Subsection \ref{subsec:tr-coh-funcs}. From it we deduce that the 
homomorphisms 
\[ \opn{Hom}_{\cat{K}}(P_1[-l], I_2) \to 
\opn{Hom}_{\cat{K}}(P_1[-l], J_2) \]
are bijective for all $l$. Using the isomorphisms like (\ref{eqn:2046}) for 
$I_2$ and $J_2$ we see that 
\[ \opn{Hom}_{A, \cat{M}}(\opn{id}_{P_1}, \phi_2) : 
\opn{Hom}_{A, \cat{M}}(P_1, I_2) \to 
\opn{Hom}_{A, \cat{M}}(P_1, J_2) \]
is a quasi-isomorphism. 

According to Corollary \ref{cor:1923} the homomorphism 
$\phi_1 : Q_1 \to P_1$ is a homotopy equivalence; so it is an isomorphism in 
$\cat{K}$. Therefore for every integer $l$ the homomorphism 
\[ \opn{Hom}_{\cat{K}}(Q_1, J_2[l]) \to 
\opn{Hom}_{\cat{K}}(P_1, J_2[l]) \]
is bijective. As above we conclude that
\[ \opn{Hom}_{A, \cat{M}}(\phi_1, \opn{id}_{J_2}) : 
\opn{Hom}_{A, \cat{M}}(Q_1, J_2) \to 
\opn{Hom}_{A, \cat{M}}(P_1, J_2) \]
is a quasi-isomorphism. 
\end{proof}

\begin{rem} \label{rem:2040}
Theorem \ref{thm:2005} should be viewed as a template. It has 
commutative and noncommutative variants, that we shall study later. 
And there are geometric variants in which the source and target are categories 
of sheaves.
\end{rem}

We end this section with the connection between RHom and morphisms in the 
derived category. 

\begin{dfn} \label{dfn:2121}
Under the assumptions of Theorem \tup{\ref{thm:2005}},  
for DG modules $M_1 \in \cat{K}_1$ and $M_2 \in \cat{K}_2$,
and for an integer $i$, we write% 
\index{1-ExtiA@$\opn{Ext}^i_{A, \cat{M}}(-, -)$}
\[ \opn{Ext}^i_{A, \cat{M}}(M_1, M_2) := 
\opn{H}^i \bigl( \opn{RHom}_{A, \cat{M}}(M_1, M_2) \bigr) \in \dcat{M}(\K) . \]
\end{dfn}

\begin{exer} \label{exer:2125}
Let $A$ be a ring. Prove that for modules 
$M_1, M_2 \in \dcat{M}(A)$ the $\K$-module
$\opn{Ext}^i_{A}(M_1, M_2)$ defined above is canonically isomorphic to the 
module in the classical definition. Moreover, this is true regardless of the 
choices of the subcategories $\cat{K}_1 $ and $\cat{K}_2$,
as long as $M_i \in \cat{K}_i$. 
\end{exer}

\begin{cor} \label{cor:2120}
Under the assumptions of Theorem \tup{\ref{thm:2005}}, there is an  
isomorphism 
\[  \opn{Ext}^0_{A, \cat{M}}(-, -) \iso 
\opn{Hom}_{\dcat{D}(A, \cat{M})}(-, -) \]
of additive bifunctors 
$\cat{D}_1^{\mrm{op}} \times \cat{D}_2 \to \dcat{M}(\K)$.
\end{cor}

\begin{exer} \label{exer:2140}
Prove Corollary \ref{cor:2120}. (Hint: use Theorems \ref{thm:3135} and 
\ref{thm:3145}.)
\end{exer}




\mysubsection{The Bifunctor \texorpdfstring{$\ot^{\mrm{L}}$}{LTors}}  
\label{subsec:LTors}

Consider a DG ring $A$. Like in Example \ref{exa:1965} we get a DG bifunctor 
\[ F := (- \ot_A -) : \dcat{C}(A^{\mrm{op}}) \times \dcat{C}(A) 
\to \dcat{C}(\K) . \]
Passing to homotopy categories, and postcomposing with 
$\opn{Q} : \dcat{K}(\K) \to \dcat{D}(\K)$, 
we obtain a triangulated bifunctor 
\[ F = (- \ot_A -) : \dcat{K}(A^{\mrm{op}}) \times
\dcat{K}(A) \to \dcat{D}(\K) . \]
Next we pick full triangulated subcategories 
$\cat{K}_1 \sub \dcat{K}(A^{\mrm{op}})$
and
$\cat{K}_2 \sub \dcat{K}(A)$. 
In practice this choice would be by some boundedness conditions; for instance
$\cat{K}_1 := \dcat{C}^{-}(A^{\mrm{op}})$
or $\cat{K}_2 := \dcat{C}^{-}(A)$, cf.\ Corollary \ref{cor:1900}.
We want to construct the left derived 
bifunctor of the triangulated bifunctor 
\[ F = (- \ot_A -) : \cat{K}_1 \times
\cat{K}_2 \to \dcat{D}(\K) . \]
This is done in the next theorem. 

\begin{thm} \label{thm:2107}
Let $A$ be a DG $\K$-ring, and let 
$\cat{K}_1 \sub \dcat{K}(A^{\mrm{op}})$
and $\cat{K}_2 \sub \dcat{K}(A)$
be full triangulated subcategories.
Assume that either $\cat{K}_1$ or $\cat{K}_2$ has enough K-flat objects.
Let $\cat{D}_i$ denote the localization of $\cat{K}_i$ with respect to the 
quasi-iso\-morphisms in it. 
Then the triangulated bifunctor 
\[  (- \ot_A -)  : \cat{K}_1 \times \cat{K}_2 \to \dcat{D}(\K) \]
has a left derived bifunctor%
\index{1-LTens@$(- \ot^{\mrm{L}}_A -)$} 
\[ (- \ot^{\mrm{L}}_A -)  : \cat{D}_1 \times \cat{D}_2 \to \dcat{D}(\K) . \]
Moreover, if either $P_1 \in \cat{K}_1$ or $P_2 \in \cat{K}_2$ is K-flat, 
then the morphism 
\[ \eta^{\mrm{L}}_{P_1, P_2} : P_1 \ot^{\mrm{L}}_A P_2 \to P_1 \ot_A P_2 \]
in $\dcat{D}(\K)$ is an isomorphism. 
\end{thm}

\begin{proof}
If $\cat{K}_2$ has enough K-flats, then we can take 
$\cat{P}_2 := \cat{K}_{2, \mrm{flat}}$, the full subcategory on the 
K-flats inside $\cat{K}_2$. And we take 
$\cat{P}_1 := \cat{K}_{1}$. We claim that the conditions of Theorem 
\ref{thm:2106} are satisfied. Condition 
(b) is simply the assumption that $\cat{K}_2$ has enough K-flats. As for 
condition (a): this is Lemma \ref{lem:2109} below,
with its condition (ii). 

The other case is proved the same way (but replacing sides).  
The last assertion also follows from \ref{lem:2109}.
\end{proof}

\begin{lem}  \label{lem:2109}
Suppose 
$\phi_1 : P_1 \to Q_1$ and $\phi_2 : P_2 \to Q_2$
are quasi-isomorphisms in $\dcat{C}(A^{\mrm{op}})$
and $\dcat{C}(A)$ respectively, and either of the conditions below holds\tup{:}
\begin{enumerate}
\rmitem{i} $Q_1$ and $P_1$ are both K-flat.
\rmitem{ii} $P_2$ and $Q_2$ are both K-flat.
\end{enumerate}
Then the homomorphism
\[ \phi_1 \ot \phi_2 : P_1 \ot_A P_2 \to Q_1 \ot_A Q_2 \]
in $\dcat{C}(\K)$ is a quasi-isomorphism. 
\end{lem}

\begin{proof}
We will only prove the lemma under condition (i); the other 
case is very similar. The homomorphism in question factors as follows:
\[ \phi_1 \ot \phi_2 = (\phi_1 \ot \opn{id}_{P_2}) \circ 
(\opn{id}_{P_1} \ot \, \phi_2 ) . \]
It suffices to prove that each of the factors is a quasi-isomorphism. This can 
be done by a messy direct calculation, but we will provide an indirect proof 
that relies on properties of the DG categories  
$\dcat{C}(A^{\mrm{op}})$ and $\dcat{C}(A)$ that 
were already established. 

First we shall prove that $\opn{id}_{P_1} \ot \, \phi_2$ is a quasi-isomorphism.
Let $R_2$ be the standard cone on the strict homomorphism 
$\phi_2 : P_2 \to Q_2$. So there is a standard triangle 
\begin{equation} \label{eqn:2110}
P_2 \xar{\phi_2} Q_2 \to R_2 \to P_2[1]
\end{equation}
in $\dcat{C}_{\mrm{str}}(A)$, and $R_2$ is acyclic. 
Applying the DG functor $P_1 \ot_A (-)$ to the triangle (\ref{eqn:2110}), and 
using Theorem \ref{thm:1185}, we see that there is a standard triangle 
\begin{equation} \label{eqn:2111}
 P_1 \ot_A P_2 \xar{\opn{id}_{P_1} \ot \, \phi_2} P_1 \ot_A Q_2 \to P_1 \ot_A 
R_2 \to (P_1 \ot_A P_2)[1]
\end{equation}
in $\dcat{C}(\K)$. This becomes a distinguished triangle in the triangulated 
category $\dcat{K}(\K)$. Thus there is a long exact sequence in cohomology 
associated to (\ref{eqn:2111}). Because $P_1$ is K-flat it follows that 
$P_1 \ot_A R_2$ is acyclic. We conclude that 
$ \opn{H}^i(\opn{id}_{P_1} \ot \, \phi_2)$
is bijective for all $i$. 

Now we shall prove that $\phi_1 \ot \opn{id}_{P_2}$ is a quasi-isomorphism.
Let $R_1 \in \dcat{C}(A^{\mrm{op}})$ be the standard cone on the homomorphism 
$\phi_1 : P_1 \to Q_1$. It is both acyclic and K-flat. 
Using standard triangles like (\ref{eqn:2110}) and (\ref{eqn:2111}) we reduce 
the problem to showing that $R_1 \ot_A P_2$ is acyclic. 
According to Corollary \ref{cor:1580} and Proposition \ref{prop:1525} there is 
a quasi-isomorphism 
$\til{P}_2 \to P_2$ in $\dcat{C}(A)$ from some K-flat DG module $\til{P}_2$. 
As already proved in the previous paragraph, since $R_1$ is K-flat, the 
homomorphism
$R_1 \ot_A \til{P}_2 \to R_1 \ot_A P_2$
is a quasi-isomorphism. But $R_1$ is acyclic and $\til{P}_2$ is K-flat, and 
therefore
$R_1 \ot_A \til{P}_2$ is acyclic. We conclude that 
$R_1 \ot_A P_2$ is acyclic, as required. 
\end{proof}

\begin{rem} \label{rem:2105} 
Theorem \ref{thm:2107} should be viewed as a template. It has commutative and 
noncommutative variants, that we will talk about later. 
And there are geometric variants where the source and target are categories of 
sheaves.
\end{rem}

\begin{dfn} \label{dfn:4295} 
Under the assumptions of Theorem \tup{\ref{thm:2107}},  
for DG modules $M_1 \in \cat{K}_1$ and $M_2 \in \cat{K}_2$,
and for an integer $i$, we write%
\index{1-ToriA@$\opn{Tor}_i^{A}(-, -)$}
\[ \opn{Tor}_i^{A}(M_1, M_2) := 
\opn{H}^{-i} ( M_1 \ot^{\mrm{L}}_{A} M_2 ) \in \dcat{M}(\K) . \]
\end{dfn}

\begin{exer} \label{exer:4295} 
Let $A$ be a ring. Prove that for modules
$M_1 \in \dcat{M}(A^{\mrm{op}})$ and $M_2 \in \dcat{M}(A)$ the $\K$-module
$\opn{Tor}_i^{A}(M_1, M_2)$ defined above is canonically isomorphic to the 
module in classical definition. Moreover, this is true regardless of the 
choices of the subcategories $\cat{K}_1 $ and $\cat{K}_2$,
as long as $M_i \in \cat{K}_i$. 
\end{exer}

Recall that the category $\dcat{D}(A)$ admits infinite direct sums. 

\begin{prop} \label{prop:4300}
For every $M \in \dcat{D}(A^{\mrm{op}})$ the functor 
\[ M \ot^{\mrm{L}}_{A} (-) : \dcat{D}(A) \to \dcat{D}(\K) \]
commutes with infinite direct sums. 
\end{prop}

\begin{proof}
Fix a K-flat resolution $P \to M$ in $\dcat{C}_{\mrm{str}}(A^{\mrm{op}})$, so 
we have an isomorphism of functors 
$M \ot^{\mrm{L}}_{A} (-) \cong P \ot_A (-)$.
The functor 
\[ P \ot_A (-) : \dcat{C}_{\mrm{str}}(A) \to 
\dcat{C}_{\mrm{str}}(\K) \]
commutes with infinite direct sums. But by Theorem \ref{thm:3140}
the direct sums in $\dcat{C}_{\mrm{str}}(-)$ and in
$\dcat{D}(-)$ are the same. 
\end{proof}




\mysubsection{Cohomological Dimensions of Functors and Objects} 
\label{subsec:way-out}

The material here is a refinement of the notion of ``way-out functors'' from 
\cite[Section II.7]{RD}. Most of it is taken from 
\cite{RD} and \cite{Ye8}. As always, there is a fixed base ring $\K$, and 
Convention \ref{conv:2490} is in force. 

Integer intervals and their properties were introduced in Subsection 
\ref{subsec:bounded-revis}.

\begin{dfn} \label{dfn:2126}
Let $A, B$ be DG rings, let $\cat{M}, \cat{N}$ be abelian categories, and let 
$\cat{C}_0 \sub \cat{C} \sub \dcat{D}^{}(A, \cat{M})$
be full subcategories. 
\begin{enumerate}
\item Let  
$F : \cat{C} \to \dcat{D}(B, \cat{N})$
be an additive functor, and let $S$ be an integer interval. We say that 
{\em $F$ has cohomological displacement at 
most $S$ relative to $\cat{C}_0$}%
\index{Cohomological displacement! of covariant functor}
if 
\[ \opn{con} \bigl( \opn{H}(F(M)) \bigr) \subseteq
\opn{con} \bigl( \opn{H}(M) \bigr) + S \]
for every $M \in \cat{C}_0$.

\item Let 
$F : \cat{C}^{\mrm{op}} \to \dcat{D}(B, \cat{N})$
be an additive functor, 
and let $S$ be an integer interval. We say that {\em $F$ has 
cohomological displacement at most 
$S$ relative to $\cat{C}_0$}% 
\index{Cohomological displacement! of contravariant functor}
if 
\[ \opn{con} \bigl( \opn{H}(F(M)) \bigr) \subseteq
- \opn{con} \bigl( \opn{H}(M) \bigr) + S \]
for every $M \in \cat{C}_0$.

\item Let $F$ be as in item (1) or (2).
The {\em cohomological displacement of $F$ relative to $\cat{C}_0$} is the 
smallest integer interval $S$ for which $F$ has cohomological 
displacement at most $S$ relative to $\cat{C}_0$. 

\item Let $S$ be the cohomological displacement of $F$ relative to 
$\cat{C}_0$. The 
{\em cohomological dimension of $F$  relative to $\cat{C}_0$}%
\index{Cohomological dimension! of functor}
is defined to be the length of the integer interval $S$. 

\item In case $\cat{C}_0 = \cat{C}$, we omit the clause ``relative to 
$\cat{C}_0$'' in all items above. 
\end{enumerate}
\end{dfn}

To emphasize the most important case of item (4) of the definition: {\em The 
functor $F$ has finite cohomological dimension if its cohomological 
displacement 
is bounded}. 

\begin{exa} \label{exa:2140}
The functor $F$ is the zero functor iff it has cohomological displacement 
$\varnothing$ and cohomological dimension $-\infty$. 
\end{exa}

\begin{exa} \label{exa:4285}
Let $F : \dcat{D}(A, \cat{M}) \to \dcat{D}(B, \cat{N})$
be an additive functor of finite cohomological dimension.
Then for each boundedness indicator $\star$ we have 
$F \bigl( \dcat{D}^{\star}(A, \cat{M}) \bigr) \sub 
\dcat{D}^{\star}(B, \cat{N})$.
\end{exa}

\begin{exa} \label{exa:2125}
Consider a nonzero commutative ring $A = B$, and the abelian category
$\cat{M} = \cat{N} := \dcat{M}(\K)$. So 
$\dcat{D}(A, \cat{M}) = \dcat{D}(B, \cat{N}) = \dcat{D}(A)$.
Take $\cat{C} := \dcat{D}(A)$.
For the covariant case (item (1) in Definition \ref{dfn:2126}) take a nonzero 
projective module $P$, and let
\[ F := \opn{RHom}_A \bigl( P \oplus P[1], - \bigr) : \dcat{D}(A) \to 
\dcat{D}(A) . \]
Then $F$ has cohomological displacement $[0, 1]$. 
For the contravariant case (item (2)) take a nonzero 
injective module $I$, and let 
\[ F := \opn{RHom}_A \bigl( -, I \oplus I[1] \bigr) : \dcat{D}(A)^{\mrm{op}} 
\to \dcat{D}(A) . \]
Then $F$ has cohomological displacement $[-1, 0]$. 
In both cases the cohomological dimension of $F$ is $1$. 
\end{exa}

\begin{exa} \label{exa:2135}
Suppose $A$ and $B$ are rings and 
$F : \dcat{M}(A) \to \dcat{M}(B)$ is a left exact additive functor. 
We get a triangulated functor 
$\mrm{R} F : \dcat{D}(A) \to \dcat{D}(B)$,
and $\opn{H}^i(\mrm{R} F(M)) = \mrm{R}^i F(M)$ for all $M \in \dcat{M}(A)$. 
Taking $\cat{C} := \dcat{M}(A)$, with its canonical embedding into
$\dcat{D}(A)$, the cohomological dimension of $\mrm{R} F$ relative to 
$\dcat{M}(A)$ equals the usual right cohomological dimension of the functor 
$F$. 
\end{exa}

Next is an example about triangulated functors of finite 
cohomological dimensions. Another example, of a contravariant triangulated 
functor of finite cohomological dimension, is the duality functor $D$ 
associated to a dualizing complex; see Subsections \ref{subsec:du-cplxs},
\ref{subsec:graded-NCDC} and \ref{subsec:NC-DC}.

\begin{exa} \label{exa:4595}  
Let $A$ be a commutative ring, let $\bsym{a} = (a_1, \ldots, a_n)$ be a finite 
sequence of elements of $A$, and let $\a \sub A$ be the ideal generated by 
$\bsym{a}$. The {\em $\a$-adic completion functor} $\Lambda_{\a}$ and the 
{\em $\a$-torsion functor} $\Ga_{\a}$ have derived functors 
$\mrm{L} \Lambda_{\a}, \mrm{R} \Ga_{\a} : \dcat{D}(A) \to  \dcat{D}(A)$.

We call the sequence $\bsym{a}$ {\em weakly proregular} if  
the Koszul complexes associated to powers of $\bsym{a}$ satisfy a rather 
complicated asymptotic formula -- see \cite[Definition 4.21]{PSY}. If $A$ is 
noetherian then weak proregularity is 
automatic; but this condition is also satisfied in many important 
non-noetherian cases. 

When $\bsym{a}$ is weakly proregular, the functors 
$\mrm{L} \Lambda_{\a}$ and $\mrm{R} \Ga_{\a}$ have finite cohomological 
dimensions, bounded by the length of $\bsym{a}$.
Furthermore, the {\em commutative MGM Equivalence} holds. To explain it, we 
need to introduce two full triangulated subcategories of $\dcat{D}(A)$~: 
\begin{itemize}
\item The category $\dcat{D}(A)_{\mrm{com}}$ of {\em derived $\a$-adically 
complete complexes}. These are the complexes $M$ such that the canonical 
morphism $M \to \mrm{L} \Lambda_{\a}(M)$ is an isomorphism. 

\item The category $\dcat{D}(A)_{\mrm{tor}}$ of {\em derived $\a$-torsion
complexes}. These are the complexes $M$ such that the canonical morphism
$\mrm{R} \Ga_{\a}(M) \to M$ is an isomorphism. 
\end{itemize}
The MGM Equivalence says that the functor 
$\mrm{R} \Ga_{\a} : \dcat{D}(A)_{\mrm{com}} \to  \dcat{D}(A)_{\mrm{tor}}$ 
is an equivalence of triangulated categories, with quasi-inverse 
$\mrm{L} \Lambda_{\a}$.
The name ``MGM'' stands for ``Matlis-Greenlees-May''.
This duality was worked out in the paper \cite{PSY} by M. Porta, L. Shaul and 
A. Yekutieli; expanding earlier work by E. Matlis; A. Grothendieck; 
J. Greenlees and P. May;  L. Alonso, A. Jeremias and J. Lipman; and P. 
Schenzel. 

In Subsection \ref{subsec:NC-MGM} we give a noncommutative variant of the MGM 
Equivalence, adapted from the paper \cite{VyYe} by R. Vyas and Yekutieli. 
\end{exa}

\begin{rem} \label{rem:2126}
Assume that in Definition \ref{dfn:2126}(1) we take $A = B = \K$,  
$\cat{C} = \dcat{D}(\cat{M})$, and 
$F : \dcat{D}(\cat{M}) \to \dcat{D}(\cat{N})$
is a triangulated functor. The functor $F$ has bounded below (resp.\ above) 
cohomological displacement if and only if it is a {\em way-out right} (resp.\ 
{\em left}) functor, in the sense of \cite[Section I.7, Definition]{RD}.
For instance, if $F$ is a way-out right functor, with bounding integers 
$n_1$ and $n_2$ (as defined in \cite{RD}), then the cohomological displacement 
of $F$ is contained in the integer interval $[n_1 - n_2, \infty]$.
Conversely, if $F$ has cohomological displacement at most
$[i_0, \infty]$ for some integer $i_0$, then $F$ is way-out right, and for 
every $n_1 \in \Z$ the integer $n_2 :=  n_1 - i_0$ satisfies the condition in 
\cite{RD}. Likewise for way-out left functors.  
\end{rem}

\begin{dfn} \label{dfn:2128}
Let $\star_1, \star_2$ be boundedness indicators, and assume the right derived 
bifunctor 
\[ \opn{RHom}_{A, \cat{M}} : 
\dcat{D}^{\star_1}(A, \cat{M})^{\mrm{op}} \times 
\dcat{D}^{\star_2}(A, \cat{M}) \to \dcat{D}(\K) \]
exists. Let $S$ be an integer interval of length 
$i\in \N \cup \{\pm  \infty \}$.
\begin{enumerate}
\item Let $M \in \dcat{D}^{\star_1}(A, \cat{M})$, and let 
$\cat{C} \sub \dcat{D}^{\star_2}(A, \cat{M})$ be a full 
subcategory. We say that $M$ has {\em projective concentration} $S$
and {\em projective dimension}% 
\index{Projective! concentration of DG module}%
\index{Projective! dimension of DG module}
$i$ relative to $\cat{C}$ if the functor
\[ \opn{RHom}_{A, \cat{M}}(M, -)|_{\cat{C}} : \cat{C} \to \dcat{D}(\K) \]
has cohomological displacement $-S$.

\item Let $M \in \dcat{D}^{\star_2}(A, \cat{M})$, and let 
$\cat{C} \sub \dcat{D}^{\star_1}(A, \cat{M})$ be a full  
subcategory. We say that $M$ has {\em injective concentration} $S$ 
and {\em injective dimension}% 
\index{Injective! dimension of DG module}%
\index{Injective! concentration of DG module}
$i$ relative to $\cat{C}$ if the functor
\[ \opn{RHom}_{A, \cat{M}}(-, M)|_{\cat{C}^{\mrm{op}}} 
: \cat{C}^{\mrm{op}} \to \dcat{D}(\K) \]
has  cohomological displacement $S$.

\item If $\cat{C} = \dcat{D}^{}(A, \cat{M})$,
then we omit the clause ``relative to $\cat{C}$'' in items (1) and (2). 
\end{enumerate}
\end{dfn}

\begin{exa} \label{exa:2127}
Continuing with the setup of Example \ref{exa:2125}, the DG module
$P \oplus P[1]$ (resp.\ $I \oplus I[1]$) has projective (resp.\ injective)
concentration $[-1, 0]$.
\end{exa}

\begin{exa} \label{exa:2128}
Let $A$ be a DG ring such that $\opn{H}(A)$ is nonzero, and consider the free 
DG module $P := A \in \dcat{D}(A)$.
The functor 
$F := \opn{RHom}_{A}(P, -) : \dcat{D}(A) \to \dcat{D}(\K)$ 
is isomorphic to the forgetful functor, so it has cohomological displacement 
$[0, 0]$ and cohomological dimension $0$.
Thus the DG module $P$ has projective concentration $[0, 0]$ and projective 
dimension $0$.
Note however that the cohomology $\opn{H}(P)$ could be unbounded! 
\end{exa}

\begin{prop} \label{prop:4291} 
Let $\cat{M}$ be an abelian category with enough injectives. The following are 
equivalent for $M \in \cat{M}$~\tup{:}
\begin{enumerate} [itemsep=0.3ex]
\rmitem{i} $M$ is an injective object of $\cat{M}$.

\rmitem{ii} $\opn{Ext}^{p}_{\cat{M}}(N, M) = 0$ for every 
$N \in \cat{M}$ and every $p \geq 1$. 

\rmitem{iii}  $\opn{Ext}^1_{\cat{M}}(N, M) = 0$ for every 
$N \in \cat{M}$. 
\end{enumerate}
\end{prop}

Note that by Corollary \ref{cor:1902} and Definition \ref{dfn:2121}, the 
$\K$-modules $\opn{Ext}^p_{\cat{M}}(N, M)$ exist.

\begin{exer} \label{exer:4291}
Prove Proposition \ref{prop:4291}. (Hint: the proof is just like in the case 
$\cat{M} = \dcat{M}(A)$. It can be found in standard textbooks on homological 
algebra.)
\end{exer}

\begin{prop} \label{prop:2115}
Let $\cat{M}$ be an abelian category with enough injectives. The following are 
equivalent for $M \in \dcat{D}^+(\cat{M})$~\tup{:}
\begin{enumerate}
\rmitem{i} $M$ has finite injective dimension.

\rmitem{ii} $M$ has finite injective dimension relative to $\cat{M}$.

\rmitem{iii} There is a quasi-isomorphism $M \to I$ in 
$\dcat{C}_{\mrm{str}}(\cat{M})$ 
to a bounded complex of injective $A$-modules $I$.
\end{enumerate}
\end{prop}

Note that by Corollary \ref{cor:1902} we can apply Definition \ref{dfn:2128} 
with boundedness type $\star_2 = +$, so we can talk about the injective 
dimension of $M$. 

\begin{proof} \mbox{}

\smallskip \noindent
(i) $\Rightarrow$ (ii): This is trivial, since 
$\cat{M} \sub \dcat{D}(\cat{M})$. 

\medskip \noindent 
(ii)  $\Rightarrow$ (iii): We may assume that $\opn{H}(M)$ is nonzero. 
Let $S$ be the injective concentration of the complex $M$
relative to $\cat{M}$, as in Definition \ref{dfn:2128}; this is a bounded 
integer interval (possibly empty, a priory). 

Let $q$ be an integer such that $\opn{H}^q(M) \neq 0$.
Consider the object $N := \opn{Z}^q(M) \in \cat{M}$ and the canonical 
monomorphism $N \to M^q$. This can be viewed as a morphism 
$\phi : N[-q] \to M$ in $\dcat{C}_{\mrm{str}}(\cat{M})$, and then 
$\opn{H}^q(\phi) : N = \opn{H}^q(N[-q]) \to \opn{H}^q(M)$ 
is a nonzero morphism in $\cat{M}$. Therefore
$\opn{Q}(\phi) :  N[-q] \to M$
is a nonzero morphism in $\dcat{D}(\cat{M})$, so this is a nonzero element of 
$\opn{H}^q \bigl( \opn{RHom}_{\cat{M}}(N, M) \bigr)$. 
Looking at Definition \ref{dfn:2128} we see that $q \in S$.
Since $S$ is nonempty and bounded, it has to be 
$S = [q_0, q_1]$ for some integers
$q_0 = \opn{inf}(\opn{H}(M))$ and $q_1 = \opn{sup}(\opn{H}(M))$.

According to Corollary \ref{cor:1902} there is 
quasi-isomorphism $M \to J$, where $J$ is a complex of injective objects of 
$\cat{M}$, and $\opn{inf}(J) = q_0$. 
Take $I := \opn{smt}^{\leq q_1}(J)$, the smart truncation from Definition 
\ref{dfn:2320}. 
Then the canonical homomorphism 
$I \to J$ is a quasi-isomorphism. The complex $I$ is concentrated in the 
integer interval $[q_0, q_1]$, and
$I^q = J^q$ is an injective object for all $q < q_1$. 

Let us prove that $I^{q_1} = \opn{Z}^{q_1}(J)$ is also an injective object of 
$\cat{M}$. Classically we would use a cosyzygy argument. Here we use another 
trick. Define $I' := \opn{stt}^{\leq q_1 - 1}(I)$, 
the stupid truncation of $I$ from Definition \ref{dfn:3066}. so
\[ I' = \bigl( \cdots 0 \to I^{q_0} \to \cdots \to I^{q_1 - 1} \to 0 \to 
\cdots \bigr) . \]
This is a bounded complex of injective objects.
Consider the short exact sequence 
\[ 0 \to I^{q_1}[-q_1] \to I \to I' \to 0 \]
in $\dcat{C}^+_{\mrm{str}}(\cat{M})$.   
According to Proposition \ref{prop:2165} this gives a distinguished triangle 
\begin{equation} \label{eqn:2129}
I^{q_1}[-q_1] \to I \to I' \xar{ \ \triangle \ }  
\end{equation}
in $\dcat{D}^+(\cat{M})$. Take any object $L \in \cat{M}$. Applying 
$\opn{RHom}_{\cat{M}}(L, -)$ to the distinguished triangle (\ref{eqn:2129}) and 
then taking cohomologies, we get a long exact sequence
\[ \cdots \to \opn{Ext}_{\cat{M}}^{q + q_1 - 1}(L, I') \to 
\opn{Ext}_{\cat{M}}^{q}(L, I^{q_1}) \to 
\opn{Ext}_{\cat{M}}^{q + q_1}(L, I) \to \cdots \]
in $\dcat{M}(\K)$. For every $q > 0$ the $\K$-module 
$\opn{Ext}_{\cat{M}}^{q + q_1 - 1}(L, I')$ 
vanishes trivially. By the definition of the interval $[q_0, q_1]$, and the 
existence of an isomorphism $M \cong I$ in $\dcat{D}(\cat{M})$, for every 
$q > 0$ the $\K$-module $\opn{Ext}_{\cat{M}}^{q + q_1}(L, I)$ is zero. Hence 
$\opn{Ext}_{\cat{M}}^{q}(L, I^{q_1}) = 0$ for all $q > 0$. 
This proves that the object $I^{q_1}$ is injective (see Proposition 
\ref{prop:4291}). 

We have quasi-isomorphisms $M \to J$ and $I \to J$. Since $I$ is 
K-injective, there is a quasi-isomorphism $M \to I$. 

\medskip \noindent 
(iii) $\Rightarrow$ (i): This is also trivial.
\end{proof}

\begin{prop} \label{prop:4293} 
Let $\cat{M}$ be an abelian category with enough projectives. The following are 
equivalent for $M \in \cat{M}$~\tup{:}
\begin{enumerate} [itemsep=0.3ex]
\rmitem{i} $M$ is a projective object of $\cat{M}$.

\rmitem{ii} $\opn{Ext}^i_{\cat{M}}(M, N) = 0$ for every 
$N \in \cat{M}$ and every $i \geq 1$. 

\rmitem{iii}  $\opn{Ext}^1_{\cat{M}}(M, N) = 0$ for every 
$N \in \cat{M}$. 
\end{enumerate}
\end{prop}

Note that by Corollary \ref{cor:1750} and Definition \ref{dfn:2121}, the 
$\K$-modules $\opn{Ext}^i_{\cat{M}}(M, N)$ exist.

\begin{prop} \label{prop:4290}
Let $\cat{M}$ be an abelian category with enough projectives. The following are 
equivalent for $M \in \dcat{D}^-(\cat{M})$~\tup{:}
\begin{enumerate}
\rmitem{i} $M$ has finite projective dimension.

\rmitem{ii} $M$ has finite projective dimension relative to 
$\cat{M}$.

\rmitem{iii} There is a quasi-isomorphism $P \to M$ in 
$\dcat{C}_{\mrm{str}}(\cat{M})$ from a bounded complex of projective objects 
$P$.
\end{enumerate}
\end{prop}

Note that by Corollary \ref{cor:1750} we can apply Definition \ref{dfn:2128} 
with boundedness type $\star_1 = -$, so we can talk about the projective 
dimension of $M$. 

\begin{exer} \label{exer:2390}
Prove Propositions \ref{prop:4293} and \ref{prop:4290}. (Hint: compare to 
Propositions \ref{prop:4291} and \ref{prop:2115}.
Or look in any standard textbook on homological algebra.)
\end{exer}

\begin{dfn} \label{dfn:2390} 
Let $\star_1, \star_2$ be boundedness indicators, and assume the left 
derived bifunctor 
\[ (- \ot^{\mrm{L}}_A -) : 
\dcat{D}^{\star_1}(A^{\mrm{op}}) \times 
\dcat{D}^{\star_2}(A) \to \dcat{D}(\K) \]
exists. Let $S$ be an integer interval of length 
$i\in \N \cup \{\pm  \infty \}$.
\begin{enumerate}
\item Let $M \in \dcat{D}^{\star_2}(A)$, and let 
$\cat{C} \sub \dcat{D}^{\star_1}(A^{\mrm{op}})$ be a full  
subcategory. We say that $M$ has {\em flat concentration} $S$
and {\em flat dimension}% 
\index{Flat! dimension of DG module}%
\index{Flat! concentration of DG module}
$i$ relative to $\cat{C}$ if the functor
$(- \ot^{\mrm{L}}_A M) |_{\cat{C}} : \cat{C} \to \dcat{D}(\K)$
has cohomological displacement $S$.

\item If $\cat{C} = \dcat{D}^{}(A^{\mrm{op}})$,
then we omit the clause ``relative to $\cat{C}$''. 
\end{enumerate}
\end{dfn}

\begin{prop} \label{prop:2390}
Let $A$ be a ring. The following are equivalent for $M \in \dcat{D}(A)$\tup{:}
\begin{enumerate}
\rmitem{i} $M$ has finite flat dimension.

\rmitem{ii} $M$ has finite flat dimension relative to 
$\dcat{M}(A^{\mrm{op}})$.

\rmitem{iii} There is an isomorphism $P \cong M$ in $\dcat{D}(A)$, 
where $P$ is a bounded complex of flat $A$-modules.
\end{enumerate}
\end{prop}

\begin{exer} \label{exer:2391}
Prove Proposition \ref{prop:2390}. (The proof is similar to that of Proposition 
\ref{prop:4290}. It can be found in many standard books on homological algebra.)
\end{exer}

\begin{rem} \label{rem:4785}
Propositions \ref{prop:2115}, \ref{prop:4290} and \ref{prop:2390} do not seem 
to have analogues for $M \in \dcat{D}(A)$, when $A$ is a DG ring that is not 
concentrated in degree $0$. 
\end{rem}

 




\mysubsection{Theorems on Functors Satisfying Finiteness Conditions} 
\label{subsec:thm-funcs-fin}

In this subsection we apply the definitions from the previous subsection 
(on cohomological dimensions of functors), combined with noetherian conditions,
to prove Theorems \ref{thm:2135}, \ref{thm:2160} and \ref{thm:4600}. 
These results (or their analogues) were originally proved in 
\cite[Section II.7]{RD}. 

As always, there is a fixed base ring $\K$, and 
Convention \ref{conv:2490} is in force.

\begin{dfn} \label{dfn:2135}
Suppose $A$ is a left noetherian ring. 
\begin{enumerate}
\item We denote by $\dcat{M}_{\mrm{f}}(A)$ the full subcategory of 
$\dcat{M}(A) = \cat{Mod}A$ on the finitely generated modules. 

\item  We denote by $\dcat{D}_{\mrm{f}}(A)$ the full subcategory of 
$\dcat{D}(A) = \dcat{D}(\cat{Mod} A)$ on the complexes with cohomology modules
in $\dcat{M}_{\mrm{f}}(A)$.
\end{enumerate}
\end{dfn}

Because $A$ is left noetherian, the category $\dcat{M}_{\mrm{f}}(A)$ is a thick
abelian subcategory of $\dcat{M}(A)$, and the category $\dcat{D}_{\mrm{f}}(A)$ 
is a full triangulated subcategory of $\dcat{D}(A)$.  
When viewed as a left module, 
$A \in \dcat{M}_{\mrm{f}}(A) \sub \dcat{D}_{\mrm{f}}^{\mrm{b}}(A)$.
Note the similarity to Definition \ref{dfn:3180}. 

\begin{thm} \label{thm:2135}
Let $A$ be a left noetherian ring, let $\cat{N}$ be an abelian category, 
let $\star$ be a boundedness indicator, let 
$F, G : \dcat{D}^{\star}_{\mrm{f}}(A) \to \dcat{D}(\cat{N})$ 
be triangulated functors, and let 
$\ze : F \to G$ be a morphism of triangulated functors. Assume that 
the morphism 
$\ze_A : F(A) \to G(A)$
in $\dcat{D}(\cat{N})$ is an isomorphism. 
\begin{enumerate}
\item If $\star = -$, and if $F$ and $G$ have bounded above cohomological 
displacements, then 
$\ze_M : F(M) \to G(M)$
is an isomorphism for every $M \in \dcat{D}^-_{\mrm{f}}(A)$.

\item If $\star = \bra{\mrm{empty}}$, and if $F$ and $G$ have bounded 
cohomological displacements, then $\ze_M$ is an isomorphism for every 
$M \in \dcat{D}_{\mrm{f}}(A)$.
\end{enumerate} 
\end{thm}

We shall require the next lemmas for the proof of the theorem. 

\begin{lem} \label{lem:2145}
Let $\cat{D}$ be a triangulated category, let 
$F, G : \cat{D} \to \dcat{D}(\cat{N})$ be triangulated functors, let 
$\ze : F \to G$ be a morphism of triangulated functors, and let  
$L \xar{\phi} M \to N \xar{ \ \triangle \ }$
be a distinguished triangle in $\cat{D}$.
\begin{enumerate}
\item If the morphisms $\ze_L$ and $\ze_M$ are both isomorphisms, then 
$\ze_N$ is an isomorphism. 

\item Let $j$ be an integer. If $\opn{H}^{j - 1}(F(N))$,
$\opn{H}^{j - 1}(G(N))$, $\opn{H}^{j}(F(N))$ and 
$\opn{H}^{j}(G(N))$ are all zero, and if  $\opn{H}^{j}(\ze_L)$
is an isomorphism, then $\opn{H}^{j}(\ze_M)$
is an isomorphism.
\end{enumerate} 
\end{lem}

\begin{proof}
(1) In $\dcat{D}(\cat{N})$ we get the commutative diagram 
\begin{equation} \label{eqn:2148}
\UseTips \xymatrix @C=8ex @R=6ex {
F(L)
\ar[r]
\ar[d]^{\ze_L}
&
F(M)
\ar[r]
\ar[d]^{\ze_M}
&
F(N)
\ar[r]
\ar[d]^{\ze_N}
&
F(L)[1]
\ar[d]^{\ze_L[1]}
\\
G(L)
\ar[r]
&
G(M)
\ar[r]
&
G(N)
\ar[r]
&
G(L)[1]
} 
\end{equation}
with horizontal distinguished triangles. According to 
Proposition \ref{prop:1283}, $\ze_N$ is an isomorphism. 

\medskip \noindent 
(2) Passing to cohomologies in (\ref{eqn:2148})
we have a commutative diagram with exact rows
\[ \UseTips \xymatrix @C=8ex @R=6ex {
\opn{H}^{j - 1}(F(N))
\ar[r]
\ar[d]^{\opn{H}^{j - 1}(\ze_{N})}
&
\opn{H}^{j}(F(L))
\ar[r]^(0.47){\opn{H}^{j}(F(\phi))}
\ar[d]^{\opn{H}^{j}(\ze_{L})}
&
\opn{H}^{j}(F(M))
\ar[r]
\ar[d]^{\opn{H}^{j}(\ze_{M})}
&
\opn{H}^{j}(F(N))
\ar[d]^{\opn{H}^{j}(\ze_{N})}
\\
\opn{H}^{j - 1}(G(N))
\ar[r]
&
\opn{H}^{j}(G(L))
\ar[r]^(0.47){\opn{H}^{j}(G(\phi))}
&
\opn{H}^{j}(G(M))
\ar[r]
&
\opn{H}^{j}(G(N)) 
} \]
in $\cat{N}$. The vanishing assumption implies that $\opn{H}^j(F(\phi))$ and
$\opn{H}^j(G(\phi))$ are isomorphisms. Hence 
$\opn{H}^{j}(\ze_{M})$ is an isomorphism.
\end{proof}

\begin{lem} \label{lem:2146}
Let $\cat{D}$ be a triangulated category, let 
$F, G : \cat{D} \to \dcat{D}(\cat{N})$ be triangulated functors, and let 
$\ze : F \to G$ be a a morphism of triangulated functors.
The following conditions are equivalent for $M \in \cat{D}$~\tup{:}
\begin{enumerate}
\rmitem{i} $\ze_M$ is an isomorphism. 

\rmitem{ii} $\ze_{M[i]}$ is an isomorphism for every integer $i$. 

\rmitem{iii} The morphism 
$\opn{H}^j(\ze_M) : \opn{H}^j(F(M)) \to \opn{H}^j(G(M))$
is an isomorphism for every integer $j$. 
\end{enumerate}
\end{lem}

\begin{proof}
The equivalence (i) $\Leftrightarrow$ (ii) is because both $F$ and $G$ are 
triangulated functors. 
The equivalence (i) $\Leftrightarrow$ (iii) is because
the functor 
$\opn{H} : \dcat{D}(\cat{N}) \to \dcat{G}_{\mrm{str}}(\cat{N})$ 
is conservative; see Corollary \ref{cor:2145}.
\end{proof}

\begin{proof}[Proof of Theorem \tup{\ref{thm:2135}}] \mbox{}

\smallskip \noindent
(1) First assume $P$ is a bounded complex of finitely generated free 
$A$-modules. Then $P$ is obtained from $A$ by finitely many standard cones and 
translations. By Lemmas \ref{lem:2145}(1) and \ref{lem:2146} it follows that 
$\ze_P$ is an isomorphism. 

Next let $P$ be a bounded above complex of finitely generated free $A$-modules.
Choose some integer $j$. 
Let $i_1$ be an integer such that the integer interval $[-\infty, i_1]$ 
contains the cohomological displacements of $F$ and $G$. 
Define $P' := \opn{stt}^{\leq j - i_1 - 2}(P)$, the stupid truncation of $P$ 
below $j - i_1 - 2$; and let $P'' := \opn{stt}^{\geq j - i_1 - 1}(P)$, the 
complementary  stupid truncation. See Definition \ref{dfn:3066}. 
According to Proposition \ref{prop:2165}, the short exact 
sequence (\ref{eqn:2147}) 
gives a distinguished triangle 
$P'' \to P \to P' \xar{\triangle}$ 
in $\dcat{D}^{-}_{\mrm{f}}(A)$. The complex $P''$ is a bounded complex of 
finitely 
generated free $A$-modules, so we already know that $\ze_{P''}$ is an 
isomorphism. Hence $\opn{H}^j(\ze_{P''})$ is an isomorphism.
On the other hand $\opn{H}(P')$ is concentrated in the degree interval 
$[-\infty, j - i_1 - 2]$. Therefore 
$\opn{H}^{k}(F(P')) = \opn{H}^k(G(P')) = 0$ for all 
$k \geq j - 1$. By Lemma \ref{lem:2145}(2), $\opn{H}^j(\ze_P)$ is an 
isomorphism. Because $j$ is arbitrary, Lemma \ref{lem:2146} says that 
$\ze_{P}$ is an isomorphism. 

Now take an arbitrary $M \in \dcat{D}^-_{\mrm{f}}(A)$.
By Theorem \ref{thm:3340} and Example \ref{exa:1930} there is a 
resolution $P \to M$, where $P$ is a bounded above complex of finitely 
generated free $A$-modules. Since $\ze_P$ is an isomorphism, so is $\ze_M$. 

\medskip \noindent
(2) Now we assume that the functors $F$ and $G$ have finite cohomological 
dimensions. Take any complex $M \in \dcat{D}_{\mrm{f}}(A)$. 
By Lemma \ref{lem:2146} it suffices to prove 
that $\opn{H}^j(\ze_{M})$ is an isomorphism for every integer $j$. 

Let $[i_0, i_1]$ be a bounded integer interval that contains the 
cohomological displacements of the functors $F$ and $G$. 
Define $M'' := \opn{smt}^{\leq j - i_0}(M)$, the smart truncation of $M$ 
below $j - i_0$; and let $M' := \opn{smt}^{\geq j - i_0 + 1}(M)$, the 
complementary smart truncation. See Definition \ref{dfn:2320}. 
According to Proposition \ref{prop:2320}
there is a distinguished triangle
$M'' \to M \to M' \xar{\triangle}$
in $\dcat{D}_{\mrm{f}}(A)$. 
The cohomologies of these complexes satisfy 
$\opn{H}^i(M'') = \opn{H}^i(M)$ and 
$\opn{H}^i(M') = 0$ for $i \leq j - i_0$; and 
$\opn{H}^i(M'') = 0$ and $\opn{H}^i(M') = \opn{H}^i(M)$ for 
$i \geq j - i_0 + 1$. 
Note that $M'' \in \dcat{D}^-_{\mrm{f}}(A)$. 

By part (1) we know that $\ze_{M''}$ is an 
isomorphism, and therefore also $\opn{H}^j(\ze_{M''})$ is an isomorphism. 
The cohomology $\opn{H}(M')$ is concentrated in the degree interval 
$[j - i_0 + 1, \infty]$, 
and therefore the cohomologies $\opn{H}(F(M'))$ and  $\opn{H}(G(M'))$
are concentrated in the interval $[j + 1, \infty]$. In particular the objects 
$\opn{H}^{j - 1}(F(M'))$, $\opn{H}^{j - 1}(G(M'))$, 
$\opn{H}^{j}(F(M'))$ and $\opn{H}^{j}(G(M'))$ are zero. 
By Lemma \ref{lem:2145}(2), $\opn{H}^j(\ze_M)$ is an isomorphism.
\end{proof}

The triangulated structure of $\dcat{D}(A)^{\mrm{op}}$ was introduced in 
Subsection \ref{subsec:opp-dercat-triang}. 
By our notational conventions, 
$\dcat{D}_{\mrm{f}}(A)^{\mrm{op}}$ is the full subcategory of 
$\dcat{D}(A)^{\mrm{op}}$ on the complexes with finitely generated cohomology 
modules, and it is triangulated, by Propositions \ref{prop:3051} and 
\ref{prop:2356}. 

\begin{exer} \label{exer:4325}
State and prove the contravariant modification of Theorem \lb \ref{thm:2135}.
(Hint: study the proofs of Theorems \ref{thm:2135} and \ref{thm:2160}.)
\end{exer}

\begin{thm} \label{thm:2160}
Let $A$ be a left noetherian ring, let $\cat{N}$ be an abelian category, let
$\cat{N}_0 \sub \cat{N}$ be a thick abelian subcategory, let $\star$ be a 
boundedness condition, and let 
$F : \dcat{D}^{\star}_{\mrm{f}}(A)^{\mrm{op}} \to \dcat{D}(\cat{N})$ 
be a triangulated functor. Assume that 
$F(A)$ belongs to $\dcat{D}_{\cat{N}_0}(\cat{N})$.
\begin{enumerate}
\item If $\star = -$, and if $F$ has bounded below cohomological displacement, 
then $F(M)$ belongs to $\dcat{D}_{\cat{N}_0}(\cat{N})$
for every $M \in \dcat{D}^-_{\mrm{f}}(A)$.

\item If $\star = \bra{\mrm{empty}}$, and if $F$ has bounded cohomological 
displacement, then $F(M)$ belongs to $\dcat{D}_{\cat{N}_0}(\cat{N})$
for every $M \in \dcat{D}_{\mrm{f}}(A)$.
\end{enumerate} 
\end{thm}

\begin{proof} \mbox{}

\smallskip \noindent
(1) First assume $P$ is a bounded complex of finitely generated free 
$A$-modules. The complex $P$ is obtained from the free module $A$ by finitely 
many extensions and translations in
$\dcat{C}_{\mrm{str}}(A)^{\mrm{op}}$. These come either from stupid 
truncations as in (\ref{eqn:2147}) or from breaking up finite direct sums. 
According to Proposition \ref{prop:3060}, $P$ is obtained from $A$ 
by finitely many standard cones and translations in 
$\dcat{D}^{\star}_{\mrm{f}}(A)^{\mrm{op}}$. 
Since $\dcat{D}_{\cat{N}_0}(\cat{N})$ is a full triangulated 
subcategory and $F$ is a triangulated functor, it follows that 
$F(P) \in \dcat{D}_{\cat{N}_0}(\cat{N})$.

Next let $P$ be a bounded above complex of finitely generated free $A$-modules.
Choose some integer $j$. We want to prove that $\opn{H}^j(F(P)) \in \cat{N}_0$. 
Let $i_0$ be an integer such that the integer interval $[i_0, \infty]$ 
contains the cohomological displacement of $F$. 
Define $P' := \opn{stt}^{\leq -j - 1 + i_0}(P)$, the stupid truncation of 
$P$ below $-j - 1 + i_0$; and let $P'' := \opn{stt}^{\geq j + i_0}(P)$, the 
complementary  stupid truncation. 
These truncations are done in the category 
$\dcat{C}_{\mrm{str}}(A)$. 
The short exact sequence (\ref{eqn:2147}) gives, upon applying $\opn{Op}$, 
a short exact sequence in 
$\dcat{C}_{\mrm{str}}(A)^{\mrm{op}}$.
According to Proposition \ref{prop:3060} 
there is a distinguished triangle 
$P' \to P \to P'' \xar{\triangle}$ 
in $\dcat{D}^-_{\mrm{f}}(A)^{\mrm{op}}$.
Since $F$ is a triangulated functor, there is a distinguished triangle 
$F(P') \to F(P) \to F(P'') \xar{\triangle}$ 
in $\dcat{D}(\cat{N})$.
The complex $P''$ is a bounded complex of finitely generated 
free $A$-modules, so we already know that
$F(P'') \in \dcat{D}_{\cat{N}_0}(\cat{N})$,
and in particular $\opn{H}^j(F(P'')) \in \cat{N}_0$. 
On the other hand $\opn{H}(P')$ is concentrated in the interval 
$[-\infty, -j - 1 + i_0]$. Therefore 
$\opn{H}(F(P'))$ is  concentrated in the interval
$[j + 1, \infty]$, and in particular 
$\opn{H}^{j - 1}(F(P')) = \opn{H}^j(F(P')) = 0$.
As we saw in the proof of Lemma \ref{lem:2145}(2), 
$\opn{H}^j(F(P'')) \to \opn{H}^j(F(P))$ is an isomorphism. 
The conclusion is that $\opn{H}^j(F(P)) \in \cat{N}_0$. 

Now take an arbitrary $M \in \dcat{D}^-_{\mrm{f}}(A)$.
There is a quasi-isomorphism $P \to M$, where $P$ is a bounded above complex of 
finitely generated free $A$-modules. So $F(M) \cong F(P)$, and thus 
$F(M) \in \dcat{D}_{\cat{N}_0}(\cat{N})$. 
 
\medskip \noindent
(2) Now we assume that the functor $F$ has finite cohomological dimension. Take 
any complex $M \in \dcat{D}_{\mrm{f}}(A)$. We want to prove 
that for every $j \in \Z$ the object $\opn{H}^j(F(M))$ lies in $\cat{N}_0$. 

Let $[i_0, i_1]$ be a bounded integer interval that contains the 
cohomological displacement of the functor $F$. 
Define $M'' := \opn{smt}^{\leq - j + 1 + i_1}(M)$, the smart truncation of $M$ 
below $-j + 1 + i_1$; and let $M' := \opn{smt}^{\geq -j + 2 + i_1}(M)$, the 
complementary smart truncation. These truncations are preformed in the category 
$\dcat{C}_{\mrm{str}}(A)$. 
By Proposition \ref{prop:3061} there is a distinguished triangle 
$M' \to M \to M'' \xar{\triangle}$ 
in $\dcat{D}_{\mrm{f}}(A)^{\mrm{op}}$.
Since $F$ is a triangulated functor, there is a distinguished triangle 
$F(M') \to F(M) \to F(M'') \xar{\triangle}$ 
in $\dcat{D}(\cat{N})$.
The cohomology of $M'$ is concentrated in the integer
interval $[- j + 2 + i_1, \infty]$, and therefore the cohomology of 
$F(M')$ is concentrated in the interval $[- \infty, j - 2]$. 
In particular the objects $\opn{H}^{j - 1}(F(M'))$ and 
$\opn{H}^{j}(F(M'))$ are zero.
By the proof of Lemma \ref{lem:2145}(2), the morphism 
$\opn{H}^j(F(M'')) \to \opn{H}^j(F(M))$ is an isomorphism.
But $M'' \in \dcat{D}^-_{\mrm{f}}(A)$, so as we proved in part (1), its 
cohomologies are inside $\cat{N}_0$. 
\end{proof}

Here is the covariant version of Theorem \ref{thm:2160}.

\begin{thm} \label{thm:4600} 
Let $A$ be a left noetherian ring, let $\cat{N}$ be an abelian category, let
$\cat{N}_0 \sub \cat{N}$ be a thick abelian subcategory, let $\star$ be a 
boundedness indicator, and let 
$F : \dcat{D}^{\star}_{\mrm{f}}(A) \to \dcat{D}(\cat{N})$
be a triangulated functor. Assume that 
$F(A)$ belongs to $\dcat{D}_{\cat{N}_0}(\cat{N})$.
\begin{enumerate}
\item If $\star = -$, and if $F$ has bounded above cohomological displacement, 
then $F(M)$ belongs to $\dcat{D}_{\cat{N}_0}(\cat{N})$
for every $M \in \dcat{D}^-_{\mrm{f}}(A)$.

\item If $\star = \bra{\mrm{empty}}$, and if $F$ has bounded cohomological 
displacement, then $F(M)$ belongs to $\dcat{D}_{\cat{N}_0}(\cat{N})$
for every $M \in \dcat{D}_{\mrm{f}}(A)$.
\end{enumerate} 
\end{thm}

\begin{exer} \label{exer:4600}
Prove Theorem \ref{thm:4600}.  (Hint: study the proofs of Theorems 
\ref{thm:2135} and \ref{thm:2160}.)
\end{exer}

\begin{rem} \label{rem:4600}
In Theorems \ref{thm:2135}, \ref{thm:2160} and \ref{thm:4600} the 
source category is $\dcat{D}^{\star}_{\mrm{f}}(A)$, for a left noetherian 
ring $A$. There are variants of these theorems in which the source category
is $\dcat{D}^{\star}_{\mrm{f}}(A)$, for a left pseudo-noetherian DG 
ring $A$; and also with source category 
$\dcat{D}^{\star}_{\cat{M}_0}(\cat{M})$
for suitable abelian categories $\cat{M}_0 \sub \cat{M}$. 
We leave it to the reader to investigate these variations. 
\end{rem}




\mysubsection{Derived Restriction, Induction and Coinduction Functors}
\label{subsec:adjs-NC}

As before, there is a nonzero commutative base ring $\K$. All DG rings are 
$\K$-central, and all homomorphisms between them are over $\K$.

\begin{dfn} \label{dfn:3086}
Let $u : A \to B$ be a DG ring homomorphism.
\begin{enumerate}
\item The {\em restriction functor}% 
\index{Restriction functor}
is the $\K$-linear DG functor%
\index{1-Rest@$\opn{Rest}_u$}
$\opn{Rest}_u = \opn{Rest}_{B / A} : \dcat{C}(B) \to \dcat{C}(A)$,
that is the identity on the underlying DG 
$\K$-modules, and the $A$-action is via $u$.

\item Since the functor $\opn{Rest}_u$  is exact, it extends to a 
triangulated functor 
$\opn{Rest}_u : \dcat{D}(B) \to \dcat{D}(A)$.
\end{enumerate}
\end{dfn}

\begin{dfn} \label{dfn:3087}
Let $u : A \to B$ be a DG ring homomorphism.
\begin{enumerate}
\item The {\em induction functor}% 
\index{Induction functor}%
\index{1-Ind@$\opn{Ind}_u$}
is the DG functor 
$\opn{Ind}_u  = \opn{Ind}_{B / A} : \dcat{C}(A) \to \dcat{C}(B)$
with formula $\opn{Ind}_u := B \ot_A (-)$.

\item The {\em derived induction functor}% 
\index{Induction functor! derived}%
\index{1-LInd@$\opn{LInd}_u$}
$\opn{LInd}_u  = \opn{LInd}_{B / A} : \dcat{D}(A) \to \dcat{D}(B)$ 
is the triangulated left derived functor of $\opn{Ind}_u$, namely 
$\opn{LInd}_u := B \ot^{\mrm{L}}_A (-)$.
\end{enumerate}
\end{dfn}

For every $M \in \dcat{D}(A)$ there is the canonical morphism 
$\eta^{\mrm{L}}_{M} : \opn{LInd}_u(M) \to \opn{Ind}_u(M)$
in $\dcat{D}(B)$, which is part of the left derived functor. 

\begin{dfn} \label{dfn:3111}
Let $u : A \to B$ be a DG ring homomorphism, and let 
$M \in \dcat{D}(A)$ and $N \in \dcat{D}(B)$ be DG modules.
A morphism
$\la : M \to \opn{Rest}_u(N)$
in $\dcat{D}(A)$ is called a 
{\em forward morphism in $\dcat{D}(A)$  over $u$}%
\index{Forward morphism! in derived category}.

Similarly there are forward morphisms over $u$ in the categories 
$\dcat{C}(A)$, $\dcat{C}_{\mrm{str}}(A)$ and $\dcat{K}(A)$.
\end{dfn}

We often omit the functor $\opn{Rest}_u$, and just talk about a 
forward morphism $\la : M \to N$ in $\dcat{D}(A)$ over $u$, etc. 

For any DG module $M \in \dcat{C}(A)$ there is a canonical forward homomorphism
\begin{equation} \label{eqn:3194}
\opn{q}_{u, M} : M \to B \ot_A M = \opn{Ind}_u(M) \, , \quad 
\opn{q}_{u, M}(m) := 1_B \ot m 
\end{equation}
in $\dcat{C}_{\mrm{str}}(A)$ over $u$. 
Now let $N \in \dcat{C}(B)$. The usual change of ring adjunction formula gives 
a 
canonical isomorphism 
\begin{equation} \label{eqn:2401}
\opn{fadj}_{u, M, N} : 
\opn{Hom}_{\dcat{C}(A)}(M, N) \iso 
\opn{Hom}_{\dcat{C}(B)} \bigl( B \ot_A M, N  \bigr) 
\end{equation}
in $\dcat{C}_{\mrm{str}}(\K)$. It is characterized by the property that 
for every forward morphism $\la : M \to N$ there is equality 
$\opn{fadj}_{u, M, N}(\la) \circ \opn{q}_{u, M} = \la$.
We refer to the isomorphism $\opn{fadj}_{u, M, N}$ as 
{\em forward adjunction}. 
Since the iso\-morphism $\opn{fadj}_{u, M, N}$ is functorial in 
$M$ and $N$, we see that the functor $\opn{Ind}_u$ is a left adjoint of 
$\opn{Rest}_u$. 

The next theorem shows that this is also true on the derived level. 

\begin{thm} \label{thm:3085} 
Let $u : A \to B$ be a homomorphism of DG $\K$-rings.
\begin{enumerate}
\item For each $M \in \dcat{D}(A)$ there is a unique forward morphism%
\index{Forward morphism! canonical}
\index{1-QLuM@$\opn{q}_{u, M}^{\mrm{L}}$}
\[ \opn{q}_{u, M}^{\mrm{L}} : M \to B \ot^{\mrm{L}}_{A} M = 
\opn{LInd}_u(M) \]
in $\dcat{D}(A)$ over $u$, called the {\em canonical forward morphism},
which is functorial in $M$, and satisfies 
$\eta^{\mrm{L}}_M \circ \opn{q}_{u, M}^{\mrm{L}} = \opn{q}_{u, M}$.

\item For each $M \in \dcat{D}(A)$ and $N \in \dcat{D}(B)$
there is a unique $\K$-linear isomorphism%
\index{Forward adjunction! derived}%
\index{1-Fadj@$\opn{fadj}^{\mrm{L}}_{u, M, N}$}
\[ \opn{fadj}_{u, M, N}^{\mrm{L}} :  
\opn{Hom}_{\dcat{D}(A)} \bigl( M, \opn{Rest}_u(N) \bigr) \iso 
\opn{Hom}_{\dcat{D}(B)} \bigl( \opn{LInd}_u(M), N \bigr) , \]
called {\em derived forward adjunction}, such that 
$\opn{fadj}_{u, M, N}^{\mrm{L}}(\la) \circ \opn{q}_{u, M}^{\mrm{L}} = \la$
for all forward morphisms $\la : M \to N$ in $\dcat{D}(A)$ over $u$.
The isomorphism $\opn{fadj}_{u, M, N}^{\mrm{L}}$ is functorial in $M$ and $N$. 

\item The functor
$\opn{LInd}_u : \dcat{D}(A) \to \dcat{D}(B)$
is left adjoint to $\opn{Rest}_u$. 
\end{enumerate}
\end{thm}

Here are the commutative diagrams in $\dcat{D}(A)$ illustrating the theorem. 
\[ \UseTips \xymatrix @C=8ex @R=6ex {
M
\ar[d]_{\opn{q}_{u, M}^{\mrm{L}}}
\ar[dr]^{\opn{q}_{u, M}} 
\\
\opn{LInd}_u(M)
\ar[r]_{\eta^{\mrm{L}}_M}
&
*+{\opn{Ind}_u(M)}
} 
\qquad \quad 
\UseTips \xymatrix @C=14ex @R=6ex {
M
\ar[d]_{\opn{q}_{u, M}^{\mrm{L}}}
\ar[dr]^{\la} 
\\
\opn{LInd}_u(M)
\ar[r]_(0.55){\opn{fadj}_{u, M, N}^{\mrm{L}}(\la)}
&
*+++{N}
} 
\]

\begin{proof} \mbox{}

\smallskip \noindent 
(1) Let us choose a system of K-projective resolutions in $\dcat{C}(A)$, as in 
Definition \ref{dfn:1524}, and use it to present the derived functor 
$\opn{LInd}_u$. 
% According to Theorem \ref{thm:2107}, 
To construct an isomorphism $\opn{q}_{u, M}^{\mrm{L}}$
which is functorial in $M$, it suffices to consider 
a K-projective DG $A$-module $M = P$. But then $\opn{LInd}_u(P) = 
\opn{Ind}_u(P)$, 
$\eta^{\mrm{L}}_P = \opn{id}_P$,
and we have no choice but to define 
$\opn{q}_{u, P}^{\mrm{L}} := \opn{q}_{u, P}$. 

\medskip \noindent
(2) We keep the system of K-projective resolutions. 
As in item (1), it is enough to say what is $\opn{fadj}_{u, P, N}^{\mrm{L}}$ 
for a K-projective DG $A$-module $P$. Then there is a canonical isomorphism 
 \[ \opn{Q}_A : \opn{Hom}_{\dcat{K}(A)}(P, N) \iso 
\opn{Hom}_{\dcat{D}(A)}(P, N) . \]
Since $B \ot_A P = \opn{LInd}_u(P)$ is a K-projective DG $B$-module, we also 
have  a canonical isomorphism 
\[ \opn{Q}_B : \opn{Hom}_{\dcat{K}(B)}(B \ot_A P, N) \iso
\opn{Hom}_{\dcat{D}(B)}(B \ot_A P , N) . \]
But the usual adjunction formula (\ref{eqn:2401}) induces, by passing to 
$0$-th cohomology on both sides, an isomorphism 
\[ \opn{H}^0(\opn{fadj}_{u, P, N})  : \opn{Hom}_{\dcat{K}(A)}(P, N) \iso 
\opn{Hom}_{\dcat{K}(B)}(B \ot_A P, N) ,  \]
and it satisfies 
$\opn{H}^0(\opn{fadj}_{u, P, N}) (\la) \circ \opn{q}_{u, P} = \la$
for any forward morphism $\la : P \to N$ in $\dcat{K}(A)$. 
We define 
\[ \opn{fadj}_{u, P, N}^{\mrm{L}} := \opn{Q}_B \circ 
\opn{H}^0(\opn{fadj}_{u, P, N}) \circ \opn{Q}_A^{-1} . \]

\medskip \noindent
(3) Since $\opn{fadj}_{u, M, N}^{\mrm{L}}$ is a bifunctorial isomorphism, 
it determines an adjunction. 
\end{proof}

\begin{dfn} \label{dfn:3191}
Let $u : A \to B$ be a DG ring homomorphism, and let 
$M \in \dcat{D}(A)$ and $N \in \dcat{D}(B)$ be DG modules.
A forward morphism $\la : M \to N$ in $\dcat{D}(A)$ over $u$ is called a 
{\em nondegenerate forward morphism}% 
\index{Forward morphism! in derived category}%
\index{Forward morphism! nondegenerate}%
\index{1-Fadj@$\opn{fadj}^{\mrm{L}}_{u, M, N}$}
if the corresponding morphism 
\[ \opn{fadj}_{u, M, N}^{\mrm{L}}(\la) : B \ot^{\mrm{L}}_A M = 
\opn{LInd}_u(M) \to N \] 
in $\dcat{D}(B)$ is an isomorphism.
\end{dfn}

\begin{exa} \label{exa:2404}
Given $M \in \dcat{D}(A)$, let $N := B \ot^{\mrm{L}}_A M \in \dcat{D}(B)$. 
The canonical derived forward morphism%
\index{1-Fadj@$\opn{fadj}^{\mrm{L}}_{u, M, N}$}%
\index{1-QLuM@$\opn{q}_{u, M}^{\mrm{L}}$} 
$\opn{q}_{u, M}^{\mrm{L}} : M \to N$ is a nondegenerate forward morphism
in $\dcat{D}(A)$ over $u$, because 
$\opn{fadj}_{u, M, N}^{\mrm{L}}(\opn{q}_{u, M}^{\mrm{L}}) = \opn{id}_N$.
\end{exa}

\begin{prop} \label{prop:3191}
Suppose $A \xar{u} B \xar{v} C$ are homomorphisms of DG rings,
$L \in \dcat{D}(A)$, $M \in \dcat{D}(B)$ and $N \in \dcat{D}(C)$ are DG 
modules, $\la : L \to M$ is a nondegenerate forward morphism
in $\dcat{D}(A)$ over $u$, and $\mu : M \to N$ is a nondegenerate forward 
morphism in $\dcat{D}(B)$ over $v$. Then 
$\mu \circ \la : L \to  N$
is a nondegenerate forward morphism in $\dcat{D}(A)$ over $v \circ u$.
\end{prop}

\begin{proof}
Since $\la$ is a nondegenerate forward morphism, it induces an isomorphism 
$N \cong B \ot^{\mrm{L}}_A M$. We can thus assume that  
$N = B \ot^{\mrm{L}}_A M$ and that 
$\la = \opn{q}_{u, M}^{\mrm{L}} : M \to B \ot^{\mrm{L}}_A M$.
Similarly, we can assume that 
$\mu = \opn{q}_{v, N}^{\mrm{L}} : N \to C \ot^{\mrm{L}}_B N$.
But now 
$\mu \circ \la  = \opn{q}_{v, N}^{\mrm{L}} \circ \opn{q}_{u, M}^{\mrm{L}} = 
\opn{q}_{v \circ u, L}^{\mrm{L}} : L \to N$,
and this is known to be a nondegenerate forward morphism (see Example 
\ref{exa:2404}). 
\end{proof}

\begin{dfn} \label{dfn:3095}
Let $u : A \to B$ be a homomorphism of DG $\K$-rings.
\begin{enumerate}
\item The {\em coinduction functor}%
\index{Coinduction functor}%
\index{1-CInd@$\opn{CInd}_u$}
is the $\K$-linear DG functor 
$\opn{CInd}_u = \opn{CInd}_{B / A} : \dcat{C}(A) \to \dcat{C}(B)$ 
with formula
$\opn{CInd}_u := \opn{Hom}_{A}(B, -)$.

\item The  {\em derived coinduction functor}%
\index{Coinduction functor! derived}% 
\index{1-RCInd@$\opn{RCInd}_u$}
$\opn{RCInd}_u = \opn{RCInd}_{B / A} : \dcat{D}(A) \to \dcat{D}(B)$  
is the triangulated right derived functor of $\opn{CInd}_u$.
\end{enumerate}
\end{dfn}

Thus for each $M \in \dcat{D}(A)$ we have 
$\opn{RCInd}_u(M) = \opn{RHom}_{A}(B, M)$.
There is the canonical morphism 
$\eta^{\mrm{R}}_{M} : \opn{CInd}_u(M) \to \opn{RCInd}_u(M)$
in $\dcat{D}(B)$, which is part of the right derived functor, 


\begin{dfn} \label{dfn:3110}
Let $u : A \to B$ be a DG ring homomorphism, and let 
$M \in \dcat{D}(A)$ and $N \in \dcat{D}(B)$ be DG modules.
A morphism
$\th : \opn{Rest}_u(N) \to M$
in $\dcat{D}(A)$ is called a {\em backward} (or {\em trace}) {\em morphism in 
$\dcat{D}(A)$ over $u$}%
\index{Backward morphism! in derived category}%
\index{Trace morphism! in derived category}.

Similarly there are backward morphisms over $u$ in the categories 
$\dcat{C}(A)$, $\dcat{C}_{\mrm{str}}(A)$ and $\dcat{K}(A)$.
\end{dfn}

We shall use the terms ``backward morphism'' and ``trace morphism'' 
synonymously. 
We often omit the functor $\opn{Rest}_u$, and just talk about a 
backward morphism $\th : N \to M$ in $\dcat{D}(A)$  over $u$. 

For any DG module $M \in \dcat{C}(A)$ there is a canonical backward homomorphism
$\opn{tr}_{u, M} : \opn{CInd}_u(M) = \opn{Hom}_A(B, M) \to M$,
$\opn{tr}_{u, M}(\phi) := \phi(1_B)$ 
in $\dcat{C}_{\mrm{str}}(A)$ over $u$. 
Now let $N \in \dcat{C}(B)$. The usual change of ring adjunction formula gives 
a canonical isomorphism 
\[ \opn{badj}_{u, M, N} : 
\opn{Hom}_{\dcat{C}(A)}(N, M) \iso 
\opn{Hom}_{\dcat{C}(B)} \bigl( N, \opn{Hom}_A(B, M) \bigr) \]
in $\dcat{C}_{\mrm{str}}(\K)$. It is characterized by the property that 
for every backward morphism $\th : N \to M$ there is equality 
$\opn{tr}_{u, M} \circ \opn{badj}_{u, M, N}(\th) = \th$.
We refer to the isomorphism  $\opn{badj}_{u, M, N}$ as 
{\em backward adjunction}. 
Since the isomorphism $\opn{badj}_{u, M, N}$ is 
functorial in $M$ and $N$, we see that the functor $\opn{CInd}_u$ is a right 
adjoint of $\opn{Rest}_u$. 

The next theorem shows that this is also true on the derived level. 

\begin{thm} \label{thm:3192} 
Let $u : A \to B$ be a homomorphism of DG $\K$-rings.
\begin{enumerate}
\item For each $M \in \dcat{D}(A)$ there is a unique backward morphism 
\[ \opn{tr}_{u, M}^{\mrm{R}} : \opn{RCInd}_u(M) = 
\opn{RHom}_{A}(B, M) \to  M \]
in $\dcat{D}(A)$ over $u$, called the {\em canonical backward morphism}% 
\index{Backward morphism! canonical}%
\index{1-TruM@$\opn{tr}_{u, M}^{\mrm{R}}$},
which is functorial in $M$, and satisfies 
$\opn{tr}_{u, M}^{\mrm{R}}  \circ \, \eta^{\mrm{R}}_M  = \opn{tr}_{u, M}$.

\item For each $M \in \dcat{D}(A)$ and $N \in \dcat{D}(B)$
there is a unique $\K$-linear isomorphism 
\[ \opn{badj}_{u, M, N}^{\mrm{R}} :  
\opn{Hom}_{\dcat{D}(A)} \bigl( \opn{Rest}_u(N), M \bigr) \iso 
\opn{Hom}_{\dcat{D}(B)} \bigl( N, \opn{RCInd}_u(M) \bigr) , \]
called {\em derived backward adjunction}%
\index{Backward adjunction! derived}%
\index{1-Badj@$\opn{badj}^{\mrm{R}}_{u, M, N}$},
such that 
$\opn{tr}_{u, M}^{\mrm{R}} \circ \opn{badj}_{u, M, N}^{\mrm{R}}(\th) = \th$
for all backward morphisms $\th : N \to M$ in $\dcat{D}(A)$ over $u$.
The isomorphism $\opn{badj}_{u, M, N}^{\mrm{R}}$ is functorial in $M$ and $N$. 

\item The functor
$\opn{RCInd}_u : \dcat{D}(A) \to \dcat{D}(B)$
is right adjoint to $\opn{Rest}_u$. 
\end{enumerate}
\end{thm}

Here are the commutative diagrams in $\dcat{D}(A)$ illustrating the theorem. 
\[ \UseTips \xymatrix @C=8ex @R=6ex {
\opn{CInd}_u(M)
\ar[r]^{\eta^{\mrm{R}}_M}
\ar[dr]_{\opn{tr}_{u, M}} 
&
\opn{RCInd}_u(M)
\ar[d]^{\opn{tr}_{u, M}^{\mrm{R}}}
\\
&
M
} 
\qquad \quad 
\UseTips \xymatrix @C=14ex @R=6ex {
N
\ar[r]^(0.4){\opn{badj}_{u, M, N}^{\mrm{R}}(\th)}
\ar[dr]_{\th} 
& 
\opn{RCInd}_u(M)
\ar[d]^{\opn{tr}_{u, M}^{\mrm{R}}}
\\
&
M
} 
\]

\begin{proof}
The proof is very similar to that of Theorem \ref{thm:3085}. The only change is 
that now we choose a system of K-injective resolutions in $\dcat{C}(A)$. 
For a K-injective DG $A$-module $I$, the DG $B$-module
$\opn{Hom}_A(B, I)$ is K-injective, and it represents
$\opn{RCInd}_u(M)$. 
\end{proof}

\begin{dfn} \label{dfn:3190} 
Let $u : A \to B$ be a DG ring homomorphism, and let 
$M \in \dcat{D}(A)$ and $N \in \dcat{D}(B)$ be DG modules.
A backward morphism $\th : N \to M$ in $\dcat{D}(A)$ over $u$ is called  
{\em a nondegenerate backward morphism}% 
\index{Backward morphism! nondegenerate}%
\index{Backward morphism! in derived category}%
\index{1-Badj@$\opn{badj}^{\mrm{R}}_{u, M, N}$}
if the corresponding morphism
\[ \opn{badj}_{u, M, N}^{\mrm{R}}(\th) : N \to \opn{RHom}_A(B, M) 
= \opn{RCInd}_u(M) \]
in $\dcat{D}(B)$ is an isomorphism.
\end{dfn}

\begin{exa} \label{exa:2402}
Given $M \in \dcat{D}(A)$, let 
$N := \opn{RHom}_{A}(B, M) \in \dcat{D}(B)$. 
The canonical backward morphism 
$\opn{tr}_{u, M}^{\mrm{R}} : N \to M$%
\index{1-Badj@$\opn{badj}^{\mrm{R}}_{u, M, N}$}%
\index{1-TruM@$\opn{tr}_{u, M}^{\mrm{R}}$} 
is a nondegenerate backward morphism
in $\dcat{D}(A)$ over $u$, because 
$\opn{badj}_{u, M, N}^{\mrm{R}}(\opn{tr}_{u, M}^{\mrm{R}}) = \opn{id}_N$.
\end{exa}

\begin{prop} \label{prop:3190}
Suppose $A \xar{u} B \xar{v} C$ are homomorphisms of DG rings,
$L \in \dcat{D}(A)$, $M \in \dcat{D}(B)$ and $N \in \dcat{D}(C)$ are DG 
modules, $\th : M \to L$ is a nondegenerate backward morphism
in $\dcat{D}(A)$ over $u$, and $\ze : N \to M$ is a nondegenerate backward 
morphism in $\dcat{D}(B)$ over $v$. Then 
$\th \circ \ze : N \to  L$
is a nondegenerate backward morphism in $\dcat{D}(A)$ over $v \circ u$.
\end{prop}

\begin{proof}
Since $\th$ is a nondegenerate backward morphism, it induces an isomorphism 
$M \cong \opn{RHom}_{A}(B, L)$. We can thus assume that  
$M = \opn{RHom}_{A}(B, L)$ and that 
$\th = \opn{tr}_{u, L}^{\mrm{R}}$.
Similarly, we can assume that 
$N = \opn{RHom}_{B}(C, M)$ and that 
$\ze = \opn{tr}_{v, M}^{\mrm{R}}$.
But now 
$\th \circ \ze  =
\opn{tr}_{u, L}^{\mrm{R}}  \circ \opn{tr}_{v, M}^{\mrm{R}} = 
\opn{tr}_{v \circ u, L}^{\mrm{R}} : N \to L$,
and this is known to be a nondegenerate backward morphism (see Example 
\ref{exa:2402}). 
\end{proof}

\begin{exa} \label{exa:2403}
If $A = B$ and $u = \opn{id}_A$, then 
backward and forward morphisms over $u$ are just morphisms in $\dcat{D}(A)$. 
Nondegenerate (backward or forward) morphisms are just isomorphisms in 
$\dcat{D}(A)$. 
\end{exa}



\mysubsection{DG Ring Quasi-Isomorphisms}
\label{subsec:DGR-quisoms}

Suppose $A$ and $B$ are DG $\K$-rings. A homomorphism of DG $\K$-rings 
$u : A \to B$ induces a homomorphism of graded $\K$-rings 
$\opn{H}(u) : \opn{H}(A) \to \opn{H}(B)$; cf.\ Example \ref{exer:1100}.

\begin{dfn} \label{dfn:3085}
A DG ring homomorphism $u : A \to B$ is called a 
{\em quasi-isomorphism of DG rings}% 
\index{Differential graded ring! quasi-isomorphism of {\indash}s}
if $\opn{H}(u)$ is an isomorphism. 
\end{dfn}

Here is a fundamental result, due to B. Keller in 1994 \cite{Kel1}, 
and independently to V. Hinich in 1997 \cite{Hi1}. 
It is the justification behind the use of DG ring resolutions
(see Remarks \ref{rem:2285}, \ref{rem:1420} and \ref{rem:3691}). 

\begin{thm} \label{thm:2363}
Let $u : A \to B$ be a quasi-isomorphism of DG $\K$-rings%
\index{Differential graded ring! quasi-isomorphism of {\indash}s}.
Then\tup{:}
\begin{enumerate}
\item The restriction functor 
$\opn{Rest}_u : \dcat{D}(B) \to \dcat{D}(A)$
is an equivalence of $\K$-linear triangulated categories, with quasi-inverse
$\opn{LInd}_u$. 

\item For every $M \in \dcat{D}(B^{\mrm{op}})$ and 
$N \in \dcat{D}(B)$ there is an  isomorphism
$M \ot^{\mrm{L}}_A N \iso M \ot^{\mrm{L}}_B N$
in $\dcat{D}(\K)$. This isomorphism is functorial in $M$ and $N$.

\item For every $M, N \in \dcat{D}(B)$ there is an isomorphism 
$\opn{RHom}_{A}(M, N) \iso \lb \opn{RHom}_{B}(M, N)$
in $\dcat{D}(\K)$. This isomorphism is functorial in $M$ and $N$.
\end{enumerate}
\end{thm}

Notice that the restriction functor $\opn{Rest}_u$ is suppressed in parts (2) 
and (3) of the theorem. The proof hinges on the next lemma. 

\begin{lem} \label{lem:4770}
Let $u : A \to B$ be a quasi-isomorphism of DG $\K$-rings, let 
$N \in \dcat{D}(B)$, and let $\rho : P \to N$ be a K-projective resolution in 
$\dcat{C}_{\mrm{str}}(A)$. 
\begin{enumerate}
\item The homomorphism
$\opn{q}_{u, P} : P \to B \ot_A P$
in $\dcat{C}_{\mrm{str}}(A)$ from equation \lb \tup{(\ref{eqn:3194})}, and the 
homomorphism 
$\psi : B \ot_A P \to N$, $\psi(b \ot p) := b \cd \rho(p)$,  
in $\dcat{C}_{\mrm{str}}(B)$, are both quasi-isomorphisms.

\item $B \ot_A P$ is a K-projective DG $B$-module.
\end{enumerate}
\end{lem}

\begin{proof} \mbox{}

\smallskip \noindent
(1) Look at the commutative diagram 
\[ \UseTips \xymatrix @C=8ex @R=6ex {
&
B \ot_A P 
\ar[r]^(0.56){\psi}
&
N
\\
P
\ar[r]^(0.4){\cong}
\ar[ur]^{\opn{q}_{u, P}}
&
A \ot_A P
\ar[u]_{u \, \ot \, \opn{id}_P}
\ar[r]^{\opn{id}_A \ot \, \rho}
&
A \ot_A N 
\ar[u]_{\cong}
} \]
in $\dcat{C}_{\mrm{str}}(A)$, in which the arrows marked $\cong$ are the 
canonical isomorphisms. 
The homomorphism $u \ot \opn{id}_P$ is a quasi-isomorphism, because $u$ is a 
quasi-isomorphism and $P$ is K-flat over $A$. 
The homomorphism $\opn{id}_A \ot \, \rho$ is a quasi-isomorphism, because 
$\rho$ is a quasi-isomorphism and $A$ is K-flat over itself. Therefore
$\opn{q}_{u, P}$ and $\psi$ are quasi-isomorphisms.

\medskip \noindent 
(2) Take an acyclic DG $B$-module $N$. There is an isomorphism 
\[ \opn{Hom}_B(B \ot_A P, N) \cong \opn{Hom}_A(P, N) \]
in $\dcat{C}_{\mrm{str}}(\K)$, coming from adjunction for the DG ring 
homomorphism $u$. By assumption the DG module $\opn{Hom}_A(P, N)$ is acyclic. 
This shows that $B \ot_A P$ is K-projective over $B$. 
\end{proof}

\begin{proof}[Proof of Theorem \tup{\ref{thm:2363}}] \mbox{}

\smallskip \noindent
(1) Take any $N \in \dcat{D}(B)$. 
Choose a K-projective resolution $\rho : P \to N$ in 
$\dcat{C}_{\mrm{str}}(A)$, so that 
$(\opn{LInd}_u \circ \opn{Rest}_u)(N) \cong B \ot_A P$
in $\dcat{D}(B)$. By the lemma we have a quasi-isomorphism 
$\psi : B \ot_A P \to N$ in $\dcat{C}_{\mrm{str}}(B)$.
This means that we have an isomorphism
\[ \opn{Q}(\psi) : (\opn{LInd}_u \circ \opn{Rest}_u)(N) \iso N \]
in $\dcat{D}(B)$, and it is functorial in $N$. 

On the other hand, starting from a DG module $M \in \dcat{D}(A)$, 
and choosing a K-projective resolution $\rho : P \to M$
in $\dcat{C}_{\mrm{str}}(A)$, we have an isomorphism 
$\opn{LInd}_u(M) \cong B \ot_A P$
in $\dcat{D}(B)$. So there is an isomorphism 
$(\opn{Rest}_u \circ \opn{LInd}_u)(M) \cong \lb B \ot_A P$
in $\dcat{D}(A)$. According to the lemma there is a quasi-isomorphism 
$\opn{q}_{u, P} : P \to B \ot_A P$
in $\dcat{C}_{\mrm{str}}(A)$.
We obtain an isomorphism 
\[ \opn{Q}(\rho) \circ \opn{Q}(\opn{q}_{u, P})^{-1} : 
(\opn{Rest}_u \circ \opn{LInd}_u)(M) \iso M \]
in $\dcat{D}(A)$, and it is functorial in $M$. 

\medskip \noindent
(2) Choose a K-projective resolution $\rho : P \to M$ in 
$\dcat{C}(A^{\mrm{op}})$. 
This produces an isomorphism 
$\phi_1 : P \ot_A N \iso M \ot^{\mrm{L}}_A N$
in $\dcat{D}(\K)$. The DG module 
$P \ot_A B \in \dcat{C}(B^{\mrm{op}})$
is K-projective, and there is a quasi-isomorphism
$\psi : P \ot_A B \to M$ in 
$\dcat{C}_{\mrm{str}}(B^{\mrm{op}})$,
by the lemma (with $B^{\mrm{op}}$ instead of $B$). 
We get an isomorphism 
\[ \phi_2 := \opn{Q}(\psi) \ot^{\mrm{L}}_B \opn{id}_N : 
(P \ot_A B) \ot_B N \iso M \ot^{\mrm{L}}_B N \]
in $\dcat{D}(\K)$.
There is also the canonical isomorphism 
\[ \phi_3 : (P \ot_A B) \ot_B N \iso P \ot_A N \]
in  $\dcat{D}(\K)$. 
The functorial isomorphism we want is 
\[ \phi_2 \circ \phi_3^{-1} \circ \phi_1^{-1} : 
M \ot^{\mrm{L}}_A N \iso M \ot^{\mrm{L}}_B N \]
in $\dcat{D}(\K)$. 

\medskip \noindent
(3) Let us choose a K-projective resolution 
$\rho : P \to M$ in $\dcat{C}(A)$. 
This produces an isomorphism 
\[ \phi_1 : \opn{RHom}_{A}(M, N) \iso \opn{Hom}_{A}(P, N)  \]
in $\dcat{D}(\K)$. The DG module 
$B \ot_A P \in \dcat{C}(B)$ is K-projective, and the homomorphism 
$\psi : B \ot_A P \to M$ 
in $\dcat{C}_{\mrm{str}}(B)$ is a quasi-isomorphism (again using the lemma). 
In this way we have an isomorphism 
\[ \phi_2 : \opn{Hom}_{B}(B \ot_A P, N) \iso \opn{RHom}_{B}(M, N) \]
in $\dcat{D}(\K)$. And there is the canonical isomorphism of adjunction 
\[ \phi_3 : \opn{Hom}_{B}(B \ot_A P, N) \iso \opn{Hom}_{A}(P, N) \]
in $\dcat{D}(\K)$.
The functorial isomorphism we want is
\[ \phi_2 \circ \phi_3^{-1} \circ \phi_1 :
\opn{RHom}_{A}(M, N) \iso \opn{RHom}_{B}(M, N) . \qedhere \]
\end{proof}







\mysubsection{Existence of DG Ring Resolutions}
\label{subsec:DG-ring-res}

In Theorem \ref{thm:2363} we saw that a DG ring quasi-isomorphism 
$A \to B$ does not change the derived categories: the restriction functor 
$\dcat{D}(B) \to \dcat{D}(A)$ is an equivalence of triangulated categories. 

Sometimes it is advantageous to replace a given DG ring by a quasi-isomorphic 
one (see Remarks \ref{rem:4605}, \ref{rem:4945}, \ref{rem:2285}, 
\ref{rem:1420} and \ref{rem:3691}). 
In this subsection we study {\em semi-free DG rings}, and 
{\em semi-free DG ring resolutions}. 

As before, we fix a nonzero commutative base ring $\K$.
All DG rings are $\K$-central, namely we work 
in the category $\catt{DGRng} \centover \K$. We use the shorthand 
``NC'' for ``noncommutative''. 

Graded sets and filtered graded sets were introduced in Definitions 
\ref{dfn:4730} and \ref{dfn:4731} respectively. 
The next definition was mentioned briefly in Example \ref{exa:4625}.

\begin{dfn} \label{dfn:4306}
Let $X$ be a graded set. The {\em NC polynomial ring on $X$} is the graded 
$\K$-ring $\K \bra{X}$. 
\end{dfn}

Perhaps we should say a few words on the structure of the ring $\K \bra{X}$. 
Given $n \geq 0$ and elements 
$x_{1}, \ldots, x_{n} \in X$, their product
$x_{1} \cdots x_{n} \in \K \bra{X}$
is called a monomial (or a word), and its degree is 
\begin{equation} \label{eqn:4305}
\opn{deg}(x_{1} \cdots x_{n}) := \opn{deg}(x_{1}) + \cdots + 
\opn{deg}(x_{n}) \in \Z .
\end{equation}
For $n = 0$ the unique monomial is denoted by $1$. 
As a graded $\K$-module, $\K \bra{X}$ is free, with basis the set of 
monomials. Multiplication in $\K \bra{X}$ is by concatenation of monomials, 
extended $\K$-bilinearly. 

\begin{dfn} \label{dfn:4308}
Let $(X, F)$ be a filtered graded set. There is an induced filtration
$F = \bigl\{ F_{j} \bigl( \K \bra{X} \bigr) \bigr\}_{j \geq -1}$ on 
$\K \bra{X}$ by $\K$-submodules, as follows:
\[ F_j \bigl( \K \bra{X} \bigr) := 
\begin{cases}
0 & \tup{if} \quad j = -1 
\\
\K \bra{F_j(X)} & \tup{if} \quad j \geq 0  .   
\end{cases} \]
\end{dfn}

For every $j \geq 0$ the filtered piece 
$F_j \bigl( \K \bra{X} \bigr)$ is a graded subring of $\K \bra{X}$. And of 
course 
$\K \bra{X} = \bigcup_{j} F_j \bigl( \K \bra{X} \bigr)$.

Recall that given a DG ring 
$A$, we denote by $A^{\natural}$ the graded ring gotten by forgetting the 
differential $\d_A$. Here is the main definition of this subsection. 

\begin{dfn} \label{dfn:4309}
A {\em NC semi-free DG $\K$-ring}%
\index{Differential graded ring! NC semi-free}
is a DG $\K$-ring $A$ that 
admits an isomorphism of graded $\K$-rings
$A^{\natural} \cong \K \bra{X}$,
for some filtered graded set $(X, F)$, such that under this isomorphism there 
is an inclusion 
$\d_A \bigl( F_j(X) \bigr) \sub F_{j - 1} \bigl( \K \bra{X} \bigr)$ 
for every $j \geq 0$. Such a filtered graded set $(X, F)$ is called a 
{\em multiplicative semi-basis}%
\index{Semi-basis! multiplicative {\indash} of DG ring}
of $A$. 
\end{dfn}

In this book we are mostly interested in semi-free DG rings because of the 
proposition below. 

\begin{prop} \label{prop:4305}
If $A$ is a NC semi-free DG $\K$-ring, then as a DG $\K$-module $A$ is 
semi-free, and hence K-flat. 
\end{prop}

\begin{proof}
Suppose we are given an isomorphism 
$A^{\natural} \cong \K \bra{X}$
for some filtered graded set $(X, F)$, as in Definition \ref{dfn:4309}.
Let $Y$ be the set of monomials in the elements of $X$. The set $Y$ is 
graded, see (\ref{eqn:4305}), and it forms a basis of $A^{\natural}$ as a 
graded $\K$-module. However, since $X$ could have elements of positive degree, 
this is not enough to make $A$ into a semi-free DG $\K$-module; we need to put 
a suitable filtration on $Y$. Cf.\ Proposition \ref{prop:4730}.

For an element $x \in X$ let 
$\opn{ord}^F(x) := \opn{min} \{ j \mid x \in F_j(X) \} \in \N$.
Thus 
$F_j(X) = \{ x \in X \mid \opn{ord}^F(x) \leq j \}$.
We extend this order function to monomials by 
$\opn{ord}^F(x_{1} \cdots x_{n}) := \sum_{i} \opn{ord}^F(x_{i})$.
The element $1$ (the monomial of length $0$) gets order $0$. 
Then we define 
$G_j(Y) := \{ y \in Y \mid \opn{ord}^F(y) \leq j \}$.
This is a filtration 
$G = \{ G_{j}(Y) \}_{j \geq -1}$ of the graded set $Y$, and it induces a 
filtration $G = \{ G_{j}(A) \}_{j \geq -1}$ on $A$ 
by $\K$-submodules, where $G_{j}(A) \sub A$ is defined to be the $\K$-linear 
span of $G_{j}(Y)$. 

The graded Leibniz rule, together with the inclusion 
$\d_A (F_j(X)) \sub \lb F_{j - 1}(\K \bra{X})$
from Definition \ref{dfn:4309}, imply that 
$\d_A (G_{j}(A)) \sub G_{j - 1}(A)$
for every $j \geq 0$. And 
$\opn{Gr}^G_j(A)$ is a free DG $\K$-module, with basis 
$\opn{Gr}_j^G(Y) := \{ y \in Y \mid \opn{ord}^F(y) = j \}$.
Thus $G$ is a semi-free filtration on the DG $\K$-module $A$, and $A$ is a 
semi-free DG $\K$-module; see Definition \ref{dfn:1575}.

Finally, according to Theorem \ref{thm:1575} and Proposition \ref{prop:1525}, 
$A$ is a K-flat DG $\K$-module. 
\end{proof}

\begin{dfn} \label{dfn:4310}
Let $A$ be a DG $\K$-ring. A {\em NC semi-free DG ring resolution} of 
$A$ relative to $\K$ is a quasi-isomorphism $u : \til{A} \to A$ 
of DG $\K$-rings, where $\til{A}$ is a NC semi-free DG $\K$-ring.
We also call $u : \til{A} \to A$ a {\em NC semi-free resolution of $A$ in 
$\catt{DGRng} \centover \K$}. 
\end{dfn}

The important result of this subsection is the next existence theorem. 

\begin{thm} \label{thm:4310}
Let $A$ be a DG $\K$-ring. There exists a NC semi-free DG ring
resolution $u : \til{A} \to A$ of $A$ relative to $\K$. 
\end{thm}

The proof of the theorem comes after two lemmas. 
If $X$ and $Y$ are graded sets, a degree $i$ function $f : X \to Y$ is a 
function such that $f(X^j) \sub Y^{j + i}$ for all $j$. Likewise,
a degree $i$ function $f : X \to B$ to a graded ring $B$ is a function 
such that $f(X^j) \sub B^{j + i}$ for all $j$.

\begin{lem} \label{lem:4315}
Let $X$ be a graded set, and consider the graded $\K$-ring 
$B := \K \bra{X}$. Suppose $\d_X : X \to B$ is a degree $1$ function.
\begin{enumerate}
\item  The function $\d_X$ extends uniquely to a $\K$-linear degree $1$ 
derivation $\d_B : B \to B$. 

\item The derivation $\d_B$ is a differential on $B$, i.e.\ 
$\d_B \circ \d_B = 0$, if and only if $\d_B(\d_X((x)) = 0$ for every 
$x \in X$. 
\end{enumerate}
\end{lem}

\begin{proof} \mbox{}

\smallskip \noindent
(1) For elements $x_1, \ldots, x_n \in X$, with 
$k_j := \opn{deg}(x_j)$, define 
\begin{equation} \label{eqn:4315}
\d_B(x_1 \cdots x_n) := \sum_{j = 1}^n \, 
(-1)^{k_1 + \cdots + k_{j - 1}} \cd x_1 \cdots x_{j - 1} \cd 
\d_X(x_j) \cd x_{j + 1} \cdots x_{n} . 
\end{equation}
This extends uniquely to a $\K$-linear homomorphism $\d_B : B \to B$ of degree 
$1$. A calculation shows that the graded Leibniz rule holds, so this is  
degree $1$ derivation. 

\medskip \noindent
(2) It is enough to verify that $\d_B \circ \d_B = 0$ on monomials. 
This is an elementary calculation, using formula (\ref{eqn:4315}). 
\end{proof}

\begin{lem} \label{lem:4316}
Let $B, C \in \catt{DGRng} \centover \K$. Assume that $B$ is NC semi-free, with 
a multiplicative semi-basis $(X, F)$.  
Let $w : X \to C$ be a degree $0$ function. 
\begin{enumerate}
\item  The function $w$ extends uniquely to a graded $\K$-ring homomorphism 
$w : B^{\natural} \to C^{\natural}$. 

\item The graded ring homomorphism $w$ becomes a DG ring homomorphism 
$w : B \to C$ if and only if 
$\d_C(w(x)) = w(\d_B(x))$ for every $x \in X$. 
\end{enumerate}
\end{lem}

\begin{proof} 
Item (1) is clear. As for item (2), it is enough to verify that $\d_C \circ w = 
w \circ \d_B$ on monomials $x_1 \cdots x_n \in B$. This is done by induction on 
$n$, using the graded Leibniz rule. 
\end{proof}

\begin{proof}[Proof of Theorem \tup{\ref{thm:4310}}]
The strategy is this: we are going to construct an increasing sequence of 
semi-free DG rings 
\begin{equation} \label{eqn:4945}
F_0(\til{A}) \sub F_1(\til{A}) \sub F_2(\til{A}) \sub \cdots ,
\end{equation}
together with a compatible system of DG ring homomorphisms 
$F_j(u) : \lb F_j(\til{A}) \to A$.
At the same time we will construct an increasing sequence of graded sets
\begin{equation} \label{eqn:4946}
\varnothing = F_{-1}(X) \sub F_0(X) \sub F_1(X) \sub F_2(X) \sub \cdots ,
\end{equation}
such that there is equality of graded rings
$F_j(\til{A})^{\natural} = \K \bra{F_j(X)}$
for every $j \geq 0$. 

Then we will define the DG ring
\begin{equation} \label{eqn:4310}
\til{A} := \lim_{j \to} \,  F_j(\til{A}) , 
\end{equation}
the DG ring homomorphism
\begin{equation} \label{eqn:4311}
u := \lim_{j \to} \,  F_j(u) : \til{A} \to  A ,
\end{equation}
and the graded set 
\begin{equation} \label{eqn:4775}
X := \lim_{j \to} \,  F_j(X) . 
\end{equation}
The graded set $X$ is filtered by 
$F := \bigl\{ F_j(X) \bigr\}_{j \geq -1}$. By the construction 
$(X, F)$ will be a multiplicative semi-basis of $\til{A}$, so that the latter 
is a semi-free DG ring. The construction will also ensure that $u$ is a 
quasi-isomorphism. All this will be done in several steps. 

\medskip \noindent
Step 1. In this step we deal with $j = 0$. 
Let's choose a collection $\{ \bar{a}_x \}_{x \in F_0(X)}$
of nonzero homogeneous elements of $\opn{H}(A)$ that generates it as a 
$\K$-ring. We make $F_0(X)$ into a graded set by declaring
$\opn{deg}(x) := \opn{deg}(\bar{a}_x)$ for $x \in F_0(X)$.
Next, for each $x$ we choose a homogeneous cocycle 
$a_x \in \opn{Z}(A)$ that represents the cohomology class $\bar{a}_x$. 
In this way we obtain a degree $0$ function
$F_0(u) : F_0(X) \to A$, $F_0(u)(x) := a_x$. 

Define the DG ring 
$F_0(\til{A}) := \K \bra{F_0(X)}$, with zero differential. 
According to Lemma \ref{lem:4316} the function 
$F_0(u) : F_0(X) \to A$ extends uniquely to a DG ring homomorphism 
$F_0(u) : F_0(\til{A}) \to A$.
The choice of the collection $\{ \bar{a}_x \}_{x \in F_0(X)}$
guarantees that the graded $\K$-ring homomorphism 
$\opn{H}(F_0(u)) : \opn{H}(F_0(\til{A})) \to \opn{H}(A)$ 
is surjective.

\medskip \noindent
Step 2. In this step and the next one we deal with the induction. Let
$j \geq 0$. The assumption is that we have a 
semi-free DG ring $F_j(\til{A})$, with a multiplicative semi-basis
\begin{equation} \label{eqn:4776}
\varnothing = F_{-1}(X) \sub F_0(X) \sub \cdots \sub F_{j - 1}(X) \sub 
F_{j}(X) \sub F_{j}(X) \sub \cdots  
\end{equation}
(constant from level $j$ onward), and a DG ring homomorphism 
$F_j(u) : F_j(\til{A}) \to A$
such that the graded ring homomorphism 
$\opn{H}(F_j(u)) : \opn{H}(F_j(\til{A})) \to \opn{H}(A)$ 
is surjective.

Define the graded two-sided ideal 
$K_j := \opn{Ker} \bigl( \opn{H}(F_j(u)) \bigr) \sub \opn{H}(F_j(\til{A}))$.
So there is an exact sequence of graded $\opn{H}(F_j(\til{A}))$-bimodules 
\begin{equation} \label{eqn:4316}
0 \to K_j \to \opn{H}(F_j(\til{A})) \xar{\opn{H}(F_j(u))}
\opn{H}(A) \to  0 . 
\end{equation}
Choose a collection $\{ \bar{a}_y \}_{y \in Y_{j + 1}}$
of nonzero homogeneous elements of $K_j$ that generates it as an 
$\opn{H}(F_j(\til{A}))$-bimodule (i.e.\ as a two-sided ideal). 
We make $Y_{j + 1}$ into a graded set by declaring
$\opn{deg}(y) := \opn{deg}(\bar{a}_y) - 1$ 
for $y \in Y_{j + 1}$.
Next, for each element $y \in Y_{j + 1}$ we choose a homogeneous cocycle 
$a_y \in \opn{Z} \bigl( F_j(\til{A}) \bigr)$ that represents the cohomology 
class $\bar{a}_y$. In this way we obtain 
a collection $\{ a_y \}_{y \in Y_{j + 1}}$ of homogeneous cocycles in 
$F_j(\til{A})$. 

Define the graded set 
$F_{j + 1}(X) := F_{j}(X) \sqcup Y_{j + 1}$
and the graded ring \lb 
$F_{j + 1}(\til{A})^{\natural} := \K \bra{ F_{j + 1}(X) }$.
By definition $F_{j}(\til{A})^{\natural}$ is a graded subring of 
$F_{j + 1}(\til{A})^{\natural}$. 
Next, define the degree $1$ function 
% \[ \begin{aligned}
% &
$\d_{F_{j + 1}(X)} : F_{j+ 1}(X) \to F_{j + 1}(\til{A})$
by the formula  
% \\
% &
\[ \d_{F_{j + 1}(X)}(x) := 
\begin{cases}
\d_{F_{j}(\til{A})}(x) & \tup{if} \ \ x \in F_{j}(X)
\\
a_x & \tup{if} \ \ x \in Y_{j + 1} .   
\end{cases} \]
% \end{aligned} \]
In this formula we are using the collection of homogeneous cocycles
$\{ a_y \}_{y \in Y_{j + 1}}$ that has been chosen above. 
According to Lemma \ref{lem:4315} the function $\d_{F_{j + 1}(X)}$ extends 
uniquely to a differential on $F_{j + 1}(\til{A})$, making it into a DG ring. 
And of course $F_{j}(\til{A}) \sub F_{j + 1}(\til{A})$
is a DG subring. 

Let us denote by 
$v_j : F_{j}(\til{A}) \to F_{j + 1}(\til{A})$
the inclusion. 
Because the cocycles $a_y \in F_{j}(\til{A})$, for $y \in Y_{j + 1}$, become 
coboundaries in $F_{j + 1}(\til{A})$, 
we see that 
\begin{equation} \label{eqn:4317}
K_j \sub \opn{Ker} (\opn{H}(v_j)) \sub \opn{H}(F_{j}(\til{A})) .
\end{equation}

Observe that the DG ring $F_{j + 1}(\til{A})$ is semi-free, with 
multiplicative semi-basis
\[ \varnothing = F_{-1}(X) \sub F_0(X) \sub \cdots \sub F_{j}(X) \sub 
F_{j + 1}(X) \sub F_{j + 1}(X) \sub \cdots , \]
constant from level $j + 1$ onward.

\medskip \noindent
Step 3. In this step we continue step 2, to construct the DG ring homomorphism 
$F_{j + 1}(u) : F_{j + 1}(\til{A}) \to A$.
For each element $y \in Y_{j + 1}$ the cohomology class 
$\bar{a}_y \in \opn{H}(F_{j}(\til{A}))$ belongs to the kernel of 
$\opn{H}(F_j(u))$. This means that the cocycle 
$F_j(u)(a_y) \in A$ is a coboundary. So we can find a homogeneous element
$b_y \in A$ such that 
$\d_A(b_y) = F_j(u)(a_y)$. 
Define the degree $0$ function 
$F_{j + 1}(u) : F_{j + 1}(X) \to A$
by the formula 
\[ F_{j + 1}(u)(x) := 
\begin{cases}
F_{j}(u)(x) & \tup{if} \ \ x \in F_{j}(X) 
\\
b_{x} & \tup{if} \ \ x \in Y_{j + 1}  .   
\end{cases} \]
By Lemma \ref{lem:4316} this function extends uniquely to a DG ring 
homomorphism 
$F_{j + 1}(u) : F_{j + 1}(\til{A}) \to A$.
Note that the restriction of $F_{j + 1}(u)$ to $F_{j}(\til{A})$
coincides with $F_{j}(u)$. 

\medskip \noindent
Step 4. Applying steps 2-3 recursively, we obtain direct systems 
$\{ F_j(\til{A}) \}_{j \geq 0}$ and $\{ F_j(u) \}_{j \geq 0}$
of DG rings  and homomorphisms, and we can define the 
DG ring $\til{A}$ by formula (\ref{eqn:4310}), and the DG ring homomorphism 
$\til{u}$ by formula (\ref{eqn:4311}). 
The graded set $X$ from formula (\ref{eqn:4775}) is filtered, and by the 
construction in steps 1 and 2 the filtered graded set $(X, F)$ is a 
multiplicative semi-basis of $\til{A}$. 

It remains to prove that the graded ring homomorphism
$\opn{H}(u) :  \opn{H}(\til{A}) \to \opn{H}(A)$
is bijective. 
Recall that for each $j \geq0$ we denoted by $v_j$ the inclusion
$F_{j}(\til{A}) \sub F_{j + 1}(\til{A})$. 
Define the graded ring
$L^j := \opn{Im} (\opn{H}(v_j)) \sub \opn{H}(F_{j + 1}(\til{A}))$. 
We get a commutative diagram 
\[ \UseTips \xymatrix @C=3.5ex @R=5.5ex {
0 
\ar[r]
&
*++{K_j}
\ar@{>->}[rr]^(0.4){\mrm{inc}}
&
&
\opn{H}(F_{j}(\til{A}))
\ar@{->>}[rr]^{\opn{H}(F_{j}(u))}
\ar@{->>}[d]_{\al_j}
\ar@(l,l)[dd]_{\opn{H}(v_j)}
&
&
*++{\opn{H}(A)}
\ar[r]
&
0
\\
&
&
&
*++{L_j}  
\ar@{>->}[d]_{\mrm{inc}}
\ar@{->>}[urr]_{\be_j}
\\
&
&
&
\opn{H}(F_{j + 1}(\til{A}))
\ar@{->>}@(r,d)[uurr]_{\opn{H}(F_{j + 1}(u))}
} \]
in $\dcat{G}_{\mrm{str}}(\K)$. The top row is an exact sequence (it is 
(\ref{eqn:4316})). 
Because $\al_j$ is surjective, there is equality 
$\opn{Ker}(\be_j) = \al_j \bigl( \opn{Ker}(\opn{H}(F_{j}(u)) \bigr) 
= \al_j(K_j)$.
But by formula (\ref{eqn:4317}) we know that 
$\al_j(K_j) = 0$. The conclusion is that 
$\be_j : L_j \to \opn{H}(A)$ is bijective. 
Hence $\{ L_j \}_{j \geq 0}$ is a constant direct system, and 
$\lim_{j \to} \be_j : \lim_{j \to} L_j \to \opn{H}(A)$
is bijective.
 
Now the direct systems 
$\{ L_j \}_{j \geq 0}$ and 
$\bigl\{ \opn{H}(F_{j}(\til{A})) \bigr\}_{j \geq 0}$
are sandwiched. Therefore the second direct system has a limit too, and it is 
the same limit; i.e.\ 
\[ \lim_{j \to} \, \opn{H}(F_{j}(u)) : 
\lim_{j \to} \, \opn{H}(F_{j}(\til{A})) \to \opn{H}(A) \]
is bijective. Because cohomology commutes with direct limits, 
and by formula (\ref{eqn:4310}), this implies that $\opn{H}(u)$
is bijective.
\end{proof}

\begin{exer} \label{exer:4315}
Modify the proof of Theorem \ref{thm:4310} to make the quasi-iso\-morphism 
$u : \til{A} \to A$ surjective. 
\end{exer}

\begin{cor} \label{cor:4605}
If $\opn{H}(A)$ is nonpositive, then $A$ has a semi-free DG ring resolution 
$\til{A} \to A$  such that $\til{A}$ is nonpositive. 
\end{cor}

\begin{proof}
Under this assumption, the graded set $X$ in the proof of Theorem \ref{thm:4310} 
is nonpositive.
\end{proof}

\begin{rem} \label{rem:4605}
NC semi-free DG rings have important lifting properties within 
$\catt{DGRng} \centover \K$.
V. Hinich, in \cite{Hi1}, calls them {\em standard cofibrant objects} of the 
Quillen model structure on $\catt{DGRng} \centover \K$.

One can also consider the full subcategory $\catt{DGRng}^{\leq 0} \centover \K$
of nonpositive NC DG rings. NC nonpositive semi-free DG rings
are studied in detail in \cite{Ye9}. This is done in a slightly more general 
context: the base ring $\K$ is allowed to be a commutative DG ring (see next 
remark), and $\K \to A$ is a central homomorphism of DG 
rings (i.e.\ the image of $\K$ is in $\opn{Cent}(A)$). 
The lifting properties are the same. 

See Remark \ref{rem:4970} for the related $\mrm{A}_{\infty}$ theory. 
\end{rem}

\begin{rem} \label{rem:4945}
There are also {\em commutative semi-free DG rings}.
Recall that a commutative DG ring in this book means nonpositive and strongly 
commutative; see Definitions \ref{dfn:3091} and \ref{dfn:3090}. The
commutative DG rings form the 
category $\catt{DGRng}^{\leq 0}_{\mrm{sc}} / \K$,
which is a full subcategory of 
$\catt{DGRng} \centover \K$.

A nonpositive graded set $S$ gives rise to the commutative polynomial ring 
$\K[S]$, see Example \ref{exa:4625}. 
A commutative DG ring $\til{A}$ is called commutative semi-free if it admits    
a graded ring isomorphism $\til{A}^{\natural} \cong \K[S]$ for some 
nonpositive graded set $S$. The graded set $S$ is automatically filtered by 
degree, namely $F_j(S) := \bigcup_{i \geq -j} S^i$.
This implies (cf.\ Proposition \ref{prop:4305})
that $\til{A}$ is semi-free as a DG $\K$-module . 

In \cite{Ye9} it is proved that every 
$A \in \catt{DGRng}^{\leq 0}_{\mrm{sc}} / \K$
admits a surjective quasi-isomorphism $\til{A} \to A$ from a 
commutative semi-free DG ring $\til{A}$. 
It is also proved that the commutative semi-free DG rings have good lifting 
properties in $\catt{DGRng}^{\leq 0}_{\mrm{sc}} / \K$. 
The situation studied in \cite{Ye9} is a bit more general: the base ring $\K$ 
is allowed to be itself a commutative DG ring; so the theory becomes 
relative. 

If in Example \ref{exa:4205} we take the topological space $X$ to be a 
single point, then the category 
$\catt{DGRng}^{\leq 0}_{\mrm{sc}} / \K_X$
discussed there coincides with 
$\catt{DGRng}^{\leq 0}_{\mrm{sc}} / \K$,
and the semi-pseudo-free DG $\K_X$-rings are just the semi-free commutative DG 
rings mentioned here. 

Unlike the geometric setting, the category 
$\catt{DGRng}^{\leq 0}_{\mrm{sc}} / \K$
admits a Quillen model structure (at least when $\K$ contains $\Q$),
and the calculus of fractions for the derived category 
$\dcat{D}(\catt{DGRng}^{\leq 0}_{\mrm{sc}} / \K)$
can be deduced from it. 

The fully faithful embedding 
\[ \catt{DGRng}^{\leq 0}_{\mrm{sc}} / \K \to 
\catt{DGRng}^{\leq 0} \centover \K \]
induces a functor 
\[ \dcat{D}(\catt{DGRng}^{\leq 0}_{\mrm{sc}} / \K) \to
\dcat{D}(\catt{DGRng}^{\leq 0} \centover \K) \]
between the derived categories, which presumably is neither full nor 
faithful (but we did not explore this question seriously). 
\end{rem}




\mysubsection{The Derived Tensor-Evaluation Morphism}
\label{subsec:der-tens-eval}

In this subsection we pre\-sent an important theorem (Theorem 
\ref{thm:4320}), that will reappear -- with modifications -- several times in 
the book. Before presenting it, we need a few finiteness and boundedness 
conditions, some old and some new. 

Pseudo-finite semi-free DG modules were introduced in Definition 
\ref{dfn:1925}. The saturated full triangulated subcategory generated by a set 
of objects was defined in Definition \ref{dfn:4916}. 

\begin{dfn} \label{dfn:4470}
Let $A$ be a DG ring. A DG $A$-module $L$ is called 
{\em derived pseudo-finite}% 
\index{Differential graded module! derived pseudo-finite}
if it belongs to the saturated full triangulated subcategory of $\dcat{D}(A)$ 
generated by the pseudo-finite semi-free DG $A$-modules. 
\end{dfn}

\begin{exa} \label{exa:4471}
Suppose $A$ is a nonpositive cohomologically left pseudo-noe\-ther\-ian DG 
ring. Then the derived pseudo-finite DG $A$-modules are precisely the 
objects of $\dcat{D}_{\mrm{f}}^{-}(A)$. One direction is by 
Theorem \ref{thm:3340}, and the other direction is trivial.  
\end{exa}

\begin{exa} \label{exa:4470}
If $L$ is an algebraically perfect DG $A$-module (to be defined later, 
see Definition \ref{dfn:3401}) then it is derived pseudo-finite.
\end{exa}

\begin{rem} \label{rem:4470}
Let $L$ be a DG $A$-module. Following \cite{SGA6}, let us say that 
$L$ is {\em pseudo-coherent}% 
\index{Complex in abelian category! pseudo-coherent}%
\index{Differential graded module! pseudo-coherent}
if $L$ is isomorphic in $\dcat{D}(A)$ to a 
pseudo-finite semi-free DG $A$-module $P$.
Clearly pseudo-coherent implies derived pseudo-finite. According to 
\cite[Lemma tag=064X]{SP} the converse is true when $A$ is a ring; but we do 
not know if the converse is true in general. 

Anyhow, the attribute ``derived pseudo-finite'' is very useful, 
and it can replace ``pseudo-coherent'' for most applications. For instance, 
see  condition (i) in Theorem \ref{thm:4320} below. 
\end{rem}

In Definition \ref{dfn:2390} we saw the notion of bounded below flat 
concentration of a DG module. Here is a refinement of it. 

\begin{dfn} \label{dfn:4341}
Let $A$ be a DG ring. 
\begin{enumerate}
\item We say that a DG $A$-module $P$ has 
{\em tensor displacement}%
\index{Tensor displacement}
inside an integer interval $[i_0, i_1]$ if for every 
$N \in \dcat{C}(A^{\mrm{op}})$ this inclusion of integer intervals holds:
\[ \opn{con}(N \ot_A P) \sub \opn{con}(N) + [i_0, i_1] . \]

\item We say that a DG $A$-module $P$ has
{\em bounded below tensor displacement}% 
\index{Tensor displacement! bounded below}
if it has tensor displacement inside 
the integer interval $[i_0, \infty]$ for some $i_0 \in \Z$. 

\item A DG $A$-module $M$ is said to have 
{\em derived bounded below tensor displacement}%
\index{Tensor displacement! derived bounded below}
if $M$ belongs to the saturated full triangulated subcategory of 
$\dcat{D}(A)$ generated by the K-flat DG $A$-modules $P$ that have bounded 
below tensor displacement.
\end{enumerate}
\end{dfn}


\begin{exa} \label{exa:4601}
The DG $A$-module $M := A$ has derived bounded below tensor displacement, 
regardless of the boundedness of $A$ or $\opn{H}(A)$.

More generally, if $M$ is an algebraically perfect DG $A$-module (see 
Definition \ref{dfn:3401}) then it has derived bounded below tensor 
displacement.
\end{exa}

\begin{exa}  \label{exa:4473}
If $A$ is a ring, and $M$ is a complex of $A$-modules of finite flat dimension, 
then $M$ has derived bounded below tensor displacement.
See Proposition \ref{prop:2390}. 
\end{exa}

\begin{rem} \label{rem:4471} 
Clearly, if $M$ has derived bounded below tensor displacement, then it has 
bounded below flat concentration.
We don't know if the converse is true in general. 
\end{rem}

\begin{rem} \label{rem:4602}
Definition \ref{dfn:4341} can be adapted to other kinds of boundedness 
conditions. It can also be adapted to Hom displacement, in the first or second 
argument, to give refined notions of projective and injective displacements.
\end{rem}
 
\begin{thm}[Tensor-Evaluation] \label{thm:4320}
Let $A$ and $B$ be DG rings, and let $L \in \dcat{D}(A)$, 
$M \in \dcat{D}(A \ot B^{\mrm{op}})$ and $N \in \dcat{D}(B)$
be DG modules. There is a morphism 
\[ \opn{ev}^{\mrm{R, L}}_{L, M, N} : \opn{RHom}_{A}(L, M) \ot^{\mrm{L}}_{B} N 
\to \opn{RHom}_{A}(L, M \ot^{\mrm{L}}_{B} N) \]
in $\dcat{D}(\K)$, called {\em derived tensor-evaluation}%
\index{Tensor-evaluation morphism! derived},
that enjoys the properties below. 
\begin{enumerate}
\item The morphism $\opn{ev}^{\mrm{R, L}}_{L, M, N}$ is functorial in the 
objects $L, M, N$.

\item If all conditions \tup{(i)-(iii)} below hold, then 
$\opn{ev}^{\mrm{R, L}}_{L, M, N}$ is an isomorphism. 
\begin{enumerate}
\rmitem{i} The DG $A$-module $L$ is derived pseudo-finite 
\tup{(}Definition \tup{\ref{dfn:4470})}.

\rmitem{ii} The DG $A$-$B$-bimodule $M$ is cohomologically bounded below
\tup{(}Definition \tup{\ref{dfn:4785}(2))}.

\rmitem{iii} The DG $B$-module $N$ has derived bounded below tensor 
displacement \tup{(}Definition \tup{\ref{dfn:4341}(3))}.
\end{enumerate}
\end{enumerate}
\end{thm}

We need a lemma first. 

\begin{lem} \label{lem:4735}
Suppose we are given\tup{:}
\begin{itemize}
\item A pseudo-finite semi-free DG $A$-module $P$, with pseudo-finite 
semi-free filtration $\{ F_j(P) \}_{j \geq -1}$.

\item A bounded below DG $A$-$B$-bimodule $M$.

\item A DG $B$-module $Q$ with bounded below tensor displacement.
\end{itemize}
Then for every degree $l$ there is an index $j$, depending on $l$, such that 
the canonical homomorphisms  
\[ \bigl( \opn{Hom}_{A}(P, M) \ot_{B} Q \bigr)^l \to
\bigl( \opn{Hom}_{A}(F_{j}(P), M) \ot_{B} Q \bigr)^l \]
and 
\[ \bigl( \opn{Hom}_{A}(P, M \ot_{B} Q) \bigr)^l \to
\bigl( \opn{Hom}_{A}(F_{j}(P), M \ot_{B} Q) \bigr)^l  \]
in $\dcat{M}(\K)$ are bijective. 
\end{lem}

\begin{proof}
In the notation of Definition \ref{dfn:1925}, there are numbers $i_1 \in \Z$ 
and $r_k \in \N$ such that
\begin{equation} \label{eqn:4735}
\bigl( P / F_j(P) \bigr)^{\natural} \cong 
\bigoplus_{k \geq j + 1} \opn{T}^{-i_1 + k}(A^{\natural})^{\oplus r_k}
\end{equation}
in $\dcat{G}_{\mrm{str}}(A^{\natural})$ for all $j \geq 0$. 
In terms of bases, the graded free $A^{\natural}$-module 
$\bigl( P / F_j(P) \bigr)^{\natural}$
has a basis concentrated in degrees $\leq i_1 - j - 1$.

Say $M$ is concentrated in the degree interval $[s_0, \infty]$,
and $Q$ has tensor displacement inside the degree interval $[t_0, \infty]$,
for some $s_0, t_0 \in \Z$. From equation (\ref{eqn:4735}) we see that for 
every $j$, both graded $\K$-modules
\[  \opn{Hom}_{A} \bigl(  P / F_j(P), M  \bigr) \ot_{B} Q 
\quad \tup{and} \quad 
\opn{Hom}_{A} \bigl(  P / F_j(P), M  \ot_{B} Q \bigr) \]
are concentrated in the degree interval 
\[ [-i_1 + j + 1  + s_0 + t_0, \infty] \sub \Z . \]
The split short  exact sequence 
\[ 0 \to F_j(P)^{\natural} \to P^{\natural} \to 
\bigl( P / F_j(P) \bigr)^{\natural} \to 0 \]
in $\dcat{G}_{\mrm{str}}(A)$ tells us that for every 
$j \geq l + i_1 - s_0 - t_0$
the homomorphisms in question are bijective. 
\end{proof}

\begin{proof}[Proof of Theorem \tup{\ref{thm:4320}}] \mbox{} 

\smallskip \noindent
(1) Choose a K-projective resolution 
$P \to L$ in $\dcat{C}_{\mrm{str}}(A)$ and a K-projective resolution 
$Q \to N$ in $\dcat{C}_{\mrm{str}}(B)$. These choices are unique up to 
homotopy. Then 
$\opn{ev}^{\mrm{R, L}}_{L, M, N}$ 
is represented by the obvious homomorphism 
\[ \opn{ev}_{P, M, Q} : \opn{Hom}_{A}(P, M) \ot_{B} Q \to 
\opn{Hom}_{A}(P, M \ot_{B} Q) \]
in $\dcat{C}_{\mrm{str}}(\K)$. 
The morphism $\opn{ev}^{\mrm{R, L}}_{L, M, N}$ in $\dcat{D}(\K)$ does not 
depend on the choices, and hence it is functorial. 

\medskip \noindent 
(2) This is done in several steps. 

\smallskip \noindent 
Step 1. Let $P$ be a finite semi-free DG $A$-module; namely $P$ admits a 
pseudo-finite semi-free filtration $F = \{ F_j(P) \}_{j \geq -1}$
(see Definition \ref{dfn:1925}) such that $P = F_{j_1}(P)$ for some 
$j_1 \in \N$. Then the graded $A^{\natural}$-module $P^{\natural}$ is finite 
free. This implies that 
\[ \opn{ev}_{P, M, Q} : \opn{Hom}_{A}(P, M) \ot_{B} Q \to 
\opn{Hom}_{A}(P, M \ot_{B} Q) \]
is an isomorphism in $\dcat{C}_{\mrm{str}}(\K)$ for all 
$M \in \dcat{C}(A \ot B^{\mrm{op}})$ and $Q \in \dcat{C}(B)$. 

\medskip \noindent 
Step 2. Assume that $M$ is bounded below. 
Let $Q$ be a K-flat DG $B$-module of bounded below tensor displacement, and let 
$P$ be a pseudo-finite semi-free DG $A$-module, with pseudo-finite semi-free 
filtration $F = \{ F_j(P) \}_{j \geq -1}$.
By Lemma \ref{lem:4735}, for every $l \in \Z$ there is an index $j$ such that  
\[ \bigl( \opn{Hom}_{A}(P, M) \ot_{B} Q \bigr)^l =
\bigl( \opn{Hom}_{A}(F_{j}(P), M) \ot_{B} Q \bigr)^l \]
and 
\[ \bigl( \opn{Hom}_{A}(P, M \ot_{B} Q) \bigr)^l  =
\bigl( \opn{Hom}_{A}(F_{j}(P), M \ot_{B} Q) \bigr)^l . \]
By step 1 the homomorphism 
\[ \opn{ev}_{F_{j}(P), M, Q} : 
\opn{Hom}_{A} \bigl( F_{j}(P), M \bigr) \ot_{B} Q \to 
\opn{Hom}_{A} \bigl( F_{j}(P), M \ot_{B} Q \bigr) \]
is bijective. We conclude that 
\[ \opn{ev}_{P, M, Q} : \bigl( \opn{Hom}_{A}(P, M) \ot_{B} Q \bigr)^l \to 
\opn{Hom}_{A}(P, M \ot_{B} Q)^l \]
is bijective. Since $l$ is arbitrary, it follows that 
$\opn{ev}_{P, M, Q}$ is an isomorphism in $\dcat{C}_{\mrm{str}}(\K)$.

\medskip \noindent 
Step 3. In step 2 we proved that the morphism 
\[ \opn{ev}^{\mrm{R, L}}_{L, M, N} : \opn{RHom}_{A}(L, M) \ot^{\mrm{L}}_{B} N 
\to \opn{RHom}_{A}(L, M \ot^{\mrm{L}}_{B} N) \]
in $\dcat{D}(\K)$ is an isomorphism if $M$ is bounded below, 
$L = P$ is a pseudo-finite semi-free DG 
$A$-module, and $N = Q$ is K-flat DG $B$-module of bounded below tensor 
displacement. If $M$ is only assumed to satisfy condition (ii), then it is 
quasi-isomorphic to a suitable smart truncation of it which is bounded below; 
therefore $\opn{ev}^{\mrm{R, L}}_{L, M, N}$ is an isomorphism under this 
assumption too. 

Fixing $M$ and $N$, the full subcategory of $\dcat{D}(A)$
on the objects $L$ for which $\opn{ev}^{\mrm{R, L}}_{L, M, N}$ is an 
isomorphism is a saturated full triangulated subcategory (by Proposition 
\ref{prop:4603}). Therefore  
$\opn{ev}^{\mrm{R, L}}_{L, M, N}$ is an isomorphism for every $L, M, N$ 
such that $L$ satisfies condition (i), $M$ satisfies condition (ii), and 
$N = Q$ is K-flat of bounded below tensor displacement.

Now we fix $L$ and $M$ that satisfy condition (i) and (ii) respectively. 
The full subcategory of $\dcat{D}(B)$ on the objects $N$ for which 
$\opn{ev}^{\mrm{R, L}}_{L, M, N}$ is an 
isomorphism is a saturated full triangulated subcategory. By the previous 
paragraph, this subcategory contains the K-flat complexes $Q$ of bounded below 
tensor displacement. Therefore it contains all the complexes $N$ satisfying 
condition (iii).
\end{proof}


\begin{rem} \label{rem:4677}
Let us explain the syntax of the expression 
$\opn{ev}^{\mrm{R, L}}_{L, M, N}$
in Theorem \ref{thm:4320}. The letters ``ev'' stand for ``evaluation''.
The letters ``R'' and ``L'' in the superscript refer to the 
left and right derived functors, whereas the letters ``L'', ``M'' and ``N'' 
in the subscript refer to the DG modules involved. Note the different fonts. 
This sort of syntax is used frequently in the book; cf.\ the expressions 
$\opn{hm}^{\mrm{R}}_{M}$ in Definition \ref{dfn:2394},
$\opn{ev}^{\mrm{R, R}}$ in Lemma \ref{lem:2156},
$\ttev^{\mrm{R, L}}_{\m, (-)}$ in Theorem \ref{thm:3750} and
$\opn{fadj}_{u, M, N}^{\mrm{L}}$ in Theorem \ref{thm:3085}.
\end{rem}



\mysubsection{Hom-Tensor Formulas for Weakly Commutative DG Rings}
\label{subsec:adjs-comm}

In this subsection we assume that all DG rings are weakly commutative  
(see Definition \ref{dfn:3090}).
For such a DG ring $A$, its $0$-th cohomology 
$\opn{H}^0(A)$ is a commutative ring.
 
\begin{prop} \label{prop:3107}
For a weakly commutative DG ring $A$, the category $\dcat{D}(A)$ is 
$\opn{H}^0(A)$-linear. 
\end{prop}

\begin{proof}
This is true already for the homotopy category. 
Indeed, for any $M, N \in \dcat{C}(A)$ the 
$\K$-module 
$\opn{H}^0 \bigl( \opn{Hom}_A(M, N) \bigr)$
is an $\opn{H}^0(A)$-module; and composition in $\dcat{K}(A)$ is 
$\opn{H}^0(A)$-bilinear. 
\end{proof}

Let $A$ be a DG ring. There is a canonical isomorphism 
 $u : A \iso A^{\mrm{op}}$ to its opposite. The formula is 
$u(a) := (-1)^i \cd a$
for $a \in A^i$. This implies that every left DG $A$-module can be made into a 
right DG $A$-module, and vice-versa. The formula relating the left and right 
actions is 
$m \cd a = (-1)^{i \cd j} \cd a \cd m$
for $a \in A^i$ and $m \in M^j$. 
On the level of categories we obtain an isomorphism of DG categories 
$\dcat{C}(A) \cong \dcat{C}(A^{\mrm{op}})$,
which is the identity functor on the underlying DG $\K$-modules. 

Weak commutativity makes the tensor and Hom functors $A$-bilinear 
(in the graded sense), and therefore their derived functors have more 
structure: they are $\opn{H}^0(A)$-bilinear triangulated bifunctors 
\begin{equation} \label{eqn:2471}
(- \ot^{\mrm{L}}_A -) : \dcat{D}(A) \times \dcat{D}(A) \to \dcat{D}(A) 
\end{equation}
and 
\begin{equation} \label{eqn:2472}
\opn{RHom}_{A}(-, -) : \dcat{D}(A)^{\mrm{op}} \times \dcat{D}(A) \to 
\dcat{D}(A) . 
\end{equation}

When the DG rings in Theorem \ref{thm:2363} are weakly commutative, 
this result can be amplified: 

\begin{prop} \label{prop:3150} 
In the situation of Theorem \tup{\ref{thm:2363}}, assume that $A$ and $B$ 
are weakly commutative. 
\begin{enumerate}
\item The restriction functor 
$\opn{Rest}_u : \dcat{D}(B) \to \dcat{D}(A)$
is an equivalence of $\opn{H}^0(A)$-linear triangulated categories, with 
quasi-inverse $\opn{LInd}_u$. 

\item For every $M, N \in \dcat{D}(B)$ there is an isomorphism
$M \ot^{\mrm{L}}_A N \iso M \ot^{\mrm{L}}_B N$
in $\dcat{D}(A)$. This isomorphism is functorial in $M$ and $N$.

\item For every $M, N \in \dcat{D}(B)$ there is an isomorphism 
$\opn{RHom}_{A}(M, N) \iso \lb \opn{RHom}_{B}(M, N)$
in $\dcat{D}(A)$. This isomorphism is functorial in $M$ and $N$.
\end{enumerate}
\end{prop}

\begin{proof} \mbox{}

\smallskip \noindent 
(1) The ring homomorphism 
$\opn{H}^0(u) : \opn{H}^0(A) \to \opn{H}^0(B)$
makes $\dcat{D}(B)$ into an $\opn{H}^0(A)$-linear category, and then 
$\opn{Rest}_u : \dcat{D}(B) \to \dcat{D}(A)$
becomes an $\opn{H}^0(A)$-linear functor. 

\medskip \noindent 
(2,3) Going over the steps in the proof of Theorem \ref{thm:2363}, we see that 
all the moves are $A$-linear (in the graded sense). 
\end{proof}

Suppose $A \to B$ is a DG ring homomorphism. For every 
$M \in \dcat{C}(A)$ and $N \in \dcat{C}(B)$, their tensor product 
$M \ot_A N$ is a DG $B$-module, via the action on $N$. 
Thus we get a DG bifunctor 
\begin{equation} \label{eqn:4955}
(- \ot_A -) : \dcat{C}(A) \times \dcat{C}(B) \to \dcat{C}(B) . 
\end{equation}
There is an isomorphism 
\begin{equation} \label{eqn:3101}
M \ot_A N \iso N \ot_A M 
\end{equation}
in $\dcat{C}_{\mrm{str}}(B)$, that uses the Koszul sign rule
$m \ot n \mapsto (-1)^{i \cd j} \cd n \ot m$
for homo\-geneous elements $m \in M^i$ and $n \in N^j$. 

\begin{prop} \label{prop:3103}
Let $A \to B$ be a homomorphism of weakly commutative DG rings. 
The DG bifunctor \tup{(\ref{eqn:4955})}
has a left derived bifunctor 
\[ (- \ot^{\mrm{L}}_{A} -) : \dcat{D}(A) \times \dcat{D}(B) \to \dcat{D}(B) . \]
This is an $\opn{H}^0(A)$-bilinear triangulated bifunctor.
\end{prop}

\begin{proof}
The left derived bifunctor exists by Theorem \ref{thm:2106}, using K-flat 
resolutions in the first argument (i.e.\ in $\dcat{C}_{\mrm{str}}(A)$). 
As for the $\opn{H}^0(A)$-bilinearity: this is already comes in the homotopy 
category level, namely for the bifunctor 
\[ (- \ot_{A} -) : \dcat{K}(A) \times \dcat{K}(B) \to \dcat{K}(B) . \qedhere \]
\end{proof}

Similarly we have an $\opn{H}^0(A)$-bilinear triangulated bifunctor
\[ (- \ot^{\mrm{L}}_{A} -) : \dcat{D}(B) \times \dcat{D}(A) \to \dcat{D}(B) . \]
When $A = B$ these operations coincide with (\ref{eqn:2471}). 

\begin{prop}[Symmetry] \label{prop:3102}
Let $A \to B$ be a homomorphism of weakly commutative DG rings. For every 
$M \in \dcat{D}(A)$ and $N \in \dcat{D}(B)$ there is an isomorphism 
$M \ot^{\mrm{L}}_{A} N \cong N \ot^{\mrm{L}}_{A} M$
in $\dcat{D}(B)$. This isomorphisms is functorial in $M$ and $N$.  
\end{prop}

\begin{proof}
Use the isomorphism (\ref{eqn:3101}). 
\end{proof}

\begin{prop}[Associativity] \label{prop:3100}
Let $A_1 \to A_2 \to A_3$ be homomorphisms between weakly commutative DG rings.
For each $i$ let $M_i \in \dcat{D}(A_i)$.
There is an isomorphism 
\[ (M_1 \ot^{\mrm{L}}_{A_1} M_2) \ot^{\mrm{L}}_{A_2} M_3 \cong 
M_1 \ot^{\mrm{L}}_{A_1} (M_2 \ot^{\mrm{L}}_{A_2} M_3) \]
in $\dcat{D}(A_3)$.
This isomorphisms is functorial in the $M_i$. 
\end{prop}

\begin{proof}
Choose K-flat resolutions $P_i \to M_i$ in $\dcat{C}(A_i)$. Then we have 
isomorphisms
\[ \begin{aligned}
& (M_1 \ot^{\mrm{L}}_{A_1} M_2) \ot^{\mrm{L}}_{A_2} M_3 
\cong (P_1 \ot_{A_1} P_2) \ot_{A_2} P_3
\\
& \qquad 
\cong^{\dag} P_1 \ot_{A_1} (P_2 \ot_{A_2} P_3)
\cong 
M_1 \ot^{\mrm{L}}_{A_1} (M_2 \ot^{\mrm{L}}_{A_2} M_3)  
\end{aligned} \]
in $\dcat{D}(A_3)$. The isomorphism $\cong^{\dag}$ is due to the usual 
associativity of $\ot$.
\end{proof}

Suppose $A \to B$ is a DG ring homomorphism. For every 
$M \in \dcat{C}(A)$ and $N \in \dcat{C}(B)$, the DG $\K$-module
$\opn{Hom}_{A}(N, M)$ has a DG $B$-module structure, via the action on $N$. 
Thus we get a DG bifunctor 
\begin{equation} \label{eqn:4956}
\opn{Hom}_{A}(-, -) : \dcat{C}(B)^{\mrm{op}} 
\times \dcat{C}(A) \to \dcat{C}(B) .
\end{equation}

\begin{prop} \label{prop:3105}
Let $A \to B$ be a homomorphism of weakly commutative DG rings. 
The DG bifunctor \tup{(\ref{eqn:4956})}
has a right derived bifunctor
\[ \opn{RHom}_{A}(-, -) : \dcat{D}(B)^{\mrm{op}} 
\times \dcat{D}(A) \to \dcat{D}(B) . \]
This is an $\opn{H}^0(A)$-bilinear triangulated bifunctor.
\end{prop}

\begin{proof}
The right derived bifunctor exists by Theorem \ref{thm:2000}, using K-injective
resolutions in the second argument (i.e.\ in $\dcat{C}_{\mrm{str}}(A)$). 
\end{proof}

\begin{prop}[Hom-Tensor Adjunction] \label{prop:3106}
Let $A_1 \to A_2 \to A_3$ be homomorphisms between weakly commutative DG rings.
For each $i$ let $M_i \in \dcat{D}(A_i)$.
There is an isomorphism 
\[ \opn{RHom}_{A_2} \bigl( M_3, \opn{RHom}_{A_1}(M_2, M_1) \bigr) \cong
\opn{RHom}_{A_1} \bigl( M_2 \ot^{\mrm{L}}_{A_2} M_3, M_1) \]
in $\dcat{D}(A_3)$.
This isomorphisms is functorial in the $M_i$. 
\end{prop}

\begin{proof}
Choose a K-flat resolution $P_2 \to M_2$ in $\dcat{C}(A_2)$,
and a K-injective resolution $M_1 \to I_1$ in $\dcat{C}(A_1)$. 
An easy adjunction calculation show that 
$\opn{Hom}_{A_1}(P_2, I_1)$ is K-injective in $\dcat{C}(A_2)$. 
Then we have isomorphisms
\[ \begin{aligned}
& \opn{RHom}_{A_2} \bigl( M_3, \opn{RHom}_{A_1}(M_2, M_1) \bigr)
\cong
\opn{Hom}_{A_2} \bigl( M_3, \opn{Hom}_{A_1}(P_2, I_1) \bigr)
\\
& \qquad 
\cong^{\dag} 
\opn{Hom}_{A_1} (P_2 \ot_{A_2} M_3, I_1) \cong 
\opn{RHom}_{A_1} \bigl( M_2 \ot^{\mrm{L}}_{A_2} M_3, M_1)
\end{aligned} \]
in $\dcat{D}(A_3)$. The isomorphism $\cong^{\dag}$ is due to the usual 
Hom-tensor adjunction. 
\end{proof}

\begin{rem} \label{rem:3100}
It is not hard to show that $\dcat{D}(A)$ is a closed symmetric monoidal 
category, with monoidal operation $(- \ot^{\mrm{L}}_{A} -)$ and monoidal unit 
object $A$. It needs just a bit more than Propositions \ref{prop:3100}, 
\ref{prop:3102} and \ref{prop:3106}. 

For a noncommutative version of this monoidal structure see 
Remark \ref{rem:3450} and Definition \ref{dfn:4530}.
\end{rem}

We end this subsection with a weakly-commutative variant of Theorem 
\ref{thm:4320}. Cohomologically pseudo-noether\-ian nonpositive DG rings were 
introduced in Definition \ref{dfn:3180}. 
DG modules with derived bounded below tensor displacement were defined in
Definition \ref{dfn:4341}.

\begin{thm} \label{thm:2700} 
Let $A \to B$ be a homomorphism between weakly commutative DG rings, and let
$L \in \dcat{D}(A)$ and $M, N \in \dcat{D}(B)$ be DG modules. 
There is a morphism 
\[ \opn{ev}^{\mrm{R, L}}_{L, M, N} : \opn{RHom}_{A}(L, M) \ot^{\mrm{L}}_B N \to 
\opn{RHom}_{A}(L, M \ot^{\mrm{L}}_B N) \]
in $\dcat{D}(B)$, called {\em derived tensor-evaluation}% 
\index{Tensor-evaluation morphism! derived},
which is functorial in these DG modules. Moreover, if conditions 
\tup{(i)-(iii)} below hold, then $\opn{ev}^{\mrm{R, L}}_{L, M, N}$ is an 
isomorphism. 
\begin{enumerate}
\rmitem{i} The DG ring $A$ is nonpositive and cohomologically 
pseudo-noetherian, and $L \in \dcat{D}^{-}_{\mrm{f}}(A)$.

\rmitem{ii} The DG $B$-module $M$ is in $\dcat{D}^{+}(B)$.

\rmitem{iii}  The DG $B$-module $N$ has derived bounded below tensor 
displacement.
\end{enumerate}
\end{thm}

\begin{proof}
Choose a K-projective resolution 
$P \to L$ in $\dcat{C}_{\mrm{str}}(A)$ and a K-flat resolution 
$Q \to N$ in $\dcat{C}_{\mrm{str}}(B)$. 
Consider the obvious homomorphism 
\[ \opn{ev}_{P, M, Q} : \opn{Hom}_{A}(P, M) \ot_{B} Q \to 
\opn{Hom}_{A}(P, M \ot_{B} Q) \]
in $\dcat{C}_{\mrm{str}}(B)$. 
We take 
$\opn{ev}^{\mrm{R, L}}_{L, M, N} := \opn{Q}(\opn{ev}_{P, M, Q})$
in $\dcat{D}(B)$. The morphism $\opn{ev}^{\mrm{R, L}}_{L, M, N}$ does not 
depend on the choices, and hence it is functorial. 

Now assume that conditions (i)-(iii) hold. Because the restriction functor \lb 
$\dcat{D}(B) \to \dcat{D}(\K)$ is conservative, it is enough to show that 
$\opn{ev}^{\mrm{R, L}}_{L, M, N}$ is an isomorphism in $\dcat{D}(\K)$.
According to Theorem \ref{thm:3340}, $L$ is a derived pseudo-finite DG 
$A$-module (cf.\ Example \ref{exa:4471}). Now we can use Theorem 
\ref{thm:4320}.
\end{proof}











